\documentclass[11pt,twoside,openany]{book}
\pdfvariable suppressoptionalinfo 767
\usepackage{graphicx}
\usepackage[a4paper,inner=28mm,outer=25mm,top=25mm,bottom=28mm,marginparwidth=15mm]{geometry}
\usepackage{fontspec}
\defaultfontfeatures{Ligatures=TeX,Scale=MatchLowercase}
\setmainfont{TeX Gyre Pagella}
\setsansfont{TeX Gyre Heros}
\setmonofont{TeX Gyre Cursor}
\usepackage{polyglossia}
\setdefaultlanguage{vietnamese}
\usepackage{microtype}
\usepackage{mathtools}
\usepackage{amsthm}
\usepackage{unicode-math}
\setmathfont{Latin Modern Math}
\usepackage{bookmark}
\usepackage{hyperref}
\hypersetup{hidelinks,unicode=true,hypertexnames=false}

\setcounter{secnumdepth}{3}
\setcounter{tocdepth}{2}
\numberwithin{equation}{section}

\newcommand{\src}[1]{%
  \phantomsection\hypertarget{src:\fgasourceid:#1}{}%
  \marginpar{\raggedleft\footnotesize\sffamily #1}%
}
\newcommand{\fgasourceid}{unset}
\newcommand{\sid}[1]{\gdef\fgasourceid{#1}\phantomsection\hypertarget{fga:#1}{}}
\newcommand{\Osh}{\underline{O}}
\newcommand{\Zsh}{\underline{\mathbf Z}}
\newcommand{\Ccat}{\mathcal C}
\newcommand{\Homsh}{\underline{\operatorname{Hom}}}
\newcommand{\Extsh}{\underline{\operatorname{Ext}}}
\newcommand{\hsh}{\underline{h}}
\newcommand{\fquote}[1]{\textquotedbl#1\textquotedbl}
\newcommand{\seminar}[1]{%
  \begin{flushleft}\sffamily HỘI THẢO BOURBAKI\\#1\end{flushleft}%
}
\newcommand{\paperhead}[2]{%
  \begin{center}
    {\Large\bfseries #1\par}
    \vspace{0.6\baselineskip}
    #2
  \end{center}
  \vspace{0.8\baselineskip}
}
\newcommand{\unit}[1]{%
  \chapter{#1}%
  \markboth{#1}{#1}%
}
\newcommand{\fsection}[2]{%
  \section*{#1. #2}%
  \addcontentsline{toc}{section}{#1. #2}%
}

\newcommand{\result}[2]{\par\medskip\noindent\textbf{#1\if\relax\detokenize{#2}\relax\else\ #2\fi.}\quad}
\setlength{\emergencystretch}{3em}
\hypersetup{
  pdftitle={Nền tảng hình học đại số -- Bản dịch tiếng Việt theo bản in},
  pdfauthor={Alexander Grothendieck}
}
\newcommand{\editiondescription}{Bản dịch theo bản in giữ các cách đọc của nguồn và tách biệt với bản hiệu đính.}
\begin{document}
\frontmatter
\begin{titlepage}
\centering\vspace*{2cm}
{\Large Alexandre Grothendieck\par}\vspace{1.5cm}
{\Huge\bfseries Nền tảng hình học đại số\par}\vspace{1.2cm}
{\Large Bản dịch tiếng Việt theo bản in\par}
\vfill
{Bourbaki 149, 182, 190, 195, 212, 221, 232, 236\par}
{Đính chính, bổ sung và bình luận năm 1962\par}
\end{titlepage}
\chapter*{Tác phẩm, nguồn và quyền}
\addcontentsline{toc}{chapter}{Tác phẩm, nguồn và quyền}
\markboth{Tác phẩm, nguồn và quyền}{Tác phẩm, nguồn và quyền}
\textit{Nền tảng hình học đại số} của Alexandre Grothendieck gồm các bài trình bày
Bourbaki số 149, 182, 190, 195, 212, 221, 232, 236, các đính chính và bổ sung
của tám bài, cùng bình luận năm 1962. Số trang của nguồn được ghi ở lề.
\editiondescription

Bản dịch tiếng Việt do AI thực hiện; hiệu đính, biên tập và dàn trang cũng do AI
thực hiện. Người xuất bản không dịch văn bản. Bản tiếng Pháp là nguồn chính,
có thể tra cứu tại \href{https://www.numdam.org/}{NUMDAM}. Bản
\href{https://github.com/KokunoYumeto/fga-en}{tiếng Anh} được dùng để đối chiếu.
Nguồn tiếng Anh có lời văn kế thừa từ bản dịch của \href{https://github.com/thosgood/fga}{Tim Hosgood}.
Phiên bản nguồn được xác định là
\href{https://github.com/thosgood/fga/tree/9976dd5f8e2ddcb388f0a07f3aa334f120d1019f}{\nolinkurl{9976dd5f8e2ddcb388f0a07f3aa334f120d1019f}}.
Tim Hosgood cũng ghi nhận Steve Hnizdur về các sửa chữa và Jon Sterling về hỗ trợ kỹ thuật.
Các đóng góp này thuộc bản tham khảo tiếng Anh; bản dịch tiếng Việt và những
hiệu đính riêng do AI thực hiện không được gán cho các tác giả ấy.

Quyền đối với nguyên tác và mọi đóng góp được bảo hộ trong các bản dịch, bản
chép lại phái sinh vẫn thuộc về các chủ thể quyền tương ứng. Tác phẩm không được
đưa vào phạm vi công cộng và không được cấp giấy phép mới ở đây. Các thông báo
quyền và những sự cho phép đã có vẫn có hiệu lực. Ghi nguồn không đồng nghĩa với
được cho phép hay được chứng thực.

Việc cung cấp miễn phí cho mục đích giáo dục, phi thương mại dựa trên giả định
rằng các chủ thể quyền không phản đối. Đây không phải tuyên bố đã được cho phép;
quyền truy cập miễn phí bản gốc không tự nó cho phép dịch hoặc tái xuất bản.
Chủ thể quyền phản đối có thể \href{https://github.com/KokunoYumeto/ega-fga-sga-vi/issues}{mở một issue trên GitHub}.
Người xuất bản cam kết gỡ bỏ tài liệu bị phản đối theo yêu cầu của chủ thể quyền.
Không hàm ý sự chứng thực của tác giả, nhà xuất bản nguồn hay những người được dẫn.
\clearpage

\tableofcontents
\mainmatter
\unit{Bài trình bày 149 — Các định lý đối ngẫu cho các bó đại số coherent}
\sid{149}
\markboth{Thuyết trình 149}{Thuyết trình 149}
\src{169}
\seminar{(Tháng 5 năm 1957)}
\paperhead{CÁC ĐỊNH LÝ ĐỐI NGẪU CHO CÁC BÓ ĐẠI SỐ COHERENT}%
  {do Alexandre GROTHENDIECK}

Những kết quả sau đây, được gợi ý bởi \fquote{định lý đối ngẫu đại số}
của Serre, đã được tìm ra vào mùa đông năm 1955 và mùa đông năm 1956. Chúng được chứng minh
rất đơn giản nhờ những kết quả khá sơ cấp về đối đồng điều của các không gian
xạ ảnh [3], cùng với việc sử dụng rộng rãi đại số đồng điều của
Cartan--Eilenberg theo dạng trình bày trong [2].

\fsection{1}{Các Ext của các bó môđun ({[2]}, Chương 3 và 4)}

Cho $X$ là một không gian tôpô được trang bị một bó $\Osh$ các vành có đơn vị
(không nhất thiết giao hoán). Ta xét phạm trù Abel
$\Ccat^{\Osh}$ của các bó $\Osh$-môđun, cũng được gọi là $\Osh$-Môđun.
Mọi vật của phạm trù này đều có một phân giải nội xạ; nhờ đó
có thể định nghĩa trong phạm trù này các hàm tử $\operatorname{Ext}$ với những tính chất
hình thức quen thuộc. Cụ thể hơn, để tránh nhầm lẫn, ta
ký hiệu $\operatorname{Hom}_{\Osh}(X;A,B)$, hoặc đơn giản là
$\operatorname{Hom}(X;A,B)$, là nhóm Abel các $\Osh$-đồng cấu từ $A$
vào $B$, trong khi $\Homsh_{\Osh}(A,B)$ ký hiệu bó các mầm
đồng cấu từ $A$ vào $B$ ($A,B\in\Ccat^{\Osh}$). Với
$A\in\Ccat^{\Osh}$ cố định, ta định nghĩa các hàm tử $h_A$ và $\hsh_A$ nhận giá trị tương ứng
trong phạm trù $\Ccat$ các nhóm Abel và phạm trù
$\Ccat^{\Zsh}$ các bó nhóm Abel trên $X$, bằng các công thức:
\begin{equation}\tag{1.1}
  h_A(B)=\operatorname{Hom}_{\Osh}(X;A,B),
  \qquad
  \hsh_A(B)=\Homsh_{\Osh}(A,B).
\end{equation}

Các hàm tử $h_A$ và $\hsh_A$ là những hàm tử hiệp biến khớp trái;
ta xét các hàm tử dẫn xuất phải của chúng, được ký hiệu tương ứng bởi
$\operatorname{Ext}_{\Osh}^{p}(X;A,B)$ và $\Extsh_{\Osh}^{p}(A,B)$. Do đó,
theo định nghĩa,
\begin{equation}\tag{1.2}
\left\{
\begin{aligned}
\operatorname{Ext}_{\Osh}^{p}(X;A,B)
  &= (R^{p}h_A)(B)
   = H^{p}\bigl(\operatorname{Hom}_{\Osh}(X;A,C(B))\bigr),\\
\Extsh_{\Osh}^{p}(A,B)
  &= (R^{p}\hsh_A)(B)
   = H^{p}\bigl(\Homsh_{\Osh}(A,C(B))\bigr),
\end{aligned}
\right.
\end{equation}

trong đó ký hiệu $R^{p}$ chỉ việc chuyển sang các hàm tử dẫn xuất phải, còn
$C(B)$ ký hiệu một phân giải nội xạ tùy ý của $B$ trong $\Ccat^{\Osh}$.
Ký hiệu
\[
  \Gamma:\Ccat^{\Zsh}\longrightarrow\Ccat
\]
là hàm tử \fquote{tiết diện}. Nhắc lại rằng các hàm tử dẫn xuất phải của nó được
ký hiệu bởi
\[
  B\longmapsto H^{p}(X,B).
\]


\src{170}
\begin{equation}\tag{1.3}
  H^{p}(X,B)=(R^{p}\Gamma)(B)=H^{p}\bigl(\Gamma(C(B))\bigr).
\end{equation}
Hiển nhiên, ta có $h_A=\Gamma\hsh_A$; mặt khác, ta kiểm tra được rằng $\hsh_A$
biến các vật nội xạ thành các vật $\Gamma$-acyclic. Từ đó, theo cách quen thuộc,
ta suy ra:

\result{MỆNH ĐỀ}{1} Với mọi $\Osh$-môđun $A$, tồn tại một hàm tử
phổ đối đồng điều trên $\Ccat^{\Osh}$, hội tụ tới hàm tử phân bậc
$\bigl(\operatorname{Ext}_{\Osh}^{*}(X;A,B)\bigr)$ và có số hạng ban đầu là
\begin{equation}\tag{1.4}
  E_{2}^{p,q}(A,B)=H^{p}\bigl(X,\Extsh_{\Osh}^{q}(A,B)\bigr).
\end{equation}
Từ đó, ta thu được các \fquote{đồng cấu biên} và một dãy khớp năm số hạng
mà ta sẽ không viết ra.

\result{HỆ QUẢ}{1} Nếu $A$ đẳng cấu địa phương với $\Osh^{n}$, thì ta có
các đẳng cấu chính tắc
\begin{equation}\tag{1.5}
  \operatorname{Ext}_{\Osh}^{p}(X;A,B)
  \overset{\sim}{\longleftarrow}
  H^{p}\bigl(X,\Homsh_{\Osh}(A,B)\bigr)
\end{equation}
(do các đồng cấu biên của dãy phổ cho). Đặc biệt, ta có
một đẳng cấu chính tắc
\begin{equation}\tag{1.6}
  \operatorname{Ext}_{\Osh}^{p}(X;\Osh,B)=H^{p}(X,B).
\end{equation}

Để sử dụng các kết quả này, cần mô tả tường minh các
$\Extsh_{\Osh}^{p}(A,B)$. Các hàm tử này được tính địa phương;
nghĩa là, nếu $U$ là một tập mở của $X$, thì
\[
  \Extsh_{\Osh}^{p}(A,B)|U
  =\Extsh_{\Osh|U}^{p}(A|U,B|U),
\]
vì phép hạn chế lên $U$ của một $\Osh$-môđun nội xạ
là một $(\Osh|U)$-môđun nội xạ. Hơn nữa, với $x\in X$ cố định, ta có các đồng cấu có tính hàm tử
\begin{equation}\tag{1.7}
  \Homsh_{\Osh}(A,B)_x
  \longrightarrow
  \operatorname{Hom}_{\Osh_x}(A_x,B_x),
\end{equation}
chúng mở rộng duy nhất thành một đồng cấu giữa các hàm tử đối đồng điều
theo $B$:
\begin{equation}\tag{1.8}
  \Extsh_{\Osh}^{p}(A,B)_x
  \longrightarrow
  \operatorname{Ext}_{\Osh_x}^{p}(A_x,B_x).
\end{equation}

\result{MỆNH ĐỀ}{2} Nếu $A$, trong một lân cận của $x$, đẳng cấu với
đối hạch của một đồng cấu $\Osh^{m}\longrightarrow\Osh^{n}$, thì (1.7) là một
đẳng cấu với mọi $p$.
Điều này đúng,
nói riêng, nếu $A$ là một $\Osh$-môđun coherent [3].

\src{171}
\result{MỆNH ĐỀ}{3} Cho $L_{*}=(L_i)$ là một phân giải trái của
$\Osh$-môđun $A$ bởi các $\Osh$-môđun, mỗi môđun đẳng cấu địa phương với
một $\Osh^{n}$. Khi đó $\Extsh_{\Osh}^{*}(A,B)$ được đồng nhất với
$H^{*}\bigl(\Homsh_{\Osh}(L_{*},B)\bigr)$, và
$\operatorname{Ext}_{\Osh}^{*}(X;A,B)$ được đồng nhất với siêu đối đồng điều của $X$
đối với phức $\Homsh_{\Osh}(L_{*},B)$.

Chứng minh là chuẩn: ta xét song phức
$\Homsh_{\Osh}(L_{*},C(B))$, trong đó $C(B)$ là một phân giải nội xạ của $B$, cùng các
đồng cấu tự nhiên từ $\Homsh_{\Osh}(L_{*},B)$ và từ
$\Homsh_{\Osh}(A,C(B))$ vào song phức này.

Cuối cùng, lưu ý rằng hai hàm tử $\operatorname{Ext}$ được đưa vào
trong (1.2) không chỉ là các hàm tử đối đồng điều theo $B$, mà còn là các
song hàm tử đối đồng điều, hiệp biến theo $B$ và phản biến theo $A$.

\fsection{2}{Luật hợp thành trong các Ext}

Các kết quả của mục này do CARTIER và YONEDA tìm ra độc lập; xem
một bài trình bày của CARTIER [1] để biết chi tiết. Cho $\Ccat$ là một phạm trù Abel.
Cho $K$ và $L$ là hai vật phân bậc trong $\Ccat$. Ta ký hiệu
$\operatorname{Hom}(K,L)$ là nhóm Abel phân bậc mà thành phần bậc
$n$ gồm các đồng cấu thuần nhất bậc $n$ từ $K$ vào $L$ (nghĩa là
các hệ $(u_i)$ đồng cấu $K^{i}\longrightarrow L^{i+n}$).

Nếu $K$ và $L$ là các phức (với các toán tử vi phân bậc $+1$, để
cụ thể hóa), ta trang bị cho $\operatorname{Hom}(K,L)$ toán tử
vi phân cho bởi
\begin{equation}\tag{2.1}
  \delta(u)=du+(-1)^{n+1}ud,\qquad\text{trong đó }n=\deg(u),
\end{equation}
qua đó biến nó thành một phức với toán tử vi phân bậc $+1$. Các chu trình
bậc $n$ là những toán tử bậc $n$ phản giao hoán với $u$ (xét như
các toán tử thuần nhất). Khi đó ta có thể xét
$H^{*}\bigl(\operatorname{Hom}(K,L)\bigr)$; đây là một bất biến của các kiểu
đồng luân của $K$ và $L$, và ta có thể ký hiệu nó là $H^{*}(K,L)$.

Nếu có thêm một phức thứ ba $M$, thì phép hợp thành các đồng cấu
xác định một phép ghép cặp
\[
  \operatorname{Hom}(K,L)\times\operatorname{Hom}(L,M)
  \longrightarrow\operatorname{Hom}(K,M)
\]
tương thích với các toán tử vi phân; do đó, khi chuyển sang đối đồng điều,
ta được các phép ghép cặp
\begin{equation}\tag{2.2}
  H^{*}(K,L)\times H^{*}(L,M)\longrightarrow H^{*}(K,M),
\end{equation}
được ký hiệu $(u,v)\longmapsto vu$. Các phép ghép cặp này thỏa mãn tính
kết hợp hiển nhiên. Đặc biệt, $H^{*}(K,K)$ là một vành phân bậc kết hợp
có đơn vị, và $H^{*}(K,L)$ (tương ứng $H^{*}(L,K)$) là một môđun phân bậc phải
(tương ứng trái) trên vành này, v.v. Ở bậc $0$, (2.2) thu về
phép hợp thành các đồng cấu của phức. Cuối cùng, một dãy khớp các
phức
\[
  0\longrightarrow K'\longrightarrow K\longrightarrow K''\longrightarrow 0,
\]

\src{172}
sao cho $K'^{i}$ được đồng nhất với một hạng tử trực tiếp của $K^{i}$ với mọi $i$,
dẫn đến một dãy khớp các phức nhóm $\operatorname{Hom}(K'',L)$, v.v., và do đó dẫn đến một
toán tử nối
\[
  H^{i}(K',L)\longrightarrow H^{i+1}(K'',L).
\]
Các toán tử nối tương ứng với một dãy khớp trong $L$ cũng được định nghĩa tương tự. Các
phép ghép cặp (2.2) tương thích, theo nghĩa thông thường, với các toán tử nối này.

Giả sử bây giờ $\Ccat$ là một phạm trù trong đó mọi vật $A$ của $\Ccat$
đều có một phân giải nội xạ $C(A)$. Khi đó, dùng một trong nhiều biến thể của định lý song phức, ta có
\[
  H^{*}\bigl(C(A),C(B)\bigr)
  =H^{*}\bigl(\operatorname{Hom}(C(A),C(B))\bigr)
\]
là đẳng cấu chính tắc với
\[
  H^{*}\bigl(\operatorname{Hom}(A,C(B))\bigr)
  =\operatorname{Ext}^{*}(A,B).
\]
Các toán tử nối được xét ở trên cho các đồng cấu nối của
$\operatorname{Ext}$. Hơn nữa, các phép ghép cặp (2.2) ở đây cho
các phép ghép cặp kết hợp, tương thích với các đồng cấu nối:
\begin{equation}\tag{2.3}
  \operatorname{Ext}^{*}(A,B)\times\operatorname{Ext}^{*}(B,C)
  \longrightarrow\operatorname{Ext}^{*}(A,C).
\end{equation}
Đặc biệt, $\operatorname{Ext}^{*}(A,A)$ là một vành phân bậc kết hợp
có đơn vị, v.v. (Tương tự, người ta chứng minh rằng các hàm tử
$\operatorname{Ext}$ tác động trên các hàm tử dẫn xuất của một hàm tử
bất kỳ; ở đây chúng ta sẽ không dùng đến sự kiện này.)

Trong trường hợp phạm trù đang xét là phạm trù $\Ccat^{\Osh}$ gồm các
$\Osh$-môđun trên $X$, do đó ta thu được các phép ghép cặp
\begin{equation}\tag{2.4}
  \operatorname{Ext}_{\Osh}^{p}(X;A,B)
  \times\operatorname{Ext}_{\Osh}^{q}(X;B,C)
  \longrightarrow\operatorname{Ext}_{\Osh}^{p+q}(X;A,C),
\end{equation}
có thể được viết tường minh như đã chỉ ra ở trên. Hơn nữa, bằng những triển khai tương tự,
trong đó thay phạm trù các nhóm Abel bằng phạm trù các bó nhóm Abel trên $X$, và thay các hàm tử
$\operatorname{Hom}$ bằng các hàm tử $\Homsh$, ta cũng định nghĩa được các phép ghép cặp có cùng
các tính chất hình thức và lần này mang \fquote{tính địa phương}:
\begin{equation}\tag{2.5}
  \Extsh_{\Osh}^{p}(A,B)\times\Extsh_{\Osh}^{q}(B,C)
  \longrightarrow\Extsh_{\Osh}^{p+q}(A,C).
\end{equation}
Các phép ghép cặp sau được làm rõ hơn nếu nhận thấy rằng các đồng cấu (1.8) tương thích
với các phép ghép cặp giữa các nhóm $\operatorname{Ext}$.

\src{173}
Cuối cùng, ta nhắc lại rằng còn có một cấu trúc nhân giữa các hàm tử
$H^{p}(X,A)$ (tích cup). Khi đó, ta thấy rằng các dãy phổ của
Mệnh đề 1 tương thích với các cấu trúc nhân; cụ thể hơn,
ta có một phép ghép cặp từ dãy phổ $E(A,B)$ và dãy phổ
$E(B,C)$ vào dãy phổ $E(A,C)$; tại đích hội tụ, phép này
thu về phép ghép cặp giữa các nhóm $\operatorname{Ext}$ toàn cục, còn tại
số hạng đầu thì đó là phép ghép cặp được suy ra, ở vế phải của (1.4), từ
tích cup và các phép ghép cặp của những nhóm $\operatorname{Ext}$ địa phương. Đặc biệt, từ đó suy ra
rằng các \fquote{đồng cấu biên}
\begin{equation}\tag{2.6}
  \operatorname{Ext}_{\Osh}^{n}(X;A,B)
  \longrightarrow H^{0}\bigl(X,\Extsh_{\Osh}^{n}(A,B)\bigr)
\end{equation}
và
\begin{equation}\tag{2.7}
  H^{n}\bigl(X,\Homsh_{\Osh}(A,B)\bigr)
  \longrightarrow\operatorname{Ext}_{\Osh}^{n}(X;A,B)
\end{equation}
tương thích với các cấu trúc nhân. Vì vậy, nếu chỉ xét
các bó đẳng cấu địa phương với một $\Osh^{m}$, thì điều này mô tả hoàn toàn
phép hợp thành giữa các nhóm $\operatorname{Ext}$ toàn cục thông qua tích cup, theo
đẳng cấu (1.5).

\fsection{3}{Các kết quả về đối đồng điều địa phương}

Cho $A$ là một vành giao hoán có đơn vị, cùng với một iđêan $J$. Với
$A$-môđun biến thiên $M$, ta sẽ định nghĩa các đồng cấu có tính hàm tử
\begin{equation}\tag{3.1}
\left\{
\begin{aligned}
\operatorname{Ext}_{A}^{p}(A/J,M)
  &\longrightarrow
  \operatorname{Hom}_{A}\!\left(\bigwedge^{p}J/J^{2},M\otimes A/J\right),\\
\operatorname{Tor}_{p}^{A}(A/J,M)
  &\longleftarrow
  \left(\bigwedge^{p}J/J^{2}\right)
  \otimes\operatorname{Hom}_{A}(A/J,M),
\end{aligned}
\right.
\end{equation}
trong đó các tích tenxơ và tích ngoài đều được lấy trên vành $A$; cũng lưu ý
rằng $J/J^{2}$ thực ra là một $A/J$-môđun, và các lũy thừa
ngoài của nó, xét như $A$-môđun hay như $A/J$-môđun, là như nhau.
Việc định nghĩa các đồng cấu (3.1) tương đương với việc định nghĩa, cho mọi
hệ $\underline{x}=(x_1,\ldots,x_p)$ gồm các phần tử của $J$, các đồng cấu
$\varphi_{\underline{x}}$:
\begin{equation}\tag{3.2}
\left\{
\begin{aligned}
\varphi_{\underline{x}}:
  \operatorname{Ext}_{A}^{p}(A/J,M)&\longrightarrow M\otimes A/J,\\
\varphi_{\underline{x}}:
  \operatorname{Hom}_{A}(A/J,M)&\longrightarrow
  \operatorname{Tor}_{p}^{A}(A/J,M),
\end{aligned}
\right.
\end{equation}
sao cho các điều kiện sau được thỏa mãn:


\src{174}
\begin{enumerate}
\item[\textup{i.}] $\varphi_{x_1,\ldots,x_p}$ phụ thuộc theo cách
$A$-đa tuyến tính thay phiên vào hệ các $x_i\in J$.
\item[\textup{ii.}] $\varphi_{x_1,\ldots,x_p}$ bằng không khi một trong các $x_i$ thuộc
$J^{2}$.
\end{enumerate}
Thực ra, ii. suy ra từ i., vì $a\varphi_{\underline{x}}=0$ với $a\in J$,
như thấy được khi lưu ý rằng mọi môđun trong (3.2) đều bị
$J$ triệt tiêu.

Để định nghĩa $\varphi_{\underline{x}}$, ta xét phức
$K_{\underline{x}}$ có $A$-môđun nền là $\bigwedge A^{p}$, và
có toán tử vi phân là phép co $i_{\underline{x}}$ theo
$\underline{x}$ được xem như một dạng tuyến tính trên $A^{p}$ có các thành phần
$x_1,\ldots,x_p$. Toán tử vi phân có bậc $-1$, các bậc đều
không âm, và nhóm đối đồng điều của phức này ở bậc $0$ là
$A/(x_1A+\cdots+x_pA)$. Vì các $x_i$ thuộc $J$, từ đó suy ra một
đồng cấu tăng cường $K_{\underline{x},0}\longrightarrow A/J$. Do đó,
$K_{\underline{x}}$ xuất hiện như một phức tự do tăng cường, với môđun
tăng cường $A/J$. Theo cách quen thuộc, từ đó suy ra các đồng cấu
\[
\operatorname{Ext}_{A}^{*}\bigl(H_0(K_{\underline{x}}),M\bigr)
\longrightarrow
H^{*}\bigl(\operatorname{Hom}_{A}(K_{\underline{x}},M)\bigr)
\]
và
\[
\operatorname{Tor}_{*}^{A}\bigl(H_0(K_{\underline{x}}),M\bigr)
\longleftarrow
H_{*}\bigl(K_{\underline{x}}\otimes_A M\bigr),
\]
do đó, bằng cách hợp thành với các đồng cấu trên các $\operatorname{Ext}$ và các
$\operatorname{Tor}$ suy ra từ đồng cấu tăng cường
$H_0(K_{\underline{x}})\longrightarrow A/J$, ta thu được các đồng cấu
\begin{equation}\tag{3.3}
\left\{
\begin{aligned}
\psi_{\underline{x}}:
\operatorname{Ext}_{A}^{*}(A/J,M)
&\longrightarrow H^{*}\bigl(\operatorname{Hom}_{A}(K_{\underline{x}},M)\bigr),\\
\psi_{\underline{x}}:
\operatorname{Tor}_{*}^{A}(A/J,M)
&\longleftarrow H_{*}\bigl(K_{\underline{x}}\otimes_A M\bigr).
\end{aligned}
\right.
\end{equation}
Ta kiểm tra ngay rằng, ở bậc cao nhất $p$, các nhóm (đối) đồng điều ở các vế
phải lần lượt là $M/(x_1M+\cdots+x_pM)$ (tương ứng, là tập các phần tử của $M$
bị từng $x_i$ triệt tiêu). Vì các $x_i$ thuộc $J$, từ đó suy ra các
đồng cấu
\begin{equation}\tag{3.4}
\left\{
\begin{aligned}
H^{p}\bigl(\operatorname{Hom}_{A}(K_{\underline{x}},M)\bigr)
&\longrightarrow M\otimes A/J,\\
H_{p}\bigl(K_{\underline{x}}\otimes_A M\bigr)
&\longleftarrow\operatorname{Hom}_{A}(A/J,M).
\end{aligned}
\right.
\end{equation}
Bằng cách hợp thành các đồng cấu (3.3) và (3.4), ta thu được các đồng cấu (3.2)
mà ta muốn định nghĩa. Việc kiểm tra i. khá dài dòng, nhưng không
gặp khó khăn nào.


\src{175}
\result{MỆNH ĐỀ}{4} Cho $A$ là một vành giao hoán có đơn vị, và cho
$(x_1,\ldots,x_p)$ là một dãy các phần tử của $A$ sao cho, với $1\leq i\leq p$,
ảnh của $x_i$ trong vành thương của $A$ theo iđêan sinh bởi
$(x_1,\ldots,x_{i-1})$ không phải là ước của $0$. Gọi $J$ là iđêan sinh
bởi các $x_i$. Khi đó $J/J^{2}$ là một $(A/J)$-môđun tự do có cơ sở là các
ảnh chính tắc của các $x_i$, phức $K_{\underline{x}}$ là một
phân giải tự do của $A/J$, và với mọi $A$-môđun $M$, các đồng cấu (3.1) ở
bậc $p$ đều song ánh. Điều tương tự cũng đúng với các đồng cấu tương tự
được định nghĩa ở mọi bậc $i$, miễn là $J.M=0$.

(Điểm cốt yếu, mà từ đó mọi khẳng định khác đều suy ra, là việc
$K_{\underline{x}}$ không có đồng điều ở các bậc dương, một sự kiện quen thuộc dưới các giả thiết đã nêu.)

\result{HỆ QUẢ}{1} Với $A$ và $J$ như trên, giả sử thêm rằng
$A$ là một đại số afin chính quy có $\dim n$ trên một trường hoàn hảo $k$, và
$A/J$ là một đại số afin chính quy. Ký hiệu $\Omega^{1}(A)$,
$\Omega^{1}(A/J)$ là các môđun vi phân Kähler. Khi đó có một
đẳng cấu chính tắc tương thích với phép địa phương hóa,
\begin{equation}\tag{3.5}
  \operatorname{Ext}_{A}^{p}\bigl(\Omega^{n-p}(A/J),\Omega^{n}(A)\bigr)=A/J.
\end{equation}
Thật vậy, $\Omega^{n-p}(A/J)$ là một $(A/J)$-môđun tự do hạng $1$, và tương tự,
$\Omega^{n}(A)$ là một $A$-môđun tự do hạng $n$,
do đó vế trái bằng
\[
  \operatorname{Ext}_{A}^{p}(A/J,A)
  \otimes\Omega^{n-p}(A/J)'\otimes\Omega^{n}(A)
\]
(trong đó dấu « $'$ » chỉ $(A/J)$-môđun đối ngẫu). Tích tenxơ của
hai thừa số cuối được đồng nhất với $\bigwedge^{p}(J/J^{2})$,
nên toàn bộ biểu thức
được đồng nhất với
\[
  \operatorname{Ext}_{A}^{p}\!\left(A/J,\bigwedge^{p}(J/J^{2})\right),
\]
và do đó, theo mệnh đề, với
\[
  \operatorname{Hom}_{A}\!\left(\bigwedge^{p}J/J^{2},
  \bigwedge^{p}J/J^{2}\right),
\]
tức là với $A/J$. Đặc biệt,
$\operatorname{Ext}_{A}^{p}\bigl(\Omega^{n-p}(A/J),\Omega^{n}(A)\bigr)$ chứa
một phần tử phân biệt tương ứng với phần tử đơn vị của $A/J$, được gọi là
\emph{lớp cơ bản của iđêan $J$ trong $A$}. (Thật ra, nó có thể được
định nghĩa dưới những điều kiện rộng hơn đáng kể.) Hệ quả 1 có thể được viết
dưới một dạng hình học hơn và toàn cục hơn:

\result{HỆ QUẢ}{2} Cho $X$ là một đa tạp không kỳ dị trên một trường đóng
đại số $k$, cho $Y$ là một đa tạp con đóng không kỳ dị của $X$,
$\Osh_X$ là bó các vành địa phương của $X$, và $\Osh_Y$ là bó các vành
địa phương của $Y$, được xem như một bó thương của $\Osh_X$. Cho $n$ là
số chiều của $X$ và $n-p$ là số chiều của $Y$.

\src{176}
Cho $\Omega_X$ (tương ứng, $\Omega_Y$) là bó các mầm của các dạng
vi phân chính quy trên $X$ (tương ứng, trên $Y$). Khi đó ta có các đẳng cấu
chính tắc:
\begin{equation}\tag{3.6}
  \Extsh_{\Osh_X}^{p}\bigl(\Omega_Y^{n-p},\Omega_X^{n}\bigr)=\Osh_Y,
\end{equation}
hoặc tương đương
\begin{equation}\tag{3.6 bis}
  \Extsh_{\Osh_X}^{p}\bigl(\Osh_Y,\Omega_X^{n}\bigr)=\Omega_Y^{n-p}.
\end{equation}
Công thức (3.6 bis) có thể dùng làm định nghĩa của $\Omega_Y^{n-p}$ khi $Y$
là một đa tạp kỳ dị. Chính xác hơn:

\result{MỆNH ĐỀ}{5} Cho $Y$ là một tập con đại số có chiều
$q=n-p$ trong một đa tạp đại số không kỳ dị $X$ có chiều $n$. Cho
$F$ là một bó đại số nhất quán trên $X$ có giá chứa trong $Y$, và $L$
là một bó đại số tự do địa phương trên $X$. Khi đó các bó
$\Extsh_{\Osh_X}^{i}(F,L)$ bằng không khi $i<p$, còn khi $i=p$ ta có một
đẳng cấu chính tắc
\begin{equation}\tag{3.7}
  \Extsh_{\Osh_X}^{p}(F,L)
  =\Homsh_{\Osh_X}\!\left(F,
  \Extsh_{\Osh_X}^{p}(\Osh_X/J,L)\right),
\end{equation}
trong đó $J$ là một bó iđêan tùy ý trên $X$ triệt tiêu $F$ và có
tập các không điểm là $Y$. Đặc biệt, nếu $F$ là một bó đại số
nhất quán trên $Y$, thì
\begin{equation}\tag{3.7 bis}
  \Extsh_{\Osh_X}^{p}(F,L)
  =\Homsh_{\Osh_Y}\!\left(F,\Extsh_{\Osh_X}^{p}(\Osh_Y,L)\right).
\end{equation}
Sau cùng, vẫn giả sử $F$ là một bó đại số nhất quán trên $Y$, các bó
\[
  E^{i}(F)=\Extsh_{\Osh_X}^{p+i}(F,\Omega_X^{n})
\]
không phụ thuộc vào cách nhúng không gian đại số $Y$ vào một đa tạp
không kỳ dị $X$.

Vì bài toán có tính địa phương, có thể giả sử $X$ là afin và $L=\Osh_X$. Điều này quy
bài toán về đại số giao hoán, thậm chí về đại số địa phương: nếu $A$ là một
vành địa phương chính quy và $M$ là một $A$-môđun mà giá của nó có chiều
$\leq q=n-p$, cần chứng minh rằng $\operatorname{Ext}_{A}^{i}(M,A)=0$ khi $i<p$, và
\[
  \operatorname{Ext}_{A}^{p}(M,A)
  =\operatorname{Hom}_{A}\!\left(M,\operatorname{Ext}_{A}^{p}(A/J,A)\right),
\]
trong đó $J$ là một iđêan tùy ý có « số chiều » $\leq q$ và triệt tiêu $M$. Đối với
khẳng định thứ nhất, ta quy nạp theo $q$: một phép rút gọn từng lớp trực tiếp đưa
về trường hợp $M$ có dạng $A/J$; sau đó, thay $J$ bằng một iđêan nhỏ hơn
và dùng giả thiết quy nạp cùng dãy khớp của các
$\operatorname{Ext}$, ta quy về trường hợp $J$ được sinh bởi một « hệ tham số »
như trong Mệnh đề 4, và khi đó kết quả là hiển nhiên. Kết quả
trên cho thấy rằng, nếu $J$ là một iđêan


\src{177}
cố định có « số chiều » $\leq q$, thì hàm tử phản biến
$E(M)=\operatorname{Ext}_{A}^{p}(M,A)$ trên phạm trù các $A/J$-môđun là
khớp trái; hơn nữa, nó biến các tổng trực tiếp thành các tích trực tiếp,
từ đó dễ dàng suy ra rằng
\[
  E(M)=\operatorname{Hom}_{A}\bigl(M,E(A)\bigr).
\]
 Sau cùng,
khẳng định cuối của Mệnh đề 5 tinh tế hơn và suy ra từ một
đặc trưng nội tại của các $E^{i}(F)$ nhờ một định lý đối ngẫu
địa phương mà không thể trình bày ở đây.

\result{HỆ QUẢ}{} Ký hiệu $\omega_Y^{q}$ là bó
$\Extsh_{\Osh_X}^{p}(\Osh_Y,\Omega_X^{n})$. Khi đó, đối với các bó đại số
nhất quán $F$ trên $Y$, ta có một đẳng cấu có tính hàm tử:
\begin{equation}\tag{3.8}
  \Extsh_{\Osh_X}^{p}(F,\Omega_X^{n})
  =\Homsh_{\Osh_X}(F,\omega_Y^{q}).
\end{equation}

\fsection{4}{Lớp đối đồng điều gắn với một đa tạp con}

Trong toàn bộ phần sau, $X$ ký hiệu một tập đại số có chiều $n$, được xác định
trên một trường $k$, mà để đơn giản ta giả sử là đóng đại số. Ngoại trừ
ở n\textsuperscript{o} 6, $X$ được giả sử không kỳ dị. Ta ký hiệu
$\Osh_X$ là bó các vành địa phương trên $X$, và
\[
  \Omega_X^{*}=\bigcup_p\Omega_X^{p}
\]
là bó các mầm của các dạng vi phân trên $X$. Nếu $Y$ là một tập con đóng của $X$, ta đồng nhất
các bó đại số nhất quán trên $Y$ với các bó nhất quán trên $X$
bằng không ngoài $Y$; điều này đặc biệt áp dụng cho $\Osh_Y$ và $\Omega_Y$.

\result{BỔ ĐỀ}{1} Cho $F$ là một bó đại số nhất quán trên $X$ mà
giá của nó có chiều $\leq n-p$, và cho $L$ là một bó đại số nhất quán
tự do địa phương trên $X$. Khi đó $\operatorname{Ext}_{\Osh_X}^{i}(X;F,L)$ bằng không
với $i<p$, và ta có một đẳng cấu chính tắc
\begin{equation}\tag{4.1}
  \operatorname{Ext}_{\Osh_X}^{p}(X;F,L)
  =H^{0}\bigl(X,\Extsh_{\Osh_X}^{p}(F,L)\bigr).
\end{equation}
Nếu $F$ là một bó đại số nhất quán trên một tập con đóng $Y$ của $X$ có
chiều $\leq n-p$, thì ta có một đẳng cấu chính tắc
\begin{equation}\tag{4.1 bis}
  \Extsh_{\Osh_X}^{p}(F,L)
  =\Homsh_{\Osh_X}\!\left(
    F\otimes L'\otimes\Omega_X^{n},\omega_Y^{n-p}
  \right),
\end{equation}
trong đó $\omega_Y^{n-p}$ là bó trên $Y$ được định nghĩa trong hệ quả của
Mệnh đề 5 (và đồng nhất với $\Omega_Y^{n-p}$ nếu $Y$ không kỳ dị).

Công thức (4.1) là hệ quả trực tiếp của dãy phổ trong
Mệnh đề 1 và của Mệnh đề 5; theo công thức (3.8), ta có thể
viết


\src{178}
\begin{align*}
\Extsh_{\Osh_X}^{p}(F,L)
&=L\otimes(\Omega_X^{n})'\otimes
  \Extsh_{\Osh_X}^{p}(F,\Omega_X^{n})\\
&=L\otimes(\Omega_X^{n})'\otimes
  \Homsh_{\Osh_X}(F,\omega_Y^{q})\\
&=\Homsh_{\Osh_X}\!\left(
  F\otimes L'\otimes\Omega_X^{n},\omega_Y^{q}\right),
\end{align*}
trong đó $q=n-p$, từ đó suy ra công thức (4.1 bis).

Đặc biệt, đặt $F=\Osh_Y$, $L=\Omega_X^{p}$; vì
$\Omega_X^{n}\otimes(\Omega_X^{p})'=\Omega_X^{n-p}$, ta thu được một đẳng cấu
chính tắc
\begin{equation}\tag{4.2}
  \operatorname{Ext}_{\Osh_X}^{p}(X;\Osh_Y,\Omega_X^{p})
  =\operatorname{Hom}_{\Osh_X}
  (X;\Omega_X^{n-p},\omega_Y^{n-p}).
\end{equation}
Bây giờ, để đơn giản, giả sử $Y$ không kỳ dị, do đó
$\omega_Y^{n-p}=\Omega_Y^{n-p}$. Ta có một đồng cấu tự nhiên từ
$\Omega_X^{n-p}$ đến $\Omega_Y^{n-p}$, từ đó suy ra một tiết diện chính tắc $s_Y$ của
bó $\Extsh_{\Osh_X}^{p}(\Osh_Y,\Omega_X^{p})$, và nếu
mọi thành phần của $Y$ đều có chiều $n-p$, ta gọi nó là \emph{tiết diện
cơ bản} của bó $\Extsh_{\Osh_X}^{p}(\Osh_Y,\Omega_X^{p})$. Theo (4.1), tiết diện
này xác định một phần tử của
$\operatorname{Ext}_{\Osh_X}^{p}(X;\Osh_Y,\Omega_X^{p})$. Mặt khác, đồng cấu
tự nhiên $\Osh_X\longrightarrow\Osh_Y$ xác định một đồng cấu
\[
  \operatorname{Ext}_{\Osh_X}^{p}(X;\Osh_Y,\Omega_X^{p})
  \longrightarrow
  \operatorname{Ext}_{\Osh_X}^{p}(X;\Osh_X,\Omega_X^{p})
  =H^{p}(X,\Omega_X^{p}).
\]
Như vậy, ta thu được một phần tử của $H^{p}(X,\Omega_X^{p})$, ký hiệu là $P_X(Y)$ và gọi là
\emph{lớp đối đồng điều của $Y$ trong $X$}. Do đó, lớp này được suy ra từ tiết diện
$s_Y$ của $\Extsh_{\Osh_X}^{p}(\Osh_Y,\Omega_X^{p})$ thông qua sơ đồ
các đồng cấu sau đây:
\begin{equation}\tag{4.3}
\begin{aligned}
H^{p}(X,\Omega_X^{p})
&=\operatorname{Ext}_{\Osh_X}^{p}(X;\Osh_X,\Omega_X^{p})\\
&\longleftarrow\operatorname{Ext}_{\Osh_X}^{p}(X;\Osh_Y,\Omega_X^{p})
\xrightarrow{\ \sim\ }
H^{0}\!\left(X,\Extsh_{\Osh_X}^{p}(\Osh_Y,\Omega_X^{p})\right).
\end{aligned}
\end{equation}

Ta gọi một \emph{chu trình không kỳ dị chiều $n-p$} là một phần tử của nhóm
Abel tự do sinh bởi các đa tạp con bất khả quy không kỳ dị có
chiều $n-p$ trong $X$. Khi đó, ánh xạ $Y\longmapsto P(Y)$ mở rộng thành
một đồng cấu từ nhóm các chu trình không kỳ dị chiều $n-p$ trên $X$
vào nhóm $H^{p}(X,\Omega_X^{p})$.

Cho $Z^{n-p}$ và $Z'^{n-p'}$ là các chu trình không kỳ dị lần lượt có chiều $n-p$ và
$n-p'$; ta nói chúng \emph{giao hoành} nếu mọi thành phần
của $Z$ giao hoành với mọi thành phần của $Z'$. Khi đó, chu trình $Z.Z'$ được
xác định và là một chu trình không kỳ dị chiều $n-p-p'$. Với quy ước đó, ta có


\src{179}
\result{ĐỊNH LÝ}{1} Nếu $Z^{n-p}$ và $Z'^{n-p'}$ là các chu trình không kỳ dị giao hoành với nhau, thì ta có
\begin{equation}\tag{4.4}
  P_X(Z.Z')=P_X(Z).P_X(Z'),
\end{equation}
trong đó tích ở vế phải là tích cup
\[
  H^{p}(X,\Omega_X^{p})\times H^{p'}(X,\Omega_X^{p'})
  \longrightarrow H^{p+p'}(X,\Omega_X^{p+p'}).
\]
(Ta giả sử $X$ đẳng cấu với một tập con đóng địa phương của một không gian
xạ ảnh.)

Giả thiết cuối này chỉ được dùng để suy ra rằng mọi bó đại số nhất quán trên $X$
là thương của một bó đại số nhất quán tự do địa phương (SERRE),
và do đó có một phân giải trái bởi các bó tự do địa phương. Để chứng minh
Định lý 1, ta có thể giả sử $Z$ và $Z'$ là các đa tạp con bất khả quy không kỳ dị
$Y$ và $Y'$ giao hoành với nhau. Cho $L_{*}$ là một phân giải trái
của $\Osh_Y$ bởi
các bó tự do địa phương; khi đó, theo Mệnh đề 3, biểu đồ các đồng cấu
(4.3) đồng nhất với biểu đồ
\[
R^{p}\Gamma(\Omega_X^{p})
\xleftarrow{\alpha}
\mathbb R^{p}\Gamma\!\left(\Homsh_{\Osh_X}(L_{*},\Omega_X^{p})\right)
\xrightarrow[\sim]{\beta}
\Gamma\!\left(H^{p}\bigl(\Homsh_{\Osh_X}(L_{*},\Omega_X^{p})\bigr)\right),
\]
trong đó $\Gamma$ là hàm tử “nhóm các tiết diện” trên phạm trù
các bó nhóm Abel trên $X$, còn $\mathbb R^{p}\Gamma$ chỉ
siêu đối đồng điều bậc $p$ của hàm tử ấy, và $R^{p}\Gamma$ là
hàm tử dẫn xuất thứ $p$ của nó. Để đơn giản, ta giả sử $L_0=\Osh_X$ và
đồng cấu tăng cường $L_0\longrightarrow\Osh_Y$ là
đồng cấu tự nhiên (điều này được phép); khi đó $\alpha$ được suy ra
từ đồng cấu
các phức $\Osh_X\longrightarrow L_{*}$ (trong đó $\Osh_X$ được xem như một phức
tập trung ở bậc $0$), có tính đến rằng $\mathbb R^{p}\Gamma(K)=R^{p}\Gamma(K_0)$ nếu $K$
là một phức các bó tập trung ở bậc $0$. Đồng cấu $\beta$ là một đồng cấu
biên quen thuộc. Xét một biểu đồ tương tự, ứng với
một phân giải tự do địa phương
$L'_{*}$ của $\Osh_{Y'}$, và xét biểu đồ giao hoán của các phép ghép cặp:
\begin{equation}\tag{4.5}
\resizebox{\textwidth}{!}{$
\begin{array}{ccc}
R^{p}\Gamma(\Omega_X^{p})
&\longleftarrow
\mathbb R^{p}\Gamma\!\left(\Homsh_{\Osh_X}(L_{*},\Omega_X^{p})\right)
&\xrightarrow{\sim}
\Gamma H^{p}\!\left(\Homsh_{\Osh_X}(L_{*},\Omega_X^{p})\right)
\\[3pt]
\times&\times&\times\\[3pt]
R^{p'}\Gamma(\Omega_X^{p'})
&\longleftarrow
\mathbb R^{p'}\Gamma\!\left(\Homsh_{\Osh_X}(L'_{*},\Omega_X^{p'})\right)
&\xrightarrow{\sim}
\Gamma H^{p'}\!\left(\Homsh_{\Osh_X}(L'_{*},\Omega_X^{p'})\right)
\\[3pt]
\downarrow&\downarrow&\downarrow\\[3pt]
R^{p+p'}\Gamma(\Omega_X^{p+p'})
&\longleftarrow
\mathbb R^{p+p'}\Gamma\!\left(
\Homsh_{\Osh_X}(L_{*}\otimes L'_{*},\Omega_X^{p+p'})\right)
&\xrightarrow{\sim}
\Gamma H^{p+p'}\!\left(
\Homsh_{\Osh_X}(L_{*}\otimes L'_{*},\Omega_X^{p+p'})\right).
\end{array}
$}
\end{equation}


\src{180}
Các phép ghép cặp của hai cột bên phải được suy ra từ phép ghép cặp của
các phức bó
\[
\Homsh_{\Osh_X}(L_{*},\Omega_X^{p})
\times\Homsh_{\Osh_X}(L'_{*},\Omega_X^{p'})
\longrightarrow
\Homsh_{\Osh_X}(L_{*}\otimes L'_{*},\Omega_X^{p+p'}),
\]
được định nghĩa bằng tích ngoài
$\Omega_X^{p}\times\Omega_X^{p'}\longrightarrow\Omega_X^{p+p'}$.
Phép ghép cặp của cột thứ nhất là tích cup (tương ứng với tích
ngoài). Tôi khẳng định hàng cuối của (4.5) đồng nhất với biểu đồ
các đẳng cấu tương tự (4.3), trong đó $Y$ được thay bởi $Y\cap Y'$ và $p$ bởi
$p+p'$. Để làm điều này, chỉ cần chứng minh rằng $L\otimes L'$ là một phân giải
(hiển nhiên tự do địa phương) của $\Osh_{Y\cap Y'}$. Thật vậy, ta có
\[
  H_0(L\otimes L')=\Osh_Y\otimes\Osh_{Y'}=\Osh_{Y\cap Y'}
\]
và
\[
  H_i(L\otimes L')=\operatorname{Tor}_{i}^{\Osh_X}(\Osh_Y,\Osh_{Y'})=0
  \qquad(i>0),
\]
vì $Y$ và $Y'$ giao hoành với nhau. Định lý 1 giờ đây suy ra
từ công thức:
\begin{equation}\tag{4.6}
  s_Y.s_{Y'}=s_{Y.Y'}
\end{equation}
(trong đó tích ở vế trái là tích của cột cuối trong (4.5)).
Công thức (4.6) này, có bản chất thuần túy địa phương, được kiểm tra dễ dàng
bằng cách lấy $L_{*}$ và $L'_{*}$ là các phân giải đã xét trong Mệnh đề
4. Tương tự, ta chứng minh (bằng một chứng minh đơn giản hơn) rằng
$Z\longmapsto P_X(Z)$ tương thích với tích Đề-các:
\begin{equation}\tag{4.7}
  P_{X\times X'}(Z\times Z')=P_X(Z)\otimes P_{X'}(Z')
\end{equation}
(công thức có hiệu lực nếu $Z$, tương ứng $Z'$, là một chu trình không kỳ dị trên đa tạp
không kỳ dị $X$, tương ứng $X'$, và $Z\times Z'$ được xem như một chu trình không
kỳ dị trên $X\times X'$). Từ (4.4) và (4.7), suy ra rằng $P_X(Z)$ cũng
tương thích với phép “lấy ảnh ngược” theo một cấu xạ của các đa tạp không
kỳ dị $f:X\longrightarrow X'$:
\begin{equation}\tag{4.8}
  P_X\bigl(f^{-1}(Z')\bigr)=f^{*}\bigl(P_{X'}(Z')\bigr),
\end{equation}
công thức có hiệu lực nếu $Z'$ là một chu trình không kỳ dị trên $X'$ sao cho $f$
“giao hoành” với $Z'$, nghĩa là đồ thị của $f$ giao hoành với chu trình
$X\times Z'$ trong $X\times X'$.


\src{181}
\result{HỆ QUẢ}{1} Cho $X$, $X'$ là hai đa tạp không kỳ dị
đóng địa phương trong một không gian xạ ảnh, và giả sử $X'$ đầy đủ. Cho $U$
là một chu trình không kỳ dị trên $X\times X'$, và $a$, $b$ là hai điểm của $X'$ sao cho
$U$ giao hoành với các chu trình $X\times(a)$ và $X\times(b)$; cho $Z$
và $Z'$ là các chu trình không kỳ dị trên $X$ sao cho
\[
  Z\times(a)=(X\times(a)).U,
  \qquad
  Z'\times(b)=(X\times(b)).u.
\]
Khi đó ta có
\[
  P_X(Z)=P_X(Z').
\]

Thật vậy, cho $f_a:X\longrightarrow X\times X'$ được xác định bởi
$f_a(x)=(x,a)$; khi đó, theo (4.8), ta có công thức
$P(Z)=f_a^{*}(P(U))$, và tương tự $P(Z')=f_b^{*}(P(U))$. Mặt khác, dùng công thức
Künneth
\[
  H^{*}(X\times X',\Omega_{X\times X'}^{*})
  =H^{*}(X,\Omega_X^{*})\otimes H^{*}(X',\Omega_{X'}^{*})
\]
và sự kiện $H^{0}(X',\Osh_{X'})$ chỉ gồm các vô hướng, ta dễ dàng thu được
$f_a^{*}=f_b^{*}$, từ đó suy ra kết luận.

Với mọi $x\in X$, $(x)$ là một đa tạp con không kỳ dị của $X$ có
đối chiều $n$, và do đó xác định một phần tử $\varepsilon_x$ của
$H^{n}(X,\Omega_X^{n})$. Nếu $X$ là một đa tạp xạ ảnh không kỳ dị, từ
Hệ quả 1 suy ra rằng $\varepsilon_x$ không phụ thuộc vào điểm $x$ được chọn; ta ký hiệu nó
bởi $\varepsilon_X$ và gọi là \emph{lớp cơ bản} của
$H^{n}(X,\Omega_X^{n})$.

\result{NHẬN XÉT}{} Để có một lý thuyết thỏa đáng, cần định nghĩa
$P_X(Z)$ cho các chu trình $Z$ tùy ý và chứng minh Định lý 1 cho một
giao đúng của các chu trình. (Vào thời điểm viết bài trình bày này, điều đó vẫn chưa được
thực hiện trong tính tổng quát hoàn toàn). Giả sử điều đó đã được thực hiện, Hệ quả 1
trở thành: nếu $Z$ và $Z'$ là hai chu trình tương đương đại số, thì
$P_X(Z)=P_X(Z')$ (một mệnh đề dường như không suy ra từ những điều trên, ngay cả khi
$Z$ và $Z'$ không kỳ dị).

\fsection{5}{Định lý đối ngẫu}

Trong mục này, $X$ ký hiệu một đa tạp xạ ảnh không kỳ dị có chiều
$n$.

\result{ĐỊNH LÝ}{2} Lớp cơ bản $\varepsilon_X$ của
$H^{n}(X,\Omega_X^{n})$ là một cơ sở của không gian vectơ này.

(Chứng minh sẽ được trình bày ở dưới). Nhờ định lý trên, ta có thể
đồng nhất $H^{n}(X,\Omega_X^{n})$ với trường $k$. Bây giờ xét các
phép ghép cặp đã xét ở n\textsuperscript{o} 2, chúng đặc biệt cho một
phép ghép cặp


\src{182}
\[
\operatorname{Ext}_{\Osh_X}^{p}(X;\Osh_X,F)
\times\operatorname{Ext}_{\Osh_X}^{n-p}(X;F,\Omega_X^{n})
\longrightarrow
\operatorname{Ext}_{\Osh_X}^{n}(X;\Osh_X,\Omega_X^{n}),
\]
tức là,
\begin{equation}\tag{5.1}
  H^{p}(X,F)
  \times\operatorname{Ext}_{\Osh_X}^{n-p}(X;F,\Omega_X^{n})
  \longrightarrow H^{n}(X,\Omega_X^{n})
\end{equation}
Theo Định lý 2, phép ghép cặp này xác định một đồng cấu
\begin{equation}\tag{5.2}
  \operatorname{Ext}_{\Osh_X}^{n-p}(X;F,\Omega_X^{n})
  \longrightarrow\bigl(H^{p}(X,F)\bigr)'
\end{equation}
Đồng cấu này có tính hàm tử theo $F$ và giao hoán với các đồng cấu nối
tương ứng với các dãy khớp $0\longrightarrow F'\longrightarrow F
\longrightarrow F''\longrightarrow0$.

\result{ĐỊNH LÝ}{3} Đồng cấu (5.2) là một đẳng cấu.

Đặc biệt, ta thu lại kết quả của Serre:

\result{HỆ QUẢ}{} Cho $E$ là một phân thớ vectơ đại số trên $X$ và
$\Osh_X(E)$ là bó các mầm của các tiết diện chính quy của $E$. Khi đó ta có các
đẳng cấu chính tắc:
\begin{equation}\tag{5.3}
  \bigl(H^{p}(X,\Osh_X(E))\bigr)'
  =H^{n-p}\bigl(X,\Omega_X^{n}\otimes\Osh_X(E')\bigr)
\end{equation}
Chỉ cần áp dụng Định lý 3 và Hệ quả 1 cho Mệnh đề 1.

Các Định lý 2 và 3 sẽ được suy ra từ mệnh đề sau:

\noindent\textbf{(D) Đồng cấu sau đây}
\begin{equation}\tag{5.2 bis}
  \operatorname{Ext}_{\Osh_X}^{n-p}(X;F,\Omega_X^{n})
  \longrightarrow\bigl(H^{p}(X,F)\bigr)'\otimes L,
  \qquad L=H^{n}(X,\Omega_X^{n}),
\end{equation}
được cảm sinh bởi phép ghép cặp (5.1) là một đẳng cấu.

Quả vậy, ta chứng minh rằng (D) kéo theo Định lý 2. Cho
$k_x=\Osh_{(x)}$ là bó các vành địa phương của đa tạp được thu về điểm
$x\in X$. Xét đồng cấu chính tắc
$\Osh_X\longrightarrow k_x$ và đồng cấu tương ứng
\begin{equation}\tag{5.3}
  H^{0}(X,\Osh_X)\longrightarrow H^{0}(X,k_x)
\end{equation}
chuyển vị của đồng cấu này được đồng nhất với đồng cấu
\begin{equation}\tag{5.4}
  \operatorname{Ext}_{\Osh_X}^{n}(X;k_x,\Omega_X^{n})\otimes L'
  \longrightarrow
  \operatorname{Ext}_{\Osh_X}^{n}(X;\Osh_X,\Omega_X^{n})\otimes L'
\end{equation}
được cảm sinh bởi đồng cấu trên các $\operatorname{Ext}^{n}$ tương ứng với
$\Osh_X\longrightarrow k_x$:
\begin{equation}\tag{5.5}
  \operatorname{Ext}_{\Osh_X}^{n}(X;k_x,\Omega_X^{n})
  \longrightarrow
  \operatorname{Ext}_{\Osh_X}^{n}(X;\Osh_X,\Omega_X^{n})
  =H^{n}(X,\Omega_X^{n})
\end{equation}


\src{183}
Vì (5.3) là một đẳng cấu, nên (5.4) cũng là một đẳng cấu, và do đó (5.5)
cũng vậy. Vì $s_{(x)}$ là một cơ sở của
$\operatorname{Ext}_{\Osh_X}^{n}(X;k_x,\Omega_X^{n})$ theo (4.2), suy ra ảnh
$\varepsilon_X$ của nó là một cơ sở của $H^{n}(X,\Omega_X^{n})$.

Còn phải chứng minh mệnh đề (D); mệnh đề này sẽ suy ra một cách hoàn toàn hình
thức từ các sự kiện sơ cấp được tóm tắt trong những bổ đề sau. Trong đó, ta giả
sử $X$ là một phần đóng (kỳ dị hoặc không) của không gian xạ ảnh $P$ chiều
$r$. Ta dùng ký hiệu $\Osh_P(m)$ cho bó trên $P$ được ký hiệu $\Osh(m)$ trong
[3], và ký hiệu $\Osh_X(m)$ cho bó tương tự trên $X$.

\result{BỔ ĐỀ}{2} Mệnh đề (D) đúng nếu $X=P$ và $F=\Osh_P(m)$.

Bổ đề này được kiểm tra bằng một phép tính trực tiếp. Phép tính tường minh các
$H^{i}(P,\Osh_P(m))$ có trong [3], nhưng cũng có thể thực hiện một cách sơ cấp
hơn. Phép tính tích cup
\[
  H^{i}(P,\Osh_P(m))\times H^{j}(P,\Osh_P(m'))
  \longrightarrow H^{i+j}(P,\Osh_P(r+r'))
\]
cần thiết để viết tường minh phép ghép cặp (5.1) không gây khó khăn.

\result{BỔ ĐỀ}{3} Mọi bó đại số nhất quán $F$ trên $X$ đều đẳng cấu với một
bó thương của một bó $\Osh_X(-m)^{k}$, trong đó có thể giả sử $m$ lớn tùy ý.

Điều này suy ra từ sự kiện rằng $F\otimes\Osh_X(m)$ ``được sinh bởi các tiết
diện'' khi $m$ lớn; xem [3].

\result{BỔ ĐỀ}{4} Cho $i>0$. Khi đó
$H^{r-i}(P,\Osh_P(-m))=0$ với $m$ lớn; và với mọi bó đại số nhất quán $B$
trên $X$, ta có
\[
  \operatorname{Ext}_{\Osh_X}^{i}(X;\Osh_X(-m),B)=0
\]
với $m$ lớn.

Khẳng định thứ nhất suy ra từ những phép tính tường minh đã nhắc ở trên; đối
với khẳng định thứ hai, ta nhận thấy có một đẳng cấu
\[
  \operatorname{Ext}_{\Osh_X}^{i}(X;\Osh_X(-m),B)
  =H^{i}(X,B\otimes\Osh_X(m))
\]
(Mệnh đề 1, Hệ quả 1), từ đó kết luận suy ra nhờ một kết quả quen thuộc trong
[3]. Kết hợp hai bổ đề trên, ta thu được

\result{HỆ QUẢ}{} Cho $i>0$. Khi đó hàm tử
$F\longmapsto H^{r-i}(P,F)$ trên phạm trù các bó đại số nhất quán trên $P$
là xóa được, và hàm tử
$\operatorname{Ext}_{\Osh_X}^{i}(X;F,B)$ trên phạm trù các bó đại số nhất quán
trên $X$ cũng vậy.


\src{184}
\result{BỔ ĐỀ}{5} Cho $A$ và $B$ là hai bó đại số nhất quán trên
$X$, và đặt $A(m)=A\otimes\Osh_X(m)$. Khi đó, với $m$ đủ lớn,
đồng cấu chính tắc
\[
  \operatorname{Ext}_{\Osh_X}^{1}(X;A(-m),B)
  \longrightarrow H^{0}\bigl(X,\Extsh_{\Osh_X}^{1}(A(-m),B)\bigr)
  =H^{0}\bigl(X,\Extsh_{\Osh_X}^{1}(A,B)(m)\bigr)
\]
là một đẳng cấu.

Điều này suy ra ngay từ dãy phổ của Mệnh đề 1 áp dụng cho
$A(-m)$ và $B$, vì khi đó ta có
\[
  E_{2}^{p,q}(A(-m),B)
  =H^{p}\bigl(X,\Extsh_{\Osh_X}^{q}(A(-m),B)\bigr)
  =H^{p}\bigl(X,\Extsh_{\Osh_X}^{q}(A,B)(m)\bigr),
\]
và biểu thức này bằng không khi $p>0$ và $m$ lớn.

Bây giờ ta chứng minh (D) trong trường hợp $X=P$. Trước hết, ta chứng minh (5.2 bis)
là một đẳng cấu khi $p=n$; vì khi đó cả hai vế đều là các
hàm tử khớp trái (do $H^{r+1}(P,F)=0$), từ Bổ đề 3 suy ra rằng
chỉ cần chứng minh khẳng định cho $F=\Osh_P(-m)$; nhưng trong trường hợp này, khẳng định đó
đã nằm trong Bổ đề 2. Vì các đồng cấu (5.2 bis) có tính hàm tử và
tương thích với các toán tử nối, và vì khi $p<n$, cả hai vế của
(5.2 bis), xét như các hàm tử theo $F$, đều xóa được (Hệ quả của Bổ đề 4), nên
một lập luận chuẩn cho thấy (5.2 bis) là một đẳng cấu
với mọi $p$. Điều này chứng minh định lý đối ngẫu cho không gian xạ ảnh.
Bây giờ giả sử $X$ tùy ý nhưng không kỳ dị. Theo định lý
đối ngẫu cho $P$, ta có một đẳng cấu
\[
  H^{n}(X,F)'=H^{n}(P,F)'
  =\operatorname{Ext}_{\Osh_P}^{r-n}(P;F,\Omega_P^{r}).
\]
Theo Bổ đề 1 (n\textsuperscript{o} 4), vế cuối được đồng nhất với
\[
  \operatorname{Hom}_{\Osh_P}(P;F,\omega_X^{n})
  =\operatorname{Hom}_{\Osh_X}(X;F,\Omega_X^{n})
  =\operatorname{Ext}_{\Osh_X}^{0}(X;F,\Omega_X^{n}),
\]
từ đó thu được một đẳng cấu
\begin{equation}\tag{5.6}
  H^{n}(X,F)'=\operatorname{Hom}_{\Osh_X}(X;F,\Omega_X^{n})
  =\operatorname{Ext}_{\Osh_X}^{0}(X;F,\Omega_X^{n}).
\end{equation}
Cho $F=\Omega_X^{n}$, ta thu được một đẳng cấu
\begin{equation}\tag{5.7}
  \gamma:H^{n}(X,\Omega_X^{n})\xrightarrow{\sim}k.
\end{equation}


\src{185}
Ta kiểm tra rằng đẳng cấu (5.6) chính là (5.2 bis) với $p=n$, khi
đặt $L=k$ nhờ (5.7). Do đó, (5.2 bis) là một đẳng cấu với
$p=n$. Để chứng minh đây là một đẳng cấu với mọi $p$, một lần nữa chỉ cần
chứng minh rằng, với $p<n$, hai vế của (5.2 bis) là các hàm tử
xóa được theo $F$; hơn nữa, theo Bổ đề 3, chỉ cần chứng minh mệnh đề mạnh hơn
rằng cả hai vế bằng không khi đặt $F=\Osh_X(-m)$ với $m$ lớn. Đối với vế
thứ nhất, điều này đúng theo Bổ đề 4; còn đối với vế thứ hai, dùng
định lý đối ngẫu trên $P$, ta viết:
\[
  H^{p}(X,\Osh_X(-m))'
  =\operatorname{Ext}_{\Osh_P}^{r-p}
    (P;\Osh_X(-m),\Omega_P^{r}).
\]
Vế thứ hai bằng không khi $p<n$ và $m$ lớn, như suy ra từ Bổ đề 5
(trong đó ta đặt $X=P$) và từ việc $\Osh_X$, xét như một bó đại số
nhất quán trên $P$, có chiều đối đồng điều $\leq r-n$ (vì $X$ không
kỳ dị); do đó
\[
  \Extsh_{\Osh_P}^{r-p}(\Osh_X,\Omega_P^{r})=0
  \qquad\text{với }p<n.
\]

\fsection{6}{Định lý đối ngẫu cho các đa tạp kỳ dị}

Cho $X$ là một tập con đóng có chiều $n$ của không gian xạ ảnh $P$ có
chiều $r$. Ở đây, công thức (5.6) phải được viết dưới dạng
\begin{equation}\tag{6.1}
  H^{n}(X,F)'
  \simeq\operatorname{Hom}_{\Osh_X}(X;F,\omega_X^{n})
  =\operatorname{Ext}_{\Osh_X}^{0}(X;F,\omega_X^{n}),
\end{equation}
trong đó ta đặt
\begin{equation}\tag{6.2}
  \omega_X^{n}=E^{0}(\Osh_X)
  =\Extsh_{\Osh_P}^{r-n}(\Osh_X,\Omega_P^{r}).
\end{equation}
Như đã nói trong Mệnh đề 5, bó vừa được đưa vào thực ra không phụ thuộc
vào phép nhúng đã chọn của $X$ vào một đa tạp không kỳ dị $P$.
Đặt $F=\omega_X^{n}$ trong (6.1), ta được
\begin{equation}\tag{6.2}
  H^{n}(X,\omega_X^{n})'
  \simeq\operatorname{Hom}_{\Osh_X}(X;\omega_X^{n},\omega_X^{n}),
\end{equation}
 từ đó tồn tại một phần tử
phân biệt trong $H^{n}(X,\omega_X^{n})'$, tương ứng với đồng cấu
đồng nhất từ $\omega_X^{n}$ vào chính nó:
\begin{equation}\tag{6.3}
  \gamma:H^{n}(X,\omega_X^{n})\longrightarrow k.
\end{equation}

\src{186}
Khi đó, ta xét các phép ghép cặp được xác định bởi phép hợp thành các
$\operatorname{Ext}$:
\begin{equation}\tag{6.4}
  H^{p}(X,F)\times
  \operatorname{Ext}_{\Osh_X}^{n-p}(X;F,\omega_X^{n})
  \longrightarrow H^{n}(X,\omega_X^{n}),
\end{equation}
rồi hợp thành chúng với đồng cấu $\gamma$ trong (6.3); ta thu được
các đồng cấu có tính hàm tử, tương thích với các toán tử nối:
\begin{equation}\tag{6.5}
  \operatorname{Ext}_{\Osh_X}^{n-p}(X;F,\omega_X^{n})
  \longrightarrow H^{p}(X,F)'
\end{equation}
(tổng quát hóa (5.2)). Ta kiểm tra rằng khi $p=n$, cách này cho
đẳng cấu (6.1). Như vậy, ta có

\result{ĐỊNH LÝ}{3 bis} Bốn điều kiện sau đây về $X$ là tương đương đối với
một số nguyên $k\geq0$ cho trước:
\begin{enumerate}
  \item[\textit{i.}] Đồng cấu có tính hàm tử (6.5) là một đẳng cấu với
    $n-k\leq p\leq n$.
  \item[\textit{ii.}] Ta có $H^{p}(X,\Osh_X(-m))=0$ với $m$ lớn và với
    $n-k\leq p<n$.
  \item[\textit{iii.}] Hàm tử $H^{p}(X,F)$ trên phạm trù các bó đại số
    nhất quán trên $X$ là xóa được với $n-k\leq p<n$.
  \item[\textit{iv.}] Ta có
    $E^{i}(\Osh_X)=\Extsh_{\Osh_P}^{r-n+i}(\Osh_X,\Omega_P^{r})=0$
    với $0<i\leq k$.
\end{enumerate}

\noindent\textbf{CHỨNG MINH.} --- \textit{i.}$\Rightarrow$\textit{ii.} theo
Bổ đề 4, \textit{ii.}$\Rightarrow$\textit{iii.} theo Bổ đề 3, còn
\textit{iii.}$\Rightarrow$\textit{i.} bằng một lập luận chuẩn quen thuộc,
vì khi đó hai vế của (6.5), xét như các hàm tử, đều xóa được với
$n-k\leq p<n$ (vế thứ nhất là như vậy theo Bổ đề 4). Cuối cùng, ta chứng minh
\textit{ii.}$\Leftrightarrow$\textit{iv.} Điều này suy ra từ hệ quả của
mệnh đề sau đây:

\result{MỆNH ĐỀ}{6} Cho $F$ là một bó đại số nhất quán trên $X$, và cho
$i$ là một số nguyên. Khi đó, với $m$ đủ lớn, ta có một đẳng cấu:
\begin{equation}\tag{6.6}
  H^{i}(X,F(-m))'
  \simeq H^{0}\bigl(X,E^{n-i}(F)(m)\bigr),
\end{equation}
trong đó ta đặt (so sánh n\textsuperscript{o} 3, Mệnh đề 5):
\begin{equation}\tag{6.7}
  E^{j}(F)=\Extsh_{\Osh_P}^{r-n+j}(F,\Omega_P^{r}).
\end{equation}
Thật vậy, theo định lý đối ngẫu trên $P$, vế thứ nhất của (6.6)
đẳng cấu với
$\operatorname{Ext}_{\Osh_P}^{r-i}(P;F(-m),\Omega_P^{r})$, do đó (6.6) suy ra
từ Bổ đề 5.

\src{187}
\result{HỆ QUẢ}{} Ta có $H^{i}(X,F(-m))=0$ với $m$ đủ lớn nếu và chỉ nếu
$E^{n-i}(F)=0$.

Nhắc lại rằng các $E^{j}(F)$ thực ra không phụ thuộc vào phép nhúng xạ ảnh
đang xét. Điều kiện trong hệ quả có tính hoàn toàn địa phương; do đó,
nếu nó đúng đối với $F$, thì nó cũng đúng đối với mọi bó đẳng cấu địa phương
với một bó $F^{n}$. Đặc biệt, nếu điều kiện này đúng
đối với $\Osh_X$, thì nó đúng đối với mọi bó đại số nhất quán tự do địa phương.
Chẳng hạn, điều này đúng với mọi $i<n$ nếu $X$ không kỳ dị, với
$i=0$ nếu không thành phần nào của $X$ chỉ gồm một điểm duy nhất, và với $i=0,1$
nếu $S$ là chuẩn tắc và mọi thành phần của nó có chiều $>1$ (xem
[3]). Để điều này đúng với mọi
$i<k$, theo định nghĩa, cần và đủ là các vành địa phương
$\Osh_x$ ($x\in X$) có « đối chiều đồng điều » $\geq k$ (xem [4]
về khái niệm này). Nếu $k=n$, theo Định lý 3 bis điều đó có nghĩa rằng
định lý đối ngẫu đúng cho $X$, nghĩa là (6.5) là một đẳng cấu với
mọi $p$ và mọi $F$. Có thể đưa ra nhiều điều kiện tương đương trên
các vành địa phương $\Osh_x$ để điều này xảy ra (Nagata), chẳng hạn chúng
thỏa mãn định lý đẳng chiều Cohen--Macaulay. Điều này xảy ra,
chẳng hạn, nếu về địa phương $X$ là một « giao đầy đủ » trong một
đa tạp xung quanh không kỳ dị.

\fsection{7}{Đối ngẫu Poincaré}

Cho $X$ là một đa tạp xạ ảnh không kỳ dị có chiều $n$. Khi đó
\[
  H^{**}(X)=H^{*}(X,\Omega_X^{*})
\]
là một đại số phản giao hoán song phân bậc hữu hạn chiều, được phân bậc
theo bậc toàn phần, sao cho $H^{p,q}(X)=H^{p}(X,\Omega_X^{q})$ có bậc
$p+q$. Các bậc của $H^{*}(X)$ nằm giữa $0$ và $2n$. Theo
Định lý 2 và hệ quả của Định lý 3, $H^{*}(X)$ là một « đại số
Poincaré » có chiều $2n$, nghĩa là $H^{2n}(X)$ được trang bị một đẳng cấu với
trường cơ sở $k$, và tích
\[
  H^{m}(X)\times H^{2n-m}(X)\longrightarrow H^{2n}(X)=k
\]
là một phép đối ngẫu giữa $H^{m}(X)$ và $H^{2n-m}(X)$. Hơn nữa, nếu $Y$ là một
đa tạp xạ ảnh không kỳ dị khác, công thức Künneth cho các
bó đại số nhất quán cho
\begin{equation}\tag{7.1}
  H^{*}(X\times Y)=H^{*}(X)\otimes H^{*}(Y),
\end{equation}
một đẳng cấu tương thích với các cấu trúc đại số Poincaré. Hơn nữa,
$H^{*}(X)$, xét như một đại số giao hoán, là một hàm tử phản biến theo
$X$, vì một cấu xạ $f:Y\longrightarrow X$ hiển nhiên xác định một
đồng cấu của các đại số phân bậc

\src{188}
\begin{equation}\tag{7.2}
  f^{*}:H^{*}(X)\longrightarrow H^{*}(Y).
\end{equation}
Vì đây là các đại số Poincaré, phép chuyển vị cho ta một đồng cấu giữa các
không gian vectơ:
\begin{equation}\tag{7.3}
  f_{*}:H^{*}(Y)\longrightarrow H^{*}(X).
\end{equation}
Ta đã thấy ở n\textsuperscript{o} 4 rằng tác động của $f^{*}$ lên các lớp đối đồng điều
ứng với các chu trình không kỳ dị được diễn giải về mặt hình học bằng cách
lấy các lớp đối đồng điều ứng với ảnh ngược của chúng. Trong trường hợp hiện
tại, cần chứng minh rằng (7.3) cũng ứng với phép « ảnh trực tiếp » trên các
chu trình. Điều này suy ra, ít nhất dưới các điều kiện không kỳ dị thích hợp,
từ trường hợp riêng sau đây:

\result{ĐỊNH LÝ}{4} Nếu $f$ là ánh xạ đồng nhất của một đa tạp con không
kỳ dị $Y^{m}$ của $X^{n}$ vào $X^{n}$, thì, ký hiệu
phần tử đơn vị của $H(Y)$ là $1_Y$, ta có
\begin{equation}\tag{7.4}
  f_{*}(1_Y)=P_X(Y),
\end{equation}
trong đó vế thứ hai là lớp đối đồng điều trong $X$ gắn với $Y$.

Công thức này tương đương với
\begin{equation}\tag{7.4 bis}
  \bigl\langle\xi^{m,m},P_X(Y)\bigr\rangle\varepsilon_Y
  =f_{*}(\xi^{m,m})
  \qquad
  \bigl(\xi^{m,m}\in H^{m}(X,\Omega_X^{m})\bigr),
\end{equation}
 trong đó $\varepsilon_Y$ là phần tử cơ bản
trong $H^{m}(Y,\Omega_Y^{m})$. Trong trường hợp các đa tạp xạ ảnh không kỳ dị,
công thức này cho một định nghĩa mới về lớp đối đồng điều gắn với $Y$. Để chứng
minh Định lý 4, nhờ Định lý 3, ta diễn giải chuyển vị của đồng cấu
\[
  H^{m}(X,\Omega_X^{m})
  \longrightarrow H^{m}(Y,\Omega_Y^{m})
  =H^{m}(X,\Omega_Y^{m})
\]
như đồng cấu
\begin{equation}\tag{7.5}
\begin{split}
  H^{n-m}(X,\Omega_X^{n-m})
  &\simeq
  \operatorname{Ext}_{\Osh_X}^{n-m}(X;\Omega_X^{m},\Omega_X^{n})\\
  &\longleftarrow
  \operatorname{Ext}_{\Osh_X}^{n-m}(X;\Omega_Y^{m},\Omega_X^{n})
  \simeq
  \operatorname{Hom}_{\Osh_X}(X;\Omega_Y^{m},\Omega_Y^{m}).
\end{split}
\end{equation}
Ta kiểm tra rằng phần tử $1_Y$ của đối ngẫu của $H^{m}(Y,\Omega_Y^{m})$
được đồng nhất với phần tử ở vế phải tương ứng với tự đồng cấu đồng nhất của
$\Omega_Y^{m}$, và mặt khác ảnh của phần tử này trong
$H^{n-m}(X,\Omega_X^{n-m})$ chính là $P_X(Y)$. Những tính tương thích này,
vốn có thể được nêu ngay từ n\textsuperscript{o} 4, có thể được phát biểu, và thực ra đúng,
đối với mọi đa tạp không kỳ dị; tính tương thích thứ hai

\src{189}
được suy ra, chẳng hạn, từ tính giao hoán của sơ đồ các tự đồng cấu chính tắc
sau đây:
\begin{equation}\tag{7.6}
\resizebox{\textwidth}{!}{$
\left\{
\begin{array}{ccccc}
\operatorname{Ext}_{\Osh_X}^{n-m}(X;\Omega_X^{m},\Omega_X^{n})
&\longleftarrow&
\operatorname{Ext}_{\Osh_X}^{n-m}(X;\Omega_Y^{m},\Omega_X^{n})
&\simeq&
\operatorname{Hom}_{\Osh_X}(X;\Omega_Y^{m},\Omega_Y^{m})\\[3pt]
\Big\downarrow&&&&\Big\downarrow\\[3pt]
\operatorname{Ext}_{\Osh_X}^{n-m}(X;\Osh_X,\Omega_X^{n-m})
&\longleftarrow&
\operatorname{Ext}_{\Osh_X}^{n-m}(X;\Osh_Y,\Omega_X^{n-m})
&\simeq&
\operatorname{Hom}_{\Osh_X}(X;\Omega_X^{m},\Omega_Y^{m})
\end{array}
\right.
$}
\end{equation}

Như vậy ta thu được một bản tương tự chính xác của hình thức luận đối ngẫu
Poincaré trên các đa tạp compact có định hướng. Đặc biệt, Định lý 4 cho phép xác
định lớp đối đồng điều gắn với đường chéo của $X\times X$. Bằng một lập luận
quen thuộc, chẳng hạn ta suy ra một công thức Lefschetz:

\result{ĐỊNH LÝ}{5} Cho $f$ là một tự đồng cấu của một đa tạp xạ ảnh không
kỳ dị $X$, sao cho các điểm bất động của $f$ đều có bội 1. Khi đó, theo môđun
đặc số của $k$, số các điểm bất động này đồng dư với tổng luân phiên các vết
của các tự đồng cấu do $f$ cảm sinh trên các $H^{i}(X)$.

Hạn chế buộc phải đặt lên $f$ bắt nguồn từ những khó khăn được nêu trong
nhận xét ở n\textsuperscript{o} 4. Tuy nhiên, ta lưu ý rằng công thức Lefschetz
vẫn đúng nếu $f$ « đồng luân » với một tự đồng cấu mà mọi điểm bất động đều
có bội 1.

\fsection{8}{Khái quát hóa định lý đối ngẫu}

Cho $X$ là một đa tạp đại số không kỳ dị sao cho mọi bó đại số nhất quán
$F$ trên $X$ đều đẳng cấu với thương của một bó đại số nhất quán tự do địa
phương (điều này đúng nếu $X$ đóng địa phương trong một không gian xạ ảnh).
Khi đó mọi bó đại số nhất quán $F$ trên $X$ đều có một phân giải hữu hạn
$L$ bởi các bó tự do địa phương; với hai phân giải như vậy, ta luôn có thể
tìm được một phân giải thứ ba và một đồng cấu từ phân giải này vào từng phân
giải trong hai phân giải đầu, tương thích với các đồng cấu tăng cường. Tương
tự, nếu $L$ là một phân giải hữu hạn tự do địa phương của $F$, và nếu cho
một đồng cấu $F'\longrightarrow F$, thì tồn tại một phân giải hữu hạn tự do
địa phương $L'$ của $F'$ cùng một đồng cấu $L'\longrightarrow L$ tương thích
với $F'\longrightarrow F$; có thể giả sử đồng cấu này là toàn ánh nếu
$F'\longrightarrow F$ là toàn ánh. Điều này cho phép, với hai số nguyên
$r,s\geq0$ đã cho, định nghĩa hai đa hàm tử đối đồng điều theo các đối số
$A_1,\ldots,A_r$; $B_1,\ldots,B_s$,

\src{190}
biến thiên trong phạm trù các bó đại số nhất quán trên $X$, và các giá trị của
chúng lần lượt thuộc phạm trù các bó đại số nhất quán trên $X$ và phạm trù các
môđun trên $H^{0}(X,\Osh_X)$, theo các công thức:
\begin{equation}\tag{8.1}
\resizebox{\textwidth}{!}{$
\left\{
\begin{aligned}
\underline{T}_{r}^{s*}(A_1,\ldots,A_r;B_1,\ldots,B_s)
  &={}H^{*}\!\left(
    \Homsh_{\Osh_X}\!\left(
      L(A_1)\otimes\cdots\otimes L(A_r),
      L(B_1)\otimes\cdots\otimes L(B_s)
    \right)\right),\\[3pt]
T_{r}^{s*}(A_1,\ldots,A_r;B_1,\ldots,B_s)
  &={}\mathbf R^{*}\Gamma\!\left(
    \Homsh_{\Osh_X}\!\left(
      L(A_1)\otimes\cdots\otimes L(A_r),
      L(B_1)\otimes\cdots\otimes L(B_s)
    \right)\right).
\end{aligned}
\right.
$}
\end{equation}
Trong công thức này, $L(F)$ ký hiệu một phân giải hữu hạn tự do địa phương của
bó đại số nhất quán $F$, còn $\mathbf R^{*}\Gamma(K)$ ký hiệu siêu đối đồng
điều của không gian $X$ đối với phức các bó $K$. Nếu $r$ hoặc $s$ bằng 0, ta thay
tích tenxơ của các $L(A_i)$, tương ứng của các $L(B_j)$, bằng $\Osh_X$. Đặc
biệt, $\underline{T}_{0}^{0}$ và $T_{0}^{0}$ là các hàm tử phân bậc không có
đối số; $\underline{T}_{0}^{0}$ tập trung ở bậc 0, tại đó nó là bó $\Osh_X$,
còn $T_{0}^{0}$ được đồng nhất với $H^{*}(X,\Osh_X)$. Trong mọi trường hợp,
việc các vế phải của (8.1) không phụ thuộc vào các phân giải đã chọn là hiển
nhiên đối với dòng thứ nhất (vì khi đó vấn đề mang tính địa phương), còn đối
với dòng thứ hai thì suy ra từ những nhận xét tổng quát ở trên, khi xét dãy phổ
siêu đối đồng điều của phức các bó
\[
  K=\Homsh_{\Osh_X}\bigl(L(A_1)\otimes\cdots,
      L(B_1)\otimes\cdots\bigr),
\]
hội tụ tới siêu đối đồng điều của $X$ đối với $K$ và có số hạng ban đầu là
$H^{p}(X,H^{q}(K))$, tức là
\begin{equation}\tag{8.2}
  E_{2}^{p,q}
  =H^{p}\!\left(X,
    (\underline{T}_{r}^{s})^{(q)}
      (A_1,\ldots,A_r;B_1,\ldots,B_s)\right).
\end{equation}
Khi đó, ta thấy rằng chính dãy phổ này cũng không phụ thuộc vào các phân giải
đã chọn; đích hội tụ của nó là vế trái trong dòng thứ hai của (8.1). Ta dễ dàng
định nghĩa các toán tử nối ứng với từng đối số $A_i$, $B_j$, vì mọi dãy khớp
$0\longrightarrow F'\longrightarrow F\longrightarrow F''\longrightarrow0$
đều có thể được phân giải bởi một dãy khớp các phức hữu hạn tự do địa phương.

Trong hệ các hàm tử $\underline{T}_{r}^{s*}$, tương ứng $T_{r}^{s*}$, ta định
nghĩa các phép toán tương tự như các phép toán trong phép tính tenxơ; định nghĩa
của chúng suy ra ngay từ các công thức định nghĩa (8.1). Vì vậy, ta có một phép
hợp thành (khái quát hóa phép hợp thành đã xét ở n\textsuperscript{o} 2):
\begin{equation}\tag{8.3}
\resizebox{\textwidth}{!}{$
\begin{aligned}
&T_{r}^{s*}(A_1,\ldots,A_r;B_1,\ldots,B_s)
\times T_{r'}^{s'*}(A'_1,\ldots,A'_{r'};B'_1,\ldots,B'_{s'})\\
&\qquad\longrightarrow
T_{r+r'}^{s+s'*}
(A_1,\ldots,A_r,A'_1,\ldots,A'_{r'};
 B_1,\ldots,B_s,B'_1,\ldots,B'_{s'}),
\end{aligned}
$}
\end{equation}
thỏa mãn các tính chất hiển nhiên về tính kết hợp và tính tương thích với các
đồng cấu có tính hàm tử, các đồng cấu nối và các dãy phổ. Tương tự, ta có các
phép đối xứng, mà ta để độc giả tự viết tường minh. Ngoài ra, ta còn có một
phép co mỗi khi một trong các đối số $A_i$

\src{191}
trùng với một trong các đối số $B_j$:
\begin{equation}\tag{8.4}
\resizebox{\textwidth}{!}{$
\begin{aligned}
&T_{r}^{s*}(A_1,\ldots,A_{i-1},C,A_{i+1},\ldots,A_r;
 B_1,\ldots,B_{j-1},C,B_{j+1},\ldots,B_s)\\
&\qquad\longrightarrow
T_{r-1}^{s-1*}(A_1,\ldots,\widehat{A_i},\ldots,A_r;
 B_1,\ldots,\widehat{B_j},\ldots,B_s).
\end{aligned}
$}
\end{equation}
Ngoài ra, nếu một đối số $A_i$ là một bó tự do địa phương, thì có thể
lược bỏ đối số ấy, với điều kiện thay một trong các $A_j$ ($j\neq i$) bằng
$A_j\otimes A_i$, hoặc thay một trong các $B_k$ bằng $B_k\otimes A_i'$ (trong đó
$A_i'=\Homsh_{\Osh_X}(A_i,\Osh_X)$); ta có một quy tắc tính toán tương tự khi
một trong các đối số $B_i$ là một bó tự do địa phương. Đặc biệt, ta luôn có thể
lược bỏ một đối số đồng nhất với $\Osh_X$. Nếu mọi đối số đều tự do địa phương, trừ nhiều nhất
một trong các đối số $B_i$, thì quy tắc vừa nêu cho một đẳng cấu có tính hàm tử
\begin{equation}\tag{8.5}
  T_{r}^{s*}(A_1,\ldots,A_r;B_1,\ldots,B_s)
  =H^{*}\bigl(X,A_1'\otimes\cdots\otimes A_r'
    \otimes B_1\otimes\cdots\otimes B_s\bigr)
\end{equation}
(vì ta được quy về trường hợp $r=0$, $s=1$, trong đó điều này là hiển nhiên;
ta cũng có thể dùng trực tiếp dãy phổ có số hạng ban đầu (8.2)). Đối với các
$\underline{T}_{r}^{s}$, ta cũng định nghĩa các phép toán tương ứng với những
phép toán vừa nêu. Các quan hệ giữa những phép toán vừa được đưa vào cũng giống
như các quan hệ giữa những phép toán tương tự trong phép tính tenxơ.

Cho $n$ là chiều của $X$. Bằng cách áp dụng liên tiếp một phép hợp thành tenxơ
(8.3) và các phép co (8.4) trên những đối số trùng nhau, ta thu được một phép ghép cặp
\begin{equation}\tag{8.6}
\resizebox{\textwidth}{!}{$
  (T_{r}^{s})^{p}(A_1,\ldots;B_1,\ldots)
  \times
  (T_{s}^{r})^{n-p}(B_1,\ldots;A_1,\ldots,A_r\otimes\Omega_X^{n})
  \longrightarrow H^{n}(X,\Omega_X^{n}).
$}
\end{equation}

\result{ĐỊNH LÝ}{6} Nếu $X$ là một đa tạp xạ ảnh không kỳ dị, thì các phép
ghép cặp (8.6) là những phép đối ngẫu.

Điều này suy ra một cách hoàn toàn hình thức từ hệ quả của Định lý 3. Thật vậy,
từ hệ quả ấy dễ dàng suy ra rằng nếu $K$ là một phức các bó đại số nhất quán tự do
địa phương, thì siêu đối đồng điều của $X$ đối với $K$ đối ngẫu với siêu đối
đồng điều của $X$ đối với $K'\otimes\Omega_X^{n}$ qua các phép ghép cặp tự nhiên
\begin{equation}\tag{8.7}
  \mathbf R^{p}\Gamma(K)
  \times\mathbf R^{n-p}\Gamma(K'\otimes\Omega_X^{n})
  \longrightarrow\mathbf R^{n}\Gamma(\Omega_X^{n})
  =H^{n}(X,\Omega_X^{n}).
\end{equation}
Ta thấy điều này bằng cách dùng dãy phổ có số hạng ban đầu
$H^{p}(H^{q}(X,K))$ và dãy phổ tương tự cho
$K'\otimes\Omega_X^{n}$. Từ kết quả trên, Định lý 6 suy ra bằng cách dùng
định nghĩa (8.1).

\src{192}
\subsection*{NHẬN XÉT}

\noindent 1\textsuperscript{o} Đối với các định nghĩa trước Định lý 6,
không cần giả sử $X$ không kỳ dị, vì không nhất thiết chỉ làm việc với các
phân giải hữu hạn. Nhưng nếu $X$ kỳ dị, thì nói chung không còn có thể khẳng
định tiên nghiệm rằng các
$(\underline{T}_{r}^{s})^{p}(A_1,\ldots;B_1,\ldots)$ là những bó nhất quán,
bởi vì trong phức các bó
\[
  \Homsh_{\Osh_X}\bigl(L(A_1)\otimes\cdots,L(B_1)\otimes\cdots\bigr)
\]
sẽ có vô hạn thành phần có một bậc toàn phần cho trước.

\medskip
\noindent 2\textsuperscript{o} Dễ dàng kiểm tra rằng trong các công thức
(8.1), có thể thay một trong các $L(B_i)$ bằng $B_i$. Theo Mệnh đề 3, điều này
chứng tỏ rằng ta có
\begin{equation}\tag{8.8}
\left\{
\begin{aligned}
  \underline{T}_{1}^{1*}(A;B)&=\Extsh_{\Osh_X}^{*}(A,B),\\
  T_{1}^{1*}(A;B)&=\operatorname{Ext}_{\Osh_X}^{*}(X;A,B).
\end{aligned}
\right.
\end{equation}
Đặc biệt, lấy $r=s=1$ và $A_1=\Osh_X$ trong (8.6), ta thu lại Định lý 3.
Công thức (8.8) cũng kéo theo
\[
  T_{0}^{1*}(B)=H^{*}(X,B),
  \qquad
  T_{1}^{0*}(A)=\operatorname{Ext}_{\Osh_X}^{*}(X;A,\Osh_X).
\]

\medskip
\noindent 3\textsuperscript{o} Công thức (8.1) cho thấy rằng, nói chung, các
hàm tử $\underline{T}_{r}^{s*}$ và $T_{r}^{s*}$ có các thành phần ở cả bậc
dương lẫn bậc âm. Dùng Nhận xét 2, ta thấy rằng nếu $X$ có chiều $n$, thì các
thành phần khác không của $\underline{T}_{r}^{s*}$ nằm giữa $-(s-1)n$ và $rn$
khi $s>0$, và nằm giữa $0$ và $rn$ khi $s=0$; các thành phần khác không của
$T_{r}^{s*}$ nằm giữa $-(s-1)n$ và $(r+1)n$ khi $s>0$, và nằm giữa $0$ và
$(r+1)n$ khi $s=0$ (và thậm chí, nếu không nhầm, khi $r>0$, chúng nằm giữa
$-(s-1)n$ và $rn$, tương ứng giữa $0$ và $rn$).

\section*{TÀI LIỆU THAM KHẢO}

\noindent[1] Cartier (Pierre). --- Les groupes
$\operatorname{Ext}^{s}(A,B)$, Séminaire A. Grothendieck : Algèbre homologique,
tập 1, 1957, số 3.

\noindent[2] Grothendieck (Alexandre). --- Sur quelques points d'algèbre
homologique, \emph{Tohoku Math. J.}, tập 9, 1957, tr. 119--183.

\noindent[3] Serre (Jean-Pierre). --- Faisceaux algébriques cohérents,
\emph{Annals of Math.}, tập 61, 1955, tr. 197--278.

\noindent[4] Serre (Jean-Pierre). --- Sur la dimension homologique des anneaux
et des modules noethériens, \emph{Proc. Intern. Symp. on Alg. Number Theory}
[1955, Tokyo và Nikko]. --- Tokyo, Science Council of Japan, 1956;
tr. 175--189.

\clearpage

\src{193}
\section*{BỔ SUNG}

Những khó khăn được nêu trong nhận xét ở trang 13 hiện nay đã được giải quyết
hoàn toàn, nhờ việc mở rộng các định lý đối ngẫu sang các đa tạp bất kỳ
(có các kỳ dị tùy ý). Đối với lý thuyết này, vốn có thể được trình bày trong
khuôn khổ tổng quát của « lý thuyết lược đồ », chúng tôi dẫn chiếu đến bộ
« Éléments de Géométrie algébrique » của J. Dieudonné và A. Grothendieck,
hiện đang được biên soạn. Có thể tìm thấy một vài chỉ dẫn trong:

\begin{quote}
Grothendieck (A.). --- Cohomology theory of algebraic varieties,
\emph{Congrès int. des Mathématiciens}, 1958, Edinburgh (sắp xuất bản).
\end{quote}

\begin{flushright}
[Tháng 4 năm 1959]
\end{flushright}

\begin{center}\rule{36mm}{0.5pt}\end{center}

\par

\clearpage
\chapter*{Đính chính cho Bài trình bày số 149}
\addcontentsline{toc}{chapter}{Đính chính cho Bài trình bày số 149}
\setcounter{footnote}{0}
\sid{e149}
\markboth{Đính chính 149}{Đính chính 149}
\src{298}

\noindent\textbf{1. Định lý đối ngẫu cho các bó đại số nhất quán.}

\begin{center}
[Xê-mi-na Bourbaki, t. 9, 1956/57, số 149, 25 tr.]
\end{center}

\noindent\textbf{Ghi chú về trang 14.} --- Như tôi đã nêu trong bài giảng
của mình tại Đại hội Toán học Quốc tế năm 1958,\footnote{
Grothendieck (Alexander). --- The cohomology theory of abstract algebraic
varieties, \emph{Proceedings of the International Congress of Mathematicians}
[1958, Edinburgh], tr. 103--118. --- Cambridge, at the University Press, 1960.}
các vấn đề được nêu ra ở đây nay đã được giải quyết hoàn toàn.

Độc giả sẽ tìm thấy thêm thông tin về đối ngẫu của các bó nhất quán trong
\emph{loco citato}\footnotemark[1], tr. 112--115, trong EGA III, phần 2
(đang chuẩn bị), và trong khi chờ đợi, trong SGA 1962. Một cách trình bày
có hệ thống hơn sẽ được tìm thấy trong một chương sau của EGA
(chương IX theo đề cương dự kiến).

\par

\unit{Bài trình bày 182 — Hình học hình thức và hình học đại số}
\sid{182}
\markboth{Thuyết trình 182}{Thuyết trình 182}
\src{193}
\seminar{(Tháng 5 năm 1959)}
\paperhead{HÌNH HỌC HÌNH THỨC VÀ HÌNH HỌC ĐẠI SỐ}%
  {do Alexander GROTHENDIECK}

\fsection{1}{Lược đồ}

Ta biết rằng một không gian đại số affine xác định trên một trường $k$ về căn
bản được xác định bởi đại số affine $A$ của nó (vành các hàm chính quy xác định
trên $k$), vì các cấu xạ giữa các không gian đại số
$X\longrightarrow Y$ tương ứng song ánh với các đồng cấu
$k$-đại số $A(Y)\longrightarrow A(X)$. Đại số affine tương ứng với một không
gian đại số là một $k$-đại số kiểu hữu hạn và, theo quan điểm
\fquote{cổ điển}, không có phần tử lũy linh; ngược lại, mọi đại số thuộc kiểu
này đều nhận được như đại số affine của một không gian đại số xác định trên
$k$. Như vậy có một từ điển quen thuộc cho phép diễn giải các tình huống liên
quan đến không gian đại số affine bằng ngôn ngữ đại số giao hoán. Từ lâu người
ta đã nhận thấy rằng cách làm này cho những mệnh đề tổng quát hơn, vì thông
thường không còn cần giả sử các vành đang xét thuộc kiểu vừa nêu; giả thiết
Noether thường là đủ. Đặc biệt, dù có cho trước một trường cơ sở hay không,
không có lý do gì để loại trừ trường hợp các vành ấy chứa phần tử lũy linh.
Cho đến nay, các nhà hình học vẫn từ chối lưu ý đến những chỉ dẫn này và khăng
khăng chỉ xét các đại số affine không có phần tử lũy linh, nghĩa là các không
gian đại số mà bó cấu trúc không có phần tử lũy linh (và phần lớn còn chỉ xét
các không gian đại số \fquote{bất khả quy tuyệt đối}). Báo cáo viên cho rằng
quan niệm ấy đã gây trở ngại nghiêm trọng cho sự phát triển của các phương
pháp thật sự tự nhiên trong hình học đại số.

Cho $A$ là một vành giao hoán. Ta biết rằng tập
$X=\operatorname{Spec}(A)$ các iđêan nguyên tố của $A$ mang một tôpô tự nhiên,
gọi là \fquote{tôpô Zariski} hay tôpô phổ. Mặt khác, trên $X$ có một bó các
vành giao hoán $\Osh_X$, với thớ tại $\mathfrak p\in X$ là vành địa phương hóa
$A_{\mathfrak p}$ và vành các tiết diện của nó được đồng nhất với $A$. Do đó
$X$ trở thành một \emph{không gian có vành}, gọi là \emph{phổ nguyên tố} của
$A$. Một đồng cấu vành $f:A\longrightarrow B$ xác định một cấu xạ của các
không gian có vành
$f':\operatorname{Spec}(B)\longrightarrow\operatorname{Spec}(A)$; ánh xạ tập
hợp nền chính là $\mathfrak p\longmapsto f^{-1}(\mathfrak p)$. Các cấu xạ

\src{194}
của các không gian có vành từ $\operatorname{Spec}(B)$ đến
$\operatorname{Spec}(A)$ thu được theo cách này chính là những cấu xạ mà các
đồng cấu $\Osh_x\longrightarrow\Osh_y$ ($x=f'(y)$) đều là đồng cấu địa phương
(nghĩa là nghịch ảnh của iđêan cực đại là iđêan cực đại).

Ta gọi một không gian có vành đẳng cấu với một
$\operatorname{Spec}(A)$ nào đó là \emph{lược đồ affine}, và gọi một không gian
có vành affine địa phương là \emph{tiền lược đồ}; nói cách khác, mỗi điểm của
nó có một lân cận mở mà, với cấu trúc cảm sinh, là một lược đồ affine. Các
\emph{cấu xạ} của tiền lược đồ được định nghĩa theo cách hiển nhiên; ở địa
phương chúng tương ứng với các đồng cấu vành.

Khi cố định một tiền lược đồ $S$ và xét các cấu xạ tiền lược đồ
$X\longrightarrow S$, không gian $S$ đóng vai trò của một trường hay một vành
cơ sở (hay đúng hơn, của một không gian cơ sở trong một phân thớ). Khi đó ta
nói $X$ là một $S$-tiền lược đồ; nếu $S=\operatorname{Spec}(A)$, điều này cũng
có nghĩa $\Osh_X$ là một bó các $A$-đại số. Vì thế, mỗi tiền lược đồ có thể
được xem một cách duy nhất như một $\mathbf Z$-tiền lược đồ. Dĩ nhiên, các
$S$-tiền lược đồ lập thành một phạm trù; hơn nữa, người ta chứng minh rằng
trong phạm trù này tích của hai vật $X,Y$ luôn tồn tại và được ký hiệu
$X\times_S Y$. Khái niệm tích này cho phép định nghĩa phép thay đổi cơ sở của
một $S$-tiền lược đồ ứng với một cấu xạ $S'\longrightarrow S$: quả vậy,
$X\times_S S'$ có thể được xem như một $S'$-tiền lược đồ.

Ta nói $X$ \emph{tách} trên $S$ nếu đường chéo của $X$ trong
$X\times_S X$ là đóng. Một tiền lược đồ tách trên $\mathbf Z$ được gọi là một
\emph{lược đồ}; khi ấy nó tách trên mọi cơ sở. Để đơn giản, từ đây ta chỉ nói
đến các lược đồ và còn giả sử chúng là \emph{Noether}, nghĩa là hợp hữu hạn
của các tập mở affine vốn là phổ của các vành Noether. Ta nói $X$ \emph{thuộc
kiểu hữu hạn trên} $S$ nếu với mọi tập mở affine $U$ của $S$, nghịch ảnh của
$U$ trong $X$ là hợp hữu hạn của các tập mở affine có các vành là đại số kiểu
hữu hạn trên vành của $U$. Chính những $S$-lược đồ như vậy mới thích hợp cho
một nghiên cứu thực sự mang tính hình học. Đặc biệt, với mọi $s\in S$, thớ
$f^{-1}(s)$ của $X$ trên $s$ là một lược đồ đại số trên trường thặng dư
$\kappa(s)$ của vành địa phương $\Osh_s$ tại $s$ trong $S$. Bởi vậy, trong một
chừng mực nào đó, có thể xem $X$ như một họ các \fquote{không gian đại số}
$f^{-1}(s)$, trong đó tham số $s$ chạy trên $S$ (nghĩa là, theo quan điểm địa
phương, trên tập các iđêan nguyên tố của một vành đã cho). Dĩ nhiên, các
$\kappa(s)$ có thể có những đặc số khác nhau. Nếu
$S=\operatorname{Spec}(k)$, với $k$ là một trường, ta về căn bản thu lại khái
niệm thông thường về \fquote{không gian đại số}, chỉ khác ở chỗ giờ đây bó cấu
trúc có thể có các phần tử lũy linh.

\src{195}
Lấy cảm hứng từ các khái niệm quen thuộc, ta định nghĩa khái niệm
\emph{cấu xạ xạ ảnh} và, tổng quát hơn, \emph{cấu xạ proper}. Một cấu xạ như
vậy thuộc kiểu hữu hạn; hơn nữa, nó biến các tập con đóng thành các tập con
đóng và giữ nguyên tính chất này sau mọi phép thay đổi cơ sở.

Với $X$ là một lược đồ (Noether, như vẫn giả sử), bó $\Osh_X$ là một
\emph{bó vành coherent} theo nghĩa của [2]. Do đó, các bó môđun coherent trên
$X$ cũng chính là những bó đẳng cấu địa phương với đối hạch của một đồng cấu
$\Osh_X^{m}\longrightarrow\Osh_X^{n}$.

\fsection{2}{Lược đồ hình thức}

Cho $X$ là một lược đồ và $X'$ là một tập con đóng của $X$. Khi đó tồn tại một
bó con coherent $J$ của $\Osh_X$ sao cho
$X'=\operatorname{Supp}(\Osh_X/J)$ (thậm chí còn có một bó lớn nhất như vậy).
Khi được trang bị $\Osh_X/J$, không gian $X'$ trở thành một lược đồ, ký hiệu
$X_0$; lược đồ ấy được gọi là \fquote{lược đồ con đóng của $X$}. Với mọi $n$,
ta cũng có thể xét $X'$ được trang bị $\Osh_X/J^{n+1}$, ký hiệu $X_n$; đây là
một tiền lược đồ con đóng của $X$, có tập nền vẫn là $X'$ nhưng mang một bó
cấu trúc khác, cụ thể là $\Osh_{X_n}=\Osh_X/J^{n+1}$. Hiển nhiên, các
$\Osh_{X_n}$ lập thành một hệ xạ ảnh các bó vành trên $X'$; giới hạn xạ ảnh
$\overline{\Osh}_X$ của hệ ấy được gọi là \emph{hoàn thành hình thức của
$\Osh_X$ dọc theo $X'$}. Khi được trang bị bó vành này, $X'$ được gọi là
\emph{hoàn thành hình thức của $X$ dọc theo $X'$}; đó là một không gian có
vành, nhưng nói chung không phải là một lược đồ. Với mọi bó coherent $F$ trên
$X$, tương tự ta có thể xét hoàn thành hình thức
$\overline F=\varprojlim_n F_n$ của $F$ dọc theo $X'$, trong đó
$F_n=F\otimes_{\Osh_X}\Osh_X/J^{n+1}$; đây là một bó môđun trên
$\overline X$. Các tiết diện của nó được gọi là \fquote{các tiết diện hình
thức của $F$ dọc theo $X'$} và được đồng nhất với các phần tử của
$\varprojlim_n\Gamma(X',F_n)$. Với $F=\Osh_X$, ta thu được các
\fquote{hàm chỉnh hình} của $X$ dọc theo $X'$ theo nghĩa Zariski (chúng tôi
không dùng thuật ngữ ấy vì nó xung đột với thuật ngữ cổ điển).

Ta gọi một không gian tôpô $\mathfrak X$, được trang bị một bó các vành tôpô
$\Osh_{\mathfrak X}$ thỏa mãn điều kiện sau, là một \emph{lược đồ hình thức}
(ngầm hiểu: Noether): có một đẳng cấu các bó vành tôpô
$\Osh_{\mathfrak X}=\varprojlim_n\Osh_n$, trong đó các $\Osh_n$ lập thành một
hệ xạ ảnh các bó vành trên $\mathfrak X$, mỗi bó biến $\mathfrak X$ thành một
lược đồ $\mathfrak X_n$, và với $m\geq n$, đồng cấu
$\Osh_m\longrightarrow\Osh_n$ là toàn ánh, có hạch bằng $J_m^{n+1}$, trong đó
$J_m$ là

\src{196}
hạch của $\Osh_m\longrightarrow\Osh_0$. Khi đó người ta chứng minh được rằng
$\Osh_{\mathfrak X}$ là một \emph{bó coherent các vành địa phương Noether}.

Theo các định nghĩa, một hoàn thành hình thức $\overline X$ như trên là một
lược đồ hình thức; ngược lại, mọi lược đồ hình thức đều \emph{ở địa phương}
thuộc dạng này. Thật vậy, cho một lược đồ hình thức \emph{affine} (nghĩa là
$\mathfrak X_0$ affine, kéo theo mọi $\mathfrak X_n$ đều affine) tương đương
với cho một vành tôpô Noether, tách và đầy đủ đối với tôpô $J$-adic.

Các định nghĩa thông thường về cấu xạ, cấu xạ kiểu hữu hạn, cấu xạ proper,
v.v.\ của các lược đồ thông thường được mở rộng không khó khăn sang các lược
đồ hình thức.

\fsection{3}{Ba định lý cơ bản}

Cho $f:X\longrightarrow Y$ là một cấu xạ proper của các lược đồ (Noether, như
vẫn giả sử), $Y'$ là một tập con đóng của $Y$, còn $X'$ là nghịch ảnh của nó
trong $X$; xét các hoàn thành hình thức tương ứng $\overline Y$ và
$\overline X$. Khi đó $f$ cảm sinh một cấu xạ
$\overline f:\overline X\longrightarrow\overline X$ của các lược đồ hình thức,
và cấu xạ này cũng proper. Nếu $F$ là một bó coherent trên $X$, thì
$\overline F$ là một bó coherent trên $\overline X$. Trong Định lý 1, ta quên
$X,Y,F$ và chỉ xét cấu xạ proper $\overline f$ của các lược đồ hình thức cùng
bó coherent $\overline F$ trên $\overline X$. (Tuy nhiên, báo cáo viên mới chỉ
viết một chứng minh đầy đủ cho trường hợp xuất phát từ $X,Y,f,F$.)

\result{ĐỊNH LÝ}{1 (định lý hữu hạn)}
\begin{enumerate}
  \item[\textit{i.}] Các $R^{q}\overline f_{*}(\overline F)$ là những bó
    coherent trên $\overline Y$.
  \item[\textit{ii.}] Các đồng cấu tự nhiên
    \[
      R^{q}\overline f_{*}(\overline F)
      \longrightarrow \varprojlim_n R^{q}f_{n*}(F_n)
    \]
    là những đẳng cấu.
\end{enumerate}

Trong mệnh đề này, ta giả sử đã chọn một bó con coherent $J$ của $\Osh_Y$ xác
định $Y'$; lấy nghịch ảnh, ta được một bó con coherent của $\Osh_X$ xác định
$X'$, và do đó xác định $F_n,X_n,Y_n$ cùng
$f_n:X_n\longrightarrow Y_n$ như trong mục số 2. Nếu xuất phát từ một cấu xạ
proper tùy ý giữa hai lược đồ hình thức, những thay đổi nhỏ cần thực hiện khi
viết tường minh các ký hiệu này là hiển nhiên.

Định lý 1 chỉ liên quan đến \fquote{đối đồng điều hình thức}. Định lý tiếp theo

\src{197}
liên hệ nó với \fquote{đối đồng điều đại số} và tương tự một định lý quen thuộc
của Serre [4] về sự so sánh giữa đối đồng điều đại số và đối đồng điều giải
tích.

\result{ĐỊNH LÝ}{2 (Định lý so sánh thứ nhất)} Các
$R^{q}f_{*}(F)$ là những bó coherent trên $Y$ (đây là một trường hợp riêng của
Định lý 1), và các đồng cấu tự nhiên
\[
  \overline{R^{q}f_{*}(F)}
  \longrightarrow \varprojlim_n R^{q}f_{n*}(F_n)
\]
là những đẳng cấu.

\result{HỆ QUẢ}{1} Ta có các đẳng cấu chính tắc:
\[
  \overline{R^{q}f_{*}(F)}=R^{q}\overline f_{*}(\overline F).
\]

Với $q=0$, hệ quả này là một khái quát của \fquote{định lý cơ bản về các hàm
chỉnh hình} của Zariski; từ đó ta sẽ suy ra một khái quát của
\fquote{định lý liên thông} của Zariski. Cũng cần lưu ý rằng, trong khi Định
lý 1 (\textit{ii}) là tầm thường với $q=0$, điều đó hoàn toàn không đúng đối
với Định lý 2 hay với phát biểu tương đương của nó (Hệ quả 1). Thật vậy, chứng
minh tiến hành bằng quy nạp lùi theo $q$ (nó là tầm thường khi $q$ lớn, vì khi
ấy cả hai vế đều bằng không), nên trường hợp $q=0$ xuất hiện ở bước cuối của
quy nạp và, có thể nói, là trường hợp \fquote{khó nhất}.

\result{HỆ QUẢ}{2} Giả sử $Y=\operatorname{Spec}(A)$ và $Y'$ được xác định bởi
iđêan $J$ của $A$. Khi đó, với mọi bó coherent $F$ trên $X$, các
$H^{q}(X,F)$ là những $A$-môđun hữu hạn sinh, và các hoàn thành $J$-adic của
chúng là các $H^{q}(\overline X,\overline F)$.

Cuối cùng, áp dụng hệ quả này cho $H=\Homsh_{\Osh_X}(F,G)$, ta thu được:

\result{HỆ QUẢ}{3} Giả sử $Y=\operatorname{Spec}(A)$ và $Y'$ được xác định bởi
iđêan $J$ của $A$. Cho $F,G$ là hai bó coherent trên $X$; khi đó
$\operatorname{Hom}(F,G)$ là một $A$-môđun hữu hạn sinh, và hoàn thành
$J$-adic của nó được đồng nhất với
$\operatorname{Hom}(\overline F,\overline G)$.

Dĩ nhiên, ánh xạ tự nhiên
$\operatorname{Hom}(F,G)\longrightarrow
\operatorname{Hom}(\overline F,\overline G)$ là ánh xạ gán cho một đồng cấu
$u:F\longrightarrow G$ phần mở rộng \fquote{theo liên tục}
$\overline u:\overline F\longrightarrow\overline G$ của nó (nhờ đó
$\overline F$ trở thành một hàm tử theo $F$).

Bây giờ giả sử $A$ tách và đầy đủ đối với tôpô $J$-adic; khi đó Hệ quả 2 và
Hệ quả 3 ở trên cho:
\[
  H^{q}(X,F)=H^{q}(\overline X,\overline F),
  \qquad
  \operatorname{Hom}(F,G)
  =\operatorname{Hom}(\overline F,\overline G).
\]

\src{198}
Đẳng thức cuối cùng này cho thấy phạm trù các bó coherent trên $X$ được đồng
nhất với một \emph{phạm trù con} (các cấu xạ là những cấu xạ cảm sinh) của
phạm trù các bó coherent trên $\overline X$. Thật ra, ta còn có kết quả mạnh
hơn sau đây:

\result{ĐỊNH LÝ}{3} Một bó môđun trên $\overline X$ là coherent khi và chỉ khi
nó đẳng cấu với một bó có dạng $\overline F$, trong đó $F$ là một bó coherent
trên $X$ (được xác định duy nhất đến đẳng cấu chính tắc theo Định lý 2, Hệ quả
3). [Nhắc lại rằng lúc này $Y=\operatorname{Spec}(A)$, với $A$ là một vành
tôpô Noether, tách và đầy đủ đối với tôpô $J$-adic.]

\result{HỆ QUẢ}{1} Các lược đồ con đóng của $X$ tương ứng song ánh với các
lược đồ con hình thức đóng của $\overline X$.

Thật vậy, chúng lần lượt tương ứng với các bó con coherent của $\Osh_X$ và
$\overline{\Osh}_X$. Xem đồ thị của các cấu xạ như những lược đồ con đóng, từ
Hệ quả 1 ta suy ra:

\result{HỆ QUẢ}{2} Cho $X,Z$ là hai lược đồ proper trên $A$ (một vành Noether
$J$-adic, tách và đầy đủ). Khi đó ánh xạ
$g\longmapsto\overline g$ xác định một tương ứng song ánh giữa các $Y$-cấu xạ
từ $X$ đến $Z$ và các $\overline Y$-cấu xạ từ $\overline X$ đến
$\overline Z$.

Nói cách khác, các lược đồ đại số proper trên $A$ tạo thành một
\emph{phạm trù con} (các cấu xạ là những cấu xạ cảm sinh) của phạm trù các
lược đồ hình thức proper trên $\overline Y$. Tuy nhiên, cần lưu ý rằng có những
lược đồ hình thức proper trên $\overline Y$ không \fquote{đại số hóa được},
nghĩa là không đẳng cấu với một $\overline X$ nào, trong đó $X$ proper trên
$A$ (cũng như có những đa tạp giải tích phức compact không phát sinh từ các đa
tạp đại số xác định trên trường số phức). Những lược đồ hình thức như vậy xuất
hiện một cách tự nhiên trong \fquote{lý thuyết moduli}. Dẫu vậy, ta ghi nhận
một trường hợp riêng đáng chú ý trong đó một lược đồ hình thức đại số hóa được:

\result{ĐỊNH LÝ}{4} Cho $A$ là một vành địa phương Noether đầy đủ có trường
thặng dư $k$, và cho $\mathfrak X$ là một lược đồ hình thức proper trên $A$
(được trang bị tôpô $\mathfrak m(A)$-adic). Giả sử:
\begin{enumerate}
  \item[\textit{i.}] các vành địa phương của $\Osh_{\mathfrak X}$ là những
    $A$-môđun phẳng; tương đương, nếu trang bị $\Osh_{\mathfrak X}$ và $A$
    các lọc xác định bởi các lũy thừa của iđêan cực đại của $A$, thì các đối
    tượng phân bậc liên kết thỏa mãn hệ thức
    \[
      \operatorname{gr}(\Osh_{\mathfrak X})
      \simeq
      \operatorname{gr}^{0}(\Osh_{\mathfrak X})
      \otimes_k\operatorname{gr}(A).
    \]
  \item[\textit{ii.}] khi xem
    $\mathfrak X_0=\mathfrak X\otimes_A k$ như một lược đồ đại số trên $k$,
    ta có
    \[
      H^2(\mathfrak X_0,\Osh_{\mathfrak X_0})=0.
    \]
  \item[\textit{iii.}] $\mathfrak X_0$ là xạ ảnh.
\end{enumerate}
Khi đó $\mathfrak X$ đại số hóa được; chính xác hơn, nó đẳng cấu với
$\overline X$, trong đó $X$ là một $A$-lược đồ xạ ảnh.

Các điều kiện \textit{(ii)} và \textit{(iii)} đặc biệt được thỏa mãn nếu
$\mathfrak X_0$ là một đường cong trơn trên $k$; vì thế Định lý 4 có thể được
áp dụng, chẳng hạn, trong \fquote{lý thuyết moduli} của các đường cong có
giống cho trước\ldots{} Ta phác họa cách chứng minh Định lý 4. Người ta chứng
minh (xem Mệnh đề 3 dưới đây) rằng \textit{(i)} và \textit{(ii)} kéo theo:
mọi bó coherent trên $\mathfrak X_0$ đẳng cấu địa phương với bó cấu trúc đều
nhận được bằng phép rút gọn từ một bó cùng loại trên $\mathfrak X$. Bắt đầu từ
một bó \fquote{ample} trên $\mathfrak X_0$ (tồn tại theo \textit{(iii)}), ta
nâng nó thành một bó khả nghịch trên $\mathfrak X$; rồi dùng Định lý 1 để chứng
minh rằng một lũy thừa thích hợp của bó ấy xác định một phép nhúng
$\mathfrak X$ vào hoàn thành hình thức của một lược đồ $\mathbf P_A^r$
(\fquote{không gian xạ ảnh chuẩn} chiều $r$ trên $A$).

Về chứng minh các Định lý 1 đến 3, xem [1].

\src{199}
\fsection{4}{Áp dụng vào định lý liên thông và \fquote{Định lý chính} của
Zariski}

Cho $f:X\longrightarrow Y$ là một cấu xạ proper của các lược đồ. Theo định lý
hữu hạn, $f_*(\Osh_X)=\mathcal A$ là một bó coherent trên $Y$; hơn nữa, đó là
một bó các đại số giao hoán, nên tương ứng với một $Y$-lược đồ
$g:Y'\longrightarrow Y$ hữu hạn trên $Y$ (được xác định bởi các điều kiện: nó
affine trên $Y$, nghĩa là nghịch ảnh của một tập mở affine là affine, và
$g_*(\Osh_{Y'})=\mathcal A$). Hiển nhiên khi đó $f$ phân tích chính tắc thành
$f=gf'$, trong đó $f':X\longrightarrow Y'$ là một cấu xạ từ $X$ đến $Y'$ thỏa
mãn $f'_*(\Osh_X)=\Osh_{Y'}$. Phân tích này của $f$ được gọi là
\emph{phân tích Stein} của $f$. Áp dụng Định lý so sánh thứ nhất, Hệ quả 1,
cho $f'$ và cho phần của $Y'$ rút gọn về một điểm $y'$, ta thấy
$f'^{-1}(y')=X'$ liên thông (nếu không, các tiết diện hình thức của $X$ dọc
theo $X'$ sẽ không tạo thành một vành địa phương, trong khi hoàn thành của
$f'_*(\Osh_X)_{y'}=\Osh_{Y',y'}$ lại là địa phương!). Ta đã chứng minh:

\result{ĐỊNH LÝ}{5 (\fquote{Định lý liên thông} của Zariski)} Cho
$f:X\longrightarrow Y$ là một cấu xạ proper. Khi đó $f$ phân tích duy nhất
(đến đẳng cấu) thành $f=gf'$, trong đó $g:Y'\longrightarrow Y$ là hữu hạn và
$f':X\longrightarrow Y'$ thỏa mãn
$f'_*(\Osh_X)=\Osh_{Y'}$ (do đó
$g_*(\Osh_{Y'})=f_*(\Osh_X)$). Các thớ của $f'$ liên thông; nói cách khác, tập
các thành phần liên thông của một thớ $f^{-1}(y)$ của $f$ tương ứng song ánh
với tập các điểm của $Y'$ nằm trên $y$, tức là với tập các iđêan cực đại của
$f_*(\Osh_X)_y$.

Các biến thể quen thuộc của định lý liên thông được suy ra ngay lập tức. Ở đây
ta chỉ ghi nhận hệ quả sau.

\src{200}
\result{HỆ QUẢ}{1} Một điểm $x$ của $X$ cô lập trong thớ $f^{-1}(y)$ khi và
chỉ khi thớ $f'^{-1}(y')$ (trong đó $y'=f'(x)$) rút gọn về $x$; tương đương,
$f'$ cảm sinh một đẳng cấu từ một lân cận của $x$ lên một lân cận của $y'$.
Tập các điểm như vậy là một tập mở $U$, và $f'$ cảm sinh một đẳng cấu từ $U$
lên một tập mở của $Y'$.

Để chứng minh $f'$ là một đẳng cấu địa phương tại $x$, ta lưu ý rằng $f'$ cảm
sinh một đẳng cấu $\Osh_{y'}\longrightarrow\Osh_x$, như suy ra từ
$f'_*(\Osh_X)=\Osh_{Y'}$ và từ sự kiện sau: vì $f'$ là một ánh xạ đóng có thớ
trên $y'$ rút gọn về $x$, các tập $f'^{-1}(V)$ chạy qua một hệ cơ bản các lân
cận của $x$ khi $V$ chạy qua một hệ cơ bản các lân cận của $y'$. Từ đó suy ra
ngay kết quả sau, do Chevalley chứng minh trong trường hợp
\fquote{hình học}.

\result{HỆ QUẢ}{2} Một cấu xạ $f$ là hữu hạn khi và chỉ khi nó proper và các
thớ của nó hữu hạn.

Thật vậy, trong trường hợp ấy $f'$ là một đẳng cấu theo kết quả vừa chứng minh.

Cho $f:X\longrightarrow Y$ là một cấu xạ không nhất thiết proper, nhưng giả sử
$X$ nằm như một tập mở trong một $Y$-lược đồ proper
$\widetilde f:\widetilde X\longrightarrow Y$ (điều này đặc biệt đúng nếu $f$
là quasi-xạ ảnh). Áp dụng Hệ quả 1, ta thấy $\widetilde f'$ cảm sinh một đẳng
cấu từ tập $U$ các điểm của $X$ cô lập trong thớ của chúng lên một tập mở của
$Y'$ (và $U$ quả thật là tập mở). Từ đó suy ra dạng toàn cục sau đây của
\fquote{Định lý chính} của Zariski.

\result{ĐỊNH LÝ}{6 (\fquote{Định lý chính} của Zariski)} Cho
$f:X\longrightarrow Y$ là một cấu xạ kiểu hữu hạn. Khi đó tập $U$ các điểm
của $X$ cô lập trong thớ của chúng là mở; nếu $f$ quasi-xạ ảnh, thì $U$ đẳng
cấu trên $Y$ với một phần mở của một lược đồ $Y'$ hữu hạn trên $Y$.

Vì một cấu xạ kiểu hữu hạn là affine địa phương, và do đó \textit{a fortiori}
quasi-xạ ảnh địa phương, các dạng địa phương quen thuộc của Định lý chính được
suy ra ngay từ Định lý 6.

\src{201}
\result{MỆNH ĐỀ}{1} Cho $f:X\longrightarrow Y$ là một cấu xạ proper và $F$
là một bó trên $X$, coherent và phẳng trên $Y$. Cho $y\in Y$, cho $q$ là một
số nguyên, và giả sử $H^q(X_y,F_y)=0$. Khi đó $R^qf_*(F)$ triệt tiêu trong một
lân cận của $y$, và với mọi $n$, đồng cấu tự nhiên
\[
  R^{q-1}f_*(F)_y\longrightarrow
  H^{q-1}(X_y,F/\mathfrak m_y^nF)
\]
là toàn ánh.

Đặc biệt, nếu $f$ là một cấu xạ phẳng (nghĩa là $\Osh_X$ phẳng trên $Y$), thì
mọi bó coherent địa phương tự do $F$ trên $X$ đều phẳng trên $Y$. Cho $F,G$ là
hai bó như vậy; áp dụng Mệnh đề 1 cho $\Homsh_{\Osh_X}(F,G)$ với $q=1$, ta
được:

\result{ĐỊNH LÝ}{7} Cho $f$ là một cấu xạ proper và phẳng, $F,G$ là hai bó
coherent địa phương tự do trên $X$, $y\in Y$, và giả sử
\[
  H^1\bigl(X_y,\Homsh_{\Osh_{X_y}}(F_y,G_y)\bigr)=0,
\]
thì mọi đồng cấu $u_0:F_y\longrightarrow G_y$ đều được cảm sinh bởi một đồng
cấu
\[
  u:F|V\longrightarrow G|V,
  \qquad V=f^{-1}(U),
\]
trong đó $U$ là một lân cận của $y$.

\result{HỆ QUẢ}{} Nếu $u_0$ là một đẳng cấu (tương ứng, một đơn cấu, một toàn
cấu), thì $u$ cũng vậy khi $U$ đủ nhỏ.

Đặc biệt:

\src{202}
\result{HỆ QUẢ}{2} Cho $E_0$ là một bó coherent địa phương tự do trên $X_y$
thỏa mãn
\[
  H^1\bigl(X_y;\Homsh_{\Osh_{X_y}}(E_0,E_0)\bigr)=0,
\]
thì hai bó địa phương tự do mà các hạn chế xuống $X_y$ đều đẳng cấu
với $E_0$ sẽ đẳng cấu trong một lân cận của $X_y$.

Do đó:

\result{HỆ QUẢ}{3} Giả sử $H^1(X_y,\Osh_{X_y})=0$. Khi đó hai bó khả nghịch
trên $X$ (nghĩa là địa phương đẳng cấu với $\Osh_X$) mà các hạn chế xuống
$X_y$ đẳng cấu với nhau thì bản thân chúng cũng đẳng cấu.

Suy ra:

\result{MỆNH ĐỀ}{2} Cho $Y$ là một lược đồ liên thông và $E$ là một bó
coherent địa phương tự do trên $Y$; xét phân thớ không gian xạ ảnh
$X=\mathbf P(E)$ liên kết với $E$, được trang bị bó khả nghịch quen thuộc
$\Osh_X(1)$. Khi đó mọi bó khả nghịch $L$ trên $X$ đều đẳng cấu với một bó có
dạng
\[
  f^*(L')\otimes\Osh_X(n),
\]
trong đó $L'$ là một bó khả nghịch trên $Y$ và $n$ là một số nguyên. Số nguyên
vừa nêu được xác định duy nhất, còn $L'$ được xác định đến đẳng cấu.

Hệ quả 3 ở trên chứng minh rằng $L$ đẳng cấu với một $\Osh_X(n)$ nào đó trong
một lân cận của mỗi thớ. Phần còn lại gần như chỉ là hình thức.

Mệnh đề 2 cho phép xác định các $Y$-cấu xạ từ $X=\mathbf P(E)$ đến một phân
thớ xạ ảnh khác. Đặc biệt, ta được:

\result{HỆ QUẢ}{} Cho $u$ là một tự đẳng cấu của $X=\mathbf P(E)$. Khi đó tồn
tại một bó khả nghịch $L'$ trên $Y$ và một đẳng cấu $v$ từ $E$ lên
$E\otimes L'$ sao cho $u$ là đẳng cấu tương ứng
\[
  \mathbf P(E)\xrightarrow{\sim}\mathbf P(E\otimes L')=\mathbf P(E);
\]
cặp $(v,L')$ được xác định đến đẳng cấu.

Cho $\Gamma$ là tập các lớp của những bó khả nghịch $L'$ trên $Y$ sao cho
$E\otimes L'$ đẳng cấu với $E$. Các phần tử của tập này đều là phần tử xoắn:
thật vậy, nếu $n$ là hạng của $E$, thì lấy lũy thừa ngoài bậc $n$ cho thấy
nhất thiết phải có
$L'^{\otimes n}\simeq\Osh_Y$. Vì thế, hệ quả có thể được phát biểu bằng dãy
khớp các nhóm
\[
  e\longrightarrow
  \operatorname{Aut}(E)/\Gamma(Y,\Osh_Y^*)
  \longrightarrow \operatorname{Aut}_Y(X)
  \longrightarrow \Gamma\longrightarrow e
\]
(dãy này cũng suy ra từ dãy khớp đối đồng điều phát sinh từ dãy khớp sau đây
của các bó nhóm
\[
  e\longrightarrow \Osh_Y^*
  \longrightarrow \underline{\operatorname{Aut}}(E)
  \longrightarrow \underline{\operatorname{Aut}}_Y(X)
  \longrightarrow e,
\]
trong đó $\Osh_Y^*$ là bó các \fquote{đơn vị} của $\Osh_Y$, được đồng nhất với
tâm của $\underline{\operatorname{Aut}}(E)$).

\src{203}
\fsection{6}{Ứng dụng vào các định lý tồn tại và duy nhất cho bó và lược đồ
trên một vành đầy đủ theo tôpô \texorpdfstring{$J$}{J}-adic}

Định lý 7 đã cho một kết quả duy nhất đối với các bó coherent địa phương tự do,
bằng cách dùng các Định lý 1 và 2. Giờ đây, nhờ Định lý 3, ta suy ra các định
lý tồn tại cho bó, cấu xạ lược đồ hoặc lược đồ. Trong phần còn lại, $A$ ký hiệu
một vành Noether địa phương, tách và đầy đủ. Phương pháp tổng quát luôn là thực
hiện các phép dựng hình thức, về bản chất tức là làm hình học đại số trên một
vành Artin, rồi dùng ba định lý cơ bản để rút ra các kết luận có tính chất
“đại số”.

\result{MỆNH ĐỀ}{3} Cho $\mathfrak X$ là một lược đồ hình thức proper và phẳng
trên $A$, và cho $F_0$ là một bó địa phương tự do trên $X_0$ sao cho
\[
  H^2\bigl(X_0,\Homsh_{\Osh_{X_0}}(F_0,F_0)\bigr)=0.
\]
Khi đó tồn tại một bó địa phương tự do $F$ trên $\mathfrak X$ mà hạn chế xuống
$X_0$ đẳng cấu với $F_0$. (Hơn nữa, $F$ này duy nhất đến đẳng cấu nếu
$H^1(X_0,\Homsh_{\Osh_{X_0}}(F_0,F_0))=0$.)

Ta lần lượt dựng các bó địa phương tự do $F_n$ trên các $X_n$, mỗi bó cảm sinh
bó đứng trước. Cản trở đối với việc dựng $F_n$ nằm trong
\[
  H^2\bigl(X_0,\Homsh_{\Osh_{X_0}}(F_0,F_0)\bigr)
  \otimes_{A/J}(J^n/J^{n+1}),
\]
nên bằng không theo giả thiết. Bây giờ áp dụng Định lý 3, ta được:

\result{HỆ QUẢ}{1} Cho $X$ là một lược đồ proper và phẳng trên $A$, và cho
$F_0$ như trên. Khi đó tồn tại một bó địa phương tự do $F$ trên $X$ mà hạn chế
xuống $X_0$ đẳng cấu với $F_0$. Hơn nữa, $F$ này duy nhất đến đẳng cấu nếu
$H^1(X_0,\Homsh_{\Osh_{X_0}}(F_0,F_0))=0$.

Cho $X_0$ là một lược đồ kiểu hữu hạn trên trường $k$; giả sử $X_0$ trơn (ở
đây hiểu là trơn tuyệt đối) trên $k$, nhưng không nhất thiết proper trên $k$.
Cho $A$ là một vành Artin địa phương có trường thặng dư $k$. Ta muốn tìm các
lược đồ $X$ phẳng trên $A$ sao cho $X\otimes_A k=X_0$ (đây là điểm xuất phát
của “lý thuyết moduli” hoặc của “các biến thiên cấu trúc” của $X_0$). Việc cho
một $X$ như vậy tương đương với việc cho, trên không gian tôpô $X_0$, một bó
$\Osh_X$ được trang bị các cấu trúc sau:
\begin{enumerate}
  \item[\textit{i.}] Nó là một bó các $A$-đại số.
  \item[\textit{ii.}] Nó được trang bị một đồng cấu tăng cường
    $\Osh_X\longrightarrow\Osh_{X_0}$ (tương thích với các cấu trúc
    $A$-đại số); các dữ liệu này phải thỏa các điều kiện sau: phép tăng cường
    cảm sinh một đẳng cấu
    $\Osh_X\otimes_A k\xrightarrow{\sim}\Osh_{X_0}$; $\Osh_X$ phẳng trên
    $A$, nghĩa là đối tượng phân bậc liên kết với $\Osh_X$, được lọc bởi các
    lũy thừa của iđêan cực đại $\mathfrak m$ của $A$, đẳng cấu với
    $\operatorname{gr}^0(\Osh_X)\otimes_k\operatorname{gr}(A)$; tương đương,
    ta có các đẳng cấu
    \[
      \mathfrak m^n\Osh_X/\mathfrak m^{n+1}\Osh_X
      =\Osh_{X_0}\otimes_k(\mathfrak m^n/\mathfrak m^{n+1}).
    \]
\end{enumerate}
Sự kiện cơ bản là như sau:

\src{204}
\result{ĐỊNH LÝ}{8} Cho $X_0$ là một lược đồ kiểu hữu hạn và trơn trên trường
$k$, và giả sử $X_0$ affine. Cho $A$ là một vành Artin địa phương có trường
thặng dư $k$. Khi đó tồn tại một $A$-lược đồ $X$ phẳng trên $A$ sao cho
$X\otimes_A k=X_0$, và hai lược đồ bất kỳ như vậy tất yếu đẳng cấu.

Lưu ý rằng đẳng cấu nói trên không chính tắc, vì nói chung $X$ sẽ có những tự
đẳng cấu trên $A$ không tầm thường cảm sinh đồng nhất trên $X_0$. Mặt khác,
nói chung không có cách chọn “chính tắc” một $X$ thỏa các điều kiện đã nêu,
trừ khi $A$ là một $k$-đại số (trường hợp cùng đặc số); khi ấy có thể lấy
$X=X_0\otimes_k A$, nghĩa là $\Osh_X=\Osh_{X_0}\otimes_k A$ (dù $X_0$ affine
hay không). Trong trường hợp khác đặc số, khi $X_0$ không affine, nói chung ta
không biết liệu có thể “nâng” $X_0$ thành một $X$ xác định trên $A$ hay không.
Tuy nhiên, cho $n$ là một số nguyên $>0$, đặt
$A_{n-1}=A/\mathfrak m^n$, và giả sử đã nâng $X_0$ thành một
$A_{n-1}$-lược đồ phẳng $X_{n-1}$; ta muốn nâng $X_{n-1}$ thành một
$A_n$-lược đồ phẳng $X_n$. Theo Định lý 8, ta đã biết điều này thực hiện được
ở địa phương. Hơn nữa, dễ kiểm tra rằng nếu $U_n$ nâng một tập mở $U_{n-1}$
của $X_{n-1}$, thì bó nhóm các tự đẳng cấu của $U_n$ (cảm sinh đồng nhất trên
$U_{n-1}$) đẳng cấu chính tắc với
\[
  \Theta_{X_0/k}\otimes_k(\mathfrak m^n/\mathfrak m^{n+1})
\]
được hạn chế xuống $U_{n-1}$, và đặc biệt có tính giao hoán
($\Theta_{X_0/k}$ ký hiệu bó các mầm của các $k$-đạo hàm trên $X_0$). Từ đó
dễ suy ra rằng có một cản trở đối với việc dựng $X_n$ nâng $X_{n-1}$, nằm
trong
\[
  H^2(X_0,\Theta_{X_0/k})\otimes_k
  \mathfrak m^n/\mathfrak m^{n+1}.
\]
Do đó:

\result{HỆ QUẢ}{1} Cho $X_0$ là một lược đồ kiểu hữu hạn trơn trên $k$, và
giả sử $H^2(X_0,\Theta_{X_0/k})=0$. Khi đó, với mọi vành Artin địa phương
$A$ có trường thặng dư $k$, tồn tại một $A$-lược đồ phẳng $X$ sao
cho $X\otimes_A k=X_0$.

Hơn nữa, nếu đã tìm được một $X$ phẳng trên $A$ nâng $X_0$, thì theo Định lý
8, tập các lớp (đến đẳng cấu) của các $A$-lược đồ phẳng nâng $X_0$ được đồng
nhất với
\[
  H^1\bigl(X_0,\underline{\operatorname{Aut}}(X)\bigr),
\]
trong đó, dĩ nhiên, $\underline{\operatorname{Aut}}(X)$ ký hiệu bó các mầm
của những tự đẳng cấu của bó $A$-đại số $\Osh_X$ tương thích với phép tăng
cường. Lọc của $\Osh_X$ xác định một lọc của
$\underline{\operatorname{Aut}}(X)$; thương của bó này cho nhóm con thứ
$n$ trong lọc là $\underline{\operatorname{Aut}}(X_n)$. Đối tượng phân bậc
liên kết với lọc này có tính giao hoán và được đồng nhất với
$\Theta_{X_0/k}\otimes_k\operatorname{gr}(A)$. Đặc biệt, nếu
$\mathfrak m^{n+1}$ là lũy thừa đầu tiên của $\mathfrak m$ bằng không, thì
$F^n(\underline{\operatorname{Aut}}(X))$ (bậc cuối của lọc) nằm trong tâm của
$\underline{\operatorname{Aut}}(X)$ và đẳng cấu với
$\Theta_{X_0/k}\otimes_k\mathfrak m^n$; nó cũng là bó các mầm của những tự
đẳng cấu của $X$ cảm sinh đồng nhất trên
$X_{n-1}=X\otimes_A A/\mathfrak m^n$. Dùng các kết quả này, ta lập tức thu
được các mệnh đề sau:

\src{205}
\result{HỆ QUẢ}{2} Cho $X_0$ là một lược đồ kiểu hữu hạn trơn trên $k$, cho
$A$ là một vành Artin địa phương có trường thặng dư $k$ và iđêan cực đại
$\mathfrak m$, giả sử $\mathfrak m^{n+1}=0$, đặt
$A_{n-1}=A/\mathfrak m^n$, và cuối cùng cho $X_{n-1}$ là một
$A_{n-1}$-lược đồ phẳng sao cho $X_{n-1}\otimes_A k=X_0$. Khi đó tập các
lớp (đến một đẳng cấu cảm sinh đồng nhất trên $X_{n-1}$) của các $A$-lược đồ
phẳng $X_n$ sao cho $X_n\otimes_A A_{n-1}=X_{n-1}$ hoặc rỗng, hoặc là một
không gian thuần nhất chính dưới
\[
  H^1(X_0,\Theta_{X_0/k})\otimes_k\mathfrak m^n.
\]

(Lưu ý rằng nói chung không có điểm gốc ưu tiên trong không gian vừa nêu, vì
không có cách ưu tiên nào để nâng $X_{n-1}$ thành $X_n$.)

\result{HỆ QUẢ}{3} Cho $X_0$ là một lược đồ kiểu hữu hạn trơn trên $k$, và
giả sử $H^1(X_0,\Theta_{X_0/k})=0$. Khi đó, với mọi vành Artin địa phương
$A$ có trường thặng dư $k$, tồn tại nhiều nhất một (đến đẳng cấu) $A$-lược đồ
phẳng $X$ sao cho $X\otimes_A k=X_0$.

Các Hệ quả 1 và 3 lập tức kéo theo những mệnh đề bề ngoài tổng quát hơn, thu
được khi chỉ giả sử $A$ là một vành Noether địa phương đầy đủ có trường thặng
dư $k$, miễn là ta đưa $X$ vào dưới dạng một lược đồ hình thức trên $A$:

\result{ĐỊNH LÝ}{9} Cho $k$ là một trường và $X_0$ là một lược đồ kiểu hữu
hạn trơn trên $k$. Với mỗi vành Noether địa phương đầy đủ $A$ có trường thặng
dư $k$, gọi $F(A)$ là tập các lớp (đến một đẳng cấu cảm sinh đồng nhất trên
$X_0$) của các lược đồ hình thức $\mathfrak X$ trên $A$, kiểu hữu hạn và phẳng
trên $A$, sao cho $\mathfrak X\otimes_A k=X_0$. Với các ký hiệu ấy:
\begin{enumerate}
  \item[\textit{i.}] Nếu $H^1(X_0,\Theta_{X_0/k})=0$, thì với mọi $A$,
    $F(A)$ có nhiều nhất một phần tử.
  \item[\textit{ii.}] Nếu $H^2(X_0,\Theta_{X_0/k})=0$, thì với mọi $A$,
    $F(A)$ có ít nhất một phần tử.
\end{enumerate}

\src{206}
\result{HỆ QUẢ}{1} Giả sử $X_0$ proper trên $k$. Dưới điều kiện
\textit{(i)}, với mọi $A$ tồn tại nhiều nhất một (đến một đẳng cấu cảm sinh
đồng nhất trên $X_0$) lược đồ $X$ proper và phẳng trên $A$ sao cho
$X\otimes_A k=X_0$.

Ở đây ta dùng Định lý 3, Hệ quả 2. Chẳng hạn:

\result{HỆ QUẢ}{2} Nếu $X$ là một lược đồ proper và phẳng trên $A$ sao cho
$X\otimes_A k$ đẳng cấu với lược đồ xạ ảnh chuẩn $\mathbf P_k^r$ chiều $r$
trên $k$, thì $X$ đẳng cấu với $\mathbf P_A^r$.

(Kết quả này cũng có thể suy ra từ Mệnh đề 3, Hệ quả 1.)

\result{HỆ QUẢ}{3} Cho $X_0$ là một lược đồ xạ ảnh trơn trên $k$, và giả sử
\[
  H^2(X_0,\Osh_{X_0})=H^2(X_0,\Theta_{X_0/k})=0.
\]
Khi đó, với mọi $A$, tồn tại một lược đồ xạ ảnh $X$, phẳng trên $A$, sao cho
$X\otimes_A k=X_0$.

Kết hợp Định lý 9\textit{(ii)} với Định lý 4. Đặc biệt:

\result{HỆ QUẢ}{4} Cho $X_0$ là lược đồ của một đường cong đại số đầy đủ và
trơn trên $k$. Khi đó, với mọi vành Noether địa phương đầy đủ $A$ có trường
thặng dư $k$, tồn tại một “lược đồ đường cong trơn” $X$ trên $A$ sao cho
$X\otimes_A k=X_0$.

\result{NHẬN XÉT}{}

1) Các Hệ quả 3 và 4 đặc biệt đáng chú ý khi $k$ có đặc số $p\ne0$: ta lấy
$A$ là một vành định giá rời rạc có đặc số $0$ và trường thặng dư $k$; chẳng
hạn, “$A$ nhỏ nhất có thể”, tức là vành có iđêan cực đại sinh bởi $p$. (Thật
vậy, theo các định lý của Cohen, chỉ cần biết các Hệ quả 3 và 4 cho một vành
$A$ như vậy.) Về điểm này, cần lưu ý rằng theo các chuyên gia, chưa biết có
tồn tại những lược đồ trên một trường $k$ không phải là phép rút gọn modulo
$p$ của một lược đồ phẳng xác định trên một vành $A$ như vậy hay không. Ít
nhất, các kết quả của mục này cung cấp một phương pháp để khảo sát câu hỏi ấy
một cách có hệ thống. Trước hết nên xét xem cản trở đầu tiên xuất hiện trong
$H^2(X_0,\Theta_{X_0/k})$ có nhất thiết bằng không hay không.

\emph{Bổ sung năm 1962.} --- J.-P. Serre đã dựng một đa tạp xạ ảnh trơn chiều
3 trên một trường đóng đại số $k$ có đặc số $p>0$, không thu được bằng phép
rút gọn từ một lược đồ proper trên một vành địa phương nguyên có trường thặng
dư $k$ và trường phân thức đặc số không. Nghe nói Mumford đã tìm được một ví
dụ tương tự với một mặt xạ ảnh trơn.

2) Cần lưu ý rằng Định lý 3 và kỹ thuật tương ứng chỉ áp dụng khi vành cơ sở
(địa phương, để cố định ý tưởng) là đầy đủ.

Muốn chuyển từ các kết quả đã biết cho vành đầy đủ hóa của một vành địa phương
sang các kết quả tương ứng cho chính vành địa phương ấy, ta cần một “định lý
cơ bản” thứ tư, mà phát biểu dứt khoát của nó vẫn còn phải tìm.

\src{207}
3) Hãy so sánh các kết quả của mục này (đặc biệt là các Hệ quả 1 và 2 ở trên)
và của mục kế tiếp với các kết quả của Kodaira--Spencer về sự biến thiên của
các cấu trúc phức. Dùng định lý phỏng đoán vừa nói đến, dưới các điều kiện của
Hệ quả 1 nhưng không còn giả sử $A$ đầy đủ, ta hẳn có thể kết luận rằng tồn
tại một vành $A'$ chứa $A$, hữu hạn và không phân nhánh trên $A$, sao cho
$X\otimes_A A'$ và $X'\otimes_A A'$ đẳng cấu trên $A'$ (ở đây $X,X'$ là hai
$A$-lược đồ proper và phẳng đã cho, thỏa
$X\otimes_A k=X'\otimes_A k=X_0$). Ít nhất, điều này có thể chứng minh khi
$X_0=\mathbf P_k^r$, bằng cách dùng hệ quả của Mệnh đề 2. Trong mọi trường
hợp, khi $H^1(X_0,\Theta_{X_0/k})=0$, có thể chứng minh rằng các thớ của
$X,X'$ trên mọi điểm $y$ của $Y=\operatorname{Spec}(A)$ là đẳng cấu, ít nhất
sau khi chuyển sang bao đóng đại số của trường thặng dư $\kappa(y)$. (Có một
kết quả địa phương, dường như mạnh hơn, trong đó không còn cần giả sử $A$ là
địa phương.) Đối với “các biến thiên cấu trúc” của không gian xạ ảnh, ta cũng
nêu câu hỏi sau đây, do một bài toán tương ứng của Kodaira--Spencer gợi ra.
Cho $X$ là một lược đồ proper và phẳng trên vành địa phương nguyên $A$, có
trường phân thức $K$ và trường thặng dư $k$, và giả sử $X\otimes_A K$ đẳng
cấu với $\mathbf P_K^r$; liệu $X\otimes_A k=X_0$ có đẳng cấu với
$\mathbf P_k^r$ hay không (ít nhất trên bao đóng đại số của $k$)? Trong bài
toán này, có thể giả sử $A$ là một vành định giá rời rạc đầy đủ. Một bài toán
tương tự được đặt ra khi $X_0$ là một đa tạp abel.

\emph{Bổ sung năm 1962.} --- Khi ấy diễn giả đã bớt lạc quan hơn về các kết
quả phỏng đoán ở trên. Tuy nhiên, Hironaka đã giải bài toán về các biến thiên
cấu trúc của không gian xạ ảnh theo hướng khẳng định, còn Koizumi giải bài
toán tương tự cho các đa tạp abel.

\fsection{7}{Ứng dụng vào “lý thuyết moduli”}

Vì bản thân diễn giả mới chỉ vừa bắt đầu nghiên cứu lý thuyết này, ta buộc
phải giới hạn ở một vài chỉ dẫn. Để đơn giản, ta làm việc trên một trường $k$,
nghĩa là trong trường hợp cùng đặc số, mặc dù Định lý 8 dường như cũng cho phép
xử lý trường hợp tổng quát mà không có thay đổi bản chất nào. Hiện tại ta chưa
vượt quá giai đoạn “hình thức”; tuy nhiên, diễn giả hy vọng từ đó có thể dựng
được các lược đồ moduli thực sự trong một số trường hợp và, đặc biệt, với mỗi
số nguyên $g$, dựng một lược đồ trên các số nguyên đóng vai trò lược đồ moduli
phổ dụng cho các đường cong trơn có giống $g$.

\emph{Bổ sung năm 1962.} --- Mumford đã dựng các lược đồ moduli cho các đường
cong giống $g$ (xem SHMI). Hơn nữa, Định lý 10 chứng minh rằng các “lược đồ
Jacobi mức $n$” của lý thuyết moduli là không kỳ dị, và thậm chí trơn trên
$\mathbf Z$.

\src{208}
Trở lại tình huống và ký hiệu của Định lý 9, giờ đây giả sử $A$ là một đại số
địa phương hữu hạn chiều trên $k$, và để đơn giản trường này được giả sử đóng đại số.
Khi đó có thể xem $F(A)$ như một hàm tử hiệp biến theo $A$, nhận giá trị trong
phạm trù các tập hợp: một đồng cấu các $k$-đại số $A\longrightarrow B$ xác
định một ánh xạ $F(A)\longrightarrow F(B)$, vì mỗi $A$-lược đồ phẳng $X$ sao
cho $X\otimes_A k=X_0$ cho một $B$-lược đồ $X\otimes_A B$ có cùng các tính
chất. Giả sử có thể tìm được một $k$-đại số Noether địa phương đầy đủ
$\mathcal O$ và một đẳng cấu có tính hàm tử
\begin{equation}
  \tag{*}
  \operatorname{Hom}(\mathcal O,A)\xrightarrow{\sim}F(A)
\end{equation}
(trong đó vế thứ nhất ký hiệu các đồng cấu $k$-đại số). Dễ thấy một
$\mathcal O$ như vậy được xác định đến đẳng cấu chính tắc; khi ấy, ta gọi phổ
hình thức $\mathfrak Y$ của $\mathcal O$ (nghĩa là không gian tôpô chỉ gồm
một điểm và được trang bị một bó các vành tôpô thu về $\mathcal O$) là
\emph{lược đồ moduli hình thức} của $X_0$. (Lược đồ này không nhất thiết tồn
tại.) Cho $\mathfrak r$ là iđêan cực đại của $\mathcal O$; với mỗi $n$, đặt
$\mathcal O_n=\mathcal O/\mathfrak r^{n+1}$ (do đó $\mathcal O_0=k$). Khi
ấy, đồng cấu chính tắc $\mathcal O\longrightarrow\mathcal O_n$ là một phần tử
của $\operatorname{Hom}(\mathcal O,\mathcal O_n)$, nên xác định một phần tử
của $F(\mathcal O_n)$, tức là một $\mathcal O_n$-lược đồ phẳng $X_n$ có phép
rút gọn modulo $\mathfrak r$ bằng $X_0$. Các $X_n$ này suy ra từ nhau bằng
mở rộng vô hướng (ở đây tức là bằng các phép rút gọn), do đó chúng đến từ một
lược đồ hình thức $\mathfrak X$ xác định duy nhất trên lược đồ moduli hình
thức $\mathfrak Y$; $\mathfrak X$ phẳng trên $\mathfrak Y$ và
$\mathfrak X_0=X_0$. Khi đó, như thấy ngay, đẳng cấu (*) có được bằng cách
gắn với mỗi đồng cấu $k$-đại số $\mathcal O\longrightarrow A$ lớp của
$A$-lược đồ $\mathfrak X\otimes_{\mathcal O}A$ (nghĩa là, với mỗi cấu xạ các
$k$-lược đồ $Y'=\operatorname{Spec}(A)\longrightarrow\mathfrak Y$, ta gắn
$Y'$-lược đồ $\mathfrak X\otimes_{\mathfrak Y}Y'$ thu được bằng đổi cơ sở).
Hơn nữa, đẳng cấu (*) và mô tả vừa nêu vẫn còn đúng nếu chỉ giả sử $A$ là một
$k$-đại số Noether địa phương đầy đủ (không nhất thiết Artin). Dĩ nhiên, như
thường lệ, $\mathcal O$ hoàn toàn có thể có các phần tử lũy linh
\textit{a priori}, và dường như có thể tồn tại những trường hợp trong đó
$\mathcal O$ tự nó là Artin nhưng không đồng nhất với $k$. Điều này cho thấy
quan điểm Kodaira--Spencer (chỉ xét các $A$ là vành chính quy) không thích hợp
\textit{a priori} đến mức nào trong trường hợp tổng quát.

\src{209}
Ta còn phải nêu các điều kiện đủ để tồn tại một lược đồ moduli hình thức cho
$X_0$, được giả sử proper trên $k$. Nói chung, dễ nêu các điều kiện cần và đủ
đơn giản trên một hàm tử $A\longmapsto F(A)$ (từ các $k$-đại số địa phương
hữu hạn chiều đến các tập hợp) để nó có thể viết dưới dạng
$\operatorname{Hom}(\mathcal O,A)$ với một $\mathcal O$ thích hợp. Ta sẽ
không trình bày chi tiết các điều kiện ấy ở đây. Chỉ xin lưu ý rằng, trong
trường hợp đang xét, chúng áp đặt lên $X_0$ những điều kiện đối đồng điều không
tầm thường, và dường như khó có khả năng chúng luôn được thỏa mãn, dù diễn giả
chưa dựng được phản ví dụ nào. Trái lại, có vẻ hợp lý khi cho rằng điều kiện
$H^0(X_0,\Theta_{X_0/k})=0$ là điều kiện đủ (dù tuyệt nhiên không cần) cho sự
tồn tại của một lược đồ moduli hình thức. Ở đây ta chỉ phát biểu một định lý
trong một trường hợp đặc biệt đơn giản (định lý tương tự trong lý thuyết không
gian giải tích đã được biết rõ; xem Kodaira--Spencer), và định lý này được suy
ra không khó khăn từ các kết quả của mục trước:

\result{ĐỊNH LÝ}{10} Cho $X_0$ là một lược đồ proper và trơn trên trường $k$
sao cho
\[
  H^0(X_0,\Theta_{X_0/k})=H^2(X_0,\Theta_{X_0/k})=0.
\]
Khi đó tồn tại một lược đồ moduli hình thức cho $X_0$, tương ứng với một vành
địa phương chính quy $\mathcal O$ (nghĩa là một đại số chuỗi lũy thừa hình thức
trên $k$).

Như đã lưu ý, nói chung không đúng rằng lược đồ hình thức $\mathfrak X$ trên
$\mathcal O$ là đại số hóa được; tuy nhiên, ta biết điều đó đúng khi $X_0$ là
xạ ảnh và $H^2(X_0,\Osh_{X_0})=0$ (Định lý 4), chẳng hạn khi $X_0$ có chiều
1. Điều này đem lại một phần cơ sở cho hy vọng dựng được một lược đồ moduli
trên các số nguyên cho các đường cong có giống cho trước\ldots{}

Cũng cần lưu ý rằng các phương pháp như những phương pháp trình bày trong mục
này có thể áp dụng vào việc dựng và nghiên cứu các đa tạp Picard, cũng như vào
nhiều phép dựng khác. Ta sẽ sớm trở lại vấn đề này.

\fsection{8}{Ứng dụng vào nhóm cơ bản}

Các kỹ thuật đã trình bày cho phép bắt đầu nghiên cứu có hệ thống nhóm cơ bản
theo mô hình của lý thuyết tôpô. Hai định lý đầu tiên được phát biểu trong mục
này là những khái quát hóa các kết quả trong một công trình gần đây của
Lang--Serre.

Cho $X$ là một lược đồ. Một $X$-lược đồ $X'$ được gọi là một phủ không phân
nhánh của $X$ nếu
\begin{enumerate}
  \item[\textit{i.}] $X'$ hữu hạn trên $X$, nghĩa là nó được xác định bởi một
    bó coherent các đại số $\mathcal A=\mathcal A(X')$ trên $X$;
  \item[\textit{ii.}] $\mathcal A$ là một bó địa phương tự do trên $X$;
  \item[\textit{iii.}] Với mọi $x\in X$, thương
    \[
      \mathcal A_x/\mathfrak m_x\mathcal A_x
      =\mathcal A_x\otimes_{\Osh_x}\kappa(x)
    \]
    là một đại số khả phân trên $\kappa(x)$.
\end{enumerate}

\src{210}
Khái niệm phủ không phân nhánh này (do Serre và diễn giả đề xuất) có mọi tính
chất sơ cấp mà ta có thể kỳ vọng một cách hợp lý; ta sẽ không liệt kê chúng.
Chỉ xin nói rằng nó dẫn đến một lý thuyết Galois theo mô hình của lý thuyết
Galois cổ điển (và bao hàm lý thuyết ấy; các chứng minh thậm chí còn đơn giản
hơn khá nhiều so với những chứng minh thường được chấp nhận cho lý thuyết sau) và của lý thuyết
Galois về các phủ tôpô. Cụ thể hơn, ta gọi một \emph{điểm hình học} của một
lược đồ $X$ là một cấu xạ $a$ từ phổ $\xi$ của một trường đóng đại số $\Omega$
vào $X$, nghĩa là dữ liệu của một mở rộng đóng đại số của trường thặng dư
$\kappa(x)$ của một điểm $x=|a|$ thuộc $X$ (được gọi là \emph{vị trí} của
điểm hình học $a$). Nếu $X'$ là một phủ không phân nhánh của $X$, ta có thể
gắn với nó tập $E_a(X')$ các “điểm hình học của $X'$ nằm trên $a$”, nghĩa là
tập các cặp gồm một $x'\in X'$ nằm trên $x$ và một $\kappa(x)$-đồng cấu
$\kappa(x')\longrightarrow\Omega$. Như vậy (với $(X,a)$ cố định), ta thu được
một hàm tử $F(X,a)$ từ phạm trù $R(X)$ các phủ không phân nhánh $X'$ của $X$
đến phạm trù các tập hữu hạn. Nếu $X$ liên thông, cặp gồm $R(X)$ và $F(X,a)$
có các tính chất hình thức cần thiết để đẳng cấu với cặp tương tự được xác định
bởi một nhóm tôpô compact hoàn toàn không liên thông $\pi$ thích hợp (nghĩa là
một giới hạn xạ ảnh của các nhóm hữu hạn): ta lấy phạm trù $C(\pi)$ của các tập
hữu hạn $E$ trên đó $\pi$ tác động liên tục, cùng hàm tử quên
$F(\pi)(E)=E$ từ phạm trù này đến phạm trù các tập hữu hạn. Hơn nữa, nhóm
$\pi$ được xác định đến đẳng cấu chính tắc bởi điều kiện $(C(\pi),F(\pi))$
đẳng cấu với một cặp đã cho. Trong trường hợp này, $\pi$ được gọi là
\emph{nhóm cơ bản} của lược đồ liên thông $X$ tại điểm hình học $a$, và được
ký hiệu $\pi_1(X,a)$. Nếu $X$ không liên thông, ta thay nó bằng thành phần liên
thông chứa $x=|a|$. Ngược lại, nếu $X$ liên thông, thì các nhóm $\pi_1(X,a)$
ứng với hai điểm hình học $a',a''$ của $X$ là đẳng cấu (đẳng cấu ấy chỉ được
xác định đến các tự đẳng cấu trong), nên như thường lệ, ta có thể chọn $a$
theo cách thuận tiện nhất, chẳng hạn tại điểm tổng quát của $X$, với giả thiết
lược đồ này bất khả quy. Dĩ nhiên, $\pi_1(X,a)$ là một hàm tử hiệp biến theo lược
đồ có điểm gốc $(X,a)$. Khi đó, mọi phát biểu về việc phân loại các phủ bất
khả phân đều có thể chuyển sang ngôn ngữ lý thuyết nhóm theo từ điển quen thuộc
(ngoại trừ việc phải lưu ý rằng ở đây ta làm việc với các nhóm tôpô).

\src{211}
Mục đích của ta là phát triển một tương tự của dãy đồng luân chính xác của các
không gian phân thớ đối với một cấu xạ proper $f:X\longrightarrow Y$. Hiển
nhiên, vì chưa biết các nhóm đồng luân bậc cao phải là gì, ta chỉ có thể thu
được những kết quả tất yếu chưa đầy đủ. Để áp dụng các định lý cơ bản của
mục n\textsuperscript{o} 3, trước hết ta phải phát biểu tường minh vài bổ đề sơ cấp về các lược đồ
trên trường hoặc trên vành Artin (đúng theo phương pháp tổng quát!).

\result{BỔ ĐỀ}{1} Cho $(X',a')$ là một phủ không phân nhánh có điểm gốc, liên
kết với một biểu diễn có điểm gốc của $\pi_1(X,a)$ trên một tập hữu hạn $E$
(có điểm được đánh dấu $e$). Khi đó cấu xạ chính tắc
$\pi_1(X',a')\longrightarrow\pi_1(X,a)$ đồng nhất nhóm thứ nhất với nhóm ổn
định của $e$ trong $\pi_1(X,a)$ (và do đó là đơn ánh).

\result{BỔ ĐỀ}{2} Cho $X$ là một lược đồ đại số trên trường $k$, và cho $k'$
là một mở rộng thuần bất khả phân của $k$. Khi đó mọi phủ không phân nhánh của
$X\otimes_k k'$ đều thu được bằng ảnh ngược (tức là bằng mở rộng vô hướng) từ
một phủ không phân nhánh của $X$, được xác định duy nhất đến đẳng cấu chính
tắc.

Đặc biệt, từ hai bổ đề này suy ra rằng với mọi mở rộng đại số $K$ của $k$, và
mọi điểm hình học $a'$ của $X'=X\otimes_k K$ chiếu xuống điểm hình học $a$ của
$X$, đồng cấu có tính hàm tử
$\pi_1(X',a')\longrightarrow\pi_1(X,a)$ là đơn ánh.

\result{BỔ ĐỀ}{3} Cho $X$ là một lược đồ đầy đủ trên một vành Artin địa phương
$A$ sao cho $H^0(X,\Osh_X)=A$; cho $X'$ là một phủ không phân nhánh của $X$,
và đặt $A'=H^0(X',\Osh_{X'})$, vì thế là một vành hữu hạn trên $A$
(\textit{a priori} có thể phân nhánh trên $A$). Gọi $X_0,X'_0$ là các lược đồ
rút gọn liên kết với $X,X'$ (thu được bằng cách chia cho các bó phần tử lũy
linh của $\Osh_X,\Osh_{X'}$), và cho $k$ là một trường con của
$A/\mathfrak r(A)$ sao cho $A/\mathfrak r(A)$ hữu hạn trên trường con ấy (như vậy
$X_0$ là một lược đồ đại số đầy đủ trên $k$, còn $X'_0$ là một phủ không phân
nhánh của nó). Cuối cùng, cho $\Omega$ là một mở rộng đóng đại số của $k$, và
xét phủ không phân nhánh $X'_0\otimes_k\Omega$ của
$X_0\otimes_k\Omega$. Khi đó hai điều kiện sau tương đương:
\begin{enumerate}
  \item[\textit{i.}] $X'_0\otimes_k\Omega$ hoàn toàn phân rã trên
    $X_0\otimes_k\Omega$.
  \item[\textit{ii.}] Cấu xạ tự nhiên
    $X'\longrightarrow X\otimes_A A'$ là một đẳng cấu.
\end{enumerate}
Dưới các điều kiện ấy, $A'$ là một mở rộng không phân nhánh của $A$. Cuối
cùng, nếu $X'$ liên thông, điều kiện \textit{(i)} tương đương với điều kiện
sau đây, bề ngoài yếu hơn:
\begin{enumerate}
  \item[\textit{i bis.}] $X'_0\otimes_k\Omega$ có một tiết diện chính quy
    trên $X_0\otimes_k\Omega$.
\end{enumerate}
Khi điều kiện \textit{(ii)} được thỏa mãn, ta nói phủ không phân nhánh $X'$
của $X$ là \emph{tầm thường về mặt hình học}.

\src{212}
\result{BỔ ĐỀ}{4} Cho $f:X\longrightarrow Y$ là một cấu xạ proper sao cho
$f_*(\Osh_X)=\Osh_Y$. Cho $a$ là một điểm hình học của $X$ và $b$ là ảnh của
nó trên $Y$. Khi đó $\pi_1(X,a)\longrightarrow\pi_1(Y,b)$ là toàn ánh.

Thật vậy, cần chứng minh điều sau: nếu một phủ không phân nhánh $Y'$ của $Y$
(tương ứng với một bó đại số địa phương tự do $\mathcal A$) có tính chất
$X\otimes_Y Y'$ không liên thông, thì $Y'$ cũng không liên thông. Quả vậy,
$f^*(\mathcal A)$ khi ấy là tổng trực tiếp của hai bó vành khác không, nên ảnh
trực tiếp của nó cũng vậy; nhưng ảnh trực tiếp ấy chính là
$\mathcal A\otimes f_*(\Osh_X)=\mathcal A$.

\result{BỔ ĐỀ}{5} Cho $X$ là một lược đồ đầy đủ trên trường $k$; giả sử
$H^0(X,\Osh_X)$ là một vành địa phương $A$ và $A/\mathfrak r(A)$ là thuần bất
khả phân trên $k$. Cho $\Omega$ là một bao đóng đại số của $k$, và đặt
$\overline X=X\otimes_k\Omega$ (lược đồ này liên thông). Chọn một điểm hình
học $\overline a$ của $\overline X$ chiếu xuống điểm hình học $a$ của $X$;
khi đó có một dãy chính xác
\[
  e\longrightarrow\pi_1(\overline X,\overline a)
  \longrightarrow\pi_1(X,a)
  \longrightarrow\pi_1(k,b)\longrightarrow e
\]
(trong đó $\pi_1(k,b)$ là nhóm Galois của $\Omega$ trên $k$).

Tính đơn ánh của đồng cấu thứ nhất đã được chứng minh trong các Bổ đề 1 và 2;
tính chính xác ở giữa suy ra từ Bổ đề 3; cuối cùng, tính toàn ánh của đồng cấu
sau cùng (đây là chỗ duy nhất dùng việc $A/\mathfrak r(A)$ thuần bất khả phân)
suy ra từ Bổ đề 4.

\result{MỆNH ĐỀ}{4} Cho $f:X\longrightarrow Y$ là một cấu xạ proper và phẳng
sao cho, với mọi $y\in Y$, $H^0(f^{-1}(y),\Osh_{f^{-1}(y)})$ là một đại số khả
phân trên trường thặng dư $\kappa(y)$ (chẳng hạn, điều này đúng nếu
$f^{-1}(y)$ là một lược đồ khả phân trên $\kappa(y)$, tức là rút gọn và các
trường tương ứng với các thành phần bất khả quy của nó là những mở rộng khả
phân của $\kappa(y)$). Khi đó phủ $Y'$ của $Y$ liên kết với
$f_*(\Osh_X)$ là không phân nhánh.

Chứng minh suy ra dễ dàng từ Định lý 2.

\src{213}
Mệnh đề này, cùng với Bổ đề 1, trên thực tế quy việc nghiên cứu đồng luân của
các cấu xạ proper và phẳng (có thớ khả phân) về trường hợp
$f_*(\Osh_X)=\Osh_Y$ (vì khi dùng phân tích Stein, ta thay $Y$ bởi $Y'$).

\result{NHẬN XÉT}{} Một cấu xạ phẳng kiểu hữu hạn có các thớ là những lược đồ
khả phân (tương ứng, trơn) được gọi là \emph{khả phân} (tương ứng,
\emph{trơn}). Ta chứng minh được rằng nếu $f$ phẳng và $f^{-1}(y)$ khả phân
(tương ứng, trơn), thì tồn tại một lân cận của $f^{-1}(y)$ trên đó $f$ khả
phân (tương ứng, trơn). Kết quả tương tự đúng cho “chuẩn tắc tuyệt đối”
(“định lý Bertini”).

Cho $f:X\longrightarrow Y$ là một cấu xạ proper sao cho
\begin{equation}
  \tag{i}
  f_*(\Osh_X)=\Osh_Y,
\end{equation}
và cho $X'$ là một lược đồ hữu hạn trên $X$. Gọi $Y'$ là phủ của $Y$ tương ứng
với phân tích Stein của $X'\longrightarrow Y$ (xem Định lý 5). Cho $y\in Y$;
khi đó tập các thành phần liên thông của thớ $F'$ của $X'$ trên $y$ được đồng
nhất với tập các điểm $y'$ của $Y'$ nằm trên $y$ (Định lý 5). Xét cấu xạ hiển
nhiên
\begin{equation}
  \tag{*}
  X'\longrightarrow X\times_Y Y'
\end{equation}
suy ra từ các cấu xạ tự nhiên $X'\longrightarrow X$ và
$X\longrightarrow Y'$; cấu xạ này là một
đẳng cấu mỗi khi $X'$ có dạng $X\times_Y Y''$, trong đó $Y''$ là một phủ không
phân nhánh của $Y$; khi ấy $Y'$ chính là $Y''$ và (*) là đồng nhất. Ta muốn
đưa ra chính xác các điều kiện để $X'$ có dạng vừa nêu, tức là để $Y'$ không
phân nhánh và (*) là một đẳng cấu. Vì thế, xét thớ $F$ của $X$ trên $y$; đó là
một lược đồ proper trên $\kappa(y)$, còn $F'$ là một phủ của nó (không phân
nhánh nếu $X'$ không phân nhánh). Cho $F'_1$ là một thành phần liên thông của
$F'$ tương ứng với một điểm $y'_1$ của $Y'$ nằm trên $y$. Giả sử thêm rằng
\begin{enumerate}
  \item[\textit{(ii)}] $X'$ không phân nhánh trên $X$ tại các điểm của $F'_1$
    (do đó $F'_1$ là một phủ không phân nhánh của $F$), và
  \item[\textit{(iii)}] $F'_1$ là một phủ tầm thường về mặt hình học của $F$
    (xem Bổ đề 3).
\end{enumerate}

\result{ĐỊNH LÝ}{11} Dưới các điều kiện này, tồn tại một lân cận mở $U'$ của
$y'_1$ trong $Y'$ sao cho (*) là một đẳng cấu trên $U'$. Hơn nữa, $Y'$ không
phân nhánh tại $y'_1$ trên $Y$ (nhưng có thể phân nhánh tại các điểm $y'$ khác
của $Y'$ nằm trên $y$).

Dĩ nhiên, các điều kiện \textit{(ii)} và \textit{(iii)} cũng là cần thiết để
kết luận của định lý đúng. Chứng minh của định lý là dễ dàng, nhờ Bổ đề 3 và
Định lý 2.


\src{214}
\result{HỆ QUẢ}{1} Tiếp tục giả sử \textit{(i)} được thỏa mãn. Để một phủ
không phân nhánh của $X$ đẳng cấu với ảnh ngược của một phủ không phân
nhánh $Y'$ của $Y$, điều kiện cần và đủ là $X'$ cảm sinh trên mỗi thớ
$f^{-1}(y)$ một phủ tầm thường về mặt hình học.

Theo Định lý 11, tập các điểm của $Y$ tại đó điều kiện này được thỏa mãn là
mở, nên chỉ cần kiểm tra nó tại các điểm đóng $y$\ldots{} Ta ghi lại phát biểu
sau, tương đương với Hệ quả 1: hạt nhân của đồng cấu (toàn ánh theo Bổ đề 4)
$\pi_1(X)\longrightarrow\pi_1(Y)$ là nhóm con chuẩn tắc đóng sinh bởi các ảnh
trong $\pi_1(X)$ của các $\pi_1(\overline{f^{-1}(y)})$, trong đó
\[
  \overline{f^{-1}(y)}
  =f^{-1}(y)\otimes_{\kappa(y)}\overline{\kappa(y)},
\]
và $\overline{\kappa(y)}$ ký hiệu một bao đóng đại số của $\kappa(y)$. Cần lưu
ý rằng vì không thể chọn cùng một điểm gốc cho tất cả các thớ, các đồng cấu
$\pi_1(\overline{f^{-1}(y)})\longrightarrow\pi_1(X)$ trong mọi trường hợp chỉ
được xác định (sau khi đã chọn một điểm gốc cho $X$, và do đó cho $Y$) đến hợp
thành với một tự đẳng cấu trong của $\pi_1(X)$.

\result{HỆ QUẢ}{2} Dưới các giả thiết tổng quát của Định lý 11, giả sử thêm
rằng $Y,X,X'$ nguyên, và gọi $K,L,L'$ là các trường hàm tương ứng. Khi đó tồn
tại một mở rộng trung gian khả phân $K'$ của $K$ trong $L'$, độc lập tuyến
tính với $L$, sao cho $L'=LK'$ (do đó $L'=L\otimes_K K'$).

(Ta áp dụng phần cuối của Bổ đề 3 cho thớ hình học của $X$.) Trường hợp áp dụng
đáng quan tâm nhất của Định lý 11 là khi $f$ là một cấu xạ khả phân. Khi ấy
$X'$ cũng khả phân trên $Y$; do đó, theo Mệnh đề 4, $Y'$ không phân nhánh trên
$Y$, và vế thứ hai $X\times_Y Y'$ trong (*) không phân nhánh trên $X$. Ta dễ
dàng suy ra:

\result{HỆ QUẢ}{3} Ngoài \textit{(i)}, giả sử $f$ khả phân. Cho $X'$ là một
phủ không phân nhánh liên thông của $X$; để
$X$ là ảnh ngược của một phủ không phân nhánh $Y'$ của $Y$,
điều kiện cần và đủ là phủ $\overline F'$ cảm sinh trên một thớ hình học
$\overline F=f^{-1}(y)$ có một tiết diện chính quy.

Cần lưu ý rằng không phải giả sử $\overline F'$ tầm thường về mặt hình học
trên $\overline F$ (điều này sẽ đúng \textit{a posteriori}, mặc dù
\textit{a priori} điều kiện ấy mạnh hơn nhiều). Thật ra, Hệ quả 3 tương đương
với phát biểu sau:

\src{215}
\result{HỆ QUẢ}{4} Cho $f:X\longrightarrow Y$ là một cấu xạ proper và khả
phân sao cho $f_*(\Osh_X)=\Osh_Y$; cho $\overline F$ là thớ hình học của một
điểm $y\in Y$, và chọn một điểm hình học trong $\overline F$. Qua các cấu xạ
$\overline F\longrightarrow X\longrightarrow Y$, ta thu được các điểm hình
học trong $X,Y$, được lấy làm điểm gốc cho các nhóm cơ bản của
$\overline F,X,Y$. Dưới các điều kiện ấy, ta có dãy chính xác
\[
  \boxed{
    \pi_1(\overline F)\longrightarrow\pi_1(X)
    \longrightarrow\pi_1(Y)\longrightarrow e
  }.
\]

Từ đó dễ dàng suy ra hai phát biểu sau của Serre--Lang, ở đây không còn cần
bất kỳ giả thiết chuẩn tắc nào:

\result{HỆ QUẢ}{5} Cho $X,Y$ là các lược đồ liên thông trên một trường $k$;
giả sử lược đồ rút gọn $X_{\mathrm{réd}}$ khả phân trên $k$ (điều này tự động
đúng nếu $k$ hoàn hảo). Chọn một
điểm hình học $a$ (tương ứng, $b$) trong $X$ (tương ứng, $Y$), từ đó thu được
một điểm hình học $c=(a,b)$ trong $X\times_kY$ và một cấu xạ tự nhiên
\[
  \pi_1(X\times_kY,c)\longrightarrow
  \pi_1(X,a)\times\pi_1(Y,b)
\]
(suy ra từ các cấu xạ có tính hàm tử từ $\pi_1(X\times Y,c)$ vào
$\pi_1(X,a)$ và $\pi_1(Y,b)$). Cấu xạ này là đơn ánh, và thậm chí là song ánh
nếu $k$ đóng đại số.

(Tính toàn ánh trong trường hợp sau gần như hiển nhiên.) Cùng với Serre--Lang,
ta suy ra:

\result{HỆ QUẢ}{6} Cho $X$ là một lược đồ đại số liên thông trên một trường
đóng đại số $k$, và cho $K$ là một mở rộng đóng đại số của $k$;
khi đó các nhóm cơ bản của $X$ và $X\otimes_kK$ là như nhau, tức là mọi phủ
không phân nhánh của lược đồ sau đều thu được bằng mở rộng vô hướng từ một phủ
không phân nhánh của $X$, duy nhất đến đẳng cấu.

\result{NHẬN XÉT}{}

1) Dùng Mệnh đề 4, ta thấy giả thiết $f_*(\Osh_X)=\Osh_Y$ trong Hệ quả 4 không
cốt yếu. Trong trường hợp tổng quát, thay vì đặt nhóm đơn vị $e$ sau
$\pi_1(Y)$, ta phải tiếp tục bằng
\[
  \pi_0(\overline F)\longrightarrow\pi_0(X)
  \longrightarrow\pi_0(Y)\longrightarrow e.
\]
như trong tôpô đại số.

\src{216}
2) Nói chung, hiện nay chưa thể nói gì về hạt nhân của
$\pi_1(\overline F)\longrightarrow\pi_1(X)$, nơi lẽ ra phải có sự tham gia của
một $\pi_2(Y)$. Tuy nhiên, có vẻ ta phải chứng minh được rằng
$\pi_1(\overline F)\longrightarrow\pi_1(X)$ là đơn ánh nếu $Y$ là phổ của một
vành địa phương $A$, bằng cách dựa vào Định lý 12 dưới đây (định lý này khẳng
định điều đó khi $A$ đầy đủ).

Định lý 11 chỉ dùng các Định lý 1 và 2. Bây giờ ta sẽ áp dụng Định lý 3, dựa
trên bổ đề sơ cấp sau:

\result{BỔ ĐỀ}{6} Cho $X$ là một lược đồ, và $X_0$ là lược đồ rút gọn tương
ứng (tức là lược đồ trong đó các phần tử lũy linh đã bị triệt tiêu). Khi đó mọi
phủ không phân nhánh $X'_0$ của $X_0$ đều được cảm sinh bởi một phủ không phân
nhánh $X'$ của $X$, được xác định duy nhất đến đẳng cấu chính tắc.

Bổ đề này, có bản chất hoàn toàn địa phương, đóng vai trò ở đây tương tự như
Định lý 8 trong lý thuyết môđun. Kết hợp nó với định lý tồn tại (Định lý 3),
ta suy ra:

\result{ĐỊNH LÝ}{12} Cho $A$ là một vành Noether địa phương đầy đủ có trường
thặng dư $k$. Cho $X$ là một lược đồ proper trên $A$. Khi đó mọi phủ không
phân nhánh $X'_0$ của $X_0=X\otimes_A k$ đều được cảm sinh bởi một phủ không
phân nhánh $X'$ của $X$, duy nhất đến đẳng cấu.

Nói cách khác:

\result{HỆ QUẢ}{1} Chọn một điểm hình học trong $X_0$ làm điểm gốc cho các
nhóm cơ bản của $X_0$ và $X$. Khi đó đồng cấu chính tắc
$\pi_1(X_0)\longrightarrow\pi_1(X)$ là một đẳng cấu.

Bây giờ áp dụng Bổ đề 5 cho $X_0$ (để đơn giản, giả sử
$H^0(X_0,\Osh_{X_0})=k$), và lưu ý rằng vì $A$ đầy đủ, các mở rộng không phân
nhánh của nó tương ứng với các mở rộng không phân nhánh của trường thặng dư;
tức là $\pi_1(Y)=\pi_1(k)$ (trong đó $Y=\operatorname{Spec}(A)$). Ta thu được
dãy chính xác:
\[
  e\longrightarrow\pi_1(\overline X_0)
  \longrightarrow\pi_1(X)\longrightarrow\pi_1(Y)\longrightarrow e.
\]

\result{HỆ QUẢ}{2} Cho $f:X\longrightarrow Y$ là một cấu xạ proper và phẳng,
cho $y_1$ là một điểm của $Y$ và $y_0$ là một chuyên hóa của $y_1$, và xét các
thớ “hình học” tương ứng $\overline X_1$ và $\overline X_0$; giả sử
$\overline X_0$ khả phân và liên thông (điều này kéo theo
$\overline X_1$ thỏa các điều kiện tương tự). Khi đó có thể tìm được một đồng
cấu nhóm
$\pi_1(\overline X_1)\longrightarrow\pi_1(\overline X_0)$, được xác định đến
một tự đẳng cấu trong, và đồng cấu này là toàn ánh.


\src{217}
Có thể hy vọng đồng cấu này luôn song ánh; tiếc thay, nói chung điều đó không
đúng nếu $\kappa(y_0)$ có đặc số $>0$. Tuy nhiên, dưới đây ta sẽ thu được một
chặn cho hạt nhân của đồng cấu ấy (ít nhất khi $\overline X_0$ trơn), từ đó suy
ra rằng nếu $\kappa(y_0)$ có đặc số $0$ thì đồng cấu trên là song ánh (một kết
quả cũng có thể chứng minh bằng phương pháp siêu việt). Dù sao, ta đã có một
chặn cho $\pi_1$ của một thớ đặc biệt theo nhóm tương ứng của một thớ tổng quát.
Chẳng hạn, dùng sự kiện một đường cong đại số trong đặc số $p$ có thể được nâng
lên một đường cong trong đặc số $0$ (Định lý 9, Hệ quả 4), ta thu được bằng
phương pháp siêu việt:

\result{HỆ QUẢ}{3} Cho $X_0$ là lược đồ của một đường cong đầy đủ và trơn trên
một trường đóng đại số có đặc số tùy ý, và gọi $g$ là giống của $X_0$; khi đó
$\pi_1(X_0)$ có $2g$ phần tử sinh tôpô, chịu ràng buộc bởi hệ thức quen thuộc.

Từ đó, bằng kỹ thuật tiết diện siêu phẳng quen thuộc, suy ra:

\result{HỆ QUẢ}{4} Cho $X$ là một lược đồ xạ ảnh trơn trên một trường đóng đại
số có đặc số tùy ý; khi đó $\pi_1(X)$ có hữu hạn phần tử sinh tôpô.

Ta tìm hạt nhân của đồng cấu
$\pi_1(\overline X_1)\longrightarrow\pi_1(\overline X_0)$. Để làm việc này,
có thể giả sử $Y$ là phổ của một vành định giá rời rạc đầy đủ
$V=\Osh_y$ ($y=y_0$). Câu hỏi tương đương với câu hỏi sau: cho một phủ không
phân nhánh $\overline X'_1$ của $\overline X_1$ (nếu muốn, có thể giả sử nó là
Galois), xác định những điều kiện để nó đến từ một phủ không phân nhánh của
$X_0$. \textit{A priori}, phủ đã cho đến từ, bằng mở rộng vô hướng, một phủ
không phân nhánh $X'_1$ của $X_1\otimes_K K'$, trong đó $K'$ là một mở rộng
con hữu hạn của bao đóng đại số $\overline K$ của trường phân thức $K$ của
$V$; nếu $\overline X'_1$ là Galois với nhóm $G$, ta có thể chọn $X'_1$ là
Galois với nhóm $G$. Khi đó ta có: để
$X'_1\otimes_{K'}\overline K=\overline X'_1$ đến từ một phủ không phân nhánh
của $X_0$, điều kiện cần và đủ là có thể tìm được một mở rộng hữu hạn $K''$
của $K'$ trong $\overline K$ sao cho
$X''_1=X'_1\otimes_{K'}K''$ có dạng $X''\otimes_{V''}K''$, trong đó $V''$ là
bao đóng nguyên của $V$ trong $K''$, còn $X''$ là một phủ không phân nhánh của
$X\otimes_VV''$. Chẳng hạn, giả sử $X_0$ chuẩn tắc tuyệt đối; khi đó
$X\otimes_VV''$ chuẩn tắc (vì nó phẳng trên $V''$ và có thớ đặc biệt chuẩn
tắc), và trường hàm của nó đồng nhất với trường $K''(X_1)$ của
\[
  X_1\otimes_KK''
  =(X\otimes_VK)\otimes_KK''
  =X\otimes_VK''
  =(X\otimes_VV'')\otimes_{V''}K''.
\]
Cho $L''=K''(X'_1)$ là trường hàm của
$X'_1\otimes_{K'}K''$; đây là một mở rộng hữu hạn khả phân của $K''(X_1)$, và
điều kiện trên cũng có nghĩa là $L''$ là một mở rộng không phân nhánh của
trường hàm của $X\otimes_VV''$ (tức là phép chuẩn hóa của
$X\otimes_VV''$ trong $L''$ không phân nhánh trên $X\otimes_VV''$). Hơn nữa,
chỉ cần kiểm tra rằng $L''$ không phân nhánh tại các điểm của thớ đặc biệt của
$X\otimes_VV''$ (vì nó không phân nhánh trên thớ tổng quát
$X_1\otimes_KK''$). Nếu bây giờ $X_0$ trơn, “định lý thuần khiết” của
Nagata--Zariski thậm chí cho thấy chỉ cần kiểm tra rằng $L''$ không phân nhánh
trên vành địa phương $O''$ tại điểm tổng quát của thớ đặc biệt của
$X\otimes_VV''$; đây là một vành định giá rời rạc, là phép chuẩn hóa trong
$K''(X_1)$ của vành địa phương $O\subset K(X_1)$ tại điểm tổng quát của thớ
đặc biệt của $X$. Do đó ta đã thu được:

\src{218}
\result{HỆ QUẢ}{5} Dưới các giả thiết và với các ký hiệu trên, để phủ không
phân nhánh $X'_1\otimes_{K'}\overline K$ của
$\overline X_1=X_1\otimes_K\overline K$ đến từ một phủ không phân nhánh của
$\overline X_0$, điều kiện cần và đủ là tồn tại một mở rộng con hữu hạn $K''$
của $\overline K/K'$ sao cho $K''(X'_1)$ không phân nhánh trên vành định giá
rời rạc $O''\subset K''(X_1)$.

Bây giờ lưu ý rằng $O''$ là phép chuẩn hóa trong $K''(X_1)$ của vành định giá
rời rạc $O'\subset K'(X_1)$ (phép chuẩn hóa của $O$ trong $K'(X_1)$), và
$O'$ chứa bao đóng nguyên $V'$ của $V$ trong $K'$; một phần tử đều hóa $u$ của
$V'$ cũng là một phần tử đều hóa của $O'$. Giả sử bây giờ $X'_1$ là Galois với
nhóm Galois $G$, có cấp $n$ nguyên tố cùng nhau với đặc số $p$ của
$\kappa(y_0)$ (đồng thời là đặc số của trường thặng dư của $O'$). Khi đó
$K'(X'_1)$ phân nhánh thuần trên $O'$, và dễ dàng suy ra (“bổ đề Abhyankar”)
rằng sau khi ghép thêm một căn bậc $n$, ký hiệu $v$, của một phần tử đều hóa
của $O'$, mở rộng ấy trở thành không phân nhánh trên phép chuẩn hóa của $O'$
trong $K'(X_1)(v)$. Nhưng có thể lấy $v$ là một căn bậc $n$ của một phần tử
đều hóa của $V'$, điều này chứng minh rằng điều kiện của Hệ quả 5 được thỏa
mãn. (Khả năng dùng bổ đề Abhyankar và định lý thuần khiết là do Serre bán cho
tôi.) Để diễn đạt kết quả thu được, với mọi nhóm compact hoàn toàn không liên
thông $\pi$, ta đưa vào nhóm thương $\widetilde\pi$ của $\pi$ theo nhóm con
đóng sinh bởi các nhóm con Sylow $p$ của nó, tức là giới hạn xạ ảnh của các
nhóm thương rời rạc của $\pi$ có cấp nguyên tố cùng nhau với $p$. Với ký hiệu
này, ta thu được:

\result{ĐỊNH LÝ}{13} Cho $f:X\longrightarrow Y$ là một cấu xạ proper và
phẳng, cho $y_1$ là một điểm của $Y$ và $y_0$ là một chuyên hóa của $y_1$, và
giả sử $\overline X_0$ liên thông và trơn. Khi đó đồng cấu
$\widetilde\pi_1(\overline X_1)\longrightarrow
\widetilde\pi_1(\overline X_0)$ suy ra từ đồng cấu toàn ánh trong Định lý 12,
Hệ quả 2, là một đẳng cấu.

Nói cách khác:

\src{219}
\result{HỆ QUẢ}{1} Việc phân loại các phủ Galois không phân nhánh có nhóm
Galois với cấp nguyên tố cùng nhau với đặc số $p$ của $\kappa(y_0)$ là như
nhau đối với $\overline X_0$ và $\overline X_1$.

Đặc biệt, nếu $\kappa(y_0)$ có đặc số không, bằng phương pháp đại số ta thu
được rằng $\pi_1(\overline X_1)\longrightarrow\pi_1(\overline X_0)$ là song
ánh.

Cuối cùng, ta lưu ý rằng các kỹ thuật đã dùng còn cho kết quả sau, tổng quát
hơn Định lý 13:

\result{ĐỊNH LÝ}{14} Cho $f:X\longrightarrow Y$ là một cấu xạ proper và trơn,
cho $D$ là một lược đồ con đóng của $X$, trơn trên $Y$ và có đối chiều 1 tại
mọi điểm của nó. Với một thớ $Z=f^{-1}(z)$ của $f$, đặt $Z'=Z-Z\cap D$ và gọi
$\pi_1^t(Z')$ là thương của nhóm cơ bản $\pi_1(Z')$ phân loại các phủ không
phân nhánh của $Z'$ phân nhánh thuần trên $\overline{Z\cap D}$. Cho $y_0,y_1$
như trong Định lý 13. Khi đó có một đồng cấu toàn ánh (được xác định đến một tự
đẳng cấu trong)
\[
  \pi_1^t(\overline X'_1)\longrightarrow\pi_1^t(\overline X'_0),
\]
và đồng cấu tương ứng
\[
  \widetilde\pi_1^t(\overline X'_1)
  \longrightarrow\widetilde\pi_1^t(\overline X'_0)
\]
là một đẳng cấu.

Từ đó suy ra các biến thể tương ứng của các hệ quả của Định lý 13 và của Hệ
quả 4 của Định lý 12. Tương tự, dùng Định lý 9, Hệ quả 3, bằng phương pháp
siêu việt ta thu được:

\result{HỆ QUẢ}{} Cho $X_0$ là lược đồ của một đường cong đầy đủ và trơn trên
một trường đóng đại số có đặc số tùy ý, và cho
$S=(s_i)_{1\le i\le n}$ là một tập con hữu hạn gồm $n$ phần tử của $X_0$. Khi
đó $\pi_1^t(X_0-S)$ có $2g+n$ phần tử sinh tôpô
$x_i,y_i$ ($1\le i\le g$) và $\sigma_j$ ($1\le j\le n$), chịu ràng buộc bởi
hệ thức
\[
  \left(\prod_i x_i y_i x_i^{-1}y_i^{-1}\right)
  \sigma_1\cdots\sigma_n=1,
\]
trong đó các $\sigma_j$ là những phần tử sinh của các nhóm quán tính tương ứng
với các $s_j$. Với mọi nhóm hữu hạn $G$ có cấp nguyên tố cùng nhau với đặc số,
được sinh bởi các phần tử $x_i,y_i,\sigma_j$ thỏa hệ thức trên, tồn tại một
phủ Galois không phân nhánh của $X_0-S$, với nhóm $G$, sao cho các nhóm quán
tính tại các điểm $s_j$ được sinh bởi các $\sigma_j$.

Khi $X_0$ có giống 0 và $n=3$, ta có một lời giải của “bài toán ba điểm”, ít
nhất đối với các phủ Galois có cấp nguyên tố cùng nhau với đặc số. (Ở đây,
Định lý 9 thật ra không cần dùng, và hơn nữa dường như hệ quả trước có thể được
suy ra từ trường hợp riêng xét trong bài toán ba điểm.)

\src{220}
\result{NHẬN XÉT}{}

1) Một nghiên cứu đầy đủ hơn, chắc chắn phải dùng các phủ Galois tổng quát hóa
của $X,X_0,X_1$ (với nhóm Galois có thể là vô cùng bé), hẳn sẽ cho phép thu
lại hạt nhân trong Định lý 12, Hệ quả 2. Trái lại, việc nghiên cứu các phủ cho
phép phân nhánh không “thuần” dường như phải khó hơn nhiều.

2) Bổ đề 6, kết hợp với một kết quả của Grauert về phép đầy đủ hóa hình thức
của một lược đồ xạ ảnh trơn dọc theo một tiết diện siêu phẳng (hoặc với định
lý, hiện chưa được chứng minh, nêu trong Nhận xét 2 sau Định lý 11), cũng cho
phép chứng minh trong hình học đại số “trừu tượng” định lý Lefschetz cổ điển về
nhóm cơ bản.

\begin{center}
  \textbf{TÀI LIỆU THAM KHẢO}
\end{center}

\begin{description}
  \item[{[1]}] Dieudonné (J.) và Grothendieck (A.). --- \emph{Éléments de
    géométrie algébrique}, Publications mathématiques de l'Institut des Hautes
    Études Scientifiques (sắp xuất bản).
  \item[{[2]}] Grothendieck (Alexander). --- \emph{The cohomology theory of
    abstract algebraic varieties}. International Congress of Mathematicians
    [1958. Edinburgh] (sắp xuất bản).
  \item[{[3]}] Serre (Jean-Pierre). --- \emph{Faisceaux algébriques
    cohérents}, Annals of Math., Series 2, t. 61, 1955, p. 197--278.
  \item[{[4]}] Serre (Jean-Pierre). --- \emph{Géométrie algébrique et
    géométrie analytique}, Ann. Institut Fourier Grenoble, t. 6, 1955--1956,
    p. 1--42.
\end{description}

\par

\clearpage
\chapter*{Đính chính cho Bài trình bày số 182}
\addcontentsline{toc}{chapter}{Đính chính cho Bài trình bày số 182}
\setcounter{footnote}{1}
\sid{e182}
\markboth{Đính chính 182}{Đính chính 182}
\src{298}

\noindent\textbf{2. Hình học hình thức và hình học đại số.}

\begin{center}
[Hội thảo Bourbaki, t. 11, 1958/59, số 182, 28 tr.]
\end{center}

\noindent\textbf{Trang 8, Định lý 6.} --- Giả thiết «$f$ tựa xạ ảnh» có thể
được thay bằng giả thiết yếu hơn «$f$ tách», nhờ kết quả sau đây (xem SGA VIII,
6.2): mọi cấu xạ $f:X\longrightarrow Y$ tựa hữu hạn và tách đều là tựa xạ ảnh.

\src{299}
\noindent\textbf{Trang 12, và trang 14, Nhận xét 1.} --- Xin lưu ý rằng J.-P.
Serre đã xây dựng một đa tạp xạ ảnh không kỳ dị, chiều 3, trên một trường đóng
đại số $k$ có đặc số $p>0$, không thu được bằng phép rút gọn từ một lược đồ
proper trên một vành địa phương nguyên có trường thặng dư $k$ và trường phân
thức có đặc số không.\footnote{Serre (Jean-Pierre). --- Exemples de variétés
projectives en caractéristique $p$ non relevables en caractéristique zéro,
\emph{Proc. Nat. Acad. Sc. U.S.A.}, t. 47, 1961, tr. 108--109.} Mumford được
cho là đã tìm được một ví dụ tương tự với một mặt xạ ảnh không kỳ dị.

\noindent\textbf{Trang 15, các Nhận xét 2 và 3.} --- Hiện nay tôi ít lạc quan
hơn về những kết quả được phỏng đoán ở đây. Tuy nhiên, bài toán về các biến
thiên cấu trúc của không gian xạ ảnh, được nêu ở cuối Nhận xét 3, đã được
Hironaka giải quyết theo hướng khẳng định, và bài toán tương tự đối với các đa
tạp abel đã được Koizumi giải quyết.

\noindent\textbf{Trang 16, dòng 3.} --- Xin lưu ý rằng Mumford gần đây đã xây
dựng các lược đồ môđuli cho các đường cong giống $g$ (xem SHMI). Ngoài ra,
Định lý 10 chứng minh rằng các «lược đồ Jacobi mức $n$» trong lý thuyết môđuli
là không kỳ dị (và thậm chí trơn trên $\mathbf Z$).

\noindent\textbf{Trang 23, dòng 12, dòng đầu của Hệ quả 5.} --- Bổ sung:
«$X$ hoặc $Y$ proper trên $k$».

Nội dung cốt lõi của các §§ 1 đến 5 nằm trong phần EGA III đã xuất bản, còn
nội dung của các §§ 6 và 7 nằm trong SGA III. Để nghiên cứu nhóm cơ bản, xem
SGA V, IX, X, XI, cũng như SGA 1962 (các Bài trình bày X, XII và XIII), về các
định lý kiểu Lefschetz và nhiều câu hỏi mở. Chỉ riêng lý thuyết về các phủ phân
nhánh thuần (xem Định lý 14) là chưa được trình bày thành một bản hoàn chỉnh.
Hệ quả của Định lý 14, cho phép xác định hoàn toàn các phủ Galois có cấp nguyên
tố cùng nhau với đặc số của một đường cong đại số trên một trường đóng đại số,
đã được sử dụng một cách thiết yếu trong ba trường hợp:

1\textsuperscript{o} trong chứng minh của Igusa về bất đẳng thức Picard cho
các mặt xạ ảnh không kỳ dị trong đặc số tùy ý;

2\textsuperscript{o} trong việc nghiên cứu (do Ogg và Šafarevič phát triển độc
lập) nhóm các phân thớ chính thuần nhất dưới một đa tạp abel được xác định trên
một trường hàm một biến, trong đặc số tùy ý;

3\textsuperscript{o} trong chứng minh gần đây của Artin về một số định lý then
chốt của «đối đồng điều Weil» của các đa tạp đại số.

\par

\unit{Bài trình bày 190 — Kỹ thuật hạ xuống I}
\sid{190}
\markboth{Thuyết trình 190}{Thuyết trình 190}
\src{299}
\seminar{Năm thứ 12, 1959/60, số 190\\Tháng 12 năm 1959}
\paperhead{KỸ THUẬT HẠ XUỐNG VÀ CÁC ĐỊNH LÝ TỒN TẠI TRONG HÌNH HỌC ĐẠI SỐ\\
I. ĐẠI CƯƠNG. HẠ XUỐNG THEO CÁC CẤU XẠ PHẲNG TRUNG THÀNH}%
  {do Alexander GROTHENDIECK}

Xét về mặt kỹ thuật, bài trình bày này và các bài tiếp theo có thể được xem là
những biến thể của ``định lý 90'' nổi tiếng của Hilbert. Có vẻ như A.~Weil là
người đưa phương pháp hạ xuống vào hình học đại số, dưới tên gọi ``hạ xuống
trường cơ sở''. Tuy nhiên, Weil chỉ xét trường hợp các mở rộng trường hữu hạn
tách được. Sau đó P.~Cartier xử lý trường hợp các mở rộng thuần bất khả phân có
chiều cao 1. Vì chưa có ngôn ngữ lược đồ, và cụ thể hơn là vì chưa cho phép xét
các phần tử lũy linh trong các vành, Cartier đã không thể phát biểu sự đồng nhất
cốt yếu giữa hai trường hợp này.

Hiện nay, có vẻ như kỹ thuật hạ xuống tổng quát sẽ được trình bày (kết hợp, khi
thích hợp, với các định lý cơ bản của ``hình học hình thức'', xem [3]) là nền
tảng của phần lớn các định lý tồn tại trong hình học đại số. Cũng cần lưu ý rằng
kỹ thuật ấy chắc chắn có thể được chuyển sang ``hình học giải tích'', và có thể
hy vọng rằng, trong tương lai gần, các chuyên gia sẽ chứng minh được những định
lý tương tự theo nghĩa ``giải tích'' của các định lý tồn tại trong hình học
hình thức được nêu ở Bài trình bày II. Dù sao, các công trình gần đây của Kodaira--Spencer,
với những phương pháp dường như không thích hợp để định nghĩa và nghiên cứu các
``đa tạp môđuli'' ở lân cận các điểm kỳ dị của chúng, cho thấy khá rõ nhu cầu về
những phương pháp gần với lý thuyết lược đồ hơn (dĩ nhiên các kỹ thuật siêu việt
sẽ phải bổ sung cho chúng).

Trong Bài trình bày I này, ta sẽ xét trường hợp hạ xuống sơ đẳng nhất (được nêu
trong tiêu đề). Tuy vậy, các ứng dụng của các Định lý 1, 2 và 3 dưới đây đã rất
phong phú. Ta chỉ nêu một vài ứng dụng làm ví dụ, mà không cố phát biểu chúng ở
mức tổng quát tối đa có thể có tại đây.

Ta sẽ tự do sử dụng ngôn ngữ lược đồ; về ngôn ngữ này, xin xem bài trình bày đã
dẫn và [2]. Tuy nhiên, cần lưu ý rõ rằng các tiền lược đồ được xét trong bài này
không nhất thiết là Noether, và các

\src{300}
cấu xạ không nhất thiết thuộc kiểu hữu hạn. Chẳng hạn, nếu $A$ là một vành địa
phương Noether có vành đầy đủ hóa $\widehat A$, ta sẽ phải xét vành không
Noether $\widehat A\otimes_A\widehat A$, cùng các cấu xạ giữa những lược đồ
affine tương ứng với các phép nhúng giữa các vành đang xét.

\fsection{A}{Các kiến thức chuẩn bị về phạm trù}

\subsection*{1. Phạm trù phân thớ, dữ liệu hạ xuống, cấu xạ
$\mathcal F$-hạ xuống}
\addcontentsline{toc}{subsection}{1. Phạm trù phân thớ và dữ liệu hạ xuống}

\result{a. ĐỊNH NGHĨA}{1.1} Một \emph{phạm trù phân thớ} $\mathcal F$ trên cơ
sở $\mathcal C$ gồm một phạm trù $\mathcal C$; với mỗi
$X\in\mathcal C$, một phạm trù $\mathcal F_X$; với mỗi cấu xạ trong
$\mathcal C$, $f:X\longrightarrow Y$, một hàm tử
$f^*:\mathcal F_Y\longrightarrow\mathcal F_X$, cũng được ký hiệu bởi
\[
  f^*(\xi)=\xi\times_Y X
\]
với $\xi\in\mathcal F_Y$ (trong đó $X$ được xem là một ``đối tượng của
$\mathcal C$ trên $Y$'', nghĩa là được trang bị cấu xạ $f$); và cuối cùng, với
hai cấu xạ hợp thành được
$X\xrightarrow{f}Y\xrightarrow{g}Z$, một đẳng cấu hàm tử
\[
  c_{f,g}:(gf)^*\longrightarrow f^*g^*,
\]
sao cho các dữ liệu này thỏa mãn những điều kiện sau:

\begin{enumerate}
\item[(i)] $\operatorname{id}_X^*=\operatorname{id}$;
\item[(ii)] $c_{f,g}$ là đẳng cấu đồng nhất nếu $f$ hoặc $g$ là một đẳng cấu
đồng nhất;
\item[(iii)] với ba cấu xạ hợp thành được
$X\xrightarrow{f}Y\xrightarrow{g}Z\xrightarrow{h}T$, biểu đồ các đẳng cấu sau,
được tạo bởi những đẳng cấu dạng $c_{u,v}$, là giao hoán:
\[
\begin{array}{ccc}
  (h(gf))^*=((hg)f)^* & \longrightarrow & f^*(hg)^* \\
  \big\downarrow && \big\downarrow \\
  (gf)^*h^* & \longrightarrow & (f^*g^*)h^*=f^*(g^*h^*).
\end{array}
\]
\end{enumerate}

\result{VÍ DỤ}{1} Cho $\mathcal C$ là một phạm trù trong đó tồn tại các tích
thớ. Khi đó ta định nghĩa một phạm trù phân thớ $\mathcal F$ trên cơ sở
$\mathcal C$ bằng cách lấy $\mathcal F_X=$ phạm trù các đối tượng của
$\mathcal C$ trên $X$; hàm tử
$f^*:\mathcal F_Y\longrightarrow\mathcal F_X$ ứng với một cấu xạ
$f:X\longrightarrow Y$ được định nghĩa bằng tích thớ
$Z\longmapsto Z\times_YX$.

\result{VÍ DỤ}{2} Cho $\mathcal C$ là phạm trù các tiền lược đồ và, với
$X\in\mathcal C$, cho $\mathcal F_X$ là phạm trù các bó môđun tựa nhất quán
trên $X$. Nếu $f:X\longrightarrow Y$ là một cấu xạ tiền lược đồ, thì
$f^*:\mathcal F_Y\longrightarrow\mathcal F_X$ là hàm tử ảnh ngược

\src{301}
trên các bó môđun. Như vậy ta thu được một phạm trù phân thớ trên cơ sở
$\mathcal C$.

\result{b. ĐỊNH NGHĨA}{1.2} Một biểu đồ
\[
  E\xrightarrow{u}E'
  \overset{v_1}{\underset{v_2}{\rightrightarrows}}E''
\]
gồm các ánh xạ giữa những tập hợp được gọi là \emph{khớp} nếu $u$ là một song
ánh từ $E$ lên phần của $E'$ gồm các $x'\in E'$ thỏa mãn
$v_1(x')=v_2(x')$.

\result{ĐỊNH NGHĨA}{1.3} Cho $\mathcal F$ là một phạm trù phân thớ trên cơ sở
$\mathcal C$; xét một biểu đồ
\[
  \tag{+}
  S\xleftarrow{\alpha}S'
  \overset{\beta_1}{\underset{\beta_2}{\rightrightarrows}}S''
\]
gồm các cấu xạ trong $\mathcal C$ sao cho
$\alpha\beta_1=\alpha\beta_2$; biểu đồ này được gọi là
$\mathcal F$-\emph{khớp} nếu, với mọi cặp đối tượng $\xi,\eta$ của
$\mathcal F_S$, biểu đồ các ánh xạ giữa tập hợp sau là khớp:
\[
  \tag{+'}
  \operatorname{Hom}(\xi,\eta)
  \xrightarrow{\alpha^*}
  \operatorname{Hom}(\alpha^*\xi,\alpha^*\eta)
  \overset{\beta_1^*}{\underset{\beta_2^*}{\rightrightarrows}}
  \operatorname{Hom}(\gamma^*\xi,\gamma^*\eta),
\]
trong đó $\gamma=\alpha\beta_1=\alpha\beta_2$. Để đơn giản, trong biểu đồ cuối
ta đã đồng nhất $\beta_i^*\alpha^*$ với
$(\alpha\beta_i)^*=\gamma^*$ nhờ $c_{\beta_i,\alpha}$.

\result{ĐỊNH NGHĨA}{1.4} Cho $\mathcal F$ là một phạm trù phân thớ trên cơ sở
$\mathcal C$; xét hai cấu xạ $\beta_1,\beta_2$ từ $S''$ đến $S'$ trong
$\mathcal C$. Cho $\xi'\in\mathcal F_{S'}$. Một \emph{dữ liệu dán} trên
$\xi'$ (đối với cặp $(\beta_1,\beta_2)$) là một đẳng cấu từ
$\beta_1^*(\xi')$ lên $\beta_2^*(\xi')$. Nếu
$\xi',\eta'\in\mathcal F_{S'}$ đều được trang bị một dữ liệu dán, thì một cấu
xạ $u:\xi'\longrightarrow\eta'$ trong $\mathcal F_{S'}$ được gọi là
\emph{tương thích với các dữ liệu dán} nếu biểu đồ sau là giao hoán:
\[
\begin{array}{ccc}
  \beta_1^*(\xi') & \longrightarrow & \beta_2^*(\xi') \\
  \big\downarrow && \big\downarrow \\
  \beta_1^*(\eta') & \longrightarrow & \beta_2^*(\eta').
\end{array}
\]
Theo cách này, các đối tượng của $\mathcal F_{S'}$ được trang bị một dữ liệu
dán (đối với $\beta_1,\beta_2$) tạo thành một phạm trù. Nếu
$\alpha:S'\longrightarrow S$ là một cấu xạ sao cho
$\alpha\beta_1=\alpha\beta_2$, thì với mọi $\xi\in\mathcal F_S$, đối tượng
$\xi'=\alpha^*(\xi)$

\src{302}
của $\mathcal F_{S'}$ được trang bị một dữ liệu dán chính tắc, bởi vì
\[
  \beta_i^*\alpha^*(\xi)\simeq(\alpha\beta_i)^*(\xi)=\gamma^*(\xi),
\]
trong đó một lần nữa $\gamma=\alpha\beta_1=\alpha\beta_2$; hơn nữa, nếu
$u:\xi\longrightarrow\eta$ là một cấu xạ trong $\mathcal F_S$, thì
\[
  \alpha^*(u):\alpha^*(\xi)\longrightarrow\alpha^*(\eta)
\]
là một cấu xạ trong $\mathcal F_{S'}$ tương thích với các dữ liệu dán chính
tắc. Do đó ta thu được một hàm tử chính tắc từ phạm trù $\mathcal F_S$ đến
phạm trù các đối tượng của $\mathcal F_{S'}$ được trang bị một dữ liệu dán đối
với cặp $(\beta_1,\beta_2)$. Nói cách khác, Định nghĩa 1.3 cũng có thể được
phát biểu rằng biểu đồ (+) trong định nghĩa ấy là $\mathcal F$-khớp nếu hàm tử
trên là ``đầy đủ và trung thành'', nghĩa là nó xác định một tương đương giữa
phạm trù $\mathcal F_S$ và một phạm trù con của phạm trù các đối tượng
của $\mathcal F_{S'}$ được trang bị dữ liệu dán đối với
$(\beta_1,\beta_2)$.

\result{ĐỊNH NGHĨA}{1.5} Một dữ liệu dán trên
$\xi'\in\mathcal F_{S'}$ được gọi là \emph{hữu hiệu} (đối với $\alpha$) nếu
$\xi'$, cùng dữ liệu này, đẳng cấu với một $\alpha^*(\xi)$ nào đó, trong đó
$\xi\in\mathcal F_S$.

Khi biểu đồ (+) là $\mathcal F$-khớp, đối tượng $\xi$ trong định nghĩa trên vì
thế được xác định tới một đẳng cấu duy nhất, và phạm trù
$\mathcal F_S$ tương đương với phạm trù các đối tượng của
$\mathcal F_{S'}$ được trang bị một dữ liệu dán hữu hiệu.

\medskip\noindent
c. Trường hợp quan trọng nhất là khi
\[
  S''=S'\times_S S',
\]
trong đó các $\beta_i$ là hai phép chiếu $p_1$ và $p_2$ từ
$S'\times_S S'$ lên hai thừa số của nó (bây giờ ta giả sử rằng các tích thớ tồn
tại trong $\mathcal C$). Khi đó có ngay hai điều kiện cần để một dữ liệu dán
$\varphi:p_1^*(\xi')\longrightarrow p_2^*(\xi')$ trên một
$\xi'\in\mathcal F_{S'}$ là hữu hiệu. Ta có
\[
  \tag{i}
  \Delta^*(\varphi)=\operatorname{id}_{\xi'},
\]
trong đó $\Delta:S'\longrightarrow S'\times_S S'$ là cấu xạ đường chéo và ta
đã đồng nhất $\Delta^*p_i^*(\xi')$ với
$(p_i\Delta)^*(\xi')=\xi'$. Ta cũng có
\[
  \tag{ii}
  p_{31}^*(\varphi)=p_{32}^*(\varphi)\,p_{21}^*(\varphi),
\]

\src{303}
trong đó $p_{ij}$ là phép chiếu chính tắc từ
$S'\times_S S'\times_S S'$ lên tích riêng phần của thừa số thứ $i$ và thừa số
thứ $j$.

\result{ĐỊNH NGHĨA}{1.6} Một \emph{dữ liệu hạ xuống} trên
$\xi'\in\mathcal F_{S'}$, đối với cấu xạ
$\alpha:S'\longrightarrow S$, là một dữ liệu dán trên $\xi'$ đối với cặp
$(p_1,p_2)$ các phép chiếu chính tắc
$S'\times_S S'\longrightarrow S'$ thỏa mãn các điều kiện (i) và (ii) ở trên.

\result{ĐỊNH NGHĨA}{1.7} Một cấu xạ $\alpha:S'\longrightarrow S$ được gọi là
một \emph{cấu xạ $\mathcal F$-hạ xuống} nếu biểu đồ các cấu xạ
\[
  S\xleftarrow{\alpha}S'
  \overset{p_1}{\underset{p_2}{\rightrightarrows}}S'\times_S S'
\]
là $\mathcal F$-khớp (Định nghĩa 1.3); $\alpha$ được gọi là một
\emph{cấu xạ $\mathcal F$-hạ xuống nghiêm ngặt} nếu thêm vào đó, mọi dữ liệu hạ
xuống (Định nghĩa 1.6) trên một đối tượng của $\mathcal F_{S'}$ đều hữu hiệu.

Điều kiện cuối cũng có thể được diễn đạt trực quan hơn như sau: việc cho một
đối tượng của $\mathcal F_S$ là ``tương đương'' với việc cho một đối tượng của
$\mathcal F_{S'}$ được trang bị dữ liệu hạ xuống.

Lưu ý rằng nếu một cấu xạ $\alpha:S'\longrightarrow S$ có một ``tiết diện''
$s:S\longrightarrow S'$ (nghĩa là một cấu xạ $s$ sao cho
$\alpha s=\operatorname{id}_S$), thì nó là một cấu xạ
$\mathcal F$-hạ xuống nghiêm ngặt. Thật vậy, nếu
$\xi'\in\mathcal F_{S'}$ được trang bị một dữ liệu hạ xuống đối với $\alpha$,
thì nó ``đến từ'' $\xi=s^*(\xi')$.

\medskip\noindent
d. Có thể trình bày các khái niệm trên một cách trực quan hơn bằng cách đưa vào,
với một đối tượng tham số $T$ của $\mathcal C$ trên $S$, tập hợp
\[
  \operatorname{Hom}_S(T,S')=S'(T),
\]
có các phần tử được ký hiệu bởi $t,t'$, v.v. Với một đối tượng
$\xi'\in\mathcal F_{S'}$ đã cho, mỗi $t\in S'(T)$ khi đó tương ứng với một đối
tượng $t^*(\xi')$ của $\mathcal F_T$, cũng được ký hiệu là $\xi'_t$. Một dữ
liệu dán trên $\xi'$ (đối với $p_1,p_2$) khi đó được xác định bằng cách cho,
với mọi $T$ trên $S$ và mọi cặp điểm $t,t'\in S'(T)$, một đẳng cấu
\[
  \varphi_{t',t}:\xi'_t\longrightarrow\xi'_{t'}
\]
(thỏa mãn các điều kiện hiển nhiên về tính hàm tử khi $T$ biến thiên). Các điều
kiện (i) và (ii) của 1.c khi đó được viết thành
\[
  \tag{i bis}
  \varphi_{t,t}=\operatorname{id}_{\xi'_t}
  \qquad\text{với mọi $T$ và mọi $t\in S'(T)$},
\]

\src{304}
\[
  \tag{ii bis}
  \varphi_{t,t''}=\varphi_{t,t'}\varphi_{t',t''}
  \qquad\text{với mọi $T$ và mọi $t,t',t''\in S'(T)$}.
\]
Hơn nữa, lấy $t=t'=t''$ trong (ii bis), ta thấy
$\varphi_{t,t}=\varphi_{t,t}^2$; vì theo giả thiết $\varphi_{t,t}$ là một đẳng
cấu, suy ra hệ thức (i bis), nên thật ra (i bis) là hệ quả của (ii bis) (do đó
(i) là hệ quả của (ii)). Nhưng nếu không còn giả sử tiên nghiệm rằng các
$\varphi_{t,t'}$ là đẳng cấu (nghĩa là
$\varphi:p_1^*(\xi')\longrightarrow p_2^*(\xi')$ là một đẳng cấu), thì (ii bis)
không nhất thiết kéo theo (i bis); tuy vậy, kết hợp (ii bis) với (i bis) suy ra
rằng các $\varphi_{t,t'}$ là đẳng cấu (bởi ta có
$\varphi_{t,t'}\varphi_{t',t}=\varphi_{t,t}=\operatorname{id}_{\xi'_t}$).

\subsection*{2. Các biểu đồ khớp và các toàn cấu nghiêm ngặt, các cấu xạ
hạ xuống. Ví dụ}
\addcontentsline{toc}{subsection}{2. Các biểu đồ khớp và các toàn cấu nghiêm ngặt}

\result{a. ĐỊNH NGHĨA}{2.1} Cho $\mathcal C$ là một phạm trù. Một biểu đồ
\[
  T\xrightarrow{\alpha}T'
  \overset{\beta_1}{\underset{\beta_2}{\rightrightarrows}}T''
\]
gồm các cấu xạ trong $\mathcal C$ được gọi là \emph{khớp} nếu, với mọi
$Z\in\mathcal C$, biểu đồ tương ứng gồm các ánh xạ giữa tập hợp
\[
  \operatorname{Hom}(Z,T)\longrightarrow
  \operatorname{Hom}(Z,T')
  \rightrightarrows\operatorname{Hom}(Z,T'')
\]
là khớp (Định nghĩa 1.2). Khi đó ta nói $(T,\alpha)$ (hoặc, theo lối nói lạm
dụng, $T$) là một \emph{hạt nhân} của cặp cấu xạ
$(\beta_1,\beta_2)$.

Hạt nhân này hiển nhiên được xác định tới một đẳng cấu duy nhất. Nếu
$\mathcal C$ là phạm trù các tập hợp, thì định nghĩa trên tương thích với Định
nghĩa 1.2. Theo cách đối ngẫu, ta định nghĩa tính khớp của một biểu đồ
\[
  S\xleftarrow{\alpha}S'
  \overset{\beta_1}{\underset{\beta_2}{\rightrightarrows}}S''
\]
gồm các cấu xạ trong $\mathcal C$; khi đó ta nói $(S,\alpha)$ là một
\emph{đối hạt nhân} của cặp cấu xạ $(\beta_1,\beta_2)$.

\result{ĐỊNH NGHĨA}{2.2} Một cấu xạ $\alpha:S'\longrightarrow S$ được gọi là
một \emph{toàn cấu nghiêm ngặt} nếu nó là một toàn cấu và nếu, với mọi cấu xạ
$u:S'\longrightarrow Z$, điều kiện cần sau đây cũng đủ để $u$ phân tích thành
$S'\longrightarrow S\longrightarrow Z$: với mọi $S''\in\mathcal C$ và mọi
cặp cấu xạ $\beta_1,\beta_2:S''\longrightarrow S'$ sao cho
$\alpha\beta_1=\alpha\beta_2$, ta cũng có $u\beta_1=u\beta_2$.

Nếu tích thớ $S'\times_S S'$ tồn tại, điều này tương đương với việc biểu đồ
\[
  S\xleftarrow{\alpha}S'
  \overset{p_1}{\underset{p_2}{\rightrightarrows}}S'\times_S S'
\]

\src{305}
là khớp, nghĩa là $S$ là một đối hạt nhân của cặp $(p_1,p_2)$. Trong mọi
trường hợp, một cấu xạ đối hạt nhân là một toàn cấu nghiêm ngặt. Cũng lưu ý
rằng một toàn cấu nghiêm ngặt đồng thời là một đơn cấu thì là một đẳng cấu.
Ta để bạn đọc phát triển khái niệm đối ngẫu về \emph{đơn cấu nghiêm ngặt}.

Để làm rõ quan hệ giữa khái niệm cấu xạ $\mathcal F$-hạ xuống và khái niệm
toàn cấu nghiêm ngặt, ta đưa thêm vào các định nghĩa sau:

\result{ĐỊNH NGHĨA}{2.3} Một cấu xạ $\alpha:S'\longrightarrow S$ được gọi là
một \emph{toàn cấu phổ dụng} (tương ứng, một \emph{toàn cấu nghiêm ngặt phổ
dụng}) nếu, với mọi $T$ trên $S$, tích thớ $T'=S'\times_ST$ tồn tại và phép
chiếu $T'\longrightarrow T$ là một toàn cấu (tương ứng, một toàn cấu nghiêm
ngặt).

Trong các phạm trù rất tốt (như phạm trù các tập hợp, phạm trù các bó tập hợp
trên một không gian tôpô, các phạm trù abel, v.v.), bốn khái niệm về tính toàn
cấu vừa được đưa vào trùng nhau; trái lại, chúng đều khác nhau trong một phạm
trù như phạm trù các tiền lược đồ, hoặc phạm trù các tiền lược đồ trên một tiền
lược đồ không rỗng $S$ đã cho, ngay cả khi chỉ xét các $S$-lược đồ hữu hạn trên
$S$.

\result{ĐỊNH NGHĨA}{2.4} Một cấu xạ $\alpha:S'\longrightarrow S$ được gọi là
một \emph{cấu xạ hạ xuống} (tương ứng, một \emph{cấu xạ hạ xuống nghiêm
ngặt}) nếu nó là một cấu xạ $\mathcal F$-hạ xuống (tương ứng, một cấu xạ
$\mathcal F$-hạ xuống nghiêm ngặt) (xem Định nghĩa 1.7), trong đó $\mathcal F$
là phạm trù phân thớ trên cơ sở $\mathcal C$ gồm các đối tượng của
$\mathcal C$ trên các đối tượng của $\mathcal C$ (số 1, Ví dụ 1).

\result{MỆNH ĐỀ}{2.1} Nếu trong $\mathcal C$ tồn tại các tích hữu hạn và các
tích thớ (hữu hạn), thì trong $\mathcal C$, các cấu xạ hạ xuống chính là các
toàn cấu nghiêm ngặt phổ dụng.

\result{b. VÍ DỤ}{} Cho $\mathcal C$ là phạm trù các tiền lược đồ. Cho
$S\in\mathcal C$, và cho $S'$ và $S''$ là hai tiền lược đồ hữu hạn trên $S$,
nghĩa là tương ứng với các bó đại số $\mathcal A'$, $\mathcal A''$ trên $S$
mà, xét như các bó môđun, là tựa nhất quán và thuộc kiểu hữu hạn (nghĩa là
coherent nếu $S$ là Noether địa phương). Cho $\alpha:S'\longrightarrow S$ là
cấu xạ cấu trúc của $S'$, và cho $\beta_1,\beta_2$ là hai $S$-cấu xạ từ $S''$
đến $S'$, được xác định bởi các đồng cấu đại số
$\mathcal A'\longrightarrow\mathcal A''$, vẫn được ký hiệu là
$\beta_1,\beta_2$. Dùng sự kiện rằng một cấu xạ hữu hạn là đóng (định lý
Cohen--Seidenberg thứ nhất), ta dễ dàng chứng minh rằng biểu đồ
\[
  \tag{+}
  S\xleftarrow{\alpha}S'
  \overset{\beta_1}{\underset{\beta_2}{\rightrightarrows}}S''
\]
trong $\mathcal C$ là khớp khi và chỉ khi biểu đồ các bó

\src{306}
\[
  \mathcal O_S=\mathcal A\longrightarrow\mathcal A'
  \overset{\beta_1}{\underset{\beta_2}{\rightrightarrows}}\mathcal A''
\]
trên $S$ là khớp. Đặc biệt, nếu $\alpha:S'\longrightarrow S$ là một cấu xạ
hữu hạn tương ứng với một bó đại số $\mathcal A'$ trên $S$, thì $\alpha$ là
một toàn cấu nghiêm ngặt khi và chỉ khi biểu đồ các bó
\[
  \mathcal O_S=\mathcal A\longrightarrow\mathcal A'
  \overset{p_1}{\underset{p_2}{\rightrightarrows}}
  \mathcal A'\otimes_{\mathcal A}\mathcal A'
\]
là khớp (nó là một toàn cấu khi và chỉ khi
$\mathcal A\longrightarrow\mathcal A'$ là đơn ánh). Nếu $S$ là affine với
vành $A$, và do đó $S'$ là affine với vành $A'$ hữu hạn trên $A$, thì
$S'\longrightarrow S$ là một toàn cấu nghiêm ngặt khi và chỉ khi
$A\longrightarrow A'$ là một đẳng cấu từ $A$ lên vành con của $A'$ gồm những
$x'\in A'$ sao cho
\[
  1_{A'}\otimes_Ax'-x'\otimes_A1_{A'}=0
\]
(nó là một toàn cấu khi và chỉ khi $A\longrightarrow A'$ là đơn ánh). Như ta
đã lưu ý, ngay cả khi $S$ là lược đồ của một vành địa phương Artin, một cấu xạ
hữu hạn $S'\longrightarrow S$ là toàn cấu không nhất thiết là một toàn cấu
nghiêm ngặt. Tuy nhiên, có thể chứng minh rằng, nếu $S$ là một tiền lược đồ
Noether, thì mọi cấu xạ hữu hạn $S'\longrightarrow S$ là toàn cấu đều là hợp
thành của một dãy hữu hạn các toàn cấu nghiêm ngặt, mỗi cấu xạ trong dãy cũng
hữu hạn. Điều này cũng cho thấy hợp thành của hai toàn cấu nghiêm ngặt không
nhất thiết là một toàn cấu nghiêm ngặt.

\result{c}{} Nếu $(+)$ là một biểu đồ khớp gồm các cấu xạ hữu hạn, thì với mọi
cấu xạ phẳng $T\longrightarrow S$ của các tiền lược đồ, biểu đồ thu được từ
$(+)$ qua phép đổi cơ sở $T\longrightarrow S$ vẫn khớp. Suy ra rằng, nếu $X,Y$
là hai $S$-tiền lược đồ, với $X$ phẳng trên $S$, thì biểu đồ các ánh xạ giữa
tập hợp sau đây (trong đó $X',Y'$ là các ảnh ngược của $X,Y$ trên $S'$, còn
$X'',Y''$ là các ảnh ngược của chúng trên $S''$)
\[
  \operatorname{Hom}_S(X,Y)\longrightarrow
  \operatorname{Hom}_{S'}(X',Y')
  \rightrightarrows
  \operatorname{Hom}_{S''}(X'',Y'')
\]
là khớp. Đặc biệt, nếu $\mathcal F$ là phạm trù phân thớ trên cơ sở phạm trù
$\mathcal C$ các tiền lược đồ, sao cho với $X\in\mathcal C$, $\mathcal F_X$
là phạm trù các $X$-tiền lược đồ phẳng, thì biểu đồ $(+)$ là
$\mathcal F$-khớp. (Kết quả này không còn đúng nếu không đặt giả thiết phẳng;
đặc biệt, một toàn cấu nghiêm ngặt hữu hạn không nhất thiết là một cấu xạ hạ
xuống.) Tương tự, ta thấy $(+)$ là $\mathcal F$-khớp nếu $\mathcal F$ là phạm
trù phân thớ mà $\mathcal F_X$ là phạm trù các bó tựa nhất quán và phẳng trên
tiền lược đồ $X$ (ở đây giả thiết

\src{307}
phẳng cũng là thiết yếu). Trong cả hai trường hợp, vấn đề một dữ liệu dán có
hữu hiệu hay không (và cụ thể hơn, một dữ liệu hạ xuống khi
$S''=S'\times_SS'$ có hữu hiệu hay không) trên một đối tượng phẳng trên $S'$
là tinh tế (và lời giải của nó trong nhiều trường hợp riêng là một trong các
chủ đề chính của những bài trình bày này). Diễn giả không biết liệu, với mọi
toàn cấu nghiêm ngặt hữu hạn $S'\longrightarrow S$, mọi dữ liệu hạ xuống trên
một bó tựa nhất quán phẳng trên $S'$ có hữu hiệu hay không (ngay cả khi giả sử
$S$ là phổ của một vành địa phương Artin và chỉ xét các bó tự do địa phương
hạng 1). Tổng quát hơn, cho $A$ là một vành và $A'$ là một $A$-đại số (mọi
vành đều giao hoán) sao cho biểu đồ các ánh xạ
\[
  A\longrightarrow A'\rightrightarrows A'\otimes_AA'
\]
là khớp; điều này cũng tương đương với việc cấu xạ $S'\longrightarrow S$ tương
ứng giữa các phổ của $A,A'$ là một cấu xạ $\mathcal F$-hạ xuống, trong đó
$\mathcal F$ là phạm trù phân thớ các bó tựa nhất quán phẳng. Cho $M'$ là một
$A'$-môđun phẳng được trang bị một dữ liệu hạ xuống về $A$, nghĩa là một đẳng
cấu
\[
  \varphi:M'\otimes_AA'\xrightarrow{\sim}A'\otimes_AM'
\]
của các $A'\otimes_AA'$-môđun, thỏa mãn các điều kiện (i) và (ii) của 1.c (ta
để bạn đọc viết tường minh các điều kiện ấy bằng ngôn ngữ môđun). Dữ liệu này
có hữu hiệu hay không (đối với phạm trù phân thớ các bó tựa nhất quán phẳng)?
Cho $M$ là tập con của $M'$ gồm những $x'\in M'$ sao cho
\[
  \varphi(x'\otimes_A1_{A'})=1_{A'}\otimes_Ax',
\]
đó là một $A$-môđun con của $M'$. Phép nhúng chính tắc
$M\longrightarrow M'$ xác định một đồng cấu $A'$-môđun
$M\otimes_AA'\longrightarrow M'$. Khi đó tính hữu hiệu của $\varphi$ có nghĩa
là: $M$ là một $A$-môđun phẳng và đồng cấu trên là một đẳng cấu.

\result{NHẬN XÉT}{} Trong các lập luận trên, ta không đặt giả thiết phẳng nào
lên các cấu xạ của biểu đồ $(+)$; vì thế, để có một kỹ thuật hạ xuống, ta buộc
phải đặt các giả thiết phẳng lên những đối tượng trên $S,S'$ đang xét. Trong
§2, ta sẽ đặt giả thiết phẳng lên $\alpha:S'\longrightarrow S$, nhờ đó có được
một kỹ thuật hạ xuống cho những đối tượng trên $S,S'$ không còn phải chịu bất
kỳ điều kiện phẳng nào. Trong mọi trường hợp đều có một giả thiết phẳng tham
gia. Đây là một trong những lý do chính khiến khái niệm tính phẳng quan trọng
trong hình học đại số (vai trò này không thể lộ ra chừng nào người ta còn chỉ
xét các trường cơ sở, bởi trên đó, quả thực, mọi thứ đều phẳng!).

\src{308}
\subsection*{3. Ứng dụng vào các cấu xạ étale}
\addcontentsline{toc}{subsection}{3. Ứng dụng vào các cấu xạ étale}

Cho $A$ là một vành địa phương và $B$ là một đại số địa phương trên $A$ mà
iđêan cực đại của nó cảm sinh iđêan cực đại của $A$. Ta nói $B$ là
\emph{étale} trên $A$ (thay cho ``không phân nhánh'', cách gọi được dùng ở nơi
khác) nếu nó thỏa mãn các điều kiện sau:
\begin{enumerate}
\item[(i)] $B$ phẳng trên $A$;
\item[(ii)] $B/\mathfrak mB$ là một mở rộng hữu hạn khả phân của
$A/\mathfrak m=k$ (trong đó $\mathfrak m$ ký hiệu iđêan cực đại của $A$).
\end{enumerate}
Khi $A$ và $B$ là Noether và $k$ đóng đại số, điều này có nghĩa là đồng cấu
$\widehat A\longrightarrow\widehat B$ giữa các vành đầy đủ hóa, mở rộng
$A\longrightarrow B$, là một đẳng cấu. Một cấu xạ thuộc kiểu hữu hạn
$f:T\longrightarrow S$ được gọi là \emph{étale tại $x\in T$}, hay cũng nói
$T$ \emph{étale trên $S$ tại $x$}, nếu $\mathcal O_x$ étale trên
$\mathcal O_{f(x)}$; còn $f$ được gọi là \emph{étale}, hay $f$ được gọi là một
\emph{cấu xạ étale}, hay $T$ được gọi là \emph{étale trên $S$}, nếu $f$ étale tại mọi
$x\in T$. Hơn nữa, lưu ý rằng nếu $S$ là Noether địa phương, thì tập các điểm của
$T$ tại đó $f$ étale là một tập mở; việc dùng ``định lý chính'' của Zariski cho
phép mô tả chính xác cấu trúc của $T/S$ trong một lân cận của một điểm như thế
(bằng một phương trình thuộc kiểu quen thuộc).

Khi $S$ là một lược đồ thuộc kiểu hữu hạn trên trường các số phức, nó tương ứng
với một không gian giải tích $\overline S$ theo nghĩa của Serre [5], ngoại trừ
việc $\overline S$ có thể chứa các phần tử lũy linh trong bó cấu trúc của nó;
điều này không làm thay đổi gì cốt yếu trong [5]. Khi đó dễ thấy $f$ étale khi
và chỉ khi $\overline f:\overline T\longrightarrow\overline S$ étale, nghĩa là
mọi điểm của $\overline T$ có một lân cận mà trên đó $\overline f$ cảm sinh một
đẳng cấu lên một tập mở của $\overline S$. Đặc biệt, với mỗi \emph{phủ étale}
$T$ của $S$ (nghĩa là một cấu xạ hữu hạn étale $f:T\longrightarrow S$), ta có
một phủ étale $\overline T$ của $\overline S$, và phủ này liên thông khi và chỉ
khi $T$ liên thông [5]. Cũng dễ thấy rằng, nếu $T,T'$ là hai lược đồ
étale trên $S$, thì ánh xạ tự nhiên
\[
  \operatorname{Hom}_S(T,T')\longrightarrow
  \operatorname{Hom}_{\overline S}(\overline T,\overline{T'})
\]
là song ánh; nghĩa là hàm tử $T\longmapsto\overline T$ từ phạm trù các lược đồ
étale trên $S$ đến phạm trù các không gian giải tích étale trên $\overline S$
là ``đầy đủ và trung thành'', nên xác định một tương đương giữa phạm trù thứ
nhất và một phạm trù con của phạm trù thứ hai. Một định lý của
Grauert--Remmert [2] suy ra rằng, nếu $S$ chuẩn tắc, ta thu được một tương đương
giữa phạm trù các \emph{phủ étale} của $S$ và phạm trù các phủ étale (hữu hạn)
của $\overline S$; nghĩa là mọi phủ étale của $\overline S$ đều
$\overline S$-đẳng cấu với một $\overline T$, trong đó $T$ là một phủ étale của
$S$. Ta chứng minh rằng \emph{định lý Grauert--Remmert vẫn đúng mà không cần
giả thiết $S$ chuẩn tắc}. Trước hết, cho $S'\longrightarrow S$ là một toàn cấu
nghiêm ngặt hữu hạn; giả sử

\src{309}
định lý đã được chứng minh cho $S'$, và ta chứng minh nó đúng cho $S$. Thật vậy,
cho $\mathfrak T$ là một phủ étale của $\overline S$, và xét nghịch ảnh
$\mathfrak T'$ của nó trên $\overline{S'}$. Nghịch ảnh này tương ứng với một bó
giải tích coherent các đại số $\mathcal A'$ trên $\overline{S'}$, là nghịch ảnh
của bó các đại số $\mathcal A$ trên $\overline S$ xác định $\mathfrak T$. Theo giả
thiết, trên $\overline{S'}$, $\mathfrak T'$ đến từ một phủ étale $T'$ của $S'$,
nghĩa là $\mathcal A'$ đến từ một bó coherent các đại số $A'$ trên $S'$. Mặt
khác, $\mathcal A'$ được trang bị một dữ liệu hạ xuống chính tắc đối với
$\overline{S'}\longrightarrow\overline S$, nghĩa là một đẳng cấu giữa hai
nghịch ảnh của nó trên
\[
  \overline{S'}\times_{\overline S}\overline{S'}
  =\overline{(S'\times_SS')},
\]
(thỏa mãn các điều kiện (i), (ii)); theo những điều đã nói, đẳng cấu này đến từ
một đẳng cấu giữa các bó đại số tương ứng, nghĩa là từ một dữ liệu hạ xuống
trên $A'$ đối với $S'\longrightarrow S$. Dễ kiểm tra dữ liệu sau là hữu hiệu
(vì dữ liệu trên $\mathcal A'$ là hữu hiệu, và tính hữu hiệu của một dữ liệu hạ
xuống, như đã được giải thích tường minh ở mục trước, có thể được nhận biết địa
phương trên các \emph{đầy đủ hóa} của những môđun tham gia). Do đó ta thu
được một bó coherent các đại số $A$ trên $S$, xác định một phủ $T$ của $S$;
đây chính là phủ cần tìm. Kết quả trên hiển nhiên vẫn đúng nếu
$S'\longrightarrow S$ chỉ là hợp thành của một số hữu hạn các toàn cấu nghiêm
ngặt hữu hạn, nghĩa là một toàn cấu hữu hạn bất kỳ (theo kết quả phân tích thành
nhân tử được nêu trong mục 2). Suy ra định lý Grauert--Remmert vẫn đúng khi $S$
là một lược đồ \emph{rút gọn}, nghĩa là $\mathcal O_S$ không có phần tử lũy
linh, như ta thấy bằng cách đưa vào phép chuẩn hóa $S'$ của nó. Từ đó dễ dàng
chuyển sang trường hợp tổng quát.

Một chứng minh hoàn toàn tương tự, lại dùng kết quả phân tích thành nhân tử cho
các toàn cấu nghiêm ngặt hữu hạn và bản chất ``hình thức'' của tính hữu hiệu của
các dữ liệu hạ xuống, cho phép chứng minh kết quả sau: cho $S$ là một tiền lược
đồ Noether địa phương và $S'\longrightarrow S$ là một cấu xạ \emph{hữu hạn,
toàn ánh và thuần bất khả phân} (hay, tương đương, một cấu xạ thuộc kiểu hữu hạn
sao cho với mọi $T$ trên $S$, cấu xạ
$T'=S'\times_ST\longrightarrow T$ là một đồng phôi; điều này cũng được diễn đạt
bằng cách nói rằng $S'\longrightarrow S$ là một \emph{đồng phôi phổ dụng}). Với
mỗi $T$ étale trên $S$, xét nghịch ảnh $T'=T\times_SS'$ của nó, là étale trên
$S'$. Khi đó hàm tử $T\longmapsto T'$ là một \emph{tương đương phạm trù từ
phạm trù các tiền lược đồ $T$ étale trên $S$ đến phạm trù các tiền lược đồ $T'$
étale trên $S'$}. (Ta dùng tính song ánh của
\[
  \operatorname{Hom}_S(T_1,T_2)\longrightarrow
  \operatorname{Hom}_{S'}(T'_1,T'_2)
\]
đối với hai tiền lược đồ $T_1,T_2$ étale trên $S$, một sự kiện dễ kiểm tra trực
tiếp, và sự kiện rằng định lý vừa nêu đúng nếu
$S'=(S,\mathcal O_S/\mathcal J)$,

\src{310}
trong đó $\mathcal J$ là một bó coherent các iđêan lũy linh của $\mathcal O_S$
([4], Bổ đề 6)). Hơn nữa, lưu ý rằng ở đây ta không giả sử các $T$ đang xét hữu
hạn trên $S$. Đặc biệt, kết quả này suy ra cấu xạ $S'\longrightarrow S$ cảm
sinh một đẳng cấu từ nhóm cơ bản $[\,]$ của $S'$ lên nhóm cơ bản của $S$
(``\emph{tính bất biến tôpô của nhóm cơ bản của một tiền lược đồ}'').

\subsection*{4. Quan hệ với đối đồng điều bậc 1}
\addcontentsline{toc}{subsection}{4. Quan hệ với đối đồng điều bậc 1}

\medskip\noindent
a. Cho $\mathcal C$ là một phạm trù trong đó tích của hai đối tượng bất kỳ luôn
tồn tại, và cho $T\in\mathcal C$. Với mỗi tập hữu hạn
$I\ne\varnothing$, ta có thể xét $T^I$; khi $I$ biến thiên, ta thu được một hàm
tử hiệp biến từ phạm trù các tập hữu hạn khác rỗng đến $\mathcal C$, nghĩa là
một cái có thể gọi là \emph{đối tượng đơn hình} của $\mathcal C$, ký hiệu
$K_T$. Đối tượng này phụ thuộc hiệp biến vào $T$; hơn nữa, nếu $u,v$ là hai
cấu xạ $T\longrightarrow T'$, thì các cấu xạ tương ứng
$K_T\longrightarrow K_{T'}$ đồng luân với nhau. Ta nói $T$
\emph{trội hơn} $T'$ nếu $\operatorname{Hom}(T,T')\ne\varnothing$; đây là một
quan hệ tiền thứ tự lọc tăng trên $\mathcal C$. Từ những điều trên suy ra rằng,
nếu $T$ trội hơn $T'$, thì có một lớp chính tắc (sai khác một đồng luân) các
đồng cấu giữa những đối tượng đơn hình $K_T\longrightarrow K_{T'}$; đặc biệt,
nếu $K_T$ và $K_{T'}$ sao cho mỗi đối tượng trội hơn đối tượng kia, thì $K_T$
và $K_{T'}$ tương đương đồng luân. Bây giờ cho $F$ là một hàm tử (phản biến,
để cố định ý tưởng) từ $\mathcal C$ đến một phạm trù Abel $\mathcal C'$; khi đó
\[
  C^*(T,F)=F(K_T)
\]
là một đối tượng đối đơn hình của $\mathcal C'$, và do đó xác định theo cách
quen thuộc một phức (đối dây chuyền) trong $\mathcal C'$, mà ta có thể lấy đối
đồng điều:
\[
  H^*(T,F)=H^*(C^*(T,F))=H^*(F(K_T))
\]
(Ta có thể đặt một $\mathcal C$ làm chỉ số dưới của $H^*$ nếu có khả năng nhầm
lẫn.) Đây là một hàm tử đối đồng điều theo $F$; tính hiệp biến hoặc phản biến
của nó khi $T$ biến thiên suy ra từ những điều đã nói về $K_T$. Chính xác hơn,
với $F$ cố định và $T$ biến thiên trong $\mathcal C$ (được tiền thứ tự hóa bởi
quan hệ trội hơn), các $H^*(T,F)$ lập thành một hệ quy nạp các đối tượng phân
bậc của $\mathcal C'$; đặc biệt, nếu $T$ và $T'$ sao cho mỗi đối tượng trội hơn
đối tượng kia, thì $H^*(T,F)$ và $H^*(T',F)$ đẳng cấu chính tắc.

Giả sử các tích thớ tồn tại trong $\mathcal C$. Với $S\in\mathcal C$ cố định,
ta có thể áp dụng những điều trên cho phạm trù $\mathcal C_S$ gồm các đối tượng
của $\mathcal C$ trên $S$; ta sẽ viết $C^*(T/S,F)$ và $H^*(T/S,F)$ thay cho
$C^*(T,F)$ và $H^*(T,F)$ khi muốn nói rõ rằng ta đang làm việc trong phạm trù
$\mathcal C_S$; như vậy,

\src{311}
$C^*(T/S,F)$ là một phức đối dây chuyền trong $\mathcal C'$ mà số hạng ở bậc
$n$ bằng
\[
  F(T\times_ST\times_S\cdots\times_ST)
\]
(trong đó biểu thức trong ngoặc có $n+1$ thừa số).

Như thường lệ, lưu ý rằng ta có thể định nghĩa $H^0(T/S,F)$ mà không giả sử
phạm trù $\mathcal C'$ là Abel: đó là hạt nhân (Định nghĩa 2.1), nếu tồn tại,
của cặp cấu xạ $F(p_i)$ ($i=1,2$)
\[
  F(T)\rightrightarrows F(T\times_ST)
\]
tương ứng với hai phép chiếu
$p_1,p_2:T\times_ST\rightrightarrows T$. Đặc biệt, ta có một cấu xạ tự nhiên
(gọi là \emph{cấu xạ tăng cường})
\[
  F(S)\longrightarrow H^0(T/S,F)
\]
là một đẳng cấu trong các trường hợp thuận lợi (đặc biệt, nếu
$T\longrightarrow S$ là một toàn cấu nghiêm ngặt và $F$ “khớp trái”).
Tương tự, khi $F$ nhận giá trị trong phạm trù các nhóm trong một phạm trù
$\mathcal C''$, ta cũng có thể định nghĩa $H^1(T/S,F)$. Khi $\mathcal C''$ là
phạm trù các tập hợp (nghĩa là $F$ nhận giá trị trong phạm trù các nhóm thông
thường, không nhất thiết giao hoán), $H^1(T/S,F)$ là tập thương của nhóm con
$Z^1(T/S,F)$ của $C^1(T/S,F)=F(T\times_ST)$ gồm các phần tử $g$ thỏa mãn
\[
  F(p_{31})(g)=F(p_{32})(g)\,F(p_{21})(g)
\]
theo nhóm toán tử $F(T)$, tác động lên $C^1(T/S,F)$ và, đặc biệt, lên tập con
$Z^1(T/S,F)$ bởi
\[
  \rho(g')\cdot g=F(p_2)(g')\,g\,F(p_1)(g')^{-1}.
\]

\medskip\noindent
b. Chẳng hạn, cho $\mathcal F$ là một phạm trù phân thớ trên cơ sở
$\mathcal C$. Cho $\xi,\eta\in\mathcal F_S$ và, với mỗi $S'$ trên $S$, đặt
\[
  F_{\xi,\eta}(S')=
  \operatorname{Hom}(\xi\times_SS',\eta\times_SS').
\]
Như vậy $F_{\xi,\eta}$ là một hàm tử phản biến từ $\mathcal C_S$ đến phạm trù
các tập hợp. Với ký hiệu đó, \emph{nói rằng cấu xạ tăng cường}
\[
  F_{\xi,\eta}(S)\longrightarrow H^0(S'/S,F_{\xi,\eta})
\]
\emph{là một đẳng cấu với mọi cặp phần tử
$\xi,\eta\in\mathcal F_S$ có nghĩa là
$\alpha:S'\longrightarrow S$ là một cấu xạ $\mathcal F$-hạ xuống}
(Định nghĩa 1.7).

\medskip\noindent
c. Tương tự, với $\xi\in\mathcal F_S$ và mỗi đối tượng $S'$ của $\mathcal C$
trên $S$, đặt
\[
  G_\xi(S')=\operatorname{Aut}(\xi\times_SS') ;
\]
như vậy ta đã định nghĩa một hàm tử phản biến $G_\xi$ từ $\mathcal C_S$ đến
phạm trù các

\src{312}
nhóm. Khi đó ta thấy $Z^1(S'/S,G_\xi)$ \emph{được đồng nhất chính tắc với tập
các dữ liệu hạ xuống trên $\xi'=\xi\times_SS'$ đối với
$S'\longrightarrow S$} (Định nghĩa 1.6), còn $H^1(S'/S,G_\xi)$ \emph{được
đồng nhất với tập các lớp, sai khác một đẳng cấu, của những đối tượng trong
$\mathcal F_{S'}$ được trang bị một dữ liệu hạ xuống đối với
$\alpha:S'\longrightarrow S$ và, xét như các đối tượng của
$\mathcal F_{S'}$, đẳng cấu với $\xi'=\xi\times_SS'$}. Do đó, nếu
$\alpha:S'\longrightarrow S$ là một cấu xạ $\mathcal F$-hạ xuống (xem (b)),
thì $H^1(S'/S,G_\xi)$ \emph{chứa như một tập con tập các lớp, sai khác một
đẳng cấu, của những đối tượng $\eta$ trong $\mathcal F_S$ sao cho
$\eta\times_SS'$ đẳng cấu trong $\mathcal F_{S'}$ với $\xi\times_SS'$; và phép
nhúng này là ánh xạ đồng nhất khi và chỉ khi mọi dữ liệu hạ xuống trên
$\xi'=\xi\times_SS'$ đối với $\alpha:S'\longrightarrow S$ đều hữu hiệu}.
(Điều này đặc biệt đúng nếu $\alpha:S'\longrightarrow S$ là một cấu xạ
$\mathcal F$-hạ xuống nghiêm ngặt.)

\result{NHẬN XÉT}{} Các phức đối dây chuyền dạng $C^*(T/S,F)$ chứa, như những
trường hợp riêng, phần lớn các phức chuẩn quen thuộc (đối đồng điều Čech, đối
đồng điều nhóm, v.v.), và giữ một vai trò quan trọng trong hình học đại số,
đặc biệt trong “đối đồng điều Weil” của các tiền lược đồ.

\medskip\noindent
d. \result{VÍ DỤ}{1} Cho $S'$ là một đối tượng trên $S\in\mathcal C$, và cho
$\Gamma$ là một nhóm tự đẳng cấu của $S'$ sao cho $S'$ “là chính trên $S$
theo nghĩa hình thức, với nhóm $\Gamma$”, nghĩa là cấu xạ tự nhiên
\[
  \Gamma\times S'\longrightarrow S'\times_SS',
\]
trong đó $\Gamma\times S'$ chỉ tổng trực tiếp của $\Gamma$ bản sao của $S'$,
là một đẳng cấu. (Ta giả sử các tổng trực tiếp xuất hiện
ở đây tồn tại trong $\mathcal C$.) Cho $F$ là một hàm tử phản biến từ
$\mathcal C$ đến phạm trù các nhóm Abel. Khi đó $C^*(S'/S,F)$ đẳng cấu chính
tắc với nhóm đơn hình của các đối dây chuyền thuần nhất chuẩn
$C^*(\Gamma,F(S'))$; do đó $H^*(S'/S,F)$ đẳng cấu chính tắc với
$H^*(\Gamma,F(S'))$.

\medskip\noindent
e. \result{VÍ DỤ}{2} Cho $\mathcal C$ là phạm trù các tiền lược đồ. Ký hiệu
$G_a$ (“nhóm cộng”) là hàm tử phản biến từ $\mathcal C$ đến phạm trù các
nhóm Abel, được định nghĩa bởi
\[
  G_a(X)=H^0(X,\mathcal O_X).
\]
Tương tự, ta định nghĩa hàm tử $G_m$ (“nhóm nhân”) bởi
\[
  G_m(X)=H^0(X,\mathcal O_X)^*
\]
(nhóm các phần tử khả nghịch của vành $H^0(X,\mathcal O_X)$), và tổng quát hơn,
hàm tử $\operatorname{Gl}(n)$ (“nhóm tuyến tính cấp $n$”) bởi

\src{313}
\[
  \operatorname{Gl}(n)(X)=
  \operatorname{Gl}\bigl(n,H^0(X,\mathcal O_X)\bigr),
\]
là một hàm tử từ $\mathcal C$ đến phạm trù các nhóm (không nhất thiết giao
hoán nếu $n>1$; với $n=1$ ta thu lại $G_m$). Hơn nữa,
$\operatorname{Gl}(n)$ có thể được diễn giải như một hàm tử các tự đẳng cấu (xem
(c)) bằng cách xét phạm trù phân thớ $\mathcal F$ trên cơ sở $\mathcal C$ sao
cho, với $X\in\mathcal C$, $\mathcal F_X$ là phạm trù các bó tự do địa phương
trên $X$: quả vậy,
\[
  \operatorname{Gl}(n)(X)=
  \operatorname{Aut}_{\mathcal F_X}(\mathcal O_X^n).
\]
Theo (b), suy ra rằng nếu $\alpha:S'\longrightarrow S$ là một cấu xạ
$\mathcal F$-hạ xuống (xem đoạn 2.c), thì
$H^1(S'/S,\operatorname{Gl}(n))$ chứa tập các lớp, sai khác một đẳng cấu, của
những bó tự do địa phương trên $S$ mà nghịch ảnh trên $S'$ đẳng cấu với
$\mathcal O_{S'}^n$; và phép nhúng này là một đẳng thức khi và chỉ khi mọi dữ
liệu hạ xuống trên $\mathcal O_{S'}^n$ (đối với
$\alpha:S'\longrightarrow S$) đều hữu hiệu. Khi $S$ là phổ của một vành địa
phương, điều đó có nghĩa là
$H^1(S'/S,\operatorname{Gl}(n))=(e)$, vì mọi bó tự do địa phương trên $S$ khi
ấy đều tầm thường.

Lưu ý tính tương đương của các điều kiện sau đây đối với một cấu xạ
$\alpha:S'\longrightarrow S$:
\begin{enumerate}
\item[(i)] đồng cấu tăng cường
\[
  H^0(S,\mathcal O_S)=G_a(S)\longrightarrow H^0(S'/S,G_a)
\]
là một đẳng cấu;
\item[(ii)] $S'\longrightarrow S$ là một cấu xạ $\mathcal F$-hạ xuống
(trong đó $\mathcal F$ là phạm trù phân thớ trên cơ sở $\mathcal C$ đã xét ở
trên).
\end{enumerate}
Nếu $S'\longrightarrow S$ là hữu hạn, các điều kiện này còn tương đương với:
\begin{enumerate}
\item[(iii)] $S'\longrightarrow S$ là một toàn cấu nghiêm ngặt
(xem đoạn 2.c).
\end{enumerate}

Bây giờ giả sử $S=\operatorname{Spec}(A)$ và
$S'=\operatorname{Spec}(A')$. Khi đó
\[
  C^n(S'/S,G_a)=C^n(A'/A,G_a)=\bigotimes_A^{n+1}A',
\]
trong đó toán tử đối biên
$C^n(A'/A,G_a)\longrightarrow C^{n+1}(A'/A,G_a)$ là tổng luân phiên của các
toán tử mặt
\[
  \partial_i(x_0\otimes x_1\otimes\cdots\otimes x_n)
  =x_0\otimes\cdots\otimes x_{i-1}\otimes1_{A'}\otimes
   x_i\otimes\cdots\otimes x_n.
\]
Tương tự,
\[
  C^n(S'/S,G_m)=C^n(A'/A,G_m)
  \simeq\left(\bigotimes_A^{n+1}A'\right)^*,
\]
các phép toán đơn hình trong $C^*(A'/A,G_m)$ được cảm sinh bởi những phép toán
trong $C^*(S'/S,G_a)$. Các phép toán đơn hình trong
$C^*(A'/A,\operatorname{Gl}(n))$ cũng được viết tường minh theo cách ấy.
\emph{Trong mọi trường hợp mà báo cáo viên biết, ta có}
\[
  H^i(A'/A,G_a)=0\qquad\text{với $i>0$},
\]
\emph{và nếu $A$ là địa phương, ta có}
\[
  H^1(A'/A,G_m)=0
  \quad\text{và tổng quát hơn}\quad
  H^1(A'/A,\operatorname{Gl}(n))=(e)
\]
(khi $S'\longrightarrow S$ là một cấu xạ $\mathcal F$-hạ xuống, nghĩa là biểu
đồ
$A\longrightarrow A'\rightrightarrows A'\otimes_AA'$ khớp; so sánh với đoạn
2.c). Lưu ý rằng \emph{“định lý 90” của Hilbert chẳng qua là hệ thức}

\src{314}
\[
  H^1(S'/S,G_m)=0
\]
\emph{khi $A$ là một trường và $A'$ là một mở rộng Galois hữu hạn của trường ấy
(xem Ví dụ 1), và còn có thể được phát biểu rằng, trong trường hợp đang xét,
$S'\longrightarrow S$ là một cấu xạ hạ xuống nghiêm ngặt đối với phạm trù phân
thớ các bó tự do địa phương hạng 1.} Chính dưới dạng cuối này mà định lý Hilbert
nên được tổng quát hóa, bằng cách thay đổi cả các giả thiết về cấu xạ
$S'\longrightarrow S$ lẫn các giả thiết về những bó tựa nhất quán đang xét.

Cuối cùng, lưu ý tính tương đương của các tính chất sau khi $A$ là một vành
Artin địa phương có iđêan cực đại $\mathfrak m$ và $A'$ là một $A$-đại số
(trong đó, với mỗi số nguyên $k>0$, $A_k$ (tương ứng $A'_k$) ký hiệu vành
$A/\mathfrak m^{k+1}$ (tương ứng $A'/\mathfrak m^{k+1}A'$)):
\begin{enumerate}
\item[(i)] $H^1(A'_k/A_k,G_a)=0$ với mọi $k$.
\item[(ii)] $H^1(A'_k/A_k,G_m)=0$ với mọi $k$.
\item[(iii)] $H^1(A'_k/A_k,\operatorname{Gl}(n))=(e)$ với mọi $k$ và mọi
$n$.
\end{enumerate}
Nếu $S'\longrightarrow S$ là một toàn cấu nghiêm ngặt, thì các điều kiện trên
thậm chí còn kéo theo rằng đó là một cấu xạ hạ xuống nghiêm ngặt đối với các
môđun tự do trên $A'$ (thuộc kiểu hữu hạn hoặc không).

\result{NHẬN XÉT}{} Định nghĩa các nhóm $H^i(S'/S,G_m)$, khi $S,S'$ là phổ của
các trường $A,A'$, là do Amitsur. Nhóm $H^2(S'/S,G_m)$ đặc biệt đáng quan tâm
như một biến thể “toàn cục” của nhóm Brauer; về biến thể này có thể tham khảo
[1], Chương VII.

\fsection{B}{Hạ xuống theo các cấu xạ phẳng trung thành}

\subsection*{1. Phát biểu các định lý hạ xuống}
\addcontentsline{toc}{subsection}{1. Phát biểu các định lý hạ xuống}

\result{ĐỊNH NGHĨA}{1.1} Một cấu xạ
$\alpha:S'\longrightarrow S$ của các tiền lược đồ được gọi là \emph{phẳng} nếu,
với mọi $x'\in S'$, $\mathcal O_{x'}$ là một môđun phẳng trên vành
$\mathcal O_{\alpha(x')}$ (nghĩa là
$\mathcal O_{x'}\otimes_{\mathcal O_{\alpha(x')}}M$ là một hàm tử khớp theo
$\mathcal O_{\alpha(x')}$-môđun $M$). Một cấu xạ được gọi là \emph{phẳng trung
thành} nếu nó phẳng và toàn ánh.

Chẳng hạn, nếu $S=\operatorname{Spec}(A)$ và
$S'=\operatorname{Spec}(A')$, thì $S'$ phẳng trên $S$ khi và chỉ khi $A'$ là
một $A$-môđun phẳng, và $S'$ phẳng trung thành trên $S$ khi và chỉ khi $A'$ là
một $A$-môđun phẳng trung thành (nghĩa là hàm tử $A'\otimes_A M$ theo
$A$-môđun $M$ là khớp và trung thành); theo thuật ngữ của Serre [5], điều này
cũng có nghĩa là cặp $(A,A')$ phẳng. Nếu $S'$ phẳng trung thành trên $S$, thì
hàm tử ảnh ngược trên các bó tựa nhất quán trên $S$ là khớp và trung thành;
nói cách khác, để một dãy các đồng cấu của các bó

\src{315}
tựa nhất quán trên $S$ là khớp, điều kiện cần và đủ là ảnh ngược của nó trên
$S'$ là khớp (đặc biệt, để một đồng cấu của các bó tựa nhất quán trên $S$ là
đơn cấu, tương ứng là toàn cấu, tương ứng là đẳng cấu, điều kiện cần và đủ là
ảnh ngược của nó trên $S'$ có tính chất ấy). Tính chất này vẫn đúng khi hạn chế
trên một phần mở tùy ý của $S$ và, dưới dạng đó, nó đặc trưng các cấu xạ phẳng
trung thành.

\result{ĐỊNH NGHĨA}{1.2} Một cấu xạ $\alpha:S'\longrightarrow S$ được gọi là
\emph{quasi-compact} nếu ảnh ngược của mọi phần mở quasi-compact $U$ của $S$
là quasi-compact (nghĩa là hợp hữu hạn của các phần mở affine).

Hiển nhiên, chỉ cần kiểm tra tính chất này đối với các phần mở affine của $S$.
Chẳng hạn, một cấu xạ affine (nghĩa là cấu xạ sao cho ảnh ngược của một phần mở
affine là affine) là quasi-compact.

Lớp các cấu xạ phẳng, tương ứng phẳng trung thành, tương ứng quasi-compact, ổn
định qua phép hợp thành và qua ``đổi cơ sở'', và dĩ nhiên chứa các đẳng cấu.

\result{ĐỊNH LÝ}{1} Cho $\alpha:S'\longrightarrow S$ là một cấu xạ phẳng trung
thành và quasi-compact của các tiền lược đồ. Khi đó $\alpha$ là một \emph{cấu
xạ hạ xuống nghiêm ngặt} (A, Định nghĩa 1.7) đối với phạm trù phân thớ
$\mathcal F$ của các bó tựa nhất quán (A, Mục 1, Ví dụ 2).

Phát biểu này có hai ý:
\begin{enumerate}
\item[(i)] Nếu $F$ và $G$ là hai bó tựa nhất quán trên $S$, còn $F'$ và $G'$
là các ảnh ngược của chúng trên $S'$, thì đồng cấu tự nhiên
\[
  \operatorname{Hom}(F,G)\longrightarrow\operatorname{Hom}(F',G')
\]
là một song ánh từ vế thứ nhất lên nhóm con của vế thứ hai gồm các đồng cấu
$F'\longrightarrow G'$ tương thích với các dữ liệu hạ xuống chính tắc trên các
bó ấy, nghĩa là những đồng cấu mà ảnh ngược theo hai phép chiếu từ
$S''=S'\times_SS'$ lên $S'$ cho cùng một đồng cấu
$F''\longrightarrow G''$.
\item[(ii)] Mọi bó tựa nhất quán $F'$ trên $S'$, được trang bị một dữ liệu hạ
xuống đối với cấu xạ $\alpha:S'\longrightarrow S$ (A, Định nghĩa 1.6), đều
đẳng cấu (cùng với dữ liệu đó) với ảnh ngược của một bó tựa nhất quán $F$ trên
$S$.
\end{enumerate}
Lấy $F=\mathcal O_S$ trong (i), ta được:

\result{HỆ QUẢ}{1} Cho $G$ là một bó tựa nhất quán trên $S$; gọi $G'$ và $G''$
là các ảnh ngược của nó trên $S'$ và trên $S''=S'\times_SS'$, và gọi $p_1,p_2$
là hai phép chiếu từ $S''$ lên $S'$. Khi đó biểu đồ các ánh xạ tập hợp sau

\src{316}
\[
  \Gamma(G)\xrightarrow{\alpha^*}\Gamma(G')
  \underset{p_2^*}{\overset{p_1^*}{\rightrightarrows}}\Gamma(G'')
\]
là khớp (A, Định nghĩa 1.(a)).

Mặt khác, kết hợp (i), (ii) và Định nghĩa 1.1 ta được

\result{HỆ QUẢ}{2} Cho $G$ như trong Hệ quả 1. Khi đó có một tương ứng song ánh
giữa các bó con tựa nhất quán của $G$ và các bó con tựa nhất quán của $G'$ mà
ảnh ngược của chúng trên $S''$ theo hai phép chiếu $p_1,p_2$ cho cùng một bó
con của $G''$.

Dĩ nhiên có một phát biểu tương đương bằng ngôn ngữ các bó thương. Như đã thấy
ở A, Mục 4, (e), Định lý 1 phải được xem là một sự khái quát hóa ``Định lý 90''
của Hilbert, và nó bao hàm như những trường hợp riêng nhiều cách phát biểu bằng
đối đồng điều bậc 1. Để chứng minh, ta dễ dàng quy về trường hợp
$S=\operatorname{Spec}(A)$, $S'=\operatorname{Spec}(A')$; và đối với (i), ta
dễ dàng quy về việc chứng minh Hệ quả 1, nghĩa là tính khớp của biểu đồ
\[
  M=A\otimes_A M\longrightarrow A'\otimes_A M
  \rightrightarrows A'\otimes_AA'\otimes_A M
\]
đối với mọi $A$-môđun $M$, điều này suy ra từ bổ đề tổng quát hơn sau đây:

\result{BỔ ĐỀ}{1.1} Cho $A'$ là một $A$-đại số phẳng trung thành. Khi đó, với
mọi $A$-môđun $M$, phức được tăng cường bởi $M$
$C^*(A'/A,G_a)\otimes_A M$ (xem A, Mục 4, (e)) là một \emph{phân giải} của
$M$.

Chỉ cần chứng minh rằng phức tăng cường thu được từ phức trên bằng phép đổi cơ
sở từ $A$ sang $A'$ thỏa mãn các kết luận tương tự. Điều này dẫn đến việc kiểm
tra phát biểu sau khi thay $A$ bởi $A'$ và $A'$ bởi $A'\otimes_AA'$, do đó quy
ta về trường hợp có một đồng cấu $A$-đại số $A'\longrightarrow A$ (hay, bằng
ngôn ngữ hình học, trường hợp $S'$ trên $S$ có một tiết diện). Trong trường hợp
ấy, khẳng định suy ra từ các sự kiện tổng quát ở A, Mục 4, (a). Nhân tiện, ta
nêu hệ quả sau, khái quát hóa một phát biểu quen thuộc trong đối đồng điều
Galois (so sánh A, Mục 4, (e)):

\result{HỆ QUẢ}{} Nếu $A'$ phẳng trung thành trên $A$, thì
$H^0(A'/A,G_a)=A$ và
\[
  H^i(A'/A,G_a)=0\qquad\text{với $i\geq1$}.
\]

Để chứng minh phần (ii) của Định lý 1, ta tiến hành như đối với (i), quy về
trường hợp $S'$ trên $S$ có một tiết diện; khi đó khẳng định suy ra từ (i)
(xem A, Mục 1, (c)).

Hiển nhiên, có thể biến đổi tùy ý Định lý 1 và các hệ quả của nó bằng cách đưa
thêm nhiều cấu trúc khác nhau lên các bó (hoặc các hệ bó) tựa nhất quán đang
xét. Chẳng hạn, việc cho trên $S$ một bó tựa nhất quán các đại số giao hoán

\src{317}
``tương đương'' với việc cho trên $S'$ một bó cùng loại được trang bị dữ liệu
hạ xuống đối với $\alpha:S'\longrightarrow S$. Xét đến sự tương ứng hàm tử
giữa các bó tựa nhất quán như vậy trên $S$ và các tiền lược đồ affine trên $S$,
ta thu được khẳng định thứ hai của định lý sau:

\result{ĐỊNH LÝ}{2} Cho $\alpha:S'\longrightarrow S$ như trong Định lý 1. Khi
đó $\alpha$ là một \emph{cấu xạ hạ xuống} (nói chung không nghiêm ngặt)
(A, Định nghĩa 2.4), và là một \emph{cấu xạ hạ xuống nghiêm ngặt} đối với phạm
trù phân thớ gồm các lược đồ affine trên các tiền lược đồ (A, Định nghĩa 1.7).

Khẳng định thứ nhất của định lý có nghĩa như sau: cho $X,Y$ là hai tiền lược đồ
trên $S$; gọi $X',Y'$ là các ảnh ngược của chúng trên $S'$, và gọi $X'',Y''$
là các ảnh ngược của chúng trên $S''=S'\times_SS'$. Khi đó biểu đồ các ánh xạ
tự nhiên sau
\[
  \operatorname{Hom}_S(X,Y)\xrightarrow{\alpha^*}
  \operatorname{Hom}_{S'}(X',Y')
  \underset{p_2^*}{\overset{p_1^*}{\rightrightarrows}}
  \operatorname{Hom}_{S''}(X'',Y'')
\]
là khớp; nghĩa là $\alpha^*$ là một song ánh từ
$\operatorname{Hom}_S(X,Y)$ lên phần của $\operatorname{Hom}_{S'}(X',Y')$ gồm
các đồng cấu tương thích với các dữ liệu hạ xuống chính tắc trên $X',Y'$ (nghĩa
là những đồng cấu mà ảnh ngược theo hai phép chiếu từ $S''$ lên $S'$ bằng
nhau). Điều này dễ dàng suy ra từ Định lý 1, Hệ quả 1, khi hạn chế $Y$ là affine
trên $S$; trong trường hợp tổng quát, phải kết hợp Định lý 1 với kết quả sau:

\result{BỔ ĐỀ}{1.2} Cho $\alpha:S'\longrightarrow S$ là một cấu xạ phẳng trung
thành và quasi-compact. Khi đó $S$ đồng nhất với một \emph{không gian tôpô
thương} của $S'$; nghĩa là mọi phần $U$ của $S$ sao cho $\alpha^{-1}(U)$ mở thì
chính nó cũng mở.

Để hoàn tất Định lý 2, cần cho các tiêu chuẩn về tính hữu hiệu của một dữ liệu
hạ xuống trên một $S'$-tiền lược đồ $X'$ (khi không giả sử $X'$ affine trên
$S'$). Trước hết, lưu ý rằng \emph{một dữ liệu hạ xuống như vậy không nhất
thiết hữu hiệu}, ngay cả khi $S$ là phổ của một trường $k$, $S'$ là phổ của một
mở rộng bậc hai $k'$ của trường ấy, và $X'$ là một lược đồ đại số proper chiều 2 trên
$S'$ (như có thể thấy, theo Serre, bằng cách dùng mặt không xạ ảnh của Nagata).
\emph{Để một dữ liệu hạ xuống trên $X'/S'$ đối với
$\alpha:S'\longrightarrow S$ (phẳng trung thành và quasi-compact) là hữu hiệu,
điều kiện cần và đủ là $X'$ là hợp của các phần mở $X'_i$, affine trên $S'$, và
``ổn định'' dưới dữ liệu hạ xuống trên $X'$.} Điều này chắc chắn đúng (bất kể
$X'/S'$ và dữ liệu hạ xuống trên $X'$ là gì) nếu cấu xạ
$\alpha:S'\longrightarrow S$ là \emph{radicial} (nghĩa là đơn ánh và các mở
rộng trường thặng dư

\src{318}
cũng radicial). Cũng có thể chứng minh rằng điều ấy vẫn đúng nếu
$\alpha:S'\longrightarrow S$ là \emph{hữu hạn} và mọi phần hữu hạn của $X'$,
nằm trong một thớ của $X'$ trên $S$, đều nằm trong một phần mở của $X'$ affine
trên $S$ (đó là \emph{tiêu chuẩn Weil}). Đặc biệt, điều này đúng nếu $X'/S'$ là
quasi-xạ ảnh, và trong trường hợp ấy có thể chứng minh tiền lược đồ ``hạ xuống''
$X/S$ cũng quasi-xạ ảnh (và xạ ảnh nếu $X'/S'$ xạ ảnh). Tóm lại:

\result{ĐỊNH LÝ}{3} Cho $\alpha:S'\longrightarrow S$ là một cấu xạ phẳng trung
thành và quasi-compact của các tiền lược đồ. Nếu $\alpha$ là \emph{radicial},
thì nó là một \emph{cấu xạ hạ xuống nghiêm ngặt}. Nếu $\alpha$ hữu hạn, thì nó
là một cấu xạ hạ xuống nghiêm ngặt đối với phạm trù phân thớ gồm các tiền lược
đồ quasi-xạ ảnh (hoặc xạ ảnh) trên các tiền lược đồ.

\result{NHẬN XÉT}{} Tôi không biết trong khẳng định thứ hai ở trên, giả thiết
$\alpha$ là một cấu xạ \emph{hữu hạn} có thực sự cần thiết hay không; dù sao,
có thể kiểm tra một cách ``hình thức'' rằng giả thiết ấy có thể được thay bằng
giả thiết sau, có vẻ yếu hơn: \emph{với mọi điểm của $S$, tồn tại một lân cận
mở $U$, một $U'$ hữu hạn và phẳng trung thành trên $U$, cùng một $S$-cấu xạ từ
$U'$ vào $S'$.} Một trường hợp điển hình không nằm trong trường hợp trên là
$S=\operatorname{Spec}(A)$, $S'=\operatorname{Spec}(\widehat A)$, trong đó $A$
là một vành Noether địa phương và $\widehat A$ là phép đầy đủ hóa của nó; hoặc
trường hợp $S'$ quasi-hữu hạn trên $S$ (nghĩa là đẳng cấu địa phương với một
phần mở của một $S$-lược đồ hữu hạn) nhưng không hữu hạn. Trong hai trường hợp
ấy, người trình bày cũng không biết câu trả lời cho câu hỏi sau: cho $X$ là một
$S$-lược đồ sao cho $X'=X\times_SS'$ xạ ảnh trên $S'$; liệu $X$ có xạ ảnh trên
$S$ hay không?

\subsection*{2. Áp dụng vào việc hạ xuống một số tính chất của cấu xạ}
\addcontentsline{toc}{subsection}{2. Hạ xuống các tính chất của cấu xạ}

Cho $P$ là một lớp các cấu xạ của tiền lược đồ. Cho
$\alpha:S'\longrightarrow S$ là một cấu xạ của tiền lược đồ, cho
$f:X\longrightarrow Y$ là một cấu xạ của các $S$-tiền lược đồ, và cho
$f':X'\longrightarrow Y'$ là nghịch ảnh của $f$ qua $\alpha$. Khi đó ta có thể
hỏi liệu quan hệ ``$f'\in P$'' có kéo theo ``$f\in P$'' hay không. Trong nhiều
trường hợp quan trọng, câu trả lời là khẳng định nếu giả sử $\alpha$
\emph{phẳng trung thành và quasi-compact} (giả thiết sau đôi khi là thừa). Điều
này thấy được trực tiếp và không khó nếu $P$ là lớp các cấu xạ toàn ánh, tương
ứng radicial (các trường hợp này suy ra từ tính toàn ánh của $\alpha$), tương
ứng phẳng, tương ứng phẳng trung thành, tương ứng trơn (các trường hợp này suy
ra từ tính phẳng trung thành của $\alpha$), tương ứng kiểu hữu hạn. Dùng các
Định lý 1 và 2 cùng Bổ đề 1.2, ta cũng thấy điều tương tự đúng nếu $P$ là một
trong các lớp sau: các đẳng cấu, các phép nhúng mở, các phép nhúng đóng, các
phép nhúng (nếu $f$ kiểu hữu hạn và $Y$ Noether địa phương), các cấu xạ affine,
các cấu xạ hữu hạn, các cấu xạ quasi-hữu hạn, các cấu xạ mở, các cấu xạ đóng,
các đồng phôi, các cấu xạ tách,

\src{319}
các cấu xạ proper. Trường hợp quan trọng duy nhất chưa được giải quyết là các
cấu xạ xạ ảnh và quasi-xạ ảnh, đã được nêu trong nhận xét ở Mục 1.

\subsection*{3. Hạ xuống theo các cấu xạ hữu hạn phẳng trung thành}
\addcontentsline{toc}{subsection}{3. Hạ xuống theo các cấu xạ hữu hạn phẳng trung thành}

Cho $\alpha:S'\longrightarrow S$ là một cấu xạ hữu hạn tương ứng với một bó
đại số $\mathcal A'$ trên $S$, bó này, khi xem như bó môđun, là tự do địa phương
kiểu hữu hạn và là $\ne0$ ở mọi nơi; khi đó $\alpha$ là một cấu xạ phẳng trung
thành và quasi-compact, nên có thể áp dụng các kết quả trên. Việc cho một bó
tựa nhất quán $F'$ trên $S'$ tương đương với việc cho bó tựa nhất quán
$\alpha_*(F')$ trên $S$, được trang bị cấu trúc $\mathcal A'$-môđun của nó
(lưu ý rằng $\mathcal A'=\alpha_*(\mathcal O_{S'})$). Để đơn giản, bó này trên
$S$ cũng sẽ được ký hiệu là $F'$. Hai nghịch ảnh $p_i^*(F')$ của $F'$ trên
$S'\times_SS'$ cũng tương ứng với các bó tựa nhất quán gồm các
$(\mathcal A'\otimes_{\mathcal O_S}\mathcal A')$-môđun
\[
  F'\otimes_{\mathcal O_S}\mathcal A'
  \quad\text{và}\quad
  \mathcal A'\otimes_{\mathcal O_S}F'.
\]
Việc cho một $(S'\times_SS')$-đồng cấu từ bó thứ nhất vào bó thứ hai tương
đương với việc cho một đồng cấu của các
$(\mathcal A'\otimes_{\mathcal O_S}\mathcal A')$-môđun và, vì $\mathcal A'$
tự do địa phương, điều này tương đương với việc cho một đồng cấu của các
$(\mathcal A'\otimes\mathcal A')$-môđun:
\[
  \mathcal U=
  \operatorname{Hom}_{\mathcal O_S}(\mathcal A',\mathcal A')
  =\mathcal A'\otimes\check{\mathcal A}'
  \longrightarrow
  \operatorname{Hom}_{\mathcal O_S}(F',F'),
\]
nghĩa là, với mỗi tiết diện $\xi$ của $\mathcal U$ trên một tập mở $V$, cho một
đồng cấu $\mathcal O_S$-môđun
$T_\xi:F'|_V\longrightarrow F'|_V$ thỏa mãn các điều kiện
\begin{equation}
  T_{f\xi}(x)=fT_\xi(x),
  \qquad T_{\xi f}(x)=T_\xi(fx),
  \tag{3.1}
\end{equation}
trong đó $f,x$ lần lượt là các tiết diện của $\mathcal A',F'$ trên một tập mở
của $S$ nằm trong $V$. Khi ấy các điều kiện (i) và (ii) của một dữ liệu hạ xuống
(A, Mục 1, (c)) lần lượt được viết thành:
\begin{equation}
  T_{1_{\mathcal U}}x=x,
  \qquad\text{nghĩa là,}\qquad
  T_{1_{\mathcal U}}=\operatorname{id}_{F'},
  \tag{3.2}
\end{equation}
\begin{equation}
  T_{\xi\eta}=T_\xi T_\eta.
  \tag{3.3}
\end{equation}
Nói cách khác, \emph{một dữ liệu hạ xuống trên $F'$ tương đương với một biểu
diễn của bó $\mathcal O_S$-đại số
$\mathcal U=\operatorname{Hom}_{\mathcal O_S}(\mathcal A',\mathcal A')$ vào
bó $\mathcal O_S$-đại số
$\operatorname{Hom}_{\mathcal O_S}(F',F')$, thỏa mãn hai điều kiện tuyến tính
(3.1).} Nếu có một phép ghép cặp các bó tựa nhất quán trên $S'$:
\[
  F'_1\times F'_2\longrightarrow F'_3
\]
(có thể được hiểu như một phép ghép cặp các bó $\mathcal A'$-môđun trên $S$),

\src{320}
và các dữ liệu dán trên $F'_i$, được xác định bởi các đồng cấu
\[
  T_i:\mathcal U\longrightarrow
  \operatorname{Hom}_{\mathcal O_S}(F'_i,F'_i)
  \qquad (i=1,2,3),
\]
thì các dữ liệu ấy tương thích với phép ghép cặp đã cho, theo nghĩa hiển nhiên,
khi và chỉ khi điều kiện sau được thỏa mãn.

Với mọi tiết diện $\xi$ của $\mathcal U$ trên một tập mở, ký hiệu
\[
  \Delta\xi=\sum_i\xi'_i\otimes_{\mathcal A'}\xi''_i
\]
là tiết diện của $\mathcal U\otimes_{\mathcal A'}\mathcal U$ (trong đó
$\mathcal U$ được xem là một $\mathcal A'$-môđun qua cấu trúc môđun trái của
nó) được xác định bởi công thức
\[
  \xi\mathbin{\cdot}(fg)=\sum_i\xi'_i(f)\xi''_i(g)
\]
(trong đó $f$ và $g$ là hai tiết diện của $\mathcal A'$ trên một tập mở nhỏ
hơn), ta có công thức
\begin{equation}
  T^{(3)}_\xi(x\mathbin{\cdot}y)
  =\sum_iT^{(1)}_{\xi'_i}x\mathbin{\cdot}T^{(2)}_{\xi''_i}y
  \tag{3.4}
\end{equation}
với mọi cặp tiết diện $x,y$ của $\mathcal A'$ trên một tập mở nhỏ hơn. (Có
thể diễn đạt tính chất này bằng cách nói
rằng các đồng cấu $T^{(i)}$ tương thích với ánh xạ đường chéo của $\mathcal U$
đối với phép ghép cặp đã cho.) Đặc biệt, các công thức (3.1) đến (3.4) cho phép
ta diễn giải, theo ngôn ngữ các biểu diễn của đại số có ánh xạ đường chéo, các
dữ liệu hạ xuống trên một bó đại số tựa nhất quán trên $S'$, và do đó cũng diễn
giải (khi hạn chế vào các đại số giao hoán) các dữ liệu hạ xuống trên một
$S'$-lược đồ affine.

Từ đây ta chuyển sang một cách diễn giải tương tự cho các dữ liệu hạ xuống trên
một $S'$-tiền lược đồ $X'$ tùy ý: việc cho một $X'$ như vậy tương đương với
việc cho một tiền lược đồ $X'$ trên $S$, được trang bị một đồng cấu
$\mathcal O_S$-đại số
\[
  \mathcal A'\longrightarrow\mathcal O_{X'},
\]
và một dữ liệu hạ xuống trên $X'$ tương đương với việc cho một đồng cấu các bó
(tương thích với cấu xạ $h:X'\longrightarrow S'$)
\[
  \mathcal U\longrightarrow
  \operatorname{Hom}_{h^{-1}(\mathcal O_S)}
  (\mathcal O_{X'},\mathcal O_{X'}),
\]
thỏa mãn các điều kiện tương tự các điều kiện (3.1) đến (3.4) ở trên.

\textsc{Ví dụ 1 (``Weil'').} -- Giả sử $S'/S$ là một phủ étale Galois có nhóm
Galois $\Gamma$ (xem A, các Mục 3 và 4.d). Khi đó, một dữ liệu hạ xuống trên một
bó tựa nhất quán $F'$ trên $S'$ (tương ứng, trên một $S'$-tiền lược đồ $X'$)
tương đương với việc cho một biểu diễn của $\Gamma$ bằng các tự đẳng cấu của
$(S',F')$ (tương ứng, của $(S',X')$), tương thích với tác động của $\Gamma$
trên $S'$. Kết quả này

\src{321}
là ``hình thức'', nghĩa là được chứng minh theo ngôn ngữ phạm trù; nhưng theo
quan điểm của mục này, nó cũng suy ra từ cấu trúc tường minh của $\mathcal U$
(được trang bị cấu trúc vành, đồng cấu vành
$\mathcal A'\longrightarrow\mathcal U$ và ánh xạ đường chéo), cấu trúc này được
biết hoàn toàn nhờ kết quả sau: $\mathcal U$, khi xem như một
$\mathcal A'$-môđun trái, có một cơ sở gồm các tiết diện của $\mathcal U$
tương ứng với các phần tử của $\Gamma$.

\textsc{Ví dụ 2 (``Cartier'').} -- Cho $p$ là một số nguyên tố; giả sử
$p\mathcal O_S=0$ (nghĩa là $\mathcal O_S$ có đặc số $p$),
$(\mathcal A')^p\subset\mathcal O_S=\mathcal A$ (nghĩa là $S'/S$ radicial có
chiều cao 1), và bó đại số $\mathcal A'$ trên $\mathcal A$ có một $p$-cơ sở
tại địa phương, nghĩa là một họ tiết diện $(x_i)$ sao cho $\mathcal A'$ được sinh
như một đại số bởi các $x_i$, chỉ chịu các điều kiện $x_i^p=0$. Ta giả sử tập
các chỉ số $i$ là hữu hạn, có lực lượng $n$. Cho $\mathfrak g$ là bó các
$\mathcal A$-đạo hàm của $\mathcal A'$; đó là một bó $\mathcal A'$-môđun tự do
địa phương hạng $n$, và hơn nữa là một bó $p$-đại số Lie trên $\mathcal A$
(nhưng không trên $\mathcal A'$), thỏa mãn điều kiện
\begin{equation}
  [X,fY]=X(f)Y+f[X,Y].
  \tag{3.5}
\end{equation}

\textsc{Bổ đề.} -- $\mathcal U=
\operatorname{Hom}_{\mathcal O_S}(\mathcal A',\mathcal A')$ được sinh, như một
$\mathcal O_S$-đại số được trang bị một đồng cấu đại số
$\mathcal A'\longrightarrow\mathcal U$, bởi $\mathcal A'$-môđun con trái
$\mathfrak g$, với các quan hệ bổ sung
\begin{equation}
  \left\{
  \begin{aligned}
    Xf-fX&=X(f),\\
    XY-YX&=[X,Y],\\
    X^p&=X^{(p)}.
  \end{aligned}
  \right.
  \tag{3.6}
\end{equation}

Từ bổ đề trên suy ra rằng một dữ liệu hạ xuống trên bó tựa nhất quán $F'$ trên
$S'$ tương đương với việc cho, với mọi $X\in\mathfrak g$, một tự đồng cấu
$\mathcal O_S$-môđun $\overline X$ của $F'$ thỏa mãn các điều kiện
\begin{align}
  \overline{fX}&=f\overline X, \tag{3.7}\\
  \overline X(fx)&=X(f)x+f\overline X(x), \tag{3.8}\\
  \overline{[X,Y]}&=[\overline X,\overline Y], \tag{3.9}\\
  \overline{X^{(p)}}&=\overline X^p. \tag{3.10}
\end{align}
(Đó là điều có thể gọi là một liên thông tuyến tính trên $F'$, không có độ
cong và tương thích với lũy thừa bậc $p$.) Khái niệm dữ liệu hạ xuống trên một
$S'$-tiền lược đồ $X'$ cũng được làm tường minh theo cách ấy; ở đây quan hệ
(3.4) được thay bằng điều kiện rằng các $\overline X$ là các đạo hàm của
$\mathcal O_{X'}$. Vì cấu xạ $S'\longrightarrow S$ là radicial, Định lý 3 bảo
đảm rằng mọi dữ liệu hạ xuống như vậy đều

\src{322}
hữu hiệu và do đó xác định một $S$-tiền lược đồ $X$.

Ta lưu ý rằng không cần đặt bất kỳ giả thiết nào về tính phẳng, tính không kỳ
dị hay tính hữu hạn đối với $F'$, tương ứng $X'$.

\subsection*{4. Ứng dụng vào các tiêu chuẩn về tính hữu tỉ}
\addcontentsline{toc}{subsection}{4. Ứng dụng vào các tiêu chuẩn về tính hữu tỉ}

Cho $X$ là một $S$-tiền lược đồ sao cho ảnh trực tiếp của $\mathcal O_X$ trên
$S$ là $\mathcal O_S$; tính chất này vẫn đúng sau mọi mở rộng phẳng
$S'\longrightarrow S$ của cơ sở $S$. Nếu $F$ là một bó khả nghịch (nghĩa là
tự do địa phương hạng 1) trên $X$, thì các tự đẳng cấu của $F$, được đồng nhất
với các tiết diện khả nghịch của $\mathcal O_X$, tương ứng song ánh với các
tiết diện khả nghịch của $\mathcal O_S$. Bây giờ cho $s$ là một tiết diện của
$X$ trên $S$; theo cách nói gợi hình, ta sẽ gọi là một tiết diện của $F$ trên
$s$ một tiết diện của bó khả nghịch $s^*(F)$ trên $S$. Từ những điều trên suy ra
rằng, nếu $F_i$ ($i=1,2$) là hai bó khả nghịch trên $X$, mỗi bó được trang bị
một tiết diện trên $s$, và nếu $F_1$ và $F_2$ đẳng cấu, thì tồn tại duy nhất
một đẳng cấu từ $F_1$ lên $F_2$ tương thích với các tiết diện đang xét (nghĩa
là biến tiết diện thứ nhất thành tiết diện thứ hai). Mặt khác, không phụ thuộc
vào tiết diện $s$, ta quy ước coi hai bó khả nghịch $F_1$ và $F_2$ trên $X$ là
tương đương nếu mọi điểm của $S$ đều có một lân cận mở $U$ sao cho các hạn chế
của $F_1$ và $F_2$ lên $X|U$ đẳng cấu. Khi đó mọi bó khả nghịch $F$ trên $X$
tương đương với một bó khả nghịch $F_1$ được trang bị một tiết diện được đánh
dấu trên $s$ (lấy $F_1=F\,s^*(F)^{-1}$), và $F_1$ được xác định đến đẳng cấu.
Nói cách khác, việc phân loại các bó khả nghịch trên $X$ đến tương đương cũng
chính là việc phân loại, đến đẳng cấu, các bó khả nghịch được trang bị một tiết
diện được đánh dấu.

Vì các tính chất này vẫn đúng sau một mở rộng phẳng
$\alpha:S'\longrightarrow S$ của cơ sở (thay tiết diện $s$ bằng ảnh ngược $s'$
của nó qua $\alpha$), nên từ Định lý 1 ta suy ra:

\emph{Cho tiền lược đồ $X/S$ như trên và giả sử nó có một tiết diện $s$; cho
$\alpha:S'\longrightarrow S$ là một cấu xạ phẳng trung thành và quasi-compact;
cho $F'$ là một bó khả nghịch trên $X'=X\times_SS'$. Điều kiện cần và đủ để
$F'$ tương đương với ảnh ngược trên $X'$ của một bó khả nghịch $F$ trên $X$ là
hai ảnh ngược $p_1^*(F')$ và $p_2^*(F')$ của nó trên
$X'\times_XX'=X\times_S(S'\times_SS')$ tương đương. Khi đó $F$ được xác định
đến tương đương.} (Khi ấy ta nói $F'$ hữu tỉ trên $S$.)

Áp dụng nguyên lý này khi $\alpha:S'\longrightarrow S$ như trong Ví dụ 1 hoặc
Ví dụ 2 của mục trước, ta thu được các tiêu chuẩn hữu tỉ của Weil hoặc Cartier.
(Lưu ý rằng các tác giả này chỉ xét trường hợp $S$ và $S'$

\src{323}
là phổ của các trường; do đó, hiển nhiên $S$ là phổ của một vành địa phương,
và quan hệ tương đương được đưa vào ở trên chính là quan hệ đẳng cấu.) Trong
trường hợp thứ nhất, $F'$ hữu tỉ trên $S$ khi và chỉ khi các biến đổi của nó
dưới $\Gamma$ tương đương với $F'$. Để phát biểu tiêu chuẩn hữu tỉ trong trường
hợp thứ hai, nói chung ta xét cấu xạ đường chéo
\[
  X'\longrightarrow X''=X'\times_XX'
\]
của $X'/X$, bó iđêan tương ứng $\mathcal I$ trên $X'\times_XX'$, và bó
$\mathcal I/\mathcal I^2$, được đồng nhất với ảnh ngược
$\Omega^1_{X'/X}$ của nó trên $X'$. Vì các hạn chế lên đường chéo của
$F''_i=p_i^*(F')$ ($i=1,2$) đẳng cấu (do cả hai đều đẳng cấu với $F'$), nghĩa
là $F''_1(F''_2)^{-1}=F''$ có hạn chế tầm thường trên đường chéo, nên hạn chế
của $F''$ lên $(X'',\mathcal O_{X''}/\mathcal I^2)$ được cho, đến đẳng cấu,
bởi một phần tử xác định duy nhất $\xi$ của
\[
  H^1(X'',\mathcal I/\mathcal I^2)
  =H^1(X',\Omega^1_{X'/X}).
\]
Hơn nữa, trong trường hợp này ta có
\[
  \Omega^1_{X'/X}
  =\Omega^1_{S'/S}\otimes_{\mathcal O_S}\mathcal O_X,
\]
và do đó, nếu $\Omega^1_{S'/S}$ tự do địa phương trên $S$ (như trong trường
hợp Cartier), thì $\xi$ xác định một tiết diện của
\[
  R^1f'_*(\mathcal O_{X'})\otimes\Omega^1_{S'/S}
\]
trên $S'$ (gọi là lớp Atiyah--Cartier của bó khả nghịch $F'$ trên $X'/S$).
Sự triệt tiêu của lớp này là điều kiện cần và đủ để các ảnh ngược của $F'$ qua
hai phép chiếu từ
\[
  (X'',\mathcal O_{X''}/\mathcal I^2)
  =X\times_S(S'',\mathcal O_{S''}/\mathcal J^2)
\]
lên $X'$ tương đương (trong đó $\mathcal J$ là bó iđêan trên
$S''=S'\times_SS'$ được xác định bởi cấu xạ đường chéo
$S'\longrightarrow S'\times_SS'$). Vì vậy, sự triệt tiêu này hiển nhiên là
điều kiện cần để các ảnh ngược của $F'$ trên chính
$X''=X\times_SS''$ tương đương, và do đó cũng là điều kiện cần để $F'$ tương
đương với ảnh ngược của một bó khả nghịch $F$ trên $X$. Lớp Atiyah--Cartier
cũng có thể được hiểu là cản trở đối với sự tồn tại, cục bộ trên $S'$, của một
liên thông trên $F'$ tương đối với các đạo hàm của $X'/X$; hơn nữa, một liên
thông như vậy được xác định khi đã biết các đạo hàm của $F'$ tương ứng với các
mở rộng tự nhiên lên $X'$ của các đạo hàm của $S'/S$. Từ điều này và các khai
triển trong mục trước, ta dễ dàng suy ra rằng trong trường hợp Ví dụ 2 của mục
ấy, khi $X/S$ có một tiết diện, sự triệt tiêu của lớp Atiyah--Cartier cũng đủ
để $F'$ hữu tỉ trên $S$.

\subsection*{5. Ứng dụng vào phép thu hẹp cơ sở trong một lược đồ abel}
\addcontentsline{toc}{subsection}{5. Thu hẹp cơ sở trong một lược đồ abel}

Cho $S$ là một tiền lược đồ. Một lược đồ abel trên $S$ là một lược đồ $X$ trơn
và proper

\src{324}
trên $S$ sao cho các thớ tại những điểm $x\in S$ là các lược đồ của những đa
tạp abel trên $\kappa(x)$. Giả sử $S$ Noether và chính quy (nghĩa là các vành
địa phương của nó chính quy). Khi đó, dùng định lý liên thông của Murre [4]
(ít nhất trong trường hợp ``đồng đặc số'', là trường hợp mà định lý được dẫn
hiện đã được chứng minh), ta có thể chứng minh rằng mọi tiết diện hữu tỉ của
$X$ trên $S$ đều xác định khắp nơi (nghĩa là thực sự là một tiết diện); điều
này tổng quát hóa một định lý cổ điển của Weil. Tổng quát hơn, từ đó suy ra
rằng nếu $X'$ là một lược đồ trơn trên $S$, thì mọi $S$-cấu xạ hữu tỉ từ $X'$
vào $X$ đều xác định khắp nơi. Ta thu được kết quả sau, tổng quát hóa một kết
quả của Igusa--Lang: Cho $S$ Noether và chính quy, và ký hiệu $K$ là
vành các hàm hữu tỉ của nó (một tích trực tiếp các trường); cho $X$ là một lược
đồ abel trên $K$. Nếu $X$ đẳng cấu với một $K$-lược đồ có dạng
$X_0\times_S\operatorname{Spec}(K)$, trong đó $X_0$ là một lược đồ abel trên
$S$, thì $X_0$ được xác định đến một đẳng cấu duy nhất.

Dùng kết quả duy nhất ở trên, ta thấy rằng bài toán thu hẹp cơ sở trong $X$ là
cục bộ trên $S$ (và do đó chỉ cần biết cách thực hiện phép thu hẹp trên các
$\operatorname{Spec}(\mathcal O_x)$, với $x\in S$). Tương tự, nếu
$\alpha:S'\longrightarrow S$ là một cấu xạ trơn kiểu hữu hạn, nếu $K'$ là
vành các hàm hữu tỉ của $S'$, và nếu $X\otimes_KK'$ có dạng
$X'_0\times_{S'}\operatorname{Spec}(K')$, thì $X'_0$ được trang bị một dữ liệu
hạ xuống chính tắc tương đối với $\alpha$. Theo Định lý 3, ta suy ra:

\result{MỆNH ĐỀ}{5.1} Cho $S$ là một tiền lược đồ Noether, chính quy và bất khả
quy, có trường hàm hữu tỉ $K$; cho $K'$ là một mở rộng hữu hạn của $K$, không
ram hóa trên $S$; cho $S'$ là chuẩn hóa của $S$ trong $K'$ (do đó là một phủ
étale của $S$); và cho $X$ là một lược đồ abel trên $K$ sao cho
$X\otimes_KK'$ có dạng $X'_0\times_{S'}\operatorname{Spec}(K')$, trong đó
$X'_0$ là một lược đồ abel xạ ảnh trên $S'$. Khi đó $X$ có dạng
$X_0\times_S\operatorname{Spec}(K)$, trong đó $X_0$ là một lược đồ abel xạ ảnh
trên $S$.

\result{NHẬN XÉT}{} Người trình bày không biết liệu có thể thay giả thiết
$S'\longrightarrow S$ là một phủ étale toàn ánh (cho phép dùng Định lý 3) bằng
giả thiết rằng đó là một cấu xạ trơn, toàn ánh và kiểu hữu hạn (ngay cả khi giả
sử nó là étale), hoặc liệu mệnh đề có còn đúng mà không cần giả sử $X'_0$ xạ
ảnh trên $S'$ hay không (một điều kiện có lẽ tự động được thỏa mãn).

\subsection*{6. Ứng dụng vào các tiêu chuẩn tầm thường cục bộ và đẳng tầm}
\addcontentsline{toc}{subsection}{6. Các tiêu chuẩn tầm thường cục bộ và đẳng tầm}

Cho $S$ là một tiền lược đồ, $G$ là một ``tiền lược đồ nhóm'' trên $S$, và $P$
là một tiền lược đồ trên $S$ mà trên đó ``$G$ tác động'' (bên phải). Ta nói
$P$ là thuần nhất chính hình thức dưới $G$ nếu cấu xạ quen thuộc
\[
  G\times_SP\longrightarrow P\times_SP
\]

\src{325}
cảm sinh bởi tác động của $G$ lên $P$ là một đẳng cấu. Từ đây trở đi, ta giả
sử $G$ phẳng trên $S$ (và do đó phẳng trung thành trên $S$), và dành tên
phân thớ chính thuần nhất dưới $G$ cho một $P$ thuần nhất chính hình thức
mà phẳng trung thành và quasi-compact trên $S$. Điều này hiển nhiên tương đương
với việc tồn tại một mở rộng phẳng trung thành và quasi-compact
$S'\longrightarrow S$ của cơ sở $S$ sao cho phân thớ thuần nhất chính hình thức
\[
  P'=P\times_SS'
\]
dưới $G'=G\times_SS'$ là tầm thường, nghĩa là đẳng cấu với $G'$ (nghĩa là có
một tiết diện); đặc biệt, có thể lấy $S'=P$. Ngoài ra, lưu ý rằng nếu $S$ địa
phương Noether, thì giả thiết $P$ phẳng trung thành tương đương với
giả thiết
\[
  \widehat P_s=P\times_S\operatorname{Spec}(\widehat{\mathcal O}_s)
\]
phẳng trung thành trên $\widehat{\mathcal O}_s$ với mọi $s\in S$ (trong đó
$\widehat{\mathcal O}_s$ ký hiệu phép đầy đủ hóa của vành địa phương
$\mathcal O_s$), vì $\widehat{\mathcal O}_s$ phẳng trung thành trên
$\mathcal O_s$. Hơn nữa, nếu $P$ có kiểu hữu hạn trên $S$ địa phương Noether,
thì tập các điểm $s$ thỏa điều kiện trên là dựng được. Vì vậy, nếu $S$ là một
``tiền lược đồ Jacobson'' (chẳng hạn một lược đồ kiểu hữu hạn trên một trường,
hoặc tổng quát hơn là một vành Jacobson), thì chỉ cần kiểm tra điều kiện
đang xét tại các điểm đóng của $S$. Điều này quy ta về trường hợp cơ sở là phổ
của một vành địa phương đầy đủ $A$. Khi $S=\operatorname{Spec}(A)$ (với $A$ là
một vành địa phương Noether đầy đủ) và $P$ có kiểu hữu hạn trên $S$, tính phẳng
trung thành của $P/S$ cũng tương đương với sự tồn tại của một $S'$ hữu hạn và
phẳng trên $S$ sao cho $P'$ tầm thường; nếu thêm $G$ trơn trên $S$, thì có thể
giả sử $S'$ étale trên $S$. Do đó, nếu trường thặng dư của $A$ còn đóng đại số
(``trường hợp hình học''), thì $P$ phẳng trung thành trên $A$ khi và chỉ khi nó
tầm thường. Như vậy, nếu $S$ là một tiền lược đồ đại số trên một trường đóng
đại số và $G$ trơn, kiểu hữu hạn trên $S$, thì điều kiện phẳng trung thành trên
$S$ tương đương với điều kiện tầm thường giải tích (SLF) của Serre ([6],
tr.~1--12).

Ta có thể đưa vào những loại điều kiện khác, mạnh hơn, trên $P$ và mang bản chất
của một ``tính tầm thường cục bộ''. Cụ thể, ta nói $P$ là đẳng tầm
(tương ứng, đẳng tầm nghiêm ngặt) nếu với mọi $s\in S$, tồn tại một lân
cận mở $U$ của $s$ và một cấu xạ hữu hạn, phẳng trung thành (tương ứng, một phủ
étale toàn ánh) $U'\longrightarrow U$ sao cho $P'=P\times_SU'$ tầm thường.
(Thuật ngữ của chúng ta khác với Serre [6], người gọi đẳng tầm cục bộ
điều mà chúng ta gọi là đẳng tầm nghiêm ngặt.) Tính đẳng tầm nghiêm ngặt
đặc biệt hữu ích khi $G$ trơn trên $S$, nhưng ngược lại là một khái niệm không
thích hợp trong những trường hợp khác.

Nếu $G$ affine trên $S$, thì theo đoạn 2, mọi phân thớ chính thuần nhất $P$ dưới
$G$ đều affine; bởi vậy, nhờ Định lý 2, ta có thể ``hạ xuống''

\src{326}
các phân thớ như vậy dọc theo những cấu xạ phẳng trung thành và quasi-compact.
Đặc biệt, lấy $G=\operatorname{Gl}(n)_S$, được xác định bởi điều kiện rằng hàm
tử trên các $S$-tiền lược đồ
\[
  S'\longmapsto\operatorname{Hom}_S(S',G)
\]
(nhận giá trị trong phạm trù các nhóm) được đồng nhất với hàm tử
\[
  \operatorname{Gl}(n)(S')
  =\operatorname{Gl}\bigl(n,H^0(S',\mathcal O_{S'})\bigr)
\]
trong A, đoạn 4.e. Dùng các sự kiện (i) mọi phân thớ chính thuần nhất dưới $G$
(tương ứng, mọi bó tự do địa phương hạng $n$ trên $S$) trở thành đẳng cấu với
đối tượng ``tầm thường'' $G$ (tương ứng, $\mathcal O_S^n$) sau một mở rộng
phẳng trung thành và quasi-compact thích hợp của $S$; (ii) các đối tượng thuộc
loại đang xét (các phân thớ chính thuần nhất dưới $G$, tương ứng các bó tự do
địa phương hạng $n$) có thể được hạ xuống dọc theo những cấu xạ như vậy; và
cuối cùng, (iii) nhóm tự đẳng cấu của phân thớ tầm thường trên một $S'/S$ đẳng
cấu một cách hàm tử với nhóm tự đẳng cấu của bó tự do địa phương tầm thường
hạng $n$ trên $S'$, ta suy ra một cách ``hình thức'' rằng việc cho trên $S$
(hoặc trên một $S'/S$) một phân thớ chính thuần nhất có nhóm $G$ ``cũng chính
là'' việc cho tại đó một bó tự do địa phương hạng $n$. (Chính xác hơn, ta có
một tương đương giữa các phạm trù phân thớ.) Đặc biệt, ta suy ra:

\result{MỆNH ĐỀ}{6.1} Mọi phân thớ chính thuần nhất có nhóm
$\operatorname{Gl}(n)_S$ đều tầm thường cục bộ.

Bằng những lập luận quen thuộc, ta thu được cùng kết quả cho các nhóm cấu trúc
như $\operatorname{Sl}(n)_S$, $\operatorname{Sp}(n)_S$ và các tích của những
nhóm như vậy. Ta cũng suy ra rằng, nếu $F$ là một nhóm con đóng của
$G=\operatorname{Gl}(n)_S$, phẳng trên $S$, sao cho thương $G/F$ tồn tại, và
nếu $G$ là một phân thớ chính thuần nhất đẳng tầm (tương ứng, đẳng tầm nghiêm
ngặt) trên $G/F$, với nhóm $F\times_S(G/F)$, thì mọi phân thớ chính thuần nhất
có nhóm $F$ đều đẳng tầm (tương ứng, đẳng tầm nghiêm ngặt). Điều này áp dụng
cho tất cả các ``nhóm tuyến tính'' trên $S$ đã được dùng cho đến đây, và đặc
biệt khi $G=S\times_k\Gamma$, trong đó $S$ là một tiền lược đồ trên trường $k$
và $\Gamma$ là một nhóm tuyến tính theo nghĩa cổ điển, đặc biệt là một nhóm
trơn, trên $k$. Do đó, đối với các nhóm như vậy, điều này giải quyết một câu hỏi
của Serre (loc. cit.).

Ta cũng lưu ý rằng, đối với phần lớn các nhóm trơn trên $S$ đã biết, dù tuyến
tính hay không, và trong mọi trường hợp đối với các nhóm có dạng
$S\times_k\Gamma$ như trên, có thể chứng minh rằng mọi phân thớ chính thuần nhất
đẳng tầm đều đẳng tầm nghiêm ngặt. Điều này đặc biệt giải quyết một câu hỏi khác
của Serre (loc. cit., tr.~1--14), vì một phân thớ chính thuần nhất thu được bằng
phép hạ xuống kiểu Cartier (xem đoạn 3, Ví dụ 2) hiển nhiên là đẳng tầm.

\result{NHẬN XÉT}{} Một trong những khó khăn cốt yếu của các vấn đề này (ngoài
vấn đề về sự tồn tại của các lược đồ thương $G/F$) là thiếu những tiêu chuẩn về
tính hữu hiệu cho một dữ liệu hạ xuống dọc theo một cấu xạ phẳng trung thành
không hữu hạn.

\src{327}
\begin{center}
  \textbf{TÀI LIỆU THAM KHẢO}
\end{center}

\begin{description}
  \item[{[1]}] Dieudonné (J.) và Grothendieck (Alexander). --- \emph{Éléments
    de géométrie algébrique}, sắp xuất bản trong Publications mathématiques de
    l'Institut des Hautes Études Scientifiques.
  \item[{[2]}] Grauert (H.) und Remmert (R.). --- \emph{Komplexe Räume}, Math.
    Annalen, tập 136, 1958, tr.~245--318.
  \item[{[3]}] Grothendieck (Alexander). --- \emph{Géométrie formelle et
    géométrie algébrique}. Séminaire Bourbaki, tập 11, 1958/59, số 182.
  \item[{[4]}] Murre (J. P.). --- \emph{On a connectedness theorem for a
    birational transformation at a simple point}, Amer. J. Math., tập 80, 1958,
    tr.~3--15.
  \item[{[5]}] Serre (Jean-Pierre). --- \emph{Géométrie algébrique et géométrie
    analytique}, Ann. Inst. Fourier Grenoble, tập 6, 1955--56, tr.~1--42.
  \item[{[6]}] Serre (Jean-Pierre). --- \emph{Espaces fibrés algébriques},
    Séminaire Chevalley, tập 2, 1958: Anneaux de Chow et applications, số 1.
\end{description}

\par

\clearpage
\chapter*{Đính chính cho Bài trình bày số 190}
\addcontentsline{toc}{chapter}{Đính chính cho Bài trình bày số 190}
\setcounter{footnote}{0}
\sid{e190}
\markboth{Đính chính 190}{Đính chính 190}
\src{300}

\noindent\textbf{3. Kỹ thuật hạ xuống và các định lý tồn tại trong hình học
đại số [TDTE I]. I: Đại cương. Hạ xuống theo các cấu xạ phẳng trung thành.}

\begin{center}
[Hội thảo Bourbaki, tập 12, 1959/60, số 190, 29 tr.]
\end{center}

\noindent\textbf{Trang 1, dòng 10.} --- Nay xem ra nói rằng kỹ thuật hạ xuống
«là nền tảng của phần lớn các định lý tồn tại trong hình học đại số» là quá
mức. Điều này đúng phần lớn đối với các kỹ thuật không xạ ảnh được trình bày
trong hai bài đầu của loạt TDTE I đến VI, nhưng không đúng đối với các kỹ
thuật xạ ảnh (TDTE IV, V, VI).

\noindent\textbf{Trang 5, dòng 16.} --- Không cần giả sử rằng $\alpha$ là một
cấu xạ $\mathcal F$-hạ xuống.

\noindent\textbf{Trang 20, nhận xét.} --- Một cấu xạ tựa hữu hạn, étale và
toàn ánh $S'\longrightarrow S$, hoặc một cấu xạ
$\operatorname{Spec}(\widehat A)\longrightarrow\operatorname{Spec}(A)$, không
nhất thiết luôn là một cấu xạ hạ xuống nghiêm ngặt, ngay cả khi $A$ là vành
địa phương của một đường cong đại số trên một trường đóng đại số $k$ và
$S=\operatorname{Spec}(A)$. Do đó có thể tìm được một cấu xạ riêng và trơn
$f:X\longrightarrow S$, làm cho $X$ thành
một phân thớ chính trên $S$ dưới tác động của một đường cong elliptic $E$,
sao cho $f':X'\longrightarrow S'$ là xạ ảnh nhưng $f$ không xạ ảnh. Vì vậy,
đây đồng thời là một ví dụ về phân thớ chính thuần nhất không đẳng tầm dưới
tác động của một lược đồ abel.

\noindent\textbf{Trang 26, dòng 9.} --- Đọc «Chow--Iang» thay cho «Igusa--Lang».

Về các chi tiết khác nhau của lý thuyết hạ xuống, xem SGA VI, VII và VIII.

\par

\unit{Bài trình bày 195 -- Kỹ thuật hạ xuống II}
\sid{195}
\markboth{Thuyết trình 195}{Thuyết trình 195}
\src{369}
\seminar{Năm thứ 12, 1959/60, số 195\\Tháng 2 năm 1960}
\paperhead{KỸ THUẬT HẠ XUỐNG VÀ CÁC ĐỊNH LÝ TỒN TẠI TRONG HÌNH HỌC ĐẠI SỐ\\
II. ĐỊNH LÝ TỒN TẠI TRONG LÝ THUYẾT HÌNH THỨC VỀ MÔĐUN}%
  {do Alexander GROTHENDIECK}

\fsection{A}{Các hàm tử biểu diễn được và pro-biểu diễn được}

\subsection*{1. Các hàm tử biểu diễn được}
\addcontentsline{toc}{subsection}{1. Các hàm tử biểu diễn được}

Cho $\mathcal C$ là một phạm trù. Với mọi $X\in\mathcal C$, ký hiệu $h_X$ là
hàm tử phản biến từ $\mathcal C$ đến phạm trù $(\operatorname{Ens})$ của các
tập hợp,
\[
  h_X:\mathcal C\longrightarrow(\operatorname{Ens}),
\]
được xác định bởi công thức
\[
  h_X(Y)=\operatorname{Hom}(Y,X).
\]
Một cấu xạ $X\longrightarrow X'$ trong $\mathcal C$ hiển nhiên cảm sinh một
đồng cấu hàm tử $h_X\longrightarrow h_{X'}$: $h_X$ là một hàm tử hiệp biến
theo $X$; nói cách khác, ta đã xác định một hàm tử hiệp biến chính tắc
\[
  h:\mathcal C\longrightarrow
  \operatorname{Hom}\bigl(\mathcal C^o,(\operatorname{Ens})\bigr)
\]
từ $\mathcal C$ đến phạm trù các hàm tử hiệp biến từ phạm trù đối ngẫu
$\mathcal C^o$ của $\mathcal C$ đến phạm trù các tập hợp. Ta nhắc lại:

\result{MỆNH ĐỀ}{1.1} Hàm tử $h$ này \emph{đầy đủ và trung thành}; nói cách
khác, với mọi cặp đối tượng $X,X'$ của $\mathcal C$, ánh xạ tự nhiên
\[
  \operatorname{Hom}(X,X')\longrightarrow
  \operatorname{Hom}(h_X,h_{X'})
\]
là song ánh.

Đặc biệt, nếu một hàm tử
$F\in\operatorname{Hom}(\mathcal C^o,(\operatorname{Ens}))$ đẳng cấu với một
hàm tử dạng $h_X$, thì $X$ được xác định sai khác một đẳng cấu duy nhất.
Trong trường hợp đó ta nói hàm tử $F$ là \emph{biểu diễn được}. Mệnh đề trên
có nghĩa rằng hàm tử chính tắc $h$ xác lập một tương đương giữa phạm trù
$\mathcal C$ và phạm trù con đầy đủ của
$\operatorname{Hom}(\mathcal C^o,(\operatorname{Ens}))$ gồm các hàm tử biểu
diễn được. Sự kiện này là nền tảng của khái niệm «nghiệm của một bài toán phổ
quát»: một bài toán như vậy luôn quy về việc xét xem một hàm tử đã cho
(phản biến như ở đây, hoặc hiệp biến trong trường hợp đối ngẫu) từ
$\mathcal C$ đến $(\operatorname{Ens})$ có biểu diễn được hay không.

\src{370}
Ngoài ra, theo chính định nghĩa của tích trong một phạm trù [1], hàm tử
$h:X\longmapsto h_X$ giao hoán với các tích mỗi khi chúng tồn tại (và tổng
quát hơn, với các giới hạn xạ ảnh hữu hạn hoặc vô hạn; đặc biệt với các tích
thớ, với phép lập «hạt nhân» [2], v.v., mỗi khi chúng tồn tại): ta có một
đẳng cấu hàm tử
\[
  h_{X\times X'}\xrightarrow{\sim}h_X\times h_{X'}
\]
mỗi khi $X\times X'$ tồn tại; nghĩa là ta có các song ánh hàm tử theo $Y$
\[
  h_{X\times X'}(Y)\xrightarrow{\sim}h_X(Y)\times h_{X'}(Y).
\]
Đặc biệt, việc cho một cấu xạ
\[
  X\times X'\longrightarrow X''
\]
trong $\mathcal C$ (tức một «luật hợp thành» trong $\mathcal C$ giữa
$X,X',X''$) tương đương với việc cho một cấu xạ
\[
  h_{X\times X'}=h_X\times h_{X'}\longrightarrow h_{X''},
\]
tức là, với mọi $Y\in\mathcal C$, cho một luật hợp thành các tập hợp
\[
  h_X(Y)\times h_{X'}(Y)\longrightarrow h_{X''}(Y)
\]
sao cho, với mọi cấu xạ $Y\longrightarrow Y'$ trong $\mathcal C$, hệ các ánh
xạ giữa các tập hợp
\[
  h_{X^{(i)}}(Y)\longleftarrow h_{X^{(i)}}(Y')
  \qquad(\text{với }i=0,1,2)
\]
là một đồng cấu đối với hai luật hợp thành tương ứng với $Y$ và $Y'$. Do đó,
khái niệm cấu trúc «$\mathcal C$-nhóm», «$\mathcal C$-vành», v.v. trên một
đối tượng $X$ của $\mathcal C$ được diễn đạt thuận tiện nhất (cả trong lý
thuyết lẫn thực hành) bằng cách nói rằng, với mọi $Y\in\mathcal C$, tập
$h_X(Y)$ được trang bị một luật nhóm (tương ứng, luật vành, v.v.) theo nghĩa
thông thường, và các ánh xạ
$h_X(Y)\longleftarrow h_X(Y')$ tương ứng với các cấu xạ
$Y\longrightarrow Y'$ là các đồng cấu nhóm (tương ứng, đồng cấu vành, v.v.).
Chẳng hạn, đây là cách trực quan và thuận tiện nhất để xác định các nhóm cổ
điển khác nhau $G_a$, $G_m$, $\operatorname{Gl}(n)$, v.v. trên một tiền lược
đồ cơ sở tùy ý $S$ và viết các quan hệ cổ điển giữa chúng; để trang bị cho
lược đồ afin $V(\mathcal F)$ trên $S$, được xác định bởi một bó tựa kết hợp
$\mathcal F$, một cấu trúc «phân thớ vectơ»; và để xác định rồi nghiên cứu
các đa tạp cờ liên hệ khác nhau (các grassmannian, các phân thớ xạ ảnh),
v.v. Nguyên lý chung chỉ đơn giản là đồng nhất, thông qua hàm tử chính tắc
$h$, các đối tượng của $\mathcal C$ với một số hàm tử phản biến: các hàm tử
biểu diễn được,

\src{371}
từ $\mathcal C$ đến phạm trù các tập hợp.

Thủ tục quen thuộc đảo chiều các mũi tên, cần thiết chẳng hạn trong trường hợp
các lược đồ afin để chuyển từ ngôn ngữ hình học sang ngôn ngữ đại số giao
hoán, dẫn đến việc đối ngẫu hóa các xét đoán trên và, đặc biệt, cũng đưa vào
các hàm tử hiệp biến biểu diễn được
$\mathcal C\longrightarrow(\operatorname{Ens})$, tức là các hàm tử dạng
\[
  Y\longmapsto\operatorname{Hom}(X,Y)=h'_X(Y).
\]

\subsection*{2. Các hàm tử pro-biểu diễn được, các pro-đối tượng}
\addcontentsline{toc}{subsection}{2. Các hàm tử pro-biểu diễn được, các pro-đối tượng}

Cho $\widetilde X=(X_i)_{i\in I}$ là một hệ xạ ảnh các đối tượng của
$\mathcal C$; hệ này tương ứng với một hàm tử hiệp biến
\[
  h'_{\widetilde X}=\varinjlim_i h'_{X_i},
\]
hay tường minh hơn,
\[
  h'_{\widetilde X}(Y)=\varinjlim_i h'_{X_i}(Y)
  =\varinjlim_i\operatorname{Hom}(X_i,Y),
\]
từ $\mathcal C$ đến $(\operatorname{Ens})$. Một hàm tử từ $\mathcal C$ đến
$(\operatorname{Ens})$ đẳng cấu với một hàm tử loại này, với $I$ lọc, được
gọi là \emph{pro-biểu diễn được}. Theo mục trước, đó chính là các hàm tử đẳng
cấu với giới hạn quy nạp lọc của các hàm tử biểu diễn được. Cho
$\widetilde X'=(X'_j)_{j\in J}$ là một hệ xạ ảnh lọc thứ hai trong
$\mathcal C$ (được xây dựng trên một tập chỉ số tiền thứ tự lọc thứ hai
$J$); dễ kiểm tra rằng có một song ánh chính tắc (tổng quát hóa Mệnh đề 1.1)
\[
  \operatorname{Hom}(h'_{\widetilde X'},h'_{\widetilde X})
  =\varprojlim_j\varinjlim_i\operatorname{Hom}(X_i,X'_j).
\]
Điều này dẫn ta đến phạm trù $\operatorname{Pro}(\mathcal C)$ của các
pro-đối tượng của $\mathcal C$; các đối tượng của nó là các hệ xạ ảnh các đối
tượng của $\mathcal C$ (trên các tập chỉ số tiền thứ tự lọc thay đổi), và nếu
$\widetilde X=(X_i)_{i\in I}$ cùng
$\widetilde X'=(X'_j)_{j\in J}$ là hai đối tượng thuộc lớp này, ta đặt
\[
  \operatorname{ProHom}(\widetilde X,\widetilde X')
  =\varprojlim_j\varinjlim_i\operatorname{Hom}(X_i,X'_j),
\]
với phép hợp thành hiển nhiên của các pro-đồng cấu. Theo chính cách xây dựng,
$\widetilde X\longmapsto h'_{\widetilde X}$ có thể được xem như một hàm tử
phản biến theo $\widetilde X$, xác lập một tương đương giữa phạm trù đối ngẫu
của phạm trù $\operatorname{Pro}(\mathcal C)$ các pro-đối tượng của
$\mathcal C$ và phạm trù các hàm tử hiệp biến pro-biểu diễn được từ
$\mathcal C$ đến $(\operatorname{Ens})$. Dĩ nhiên, một đối tượng $X$ của
$\mathcal C$ xác định chính tắc một pro-đối tượng, cũng được ký hiệu là $X$,

\src{372}
sao cho $\mathcal C$ tương đương với một phạm trù con đầy đủ của
$\operatorname{Pro}(\mathcal C)$. Khi đó, nếu
$\widetilde X=(X_i)_{i\in I}$ là một pro-đối tượng tùy ý của $\mathcal C$, ta
có (với sự đồng nhất trên)
\[
  \widetilde X=\varprojlim_i X_i,
\]
trong đó giới hạn xạ ảnh được lấy trong $\operatorname{Pro}(\mathcal C)$
(vì $h'_{\widetilde X}=\varinjlim_i h'_{X_i}$). Cần lưu ý rằng, ngay cả khi
giới hạn xạ ảnh của các $X_i$ tồn tại trong $\mathcal C$, nói chung nó không
đẳng cấu với giới hạn xạ ảnh $\widetilde X$ trong
$\operatorname{Pro}(\mathcal C)$, như đã hiển nhiên khi $\mathcal C$ là phạm
trù các tập hợp. Theo định nghĩa,
$\varprojlim_{\mathcal C}X_i=L$ được xác định bởi điều kiện rằng hàm tử theo
$Y\in\mathcal C$
\[
  \varprojlim_i\operatorname{Hom}_{\mathcal C}(Y,X_i)
  =\operatorname{Hom}_{\operatorname{Pro}(\mathcal C)}(Y,\widetilde X)
\]
có giá trị trong $(\operatorname{Ens})$ được biểu diễn bởi $L$, tức là đẳng
cấu với $\operatorname{Hom}_{\mathcal C}(Y,L)$; do đó,
$\varprojlim_{\mathcal C}X_i$ đã được xác định theo pro-đối tượng
$\widetilde X$, và chính xác hơn, phụ thuộc hàm tử vào pro-đối tượng
$\widetilde X$ mỗi khi nó được xác định; vì vậy không có trở ngại gì khi ký
hiệu nó là $\varprojlim_{\mathcal C}(\widetilde X)$. Nếu các giới hạn xạ ảnh
luôn tồn tại trong $\mathcal C$,
$\varprojlim_{\mathcal C}(\widetilde X)$ là một hàm tử từ
$\operatorname{Pro}(\mathcal C)$ đến $\mathcal C$, và có một đồng cấu hàm tử
chính tắc
$\varprojlim_{\mathcal C}(\widetilde X)\longrightarrow\widetilde X$.

Mọi hàm tử, mà để định ý ta giả sử là hiệp biến,
\[
  F:\mathcal C\longrightarrow\mathcal C'
\]
đều được mở rộng một cách hiển nhiên thành hàm tử
\[
  \operatorname{Pro}(F):\operatorname{Pro}(\mathcal C)
  \longrightarrow\operatorname{Pro}(\mathcal C'),
\]
nên nếu các giới hạn xạ ảnh luôn tồn tại trong $\mathcal C'$, thì $F$ cũng
xác định chính tắc một hàm tử hợp thành
\[
  \overline F:\operatorname{Pro}(\mathcal C)\longrightarrow\mathcal C',
  \qquad
  \overline F=\varprojlim_{\mathcal C'}\circ\operatorname{Pro}(F),
\]
biến $\widetilde X=(X_i)_{i\in I}$ thành
$\varprojlim_{\mathcal C'}F(X_i)$.

Một pro-đối tượng $\widetilde X$ được gọi là \emph{pro-đối tượng nghiêm ngặt}
nếu nó đẳng cấu với một pro-đối tượng $(X_i)_{i\in I}$ mà các cấu xạ chuyển
$X_i\longrightarrow X_j$ đều là toàn cấu; một hàm tử được xác định bởi một
đối tượng như vậy được gọi là \emph{pro-biểu diễn được nghiêm ngặt}. Ngoài
ra, có thể đòi hỏi $I$ là tập có thứ tự lọc và mọi toàn cấu
$X_i\longrightarrow X'$ tương đương với một toàn cấu

\src{373}
$X_i\longrightarrow X_j$ với một $j\in I$ thích hợp (được điều kiện này xác
định duy nhất). Dưới các điều kiện ấy, hệ xạ ảnh $(X_i)_{i\in I}$ được xác
định sai khác một đẳng cấu duy nhất (theo nghĩa thông thường của đẳng cấu các
hệ xạ ảnh). Suy ra rằng trong $\operatorname{Pro}(\mathcal C)$ luôn tồn tại
giới hạn xạ ảnh của một hệ xạ ảnh các pro-đối tượng nghiêm ngặt
$\widetilde X^{(\alpha)}$, và với các ký hiệu trên cho $F$, $\overline F$, ta
có
\[
  \overline F\left(\varprojlim_\alpha\widetilde X^{(\alpha)}\right)
  =\varprojlim_{\alpha,\mathcal C'}
  \overline F\left(\widetilde X^{(\alpha)}\right).
\]
Đặc biệt, nếu mọi pro-đối tượng của $\mathcal C$ đều nghiêm ngặt (xem mục
sau), thì hàm tử mở rộng $\overline F$ giao hoán với các giới hạn xạ ảnh.

\subsection*{3. Đặc trưng các hàm tử pro-biểu diễn được}
\addcontentsline{toc}{subsection}{3. Đặc trưng các hàm tử pro-biểu diễn được}

Cho $\mathcal C$, $\mathcal C'$ là hai phạm trù trong đó các giới hạn xạ ảnh
hữu hạn tồn tại (tức là các giới hạn tương ứng với các tập tiền thứ tự hữu
hạn, không nhất thiết lọc), hay tương đương, trong đó các tích hữu hạn và các
tích thớ hữu hạn tồn tại (điều này đặc biệt kéo theo sự tồn tại của một «đối
tượng cuối» $e$, sao cho $\operatorname{Hom}(X,e)$ chỉ gồm một phần tử với
mọi $X$). Cho $F$ là một hàm tử hiệp biến từ $\mathcal C$ đến
$\mathcal C'$. Các điều kiện sau tương đương:

\begin{enumerate}
\item[(i)] $F$ giao hoán với các giới hạn xạ ảnh hữu hạn;
\item[(ii)] $F$ giao hoán với các tích hữu hạn và các tích thớ hữu hạn;
\item[(iii)] $F$ giao hoán với các tích hữu hạn và, với mọi biểu đồ khớp
\[
  X\longrightarrow X'\rightrightarrows X''
\]
trong $\mathcal C$ ([3], A, Định nghĩa 2.1), biểu đồ được biến đổi bởi $F$
\[
  F(X)\longrightarrow F(X')\rightrightarrows F(X'')
\]
là khớp.
\end{enumerate}

Khi đó ta nói $F$ \emph{khớp trái}.

Từ đây trở đi, ta giả sử các giới hạn xạ ảnh hữu hạn tồn tại trong
$\mathcal C$. Từ các định nghĩa, lập tức suy ra rằng một hàm tử biểu diễn
được là khớp trái và, khi chuyển sang giới hạn, một hàm tử pro-biểu diễn được
là khớp trái.

Để có một mệnh đề đảo, xét
\[
  F:\mathcal C\longrightarrow(\operatorname{Ens})
\]
là một hàm tử hiệp biến, và cho $X\in\mathcal C$, $\xi\in F(X)$. Ta nói
$\xi$ (hay cặp \src{374}$(X,\xi)$) là \emph{tối tiểu} nếu, với mọi
$X'\in\mathcal C$, $\xi'\in F(X')$ và mọi đơn cấu nghiêm ngặt ([3], A, 2)
$u:X'\longrightarrow X$ sao cho $\xi=F(u)(\xi')$, thì $u$ là một đẳng cấu.
Mặt khác, ta nói một cặp $(X,\xi)$ \emph{trội} $(X'',\xi'')$
($\xi\in F(X)$, $\xi''\in F(X'')$) nếu tồn tại một cấu xạ
$v:X\longrightarrow X''$ sao cho $\xi''=F(v)(\xi)$; nếu $\xi$ tối tiểu và
$F$ khớp trái, thì cấu xạ $v$ là duy nhất; nếu $\xi''$ tối tiểu, thì $v$ là
toàn ánh. Dễ dàng suy ra mệnh đề sau:

\result{MỆNH ĐỀ}{3.1} Để $F$ pro-biểu diễn được nghiêm ngặt, điều kiện cần và
đủ là nó thỏa mãn hai điều kiện:
\begin{enumerate}
\item[(i)] $F$ khớp trái;
\item[(ii)] mọi cặp $(X,\xi)$, với $\xi\in F(X)$, đều bị trội bởi một cặp
tối tiểu.
\end{enumerate}

Điều kiện cuối là trống nếu mọi đối tượng của $\mathcal C$ đều là Artin (lấy
một đối tượng con $X'$ của $X$ tối tiểu trong số những đối tượng mà với chúng
tồn tại $\xi'\in F(X')$ sao cho $\xi$ là ảnh của $\xi'$). Do đó:

\result{HỆ QUẢ}{} Cho $\mathcal C$ là một phạm trù mà các đối tượng đều là
Artin và trong đó các giới hạn xạ ảnh hữu hạn tồn tại. Khi đó, các hàm tử
pro-biểu diễn được từ $\mathcal C$ đến $(\operatorname{Ens})$ chính là các
hàm tử khớp trái, và thậm chí đều pro-biểu diễn được nghiêm ngặt.

Sự kiện cuối này cũng có nghĩa rằng \emph{mọi pro-đối tượng của $\mathcal C$
đều nghiêm ngặt}.

\subsection*{4. Ví dụ: các nhóm kiểu Galois, các nhóm pro-đại số}
\addcontentsline{toc}{subsection}{4. Các nhóm kiểu Galois, các nhóm pro-đại số}

Nếu $\mathcal C$ là phạm trù các nhóm hữu hạn thông thường, thì
$\operatorname{Pro}(\mathcal C)$ tương đương với phạm trù các nhóm tôpô
compact hoàn toàn không liên thông. Các nhóm loại này và các tổng quát hóa
của chúng, thu được bằng cách thay các nhóm hữu hạn thông thường bởi các lược
đồ nhóm hữu hạn trên một tiền lược đồ cơ sở đã cho (chẳng hạn, các nhóm đại
số hữu hạn trên một trường $k$), sẽ đóng vai trò các nhóm cơ bản, đồng luân
và đồng điều tuyệt đối và tương đối của các tiền lược đồ. Trong tất cả các
ví dụ này, Hệ quả của Mệnh đề 3.1 được áp dụng, và chính thông qua hàm tử liên
hệ mà ta phải xác định $\pi_1$ [2]. Điều tương tự xảy ra khi bắt đầu từ phạm
trù các nhóm đại số hoặc tựa đại số trên một trường (hay tổng quát hơn, trên
một tiền lược đồ Noether): ta thu được các «nhóm pro-đại số» của Serre [4].

\subsection*{5. Ví dụ: «các đa tạp hình thức»}
\addcontentsline{toc}{subsection}{5. Ví dụ: các đa tạp hình thức}

Cho $\Lambda$ là một vành Noether và $\mathcal C$ là phạm trù các
$\Lambda$-đại số $A$ là những môđun Artin kiểu hữu hạn trên $\Lambda$ (hay
gọi tắt là

\src{375}
các $\Lambda$-đại số Artin). Các điều kiện của Hệ quả thuộc Mệnh đề 3.1 được
thỏa mãn. Phạm trù $\operatorname{Pro}(\mathcal C)$ ở đây tương đương với
phạm trù các đại số tôpô $O$ trên $\Lambda$ đẳng cấu với các giới hạn xạ ảnh
tôpô
\[
  O=\varprojlim_i O_i
\]
của các đại số $O_i\in\mathcal C$, tức là tôpô của chúng tuyến tính, tách và
đầy đủ, đồng thời với mọi iđêan mở $J_i$ của $O$, $O/J_i$ là một đại số Artin
trên $\Lambda$. Hàm tử
$\mathcal C\longrightarrow(\operatorname{Ens})$ liên hệ với một đại số như
vậy chính là
\begin{align*}
  F(A)=h'_O(A)
  &=\{\text{các đồng cấu liên tục của các $\Lambda$-đại số tôpô }
       O\longrightarrow A\}\\
  &=\varinjlim_i
    \operatorname{Hom}_{\Lambda\text{-đại số}}(O_i,A).
\end{align*}

Ngoài ra, lưu ý rằng phạm trù $\mathcal C$ về cơ bản là tích của các phạm trù
tương tự ứng với các phép đầy đủ hóa của những vành địa phương
$\Lambda_{\mathfrak m}$, với $\mathfrak m$ chạy qua các iđêan tối đại của
$\Lambda$; do đó, nếu muốn, có thể giới hạn vào trường hợp $\Lambda$ là một
vành địa phương đầy đủ thuộc lớp này. Trong mọi trường hợp, $O$ phân rã chính
tắc thành tích tôpô của các thành phần địa phương, tương ứng với các «điểm»
của lược đồ hình thức [2] được xác định bởi $O$. Một điểm thuộc lớp này được
xác định bởi một phần tử $\xi$ của một $F(K)$ nào đó, trong đó
$K\in\mathcal C$ là một trường (chẳng hạn trường thặng dư của thành phần địa
phương đang xét); hai cặp $(\xi,K)$ và $(\xi',K')$ xác định cùng một điểm khi
và chỉ khi chúng bị trội bởi cùng một $(\xi'',K'')$, hoặc cùng trội một
$(\xi''',K''')$. (Nếu các trường $\Lambda/\mathfrak m$ đóng đại số, chỉ cần
lấy hợp rời của các $F(\Lambda/\mathfrak m)$.)

Điều quan trọng là cho các điều kiện để thành phần địa phương $O_\xi$ của
$O$ ứng với $\xi\in F(K)$ là một vành Noether. Khi $\Lambda$ là một vành địa
phương đầy đủ (nhắc lại là Noether), điều này tương đương với việc $O_\xi$
đẳng cấu với một vành thương của vành chuỗi lũy thừa hình thức
$\Lambda[[t_1,\ldots,t_n]]$ trên $\Lambda$. Để nêu một tiêu chuẩn như vậy,
với mỗi vành $A$ ta đưa vào $A$-đại số $I_A$ các «số đối ngẫu» của $A$ bằng
công thức
\[
  I_A=A[t]/t^2A[t].
\]
Cho $\varepsilon:I_A\longrightarrow A$ là đồng cấu tăng cường; đồng cấu này
xác định (nếu $A\in\mathcal C$) một ánh xạ
\[
  F(\varepsilon):F(I_A)\longrightarrow F(A),
\]

\src{376}
và, bằng cách sử dụng tính khớp trái của $F$, ta xác định một cách nội tại
một cấu trúc $K$-môđun trên phần
\[
  F(I_A,\xi)=F(\varepsilon)^{-1}(\xi)
\]
của $F(I_A)$ gồm các $\xi'\in F(I_A)$ «rút gọn thành $\xi$»; từ dạng tường
minh của $F$ theo $O$, ta kiểm tra được rằng $A$-môđun này đồng nhất với
\[
  \operatorname{Hom}_{\Lambda}(\mathfrak m_\xi/\mathfrak m_\xi^2,A),
\]
trong đó $\mathfrak m_\xi$ là hạt nhân của đồng cấu
$\xi:O\longrightarrow A$, tức là, nếu $A$ là một trường, iđêan tối đại của
thành phần địa phương $O_\xi$ của $O$. Ta lập tức suy ra mệnh đề sau:

\result{MỆNH ĐỀ}{5.1} Cho $\xi\in F(K)$, trong đó $K\in\mathcal C$ là một
trường. Để thành phần địa phương tương ứng $O_\xi$ của $O$ là một vành
Noether, điều kiện cần và đủ là tập $F(I_K,\xi)$ các phần tử của $F(I_K)$ rút
gọn thành $\xi$ là một không gian vectơ hữu hạn chiều trên $K$. Khi đó, ta có
một đẳng cấu chính tắc
\[
  F(I_K,\xi)=
  \operatorname{Hom}\!\left(
    \mathfrak m_\xi/\mathfrak m_\xi^2+\mathfrak n_\xi O_\xi,K
  \right),
\]
trong đó $\mathfrak n_\xi$ là iđêan tối đại của $\Lambda$, tức hạt nhân của
đồng cấu $\Lambda\longrightarrow K$.
Đặc biệt, chiều của $K$-không gian vectơ $F(I_K,\xi)$ bằng chiều của không
gian vectơ $\mathfrak m_\xi/\mathfrak m_\xi^2$ trên trường
$O_\xi/\mathfrak m_\xi=K(\xi)$.

Giả sử $O_\xi$ là Noether và, để đơn giản ký hiệu, $\Lambda$ là địa phương
đầy đủ và $O=O_\xi$. Ta nói $O$ \emph{trơn trên $\Lambda$} nếu $O$ là một
đại số hữu hạn và étale trên đại số đầy đủ hóa của phép địa phương hóa
$\Lambda[t_1,\ldots,t_n]$ tại một iđêan tối đại của vành này, iđêan ấy cảm
sinh iđêan tối đại của $\Lambda$; khi mở rộng thặng dư của $O$ trên
$\Lambda$ là tầm thường (chẳng hạn, nếu trường thặng dư của $\Lambda$ đóng
đại số), điều này tương đương với việc chính $O$ đẳng cấu với một đại số
chuỗi lũy thừa hình thức loại đó.
Cuối cùng, nếu không còn giả sử $O$ nhất thiết Noether, ta cũng nói $O$ trơn
trên $\Lambda$ nếu $O$ đẳng cấu với một giới hạn xạ ảnh tôpô của các
$\Lambda$-đại số thương Noether và trơn trên $\Lambda$ theo nghĩa trên. Điều
này lập tức tổng quát hóa cho trường hợp không giả sử $\Lambda$ và $O$ địa
phương. Khi đó ta có mệnh đề sau:

\result{MỆNH ĐỀ}{5.2} Để $O$ trơn trên $\Lambda$, điều kiện cần và đủ là hàm
tử liên hệ $F$ biến các toàn cấu thành các toàn cấu.

Điều này có nghĩa rằng với mọi đồng cấu toàn ánh
$A\longrightarrow A'$ trong $\mathcal C$, ánh xạ
$F(A)\longrightarrow F(A')$ là toàn ánh. Dĩ nhiên, chỉ cần kiểm tra điều kiện
này khi $A$ địa phương và, tiến hành từng bước, khi iđêan của $A$ là hạt nhân
của $A\longrightarrow A'$ bị triệt tiêu bởi iđêan tối đại của $A$. Trong
thực hành, điều này quy bài toán về việc kiểm tra rằng một cản trở nào đó,
gắn với các bất biến vô cùng bé

\src{377}
của tình huống sinh ra hàm tử $F$, bằng không; đây là một bài toán có bản
chất đối đồng điều.

Để kết thúc, ta nói đôi lời trong bối cảnh trên về các \emph{vành xác định}.
Vẫn cho $F$ là một hàm tử từ $\mathcal C$ đến $(\operatorname{Ens})$,
pro-biểu diễn được bởi một $\Lambda$-đại số tôpô $O$. Khi đó, với mọi
$A\in\mathcal C$ và mọi $\xi\in F(A)$, tồn tại một vành con nhỏ nhất $A'$ của
$A$ sao cho $\xi$ là ảnh của một phần tử $\xi'$ của $F(A')$ (khi đó phần tử
này được xác định duy nhất): thật vậy, chỉ cần diễn giải $\xi$ bằng một đồng
cấu từ $O$ đến $A$ rồi lấy $A'$ là ảnh của $O$ qua đồng cấu ấy. Khi đó ta
gọi $A'$ là vành xác định của đối tượng $\xi\in F(A)$. Nếu
$u:A\longrightarrow B$ là một đồng cấu đại số và $\eta=F(u)(\xi)$, thì vành
xác định của $\eta$ là ảnh qua $u$ của vành xác định của $\xi$. Khi bắt đầu
từ một hàm tử $F$ từ $\mathcal C$ đến $(\operatorname{Ens})$, sự tồn tại của
các vành xác định cùng các tính chất vừa nêu, với một vài dè dặt nhỏ, tương
đương với việc $F$ pro-biểu diễn được; nghĩa là, nói chung chúng còn lâu mới
tầm thường.

\fsection{B}{Hai định lý tồn tại}

Giữ các ký hiệu của A, mục 5, và xuất phát từ một hàm tử hiệp biến
\[
  F:\mathcal C\longrightarrow(\operatorname{Ens});
\]
ta tìm các tiêu chuẩn dễ sử dụng để $F$ pro-biểu diễn được, tức là có thể
diễn đạt bằng một $\Lambda$-đại số $O$ như trên. Theo Hệ quả của A, Mệnh đề
3.1, điều kiện cần và đủ là $F$ khớp trái. Ở tình trạng hiện nay của kỹ thuật
hạ xuống (xem các câu hỏi nêu trong [3], tr. 9), trong những trường hợp quan
trọng nhất không thể kiểm tra trực tiếp tiêu chuẩn này dưới dạng đó, và cần
các tiêu chuẩn có vẻ ít đòi hỏi hơn.

\result{ĐỊNH LÝ}{1} Để hàm tử $F$ pro-biểu diễn được, điều kiện cần và đủ là
nó thỏa mãn hai điều kiện sau:
\begin{enumerate}
\item[(i)] $F$ giao hoán với các tích hữu hạn;
\item[(ii)] với mọi đại số $A\in\mathcal C$ và mọi đồng cấu
$A\longrightarrow A'$ trong $\mathcal C$ sao cho biểu đồ
\[
  A\longrightarrow A'\rightrightarrows A'\otimes_A A'
\]
là khớp ([3], A, Định nghĩa 1.2), biểu đồ biến đổi
\[
  F(A)\longrightarrow F(A')\rightrightarrows F(A'\otimes_A A')
\]

\src{378}
là khớp.
\end{enumerate}
Hơn nữa, chỉ cần kiểm tra (ii) khi $A$ địa phương và thêm vào đó thuộc một
trong hai trường hợp sau:
\begin{enumerate}
\item[(a)] $A'$ là một môđun tự do trên $A$;
\item[(b)] môđun thương $A'/A$ là một $A$-môđun có độ dài 1.
\end{enumerate}

Chứng minh định lý này khá tinh tế và không thể phác thảo ở đây. Ta chỉ nêu
rằng nó chủ yếu dựa trên việc nghiên cứu các quan hệ tương đương (theo nghĩa
phạm trù) trên phổ của một đại số Artin; nghiên cứu này vẫn còn đặt ra nhiều
bài toán mà lời giải dường như thiết yếu cho sự phát triển tiếp theo của lý
thuyết.

Trong các ứng dụng, việc kiểm tra điều kiện (i) luôn tầm thường. Việc kiểm tra
(ii) tách thành hai phần: trường hợp $A'$ là một $A$-môđun tự do tương ứng
với kỹ thuật hạ xuống theo các cấu xạ phẳng ([3], các Định lý 1, 2 và 3) và
không gây khó khăn. Để xử lý trường hợp (b), ta sẽ dùng kết quả sau:

\result{ĐỊNH LÝ}{2} Cho $A$ là một vành địa phương Artin có iđêan tối đại
$\mathfrak m$, và $A'$ là một $A$-đại số chứa $A$, sao cho
$\mathfrak mA'\subset A$ và
\[
  A\longrightarrow A'\rightrightarrows A'\otimes_A A'
\]
là khớp (điều này đặc biệt xảy ra nếu $A'/A$ là một $A$-môđun có độ dài 1).
Cho $\mathcal F$ là phạm trù sợi ([3], A, Định nghĩa 1.1) của các bó tựa kết
hợp và phẳng trên các tiền lược đồ thay đổi. Khi đó cấu xạ
$\operatorname{Spec}(A')\longrightarrow\operatorname{Spec}(A)$ là một cấu
xạ $\mathcal F$-hạ xuống nghiêm ngặt ([3], A, Định nghĩa 1.7).

Nói cách khác, dữ liệu của một $A$-môđun phẳng $M$ hoàn toàn tương đương với
dữ liệu của một $A'$-môđun phẳng $M'$, được trang bị một
$A'\otimes_AA'$-đẳng cấu từ $M'\otimes_AA'$ lên $A'\otimes_AM'$ thỏa mãn
điều kiện bắc cầu thông thường của dữ liệu hạ xuống (loco citato).

Định lý 2 được chứng minh bằng cách trước hết chứng minh
\[
  H^i(A'/A,G_a)=0\qquad\text{với }i\geq 1
\]
([3], A, 4(e)), vì giả thiết $\mathfrak mA'\subset A$ cho phép dễ dàng quy về
trường hợp $A$ là một trường $k$ (cụ thể là $A/\mathfrak m$). Khi đó áp dụng
các tương đương được nêu trong [3], trang 16.

\fsection{C}{Ứng dụng vào một số trường hợp riêng}

\subsection*{1. Những nhận xét tổng quát về các hàm tử được biểu diễn bởi các tiền lược đồ}
\addcontentsline{toc}{subsection}{1. Những nhận xét tổng quát về các hàm tử được biểu diễn bởi các tiền lược đồ}

Cho $S$ là một tiền lược đồ Noether địa phương. Một tiền lược đồ $X$ trên
$S$ được gọi là \src{379}kiểu hữu hạn địa phương trên $S$ nếu, với mọi
$x\in X$ chiếu xuống $y\in Y$, tồn tại một lân cận afin của $y$, có vành $A$,
và một lân cận afin của $x$ nằm trên lân cận ấy, có vành $B$, sao cho $B$ là
một $A$-đại số kiểu hữu hạn. Có nhiều ví dụ quan trọng về các tiền lược đồ kiểu
hữu hạn địa phương trên $S$ nhưng không phải kiểu hữu hạn trên $S$, xuất hiện
như nghiệm của các bài toán phổ quát cổ điển; vì vậy cần có khả năng xem lược
đồ Picard của một đường cong như hợp của vô hạn thành phần liên thông (không
được nhầm chúng với thành phần liên thông của phần tử đơn vị, tức ``đa tạp
Picard''). Do đó đôi khi thuận tiện khi làm việc trong phạm trù $\mathcal C$
các tiền lược đồ kiểu hữu hạn địa phương trên $S$, để xét trong đó vấn đề biểu
diễn được của một hàm tử phản biến $F$. Mục tiêu chính của các bài trình bày
này là phát triển một kỹ thuật tổng quát cho phép nhận biết một hàm tử $F$ như
vậy có biểu diễn được hay không và nghiên cứu các tính chất của $S$-tiền lược
đồ $X$ tương ứng từ các tính chất của $F$. Nhân đây, trong nghiên cứu ấy xuất
hiện những ví dụ không bệnh lý về các tiền lược đồ trên $S$ không tách trên
$S$, đặc biệt dưới dạng các ``tiền lược đồ Picard'' của những $S$-lược đồ rất
tốt; vì thế không nên loại khỏi hình học đại số các tiền lược đồ không phải là
lược đồ.

Cho $X$ là một tiền lược đồ kiểu hữu hạn địa phương trên $S$, và đặt
\[
  F(Y)=\operatorname{Hom}_S(Y,X)
\]
là hàm tử phản biến liên kết. Ta có thể xét hạn chế $F_0$ của $F$ lên phạm trù
con $\mathcal C_0$ của $\mathcal C$ gồm các tiền lược đồ $Y$ trên $S$ vừa
Artin vừa hữu hạn trên $S$: khi $S=\operatorname{Spec}(\Lambda)$,
$\mathcal C_0$ chính là phạm trù đối ngẫu của phạm trù các $\Lambda$-đại số
Artin đã xét trong B. Nếu $Y=\operatorname{Spec}(A)$, với $A$ là một vành
Artin địa phương, thì $Y$ rút về một điểm duy nhất $y$ nằm trên một điểm đóng
$s$ của $S$, và một $S$-đồng cấu từ $Y$ đến $X$ (tức một phần tử của $F(Y)$)
được xác định bởi dữ liệu một điểm $x\in X$ nằm trên $s$ và một
$\mathcal O_{s,S}$-đồng cấu từ $\mathcal O_{x,X}$ đến $A$. Nếu tồn tại một
đồng cấu như vậy thì $x$ tất yếu là một điểm đóng của $X$ (vì trường thặng dư
của nó là đại số trên trường thặng dư của $s$). Điều này cho thấy hạn chế
$F_0$ của $F$ lên các ``$Y$-đại số Artin'' là pro-biểu diễn được, và được biểu
diễn bởi đại số tôpô trên $Y$ mà các thành phần địa phương là các đầy đủ hóa
$\widehat{\mathcal O}_{x,X}$ của các vành địa phương của $X$ tại những điểm
$x$ của $X$ vừa đóng vừa nằm trên các điểm đóng của $Y$. Vì thế, chỉ riêng
việc biết $F_0$ đã cung cấp thông tin chính xác về cấu trúc của $X$ (cụ thể là
cấu trúc các đầy đủ hóa của các vành địa phương tại các điểm vừa nêu). Cần lưu
ý rằng, ngay cả khi $S$

\src{380}
là phổ của một trường đóng đại số, chỉ bằng cách xét có hệ thống các ``đa
tạp'' $Y$ mà $\mathcal O_Y$ có thể chứa phần tử lũy linh, đặc biệt bằng cách
làm việc với phổ của các vành Artin địa phương, ta mới đạt được ``phát biểu
đúng'' của các bài toán phổ quát cổ điển và nắm bắt được phương diện ``vô cùng
bé'' của chúng.

Khi xuất phát từ một hàm tử $F$ cho trước và muốn quyết định nó có biểu diễn
được hay không, việc nghiên cứu hàm tử $F_0$ (nhờ các Định lý 1 và 2) sẽ cho
những chỉ dẫn gần như đầy đủ: hoặc, như thường xảy ra (chẳng hạn chỉ cần kiểm
tra bản chất của các tập $F(I_K,\xi)$ và cách chúng ứng xử theo hàm tử, xem A),
ta nhận thấy ngay $F_0$ không pro-biểu diễn được (điều này giải thích sự thất
bại của các cố gắng cho đến nay nhằm định nghĩa một cách ít nhiều tự nhiên
các đa tạp môđuli để phân loại các phân thớ vectơ hạng $>1$); hoặc ta kiểm tra
được $F_0$ biểu diễn được nhưng các không gian vectơ $F(I_K,\xi)$ không hữu
hạn chiều, khi đó phải bằng lòng với nghiệm ``hình thức''; hoặc cuối cùng,
$F_0$ được biểu diễn bởi một tích các vành địa phương đầy đủ Noether, điều này
cho những chỉ dấu rất mạnh rằng chính $F$ biểu diễn được và, kết hợp với các
tính chất tương tự nhưng mang tính toàn cục hơn mà chúng tôi dự định phát
triển sau này, chắc hẳn sẽ đủ để kéo theo điều đó. Cuối cùng, còn có những bài
toán hình học thú vị (xem các mục 4 và 5 dưới đây) trong đó ta chỉ có hàm tử
$F_0$ (không phát sinh từ một hàm tử ``toàn cục'' $F$ nào), và ta có lý do để
hài lòng nếu gắn được cho nó một ``lược đồ hình thức môđuli''.

Để kết thúc các nhận xét tổng quát này, hãy chỉ ra cách lý thuyết lược đồ giải
thích những hiện tượng có vẻ dị thường, chẳng hạn mặt Igusa $V$, có ``đa tạp
Picard'' $P$ rút về một điểm nhưng vẫn thỏa
\[
  H^1(V,\mathcal O_V)\ne 0 ;
\]
trong trường hợp này, $P$ là một nhóm ``thuần túy vô cùng bé'' không rút về
nhóm đơn vị, tức là nó được xác định bởi một đại số địa phương $O$ có hạng hữu
hạn trên trường cơ sở $k$ và được trang bị một ánh xạ đường chéo tương ứng với
cấu trúc nhân của $P$; nếu $\mathfrak m$ là iđêan cực đại của $O$, đối ngẫu của
$\mathfrak m/\mathfrak m^2$ đẳng cấu chính tắc với
$H^1(V,\mathcal O_V)$ (xem mục 3 dưới đây). Chỉ khi nhóm Picard là một nhóm
đại số theo nghĩa cổ điển, tức là trơn trên trường cơ sở $k$, thì chiều của
$H^1(V,\mathcal O_V)$ (luôn bằng chiều của $\mathfrak m/\mathfrak m^2$) mới
bằng chiều của nhóm Picard.

\subsection*{2. Các lược đồ $\underline{\operatorname{Hom}}_S(X,Y)$,
$\Pi_{X/S}Z$, $\underline{\operatorname{Aut}}_S(X)$, v.v.}
\addcontentsline{toc}{subsection}{2. Các lược đồ Hom, Pi và Aut}

Cho $X,Y$ là hai tiền lược đồ trên $S$; với mọi tiền lược đồ $T$ trên $S$,

\src{381}
đặt $X_T=X\times_ST$, $Y_T=Y\times_ST$, và xét tập
\[
  F(T)=\operatorname{Hom}_T(X_T,Y_T)
      =\operatorname{Hom}_S(X_T,Y)
      =\operatorname{Hom}_S(X\times_ST,Y)
\]
như một hàm tử phản biến theo $T$. Nếu hàm tử này biểu diễn được, ta ký hiệu
bởi $\underline{\operatorname{Hom}}_S(X,Y)$ tiền lược đồ trên $S$ biểu diễn
nó; do đó có một đẳng cấu hàm tử
\[
  \operatorname{Hom}_S
  \bigl(T,\underline{\operatorname{Hom}}_S(X,Y)\bigr)
  \xrightarrow{\sim}
  \operatorname{Hom}_S(T\times_SX,Y).
\]
Có nhiều biến thể của bài toán phổ quát này mà nghiệm quy về bài toán trên:
tiền lược đồ các $S$-tự đẳng cấu của một $S$-tiền lược đồ $X$ (đó sẽ là một
tiền lược đồ nhóm), tiền lược đồ các $S$-đồng cấu từ một $S$-tiền lược đồ nhóm
đến một tiền lược đồ nhóm khác (đó sẽ là một tiền lược đồ nhóm giao hoán nếu
lược đồ nhóm thứ hai giao hoán), v.v. Ta cũng có thể tổng quát hóa định nghĩa
của $\underline{\operatorname{Hom}}_S(X,Y)$ bằng cách xét một tiền lược đồ $Z$
trên tiền lược đồ $X$ trên $S$, cùng hàm tử
\[
  F(T)=\operatorname{Hom}_{X_T}(X_T,Z_T)
\]
(tập các ``tiết diện'' của $Z_T$ trên $X_T$); nếu hàm tử này biểu diễn được,
$S$-tiền lược đồ biểu diễn nó sẽ được ký hiệu $\Pi_{X/S}Z$; theo định nghĩa ta
có một đẳng cấu hàm tử
\[
  \operatorname{Hom}_S(T,\Pi_{X/S}Z)
  =\operatorname{Hom}_{X_T}(X_T,Z_T).
\]
Lấy $Z=Y\times_SX$, ta thu lại
$\underline{\operatorname{Hom}}_S(X,Y)$. Từ các định nghĩa này suy ra một bảng
công thức vừa hiển nhiên vừa hữu ích cho các tiền lược đồ mới đưa vào, mà ta
không trình bày ở đây (vì nó đúng trong mọi phạm trù có các tích và tích thớ).
Nghiêm túc hơn là câu hỏi về sự tồn tại của các lược đồ dạng
$\underline{\operatorname{Hom}}_S(X,Y)$. Trước hết, ta thấy rằng, với $X$ cố
định, $\underline{\operatorname{Hom}}_S(X,Y)$ chỉ có thể tồn tại với mọi $Y$
trên $S$ (hoặc $\Pi_{X/S}Z$ với mọi $Z$ trên $X$) nếu $X$ phẳng trên $S$.
Hơn nữa, chỉ hợp lý khi kỳ vọng có nghiệm với $Y$ đủ tổng quát nếu $X$ còn
riêng trên $S$. Mặt khác, dường như các điều kiện này đủ cho sự tồn tại của
$\underline{\operatorname{Hom}}_S(X,Y)$ và $\Pi_{X/S}Z$, miễn là áp đặt, khi
cần, một giả thiết kiểu ``tựa xạ ảnh'' lên $Y/S$, tương ứng lên $Z/X$; ít nhất
điều này có thể kiểm tra trong khá nhiều tình huống (chẳng hạn khi $Y$ afin
trên $S$; hoặc, bằng các phép dựng trực tiếp sơ cấp, khi $X$ hữu hạn trên $S$).
Sau đây là hệ quả của các Định lý 1 và 2:

\src{382}
\result{MỆNH ĐỀ}{2.1} Cho $\Lambda$ là một vành Noether, $X$ và $Y$ là hai
tiền lược đồ tùy ý trên $\Lambda$, và xét hàm tử
\[
  F(A)=\operatorname{Hom}_A(X_A,Y_A)
\]
trên phạm trù $\mathcal C_0$ các $\Lambda$-đại số Artin. Nếu $X$ phẳng trên
$\Lambda$, hàm tử này pro-biểu diễn được.

Hơn nữa, ta kiểm tra được rằng với mọi $A\in\mathcal C_0$ và $\xi\in F(A)$ có
một đẳng cấu chính tắc
\[
  F(I_A,\xi)=H^0\!\left(
    X_A,
    \underline{\operatorname{Hom}}_{\mathcal O_{X_A}}
    \bigl(\xi^*(\Omega^1_{Y_A/A}),\mathcal O_{X_A}\bigr)
  \right),
\]
trong đó $\Omega^1_{Y_A/A}$ là bó các $1$-vi phân Kähler của $Y_A$ đối với
$A$. Lấy $A$ là một trường rồi dùng A.5.1 và định lý hữu hạn của [2], ta được
hệ quả sau:

\result{HỆ QUẢ}{} Giả sử $X$ phẳng và riêng trên $S$, còn $Y$ kiểu hữu hạn địa
phương trên $S$. Khi đó $F$ pro-biểu diễn được và các thành phần địa phương của
đại số tôpô tương ứng trên $\Lambda$ là các vành Noether.

\result{NHẬN XÉT}{} Các bài toán được xét trong mục này, cùng nhiều bài toán
khác, theo truyền thống được nghiên cứu trong khuôn khổ hình học đại số cổ điển
bằng ``tọa độ Chow'' của các chu trình trong không gian xạ ảnh, nhờ đó có thể
xem các chu trình ấy như những điểm của các đa tạp xạ ảnh thích hợp. Thủ tục
này, và nói chung việc dùng tọa độ Chow, dường như không thể tránh khỏi sự bất
đầy đủ theo quan điểm lược đồ, vì nó tiêu hủy các phần tử lũy linh trong các đa
tạp tham số và đặc biệt không thích hợp cho việc nghiên cứu thỏa đáng các biến
phân vô cùng bé của chu trình (chưa kể tính không nội tại do phụ thuộc vào
không gian xạ ảnh). Đáng tiếc, ngôn ngữ tọa độ Chow từng là ngôn ngữ duy nhất
được nhiều nhà hình học đại số dùng để nghiên cứu các họ đa tạp hoặc các họ
chu trình, và điều đó dường như đã tạo ra một trở ngại nghiêm trọng cho việc
làm sáng tỏ những khái niệm này, bất chấp lợi ích kỹ thuật không thể phủ nhận
(có lẽ chỉ tạm thời) của chúng. Nếu muốn có cái tương tự các đa tạp Chow trong
lý thuyết lược đồ, ta đi đến bài toán phổ quát sau: cho $X$ là một tiền lược đồ
trên $S$; với mọi tiền lược đồ $T$ trên $S$, xét tập $F(T)$ các tiền lược đồ
con đóng của $X_T=X\times_ST$ phẳng trên $T$. Ta tìm cách biểu diễn hàm tử theo
$T$ này bằng một tiền lược đồ trên $S$. Tổng quát hơn, có thể cho một bó tựa
kết hợp $G$ trên $X$ và lấy $F(T)$ là tập các bó

\src{383}
thương của $G_T$ phẳng trên $T$. Dường như bài toán có nghiệm bằng một lược đồ
$C$ kiểu hữu hạn địa phương trên $S$ khi $X$ riêng trên tiền lược đồ Noether
địa phương $S$ và $G$ còn coherent. Dù sao, chỉ cần giả sử $S$ Noether địa
phương thì hạn chế của $F$ lên các ``$S$-đại số Artin'' là pro-biểu diễn được;
và nếu thêm rằng $X$ riêng trên $S$ và $G$ coherent thì các thành phần địa
phương của vành tôpô tương ứng $O$ là Noether. Dĩ nhiên, ngay cả sau khi chứng
minh sự tồn tại của ``lược đồ Chow'' của $X$ trên $S$, vẫn còn phải tìm một
phân tích thành các tập mở rời nhau $C_i$ (tương ứng với việc cố định các bất
biến liên tục như bậc và chiều của các chu trình đang biến thiên) kiểu hữu hạn
(chứ không chỉ kiểu hữu hạn địa phương) trên $S$, làm rõ quan hệ giữa lược đồ
này với các đa tạp Chow cổ điển, và xác định khi nào một $C_i$ thậm chí xạ ảnh,
hoặc ít nhất tựa xạ ảnh, trên $S$.

\subsection*{3. Các lược đồ Picard}
\addcontentsline{toc}{subsection}{3. Các lược đồ Picard}

Cho $f:X\longrightarrow S$ là một $S$-tiền lược đồ, xét bó nhân
$\mathcal O_X^*$ các đơn vị của bó cấu trúc của $X$, và nhóm
\[
  P(X/S)=H^0\bigl(S,R^1f_*(\mathcal O_X^*)\bigr),
\]
gọi là nhóm Picard tương đối của $X/S$. Như vậy, một phần tử của nhóm này được
xác định bằng cách cho một phủ mở $(U_i)$ của $S$ và một bó khả nghịch $L_i$
trên mỗi $f^{-1}(U_i)$, sao cho với mọi cặp $(i,j)$,
$L_i|f^{-1}(U_i\cap U_j)$ đẳng cấu với
$L_j|f^{-1}(U_i\cap U_j)$, ít nhất là địa phương trên $U_i\cap U_j$ (tức hai bó
này ``tương đương'' theo nghĩa của [3], B, 4). Nếu $X/S$ có một tiết diện thì
$P(X/S)$ không gì khác hơn tập các lớp bó khả nghịch trên $X/S$ modulo
``tương đương'' (loco citato). Bây giờ, với mọi $T$ trên $S$, đặt
\[
  F(T)=P(X_T/T) ;
\]
ta được một hàm tử hiệp biến theo $T$, có thể gọi là hàm tử Picard của $X/S$;
nếu hàm tử này biểu diễn được thì tiền lược đồ trên $S$ biểu diễn nó được gọi
là tiền lược đồ Picard của $X/S$ và ký hiệu $\underline P(X/S)$. Do đó có một
đẳng cấu hàm tử
\[
  \operatorname{Hom}_S\bigl(T,\underline P(X/S)\bigr)
  \xrightarrow{\sim}P(X_T/T).
\]
Phép lập tiền lược đồ Picard tương thích với mở rộng cơ sở; đặc biệt, các tiền
lược đồ Picard của các thớ của $X$ trên $S$ (là các tiền lược đồ trên các
trường thặng dư $K(s)$ của $s\in S$) chính là các thớ của
$\underline P(X/S)$. Dĩ nhiên, vì $P(X_T/T)=F(T)$ là một nhóm giao hoán, các
tiền lược đồ Picard là các tiền lược đồ nhóm. Cũng cần lưu ý rằng các Jacobian
tổng quát của Rosenlicht không gì khác hơn các thành phần liên thông của đơn
vị trong các lược đồ Picard của những đường cong đầy đủ có thể có kỳ dị; điều
này hẳn làm hiển nhiên phần lớn các tính chất của chúng (một khi đã chứng minh
sự tồn tại).

\result{NHẬN XÉT}{} Định nghĩa được chọn ở đây chỉ hợp lý khi mọi điểm của
$Y$ có một lân cận mở $U$ mà trên đó $X$ có một tiết diện. Trong trường hợp
tổng quát, phải sửa nhẹ định nghĩa hàm tử Picard để vẫn có thể thu được một
định lý tồn tại.

Ở đây, các điều kiện hợp lý cho sự tồn tại của một tiền lược đồ Picard là:
$X$ riêng và phẳng trên $S$, $f_*(\mathcal O_X)=\mathcal O_S$, và $X$ có địa
phương một tiết diện trên $S$. Điều kiện này xuất hiện tự nhiên khi áp dụng kỹ
thuật hạ xuống: để loại bỏ các tự đẳng cấu của một bó khả nghịch $L$ trên $X$,
ta trang bị cho nó một tiết diện được đánh dấu trên tiết diện $s$ ([3], B, 4).
Đặc biệt, ta thu được:

\src{384}
\result{MỆNH ĐỀ}{3.1} Giả sử $X$ phẳng trên
$S=\operatorname{Spec}(\Lambda)$, với $\Lambda$ Noether, và với mọi $T$ kiểu
hữu hạn trên $S$ ta có
\[
  f_{T*}(\mathcal O_{X_T})=\mathcal O_T
\]
(nếu $f$ riêng và có các thớ khả tách, hoặc nếu $S$ là phổ của một trường,
thì theo Künneth điều kiện cuối tương đương với
$f_*(\mathcal O_X)=\mathcal O_S$). Khi đó hàm tử Picard của $X/S$ trên phạm
trù các $\Lambda$-đại số Artin là pro-biểu diễn được.

Hơn nữa, ở đây ta có
\[
  F(I_A,\xi)=H^1(X_A,\mathcal O_{X_A}),
\]
đặc biệt:

\result{HỆ QUẢ}{} Nếu $X$ riêng trên $S$, thì các thành phần địa phương của
$\Lambda$-đại số tôpô $O$ tương ứng với hàm tử Picard là Noether.

\result{NHẬN XÉT}{} Ta có thể tổng quát hóa các định nghĩa và kết quả của mục
này cho việc phân loại các phân thớ chính trên $X$, với nhóm cấu trúc là một
lược đồ nhóm $G$ trên $S$ vừa afin, vừa phẳng trên $S$, vừa giao hoán. Khi $G$
không giao hoán, phân thớ nhóm liên hợp của một phân thớ chính (mà các tiết
diện là các tự đẳng cấu của phân thớ chính) không còn tầm thường, nên Mệnh đề
3.1 không còn đúng nguyên dạng. Tuy nhiên, ta có thể sửa bài toán phổ quát sao
cho nó vẫn có nghiệm (ít nhất, hiện tại, trong hình học hình thức). Quy tắc
vàng cần ghi nhớ, trong bối cảnh của \src{385}mục này và các mục sau, mỗi khi
ta tìm các ``lược đồ môđuli'' cho những lớp đối tượng chỉ được xác định sai
khác đẳng cấu, luôn là: loại bỏ các tự đẳng cấu khả dĩ của các đối tượng cần
phân loại bằng cách đưa vào, nếu cần, những cấu trúc phụ (các điểm hay phần tử
của tiết diện được đánh dấu, cố định các dạng vi phân, v.v.), được chọn đủ vô
hại để không làm thay đổi đáng kể bài toán ban đầu.

\subsection*{4. Môđuli hình thức của một đa tạp}
\addcontentsline{toc}{subsection}{4. Môđuli hình thức của một đa tạp}

Cho $\Lambda$ là một vành Noether địa phương có trường thặng dư $k$ (thông
thường nhất, $\Lambda$ bằng $k$, hoặc bằng một $p$-vành Cohen), và cho $X_0$
là một tiền lược đồ trên $k$. Với mọi $\Lambda$-đại số Artin địa phương $A$,
xét tập $F(A)$ các lớp, sai khác đẳng cấu, của các $A$-tiền lược đồ $X$ phẳng
trên $A$, được trang bị một đẳng cấu
\[
  (*)\qquad
  X\otimes_A\kappa(A)\xleftarrow{\sim}
  X_0\otimes_k\kappa(A),
\]
trong đó $\kappa(A)$ là trường thặng dư của $A$; dĩ nhiên, các đẳng cấu giữa
những $A$-tiền lược đồ phẳng như vậy phải tôn trọng đẳng cấu vừa nêu, vốn là
một phần của cấu trúc. Nếu $A$ là một $\Lambda$-đại số Artin không nhất thiết
địa phương, có các thành phần địa phương $A_i$, ta lấy $F(A)$ là tích các
$F(A_i)$. Như vậy, $F$ trở thành một hàm tử nhân tính theo $A$, có thể gọi là
hàm tử môđuli của $X_0$ (và $\Lambda$). Nếu hàm tử này biểu diễn được, nó tương
ứng với một $\Lambda$-đại số tôpô địa phương $O$ có trường thặng dư $k$, và
phổ hình thức của $O$ được gọi là lược đồ hình thức môđuli của $X_0$ (và
$\Lambda$); xem [2] để biết một số chi tiết về khái niệm này.

Ở đây, để áp dụng kỹ thuật hạ xuống, các tự đẳng cấu ``hữu hạn'' của $X_0$ là
vô hại, vì chúng không ảnh hưởng đến sự tồn tại của các tự đẳng cấu (theo
nghĩa đã nêu trên) của các $A$-tiền lược đồ $X$; khi $A$ không rút về một
trường, điều kiện cần và đủ để $X$ không có tự đẳng cấu $A$ không tầm thường
nào là
\[
  H^0(X_0,\Theta_{X_0/k})=0,
\]
trong đó $\Theta_{X_0/k}$ ký hiệu bó các $k$-đạo hàm (= bó tiếp tuyến) của
$X_0$. Mặt khác, ta dễ dàng thu được (ít nhất nếu $X_0$ trơn trên $k$)
\[
  F(I_A,\xi)=H^1(X_A,\Theta_{X_A/A}).
\]
Như thường lệ, suy ra:

\src{386}
\result{MỆNH ĐỀ}{4.1} Giả sử
$H^0(X_0,\Theta_{X_0/k})=0$. Khi đó tồn tại lược đồ hình thức môđuli của
$X_0$. Nếu thêm rằng $X_0$ riêng trên $k$, lược đồ hình thức môđuli là
Noether.

\result{NHẬN XÉT}{} 1\textsuperscript{o} Khi không giả sử $X_0$ trơn trên
$k$, $F(I_A,\xi)$ được đồng nhất với một $A$-môđun con của
\[
  \operatorname{Ext}^1_{\mathcal O_{P_A}}
  \bigl(P_A;\mathcal I_{X_A},\mathcal O_{X_A}\bigr),
\]
trong đó đặt $P_A=X_A\times_\Lambda X_A$, xem $\mathcal O_{X_A}$ như một bó
coherent trên $P_A$ thông qua cấu xạ đường chéo
$X_A\longrightarrow P_A$, và $\mathcal I_{X_A}$ ký hiệu bó iđêan coherent
trên $P_A$ được xác định bởi cấu xạ đường chéo. Chính xác hơn, một phép toàn
cục hóa đơn giản của lý thuyết Hochschild cho thấy nhóm
$\operatorname{Ext}^1$ trên đồng nhất với tập các lớp, sai khác đẳng cấu, của
các bó $I_A$-đại số phẳng $\mathcal O$ trên $X_A$, được trang bị một đẳng cấu
tăng cường
\[
  \mathcal O\otimes_{I_A}A\longrightarrow\mathcal O_{X_A}
\]
(nhắc lại rằng $I_A=A[t]/(t^2)$). Môđun con $F(I_A,\xi)$ là môđun tương ứng
với các bó đại số giao hoán. Vì vậy, trái với điều [2] gợi ra, các giả thiết
trơn không thiết yếu trong lý thuyết môđuli.

2\textsuperscript{o} Nhắc lại (loco citato) rằng, đặc biệt, mọi đường cong đại
số $X_0$ trơn và riêng trên $k$ đều có một lược đồ hình thức môđuli trơn trên
$\Lambda$, có chiều tương đối $3g-3$ nếu giống $g$ là $\geq2$, và bằng $g$ nếu
$g=0$ hoặc $1$. Hai trường hợp sau không còn trực tiếp thuộc Mệnh đề 4.1. Tuy
nhiên, trong trường hợp các đường cong elliptic ($g=1$), có thể quy chúng về
mệnh đề ấy bằng các xét đoán sau đây.

Dĩ nhiên, ta có thể biến đổi Mệnh đề 4.1 tùy ý bằng cách xét các hệ lược đồ
trên $k$ được trang bị nhiều cấu trúc khác nhau. Chẳng hạn, giả sử $X_0$ là
một lược đồ Abel trên $k$, có điểm gốc được đánh dấu (tức là xem $X_0$ như một
lược đồ nhóm trên $k$), và cho $F(A)$ là tập các lớp, sai khác đẳng cấu, của
các lược đồ Abel trên $A$ (tức các lược đồ nhóm riêng và trơn trên $A$) được
trang bị một đẳng cấu lược đồ Abel $(*)$. Ta kiểm tra được rằng việc áp đặt một
cấu trúc nhân (thậm chí chỉ một ``tiết diện đơn vị'') loại bỏ các tự đẳng cấu
vô cùng bé, và do đó tồn tại một lược đồ hình thức môđuli tương ứng với một
vành địa phương đầy đủ Noether $O$. Mặt khác, có thể chứng minh rằng nếu $X$
là một lược đồ riêng và trơn, có các thớ ``liên thông tuyệt đối'' trên một tiền
lược đồ Noether địa phương $S$, thì mọi cấu trúc nhân trên $X$ có một tiết
diện đơn vị đều tất yếu kết hợp và giao hoán (miễn là nó như vậy trên ít nhất
một thớ và $S$ liên thông), và hơn nữa được xác định duy nhất bởi việc biết

\src{387}
tiết diện đơn vị. Ngoài ra, giả sử $S$ là phổ của một vành Artin địa phương
$A$ có trường thặng dư $k$, $X$ riêng trên $A$ và được trang bị một tiết diện
$s$, và cuối cùng $X\otimes_Ak$ được trang bị cấu trúc lược đồ Abel trên $k$
mà phần tử trung hòa là điểm của $X\otimes_Ak$ tương ứng với $s$; một phép
tính cản trở đơn giản, kết hợp với một lập luận của Serre, cho phép chứng minh
rằng thực sự tồn tại trên $X$ một cấu trúc nhân nhận $s$ làm tiết diện đơn vị.
(Từ đó, dùng ``định lý tồn tại'' của [2] để chuyển sang trường hợp $A$ Noether
địa phương đầy đủ, rồi dùng kỹ thuật hạ xuống của [3] cho trường hợp tổng quát,
ta có thể chứng minh phát biểu tương tự với mọi $S$ Noether địa phương liên
thông.) Điều này chứng minh rằng hàm tử $F(A)$ đang xét ở đây đẳng cấu với hàm
tử tương tự được định nghĩa ở đầu mục này khi bỏ qua cấu trúc nhân trên $X_0$.
Đặc biệt, suy ra rằng nếu $\mathfrak m$ là iđêan cực đại của $O$ thì
$\mathfrak m/\mathfrak m^2$ đẳng cấu chính tắc với đối ngẫu của
$H^1(X_0,\Theta_{X_0/k})$, và do đó có chiều $n^2$, trong đó
$n=\dim X_0$. Sẽ rất thú vị nếu xác định được $O$ có trơn trên $\Lambda$ hay
không, tức có đẳng cấu với một đại số chuỗi hình thức theo $n^2$ biến trên
$\Lambda$ hay không. Mục 1, Mệnh đề 5.2, cho phép phát biểu tương đương bài
toán này như một bài toán tồn tại các lược đồ Abel rút về một lược đồ Abel cho
trước. Dù sao, các phương pháp siêu việt cho thấy câu trả lời là khẳng định nếu
$k$ có đặc số không. Trong đặc số $p\ne0$, hiển nhiên chỉ cần xét trường hợp
$\Lambda$ là vành vectơ Witt xây dựng trên một trường đóng đại số $k$. Có lẽ
đã đến lúc ``hàm tử Greenberg'' chứng tỏ giá trị của mình\ldots

\subsection*{5. Mở rộng các phủ}
\addcontentsline{toc}{subsection}{5. Mở rộng các phủ}

Cho $\mathfrak X$ là một tiền lược đồ hình thức Noether [2] và $U$ là một
phần mở của $\mathfrak X$, được xác định địa phương bởi sự ``không triệt tiêu''
của một tiết diện của $\mathcal O_{\mathfrak X}$ không là ước của không trên
$\mathfrak X$, và vì thế đủ lớn để mọi tiết diện của
$\mathcal O_{\mathfrak X}$ trên một tập mở $V$ triệt tiêu trên $U\cap V$ đều
triệt tiêu. Cho $\mathcal J$ là một ``iđêan xác định'' của $\mathfrak X$, và
đặt
\[
  X_n=(\mathfrak X,\mathcal O_{\mathfrak X}/\mathcal J^{n+1}),
\]
đây là một tiền lược đồ Noether thông thường. Nếu
$\mathfrak X',\mathfrak X''$ là hai phủ phẳng của $\mathfrak X$ (tức hai tiền
lược đồ trên $\mathfrak X$ được xác định bởi các bó đại số coherent và tự do
địa phương khi xem như các bó môđun), không phân nhánh trên $U$, thì ánh xạ
hiển nhiên
\[
  \operatorname{Hom}_{\mathfrak X}(\mathfrak X',\mathfrak X'')
  \longrightarrow
  \operatorname{Hom}_{X_0}(X'_0,X''_0)
\]
là đơn ánh; đặc biệt, một tự đẳng cấu của $\mathfrak X'$ cảm sinh đồng nhất
trên $X'_0$ chính là đồng nhất. Điều này cho phép áp dụng kỹ thuật hạ xuống vào
tình huống này. Đặc biệt, hãy xuất phát từ một phủ phẳng $X'_0$ của $X_0$,
không phân nhánh

\src{388}
trên $U_0$; ký hiệu $G(\mathfrak X)$ là tập các lớp, sai khác đẳng cấu (các
đẳng cấu cảm sinh đồng nhất trên $X'_0$), của các phủ phẳng $\mathfrak X'$ của
$\mathfrak X$ cảm sinh $X'_0$ trên $X_0$ (và do đó tất yếu không phân nhánh
trên $U$). Định nghĩa tương tự $G(V)$ với mọi phần mở $V$ của $\mathfrak X$,
và tổng quát hơn $G(\mathfrak Y)$ với mọi tiền lược đồ hình thức
$\mathfrak Y$ trên $\mathfrak X$. Khi ấy, các kết quả của [2] và [3] trước hết
kéo theo những kết quả sau:

\noindent a. Khi $V$ chạy qua các tập mở của $\mathfrak X$, các $G(V)$ lập
thành một bó trên $\mathfrak X$, ký hiệu $G_{\mathfrak X}=G$. Hạn chế của bó
này lên $U$ là bó hằng có các thớ rút về một phần tử.

Tổng quát hơn, việc xác định các thớ của $G_{\mathfrak X}$ là một vấn đề về
các vành địa phương đầy đủ; chính xác hơn:

\noindent b. Với mọi $x\in\mathfrak X$, ta có
\[
  G_x=G\bigl(\operatorname{Spec}(\widehat{\mathcal O}_{\mathfrak X,x})\bigr),
\]
(= các lớp, sai khác đẳng cấu, của các đại số $B$ hữu hạn và tự do trên
$\widehat{\mathcal O}_{\mathfrak X,x}$, được trang bị một đẳng cấu của
$B\otimes_{\widehat{\mathcal O}_{\mathfrak X,x}}\mathcal O_{X_0,x}$ với
$\mathcal O'_{0x}$, trong đó $\mathcal O'_0$ là bó đại số trên $X_0$ xác định
$X'_0$).

\noindent c. Có một đẳng cấu chính tắc
\[
  G_{\mathfrak X}=\varprojlim_nG_{X_n} ;
\]
nói cách khác, với mọi tập mở $V$ của $\mathfrak X$, ta có
\[
  G(V)=\varprojlim_nG(V_n).
\]

\noindent d. Giả sử $\mathfrak X$ thu được từ một lược đồ thông thường $X$
riêng trên một vành Noether địa phương đầy đủ $\Lambda$ có iđêan xác định
$\mathfrak m$, bằng cách lấy đầy đủ hóa $\mathcal J$-adic của $\mathcal O_X$,
trong đó $\mathcal J=\mathfrak m\mathcal O_X$. Khi đó $G(\mathfrak X)$ đẳng
cấu chính tắc với tập các lớp phủ phẳng của lược đồ thông thường $X$ ``rút về
$X'_0$''.

Nói một cách hình tượng, có thể nói rằng (a) và (b) thiết lập các quan hệ cơ
bản giữa phương diện địa phương và phương diện toàn cục của bài toán, (c) cho
các quan hệ giữa phương diện ``hữu hạn'' và phương diện ``vô cùng bé'', còn
(d) nhắc lại (dưới các điều kiện đã nêu) sự đồng nhất giữa phương diện ``hình
thức'' và phương diện ``đại số''.

Bây giờ giả sử $\mathfrak X$ được xác định trên một vành Noether địa phương
đầy đủ $\Lambda$, với $\mathcal J=\mathfrak m\mathcal O_{\mathfrak X}$ (do đó
$X_0$ là một tiền lược đồ trên $\Lambda/\mathfrak m$). Với mọi đại số $A$ hữu
hạn trên $\Lambda$, đặt
\[
  F(A)=G(\mathfrak X\times_\Lambda A).
\]
Đây là một hàm tử hiệp biến theo $A$, nhận giá trị trong phạm trù các tập hợp,
và theo (c) hàm tử này được xác định hoàn toàn khi biết nó trên các $A$ Artin;
nói rằng hàm tử này pro-biểu diễn được, tức có dạng
\[
  F(A)=\operatorname{Hom}_{\Lambda\text{-algèbres top.}}(O,A),
\]

\src{389}
trong đó $O$ là một đại số tôpô trên $\Lambda$ thuộc loại đã xét ở A, mục 5,
tương đương với việc điều ấy đúng khi hạn chế vào các $A$ Artin. Khi đó, sự kết
hợp của các Định lý 1 và 2 thực sự kéo theo:

\result{MỆNH ĐỀ}{5.1} Hàm tử trên pro-biểu diễn được.

Dĩ nhiên, theo (a), nếu $U=\mathfrak X$, $G(\mathfrak Y)$ rút về một phần tử
duy nhất với mọi $\mathfrak Y$ trên $\mathfrak X$, và khi đó hàm tử $F$ ít
đáng quan tâm (ta có $O=\Lambda$). Dường như trong hầu hết mọi trường hợp khác,
vành địa phương tôpô $O$ không Noether. Tuy nhiên, sự tồn tại của nó làm nổi
bật một cách đáng ngạc nhiên bản chất ``liên tục'' của tập nghiệm
$G(\mathfrak X)$ (tương ứng trực giác với việc có một lựa chọn ``liên tục''
trong cách sự phân nhánh lan truyền khi mở rộng $X'_0$). Hãy so sánh kết quả
này với quan điểm của J.-P. Serre [4] trong lý thuyết trường lớp địa phương,
vốn cũng làm nổi bật tính liên tục của nhóm Galois tôpô của mở rộng Abel cực
đại của một trường địa phương ``hình học''; ở đó nhóm đối ngẫu (theo nghĩa
Pontryagin) xuất hiện như một giới hạn quy nạp của các nhóm đại số (hoặc ít
nhất tựa đại số); ở đây cũng vậy, việc phân loại các mở rộng được cho bởi các
``đa tạp'' vô hạn chiều. Mặt khác, trong phần trên ta có thể lấy
$\mathfrak X$ là phổ hình thức của một vành địa phương đầy đủ (mà $\Lambda$
chẳng hạn là một vành con Cohen), và có thể hy vọng các kết quả của mục này
được dùng để nghiên cứu các mở rộng của một vành địa phương đầy đủ có chiều
$>1$. Cả trong trường hợp địa phương lẫn toàn cục, có lẽ chúng cho phép phát
biểu các quan hệ chính xác giữa những hiện tượng phân nhánh cao và những hiện
tượng ở đặc số không (có thể tiếp cận bằng phương pháp siêu việt). Dù sao,
chính phân tích sơ bộ của Mệnh đề 5.1 cho phép mở rộng các phương pháp trình
bày trong [2] để nghiên cứu nhóm cơ bản sang trường hợp ``phân nhánh vừa'' và
giải ``bài toán ba điểm'' bằng phương pháp siêu việt.

Để kết thúc, hãy lưu ý rằng tình huống trở nên đơn giản khi $X_0$ có chiều 1:
khi ấy, theo (a) và (b),
\[
  G(\mathfrak X)\simeq
  \prod_iG\bigl(\operatorname{Spec}(\widehat{\mathcal O}_{\mathfrak X,x_i})\bigr),
\]
trong đó $x_i$ là các điểm của $X_0-U$: ta có thể cho tùy ý các mở rộng ``địa
phương'' tại những điểm phân nhánh. Ngoài ra, khi $X_0$ chuẩn tắc, ta nhận thấy
lược đồ hình thức môđuli được bảo đảm bởi 5.1 là trơn trên
$\operatorname{Spec}(\Lambda)$.

\src{390}
\begin{center}
  \textbf{TÀI LIỆU THAM KHẢO}
\end{center}

\begin{description}
  \item[{[1]}] Grothendieck (Alexander). --- \emph{Sur quelques points
    d'algèbre homologique}, Tohoku Math. J., t. 9, 1957, tr. 119--221.
  \item[{[2]}] Grothendieck (Alexander). --- \emph{Géométrie formelle et
    géométrie algébrique}, Séminaire Bourbaki, t. 11, 1958/59, fasc. 3,
    số 182.
  \item[{[3]}] Grothendieck (Alexander). --- \emph{Technique de descente et
    théorèmes d'existence en géométrie algébrique, I : Généralités, Descente
    par morphismes fidèlement plats}, Séminaire Bourbaki, t. 12, 1959/60,
    fasc. 1, số 190.
  \item[{[4]}] Serre (Jean-Pierre). --- \emph{Corps locaux et isogénies},
    Séminaire Bourbaki, t. 11, 1958/59, fasc. 3, số 185.
\end{description}

\par

\clearpage
\chapter*{Đính chính cho Bài trình bày số 195}
\addcontentsline{toc}{chapter}{Đính chính cho Bài trình bày số 195}
\setcounter{footnote}{0}
\sid{e195}
\markboth{Đính chính 195}{Đính chính 195}
\src{300}

\noindent\textbf{4. TDTE II: Định lý tồn tại trong lý thuyết hình thức về
môđuli.}

\begin{center}
[Hội thảo Bourbaki, tập 12, 1959/60, số 195, 22 tr.]
\end{center}

\noindent\textbf{Trang 8.} --- Công thức viết trong Mệnh đề 5.1 chỉ đúng khi
$\Lambda$ là một trường; trong trường hợp tổng quát, phải thay
$\mathfrak m_\xi/\mathfrak m_\xi^2$ bằng thương của không gian này theo ảnh
của $\mathfrak n_\xi/\mathfrak n_\xi^2$, trong đó $\mathfrak n$ là iđêan cực
đại của $\Lambda$. Hơn nữa, định nghĩa đã cho về việc $O$ trơn trên
$\Lambda$ chỉ đúng khi mở rộng trường thặng dư $k'/k$ khả phân. Trong trường
hợp tổng quát, xem SGA III, 1.1.

\noindent\textbf{Trang 14, nhận xét.} --- Các bài toán nêu ở đây được giải
quyết hoàn toàn trong trường hợp xạ ảnh bằng các ``lược đồ Hilbert'' (xem
TDTE IV). Trái lại, những ví dụ của Nagata và Hironaka cho thấy các hàm tử
được xét không nhất thiết biểu diễn được nếu không áp đặt giả thiết xạ ảnh,
ngay cả khi chỉ phân loại các đa tạp con chiều không của một đa tạp đầy đủ,
không kỳ dị, chiều 3; điều này liên hệ với sự kiện rằng bình phương đối xứng
của một đa tạp như thế có thể không tồn tại.

\src{301}
\noindent\textbf{Trang 15, số 3.} --- Để có một nghiên cứu đầy đủ hơn, xem
TDTE V.

\noindent\textbf{Trang 16, Mệnh đề 3.1.} --- Thay vì ``nếu $f$ riêng'',
hãy đọc ``nếu $f$ riêng và khả tách''.

\noindent\textbf{Trang 18 và 19, các nhận xét, số 2.} --- Gần đây tôi đã
chứng minh rằng lược đồ hình thức môđuli của một đa tạp abel trên một trường
thực sự trơn trên vành Witt; nói cách khác, mọi lược đồ abel trên một vành
Artin địa phương vốn là thương của một vành khác đều thu được, bằng phép rút
gọn, từ một lược đồ abel trên vành sau. Chứng minh chỉ sử dụng các tính chất
biến thiên của lớp cản trở đối với phép nâng đã được đưa vào trong Bài trình
bày số 182, trang 12. Cũng xin nhắc rằng các lược đồ môđuli của các đường cong
giống $g$ và của các lược đồ abel phân cực đã được Mumford xây dựng (xem
SHMT).

\noindent\textbf{Trang 19, đoạn 5.} --- Phải giả sử tiết diện được xét để định
nghĩa $U$ không là ước của không chỉ trên $\mathfrak X$ mà còn trên mọi $X_n$.

\noindent\textbf{Trang 20, dòng 12 và 13.} --- Thay vì
\[
  G_x=G\bigl(\operatorname{Spec}(\widehat{\mathcal O}_{\mathfrak X,x})\bigr),
\]
hãy đọc
\[
  G_x=G\bigl(\operatorname{Spec}(\mathcal O_{\mathfrak X,x})\bigr)
  \subset
  G\bigl(\operatorname{Spec}(\widehat{\mathcal O}_{\mathfrak X,x})\bigr),
\]
và, thay vì ``$\mathcal O'_{0x}$'', hãy đọc
``$\widehat{\mathcal O}'_{0x}$''.

\par

\unit{Thuyết trình 212 -- Các tiền lược đồ thương}
\sid{212}
\markboth{Thuyết trình 212}{Thuyết trình 212}
\src{99}
\seminar{Năm thứ 13, 1960/61, số 212\\Tháng 2 năm 1961}
\paperhead{CÁC KỸ THUẬT XÂY DỰNG VÀ CÁC ĐỊNH LÝ TỒN TẠI TRONG HÌNH HỌC ĐẠI SỐ\\
III. CÁC TIỀN LƯỢC ĐỒ THƯƠNG}%
  {Alexander GROTHENDIECK}

\subsection*{Mở đầu}
\addcontentsline{toc}{subsection}{Mở đầu}

Các vấn đề được xét trong exposé này khác với các vấn đề của hai exposé trước
ở chỗ ta tìm cách biểu diễn một số hàm tử hiệp biến, chứ không còn là phản
biến, theo các lược đồ biến thiên. Tuy nhiên, phép lấy thương là thiết yếu
trong nhiều bài toán xây dựng của hình học đại số, kể cả các bài toán trong
Exposé I và II ([1], [2]). Chẳng hạn, câu hỏi về \emph{tính hữu hiệu của một
dữ liệu hạ xuống} trên một $T$-tiền lược đồ $X$, đối với một cấu xạ phẳng
trung thành và quasi-compact $T\longrightarrow S$, tương đương với câu hỏi về
sự tồn tại của một thương của $X$ (thỏa các tính chất hợp lý được xét dưới
đây) theo quan hệ tương đương phẳng trên $X$ do dữ liệu hạ xuống xác định; các
câu hỏi nêu trong [1], A, \S~2 c chắc chắn sẽ được giải quyết đồng thời với
các câu hỏi đặt ra trong đoạn 2 của exposé này. Tương tự, \emph{lược đồ
Picard} của một $S$-lược đồ $X$ (xem định nghĩa trong [2], C, \S~3) có thể
được định nghĩa theo nhiều cách như thương của một số lược đồ khác (các lược
đồ của ước hiệu dụng, hoặc của các phép nhúng vào một lược đồ xạ ảnh) theo
các quan hệ tương đương phẳng; trong khi đó, việc định nghĩa và xây dựng các
lược đồ phụ trợ này đơn giản hơn về mặt kỹ thuật: thật vậy, chúng là các lược
đồ dạng $\underline{\operatorname{Hom}}_S(X,Y)$ và các biến thể được định
nghĩa trong [2], C, \S~2, và việc xây dựng chúng sẽ là nội dung của exposé
tiếp theo (dưới các giả thiết thích hợp về tính xạ ảnh). Vì thế, bằng cách kết
hợp các kết quả của exposé này và exposé tiếp theo, ta thu được phép xây dựng
các lược đồ Picard dưới những giả thiết thích hợp.

Bài toán lấy thương trong phạm trù các tiền lược đồ vẫn còn đặt ra một số câu
hỏi chưa được giải quyết. Câu hỏi quan trọng nhất được nêu trong đoạn 8. Hiện
nay nó là trở ngại duy nhất đối với việc xây dựng \emph{các lược đồ moduli
trên các số nguyên cho các đường cong có giống tùy ý, các đa tạp abel phân
cực}, v.v. Điều đó cho thấy lời giải của nó xứng đáng với nỗ lực của các
chuyên gia về nhóm đại số.
\src{100}
\subsection*{1. Quan hệ tương đương, quan hệ tương đương hữu hiệu}
\addcontentsline{toc}{subsection}{1. Quan hệ tương đương, quan hệ tương đương hữu hiệu}

Cho $\mathcal C$ là một phạm trù và $X$ là một vật của $\mathcal C$. Một cặp
cấu xạ
\[
  p_1,p_2:R\rightrightarrows X
\]
được gọi là một \emph{cặp tương đương} trong $\mathcal C$, có \emph{đích}
$X$ và \emph{nguồn} $R$, nếu với mọi vật $T$ của $\mathcal C$, hai ánh xạ
tương ứng
\[
  p_1(T),p_2(T):R(T)\rightrightarrows X(T)
\]
(trong đó, với mỗi vật $Y$ của $\mathcal C$, ta đặt
$Y(T)=\operatorname{Hom}(T,Y)$) xác định một ánh xạ
\[
  R(T)\longrightarrow X(T)\times X(T)
\]
cảm sinh một song ánh từ $R(T)$ lên đồ thị của một quan hệ tương đương trên
tập $X(T)$. Ta đưa vào một quan hệ tương đương hiển nhiên giữa các cặp tương
đương có đích $X$; một lớp tương đương đối với quan hệ sau này được gọi là
một \emph{quan hệ tương đương của $\mathcal C$ trên $X$}, hoặc đơn giản là
một quan hệ tương đương nếu không thể gây nhầm lẫn.

Nếu $X\times X$ tồn tại, việc cho một quan hệ tương đương trên $X$ tương
đương với việc cho một vật con $R$ của $X\times X$ sao cho, với mọi vật $T$
của $\mathcal C$, tập con của
\[
  (X\times X)(T)=X(T)\times X(T)
\]
tương ứng với $R(T)$ là đồ thị của một quan hệ tương đương trên $X(T)$. Ký
hiệu $p_1$ và $p_2$ là các cấu xạ từ $R$ đến $X$ cảm sinh bởi các phép chiếu
$\operatorname{pr}_1$ và $\operatorname{pr}_2$, điều kiện trên nói rằng
$(p_1,p_2)$ là một cặp tương đương. Ta cũng có thể biểu diễn bằng biểu đồ
trong $\mathcal C$ (với điều kiện tồn tại $X\times X$ và tích thớ
$(R,p_2)\times_X(R,p_1)$) các tiên đề theo nghĩa tập hợp của một quan hệ
tương đương đối với các $R(T)$ trong các $X(T)$, phù hợp với nguyên lý tổng
quát được giải thích trong [2], A, \S~1. Ta sẽ không cần đến điều này.

Mỗi khi có một cặp cấu xạ $(p_1,p_2)$ cùng nguồn $R$ và cùng đích $X$, ta có
thể định nghĩa \emph{đối hạt nhân} của cặp ấy là vật $Y$ của $\mathcal C$
biểu diễn hàm tử hiệp biến theo $Z$
\[
  \operatorname{Hom}_{p_1,p_2}(X,Z),
\]
tức tập các cấu xạ $u$ từ $X$ đến $Z$ sao cho $up_1=up_2$. Nếu $Y$ tồn tại,
nó được xác định duy nhất đến một đẳng cấu duy nhất. Ta ký hiệu nó là
$X/(p_1,p_2)$ hoặc, theo lạm dụng ký hiệu, là $X/R$; ký hiệu sau chủ yếu được
dùng khi
\src{101}$(p_1,p_2)$ là một cặp tương đương: khi đó, theo thông lệ, ta đồng
nhất trong ký hiệu quan hệ tương đương do cặp ấy xác định với $R$. Lưu ý rằng
nếu xem $Y$ là một thương của $X$, thì quả thật nó chỉ phụ thuộc vào quan hệ
tương đương do cặp $(p_1,p_2)$ xác định.

Bây giờ, xét một cấu xạ
\[
  f:X\longrightarrow Y
\]
cho phép xem $X$ như một «vật trên $Y$», và giả sử tích thớ
\[
  R(f)=X\times_Y X
\]
tồn tại. Gọi $p_1$ và $p_2$ là hai phép chiếu của nó. Khi đó $(p_1,p_2)$ là
một cặp tương đương, gọi là cặp \emph{liên kết} với cấu xạ $f$. Do đó, cặp
này xác định một quan hệ tương đương, gọi là quan hệ liên kết với $f$.

Ta nói một cặp cấu xạ $(p_1,p_2)$ có đích $X$ và cùng nguồn $R$ là một
\emph{cặp tương đương hữu hiệu} nếu
\begin{enumerate}
  \item[(i)] đối hạt nhân $X/(p_1,p_2)=Y$ tồn tại;
  \item[(ii)] tích thớ $X\times_Y X$ tồn tại;
  \item[(iii)] cấu xạ $R\longrightarrow X\times_Y X$ có các thành phần $p_1$ và
  $p_2$ là một đẳng cấu.
\end{enumerate}
Khi đó, $(p_1,p_2)$ quả là một cặp tương đương. Ta cũng nói quan hệ tương
đương do nó xác định là một \emph{quan hệ tương đương hữu hiệu}.

Ta nói một cấu xạ $f:X\longrightarrow Y$ là một \emph{toàn cấu hữu hiệu}
nếu
\begin{enumerate}
  \item[(i)] tích thớ $X\times_Y X=R$ tồn tại;
  \item[(ii)] thương $X/(p_1,p_2)$ tồn tại, trong đó $p_1$ và $p_2$ là hai phép
  chiếu từ $R$ đến $X$;
  \item[(iii)] cấu xạ $X/(p_1,p_2)\longrightarrow Y$ do $f$ cảm sinh là một
  đẳng cấu.
\end{enumerate}
Khi đó, $f$ quả là một toàn cấu, và thậm chí là một toàn cấu chặt (xem [1],
A, \S~2.3); điều đảo lại đúng nếu tích thớ $X\times_Y X$ tồn tại. Ta cũng
gọi vật thương của $X$ do toàn cấu $f$ xác định là một \emph{thương hữu hiệu}
của $X$.

Các định nghĩa trên dẫn đến «tương ứng Galois» sau đây:

\src{102}
\result{MỆNH ĐỀ}{1.1} Có một tương ứng song ánh, bảo toàn các thứ tự tự
nhiên, giữa tập các quan hệ tương đương hữu hiệu $R$ trên $X$ và tập các
thương hữu hiệu $Y$ của $X$: quan hệ $R$ ứng với thương hữu hiệu $X/R$, còn
thương $Y$ ứng với quan hệ tương đương hữu hiệu do phép chiếu chính tắc
$X\longrightarrow Y$ xác định (quan hệ này được cho bởi tích thớ $X\times_Y X$
cùng với hai phép chiếu của nó).

Trong các phạm trù rất tốt (các tập hợp, các bó tập hợp, v.v.), mọi thương đều
hữu hiệu và mọi quan hệ tương đương đều hữu hiệu. Điều này không còn đúng
trong những phạm trù như phạm trù các tiền lược đồ trên một tiền lược đồ $S$
cho trước, ngay cả khi $S$ là phổ của một trường, và ngay cả khi chỉ xét các
lược đồ hữu hạn trên $S$. Những vấn đề về tính hữu hiệu, và thậm chí (đối với
các tiền lược đồ không hữu hạn trên $S$) những vấn đề về sự tồn tại của các
thương, thường tỏ ra tế nhị.

\subsection*{2. Ví dụ: các tiền lược đồ hữu hạn trên $S$}
\addcontentsline{toc}{subsection}{2. Ví dụ: các tiền lược đồ hữu hạn trên S}

Cho $\mathcal C$ là phạm trù các tiền lược đồ hữu hạn trên $S$, với giả thiết
cơ sở ấy là Noether địa phương. Phạm trù này tương đương với phạm trù đối của phạm
trù các bó đại số giao hoán coherent trên $S$, hoặc, nếu $S$ là affine có
vành $A$, với phạm trù đối của phạm trù các $A$-đại số hữu hạn trên $A$ (tức
là các môđun hữu hạn sinh trên $A$). Từ đó suy ra ngay rằng trong $\mathcal C$
tồn tại các giới hạn xạ ảnh hữu hạn và các giới hạn quy nạp hữu hạn. Đối với
loại thứ nhất, điều này đã được biết rõ (không cần bất kỳ giả thiết hữu hạn
nào\ldots). Như vậy, tích thớ của các tiền lược đồ $X,Y$ trên $S$ tương ứng
với tích tensor $B\otimes_A C$ của các đại số tương ứng; hạt nhân của hai cấu
xạ $X\rightrightarrows Y$, được xác định bởi hai đồng cấu $A$-đại số
$u,v:C\rightrightarrows B$, tương ứng với thương của $B$ theo iđêan sinh bởi
các phần tử $u(c)-v(c)$, v.v. Đối với các giới hạn quy nạp hữu hạn, chỉ cần
xét, một mặt, các tổng hữu hạn, tương ứng với tích thông thường của các
$A$-đại số, và mặt khác, các đối hạt nhân của các cặp cấu xạ
$X\rightrightarrows Y$; quả thật, chúng tương ứng (như thấy ngay) với vành con
của $C$ gồm những phần tử mà trên đó các đồng cấu
$u,v:C\rightrightarrows B$ trùng nhau (vành con này hữu hạn trên $A$ nhờ giả
thiết Noether). Hơn nữa, lưu ý rằng nhờ giả thiết Noether, ta có thể chứng
minh rằng các giới hạn quy nạp hữu hạn, và đặc biệt là các thương, được xây
dựng như trên trong phạm trù $\mathcal C$ các tiền lược đồ hữu hạn trên $S$,
thực ra là các thương trong phạm trù của \emph{tất cả} các tiền lược đồ.
\src{103}
Như đã nói trong [1], trong $\mathcal C$ có những toàn cấu không hữu hiệu
(hoặc chặt; hai điều này là như nhau vì
các tích thớ tồn tại). \emph{Tôi không biết liệu các quan hệ tương đương có
luôn hữu hiệu hay không} khi không đặt giả thiết nào về tính phẳng. Theo hướng
này, tôi chỉ thu được những kết quả khẳng định rất bộ phận, nhưng không thể
thiếu cho việc chứng minh định lý cơ bản trong lý thuyết hình thức về moduli
(xem [2], B, đl. 1). Lưu ý rằng trong bài toán đã nêu, ta dễ dàng quy về
trường hợp $S$ là phổ của một vành Artin địa phương $A$ có trường thặng dư
đóng đại số. Nhưng ngay cả khi $A$ là một trường, câu trả lời vẫn chưa được
biết.

Ta cũng có thể xét trường hợp một tiền lược đồ $X$ trên $S$ không còn được
giả sử là hữu hạn trên $S$, nhưng được trang bị một quan hệ tương đương $R$
trên $X$ sao cho $p_1:R\longrightarrow X$ là một cấu xạ hữu hạn. Khi đó, ta
nói $R$ là một quan hệ tương đương hữu hạn. Để đơn giản, giả sử $S$ và $X$
đều affine (do đó $R$ affine, nên tình huống được quy về một tình huống đại số
giao hoán thuần túy); \emph{ngay cả trong trường hợp này, ta cũng chưa biết có
tồn tại thương $X/R=Y$ hay không và cấu xạ chính tắc
$X\longrightarrow Y$ có hữu hạn hay không}. (Trường hợp đơn giản nhất là khi
$S$ là phổ của một trường $k$, còn $X$ là phổ của $k[t]$, tức đường thẳng
affine.) Dĩ nhiên, nếu hai bài toán trên đều có lời giải khẳng định, thì trong
tình huống hiện tại ta có thể kết luận rằng $R$ là hữu hiệu. Lưu ý rằng bài
toán về \emph{sự tồn tại} của một thương $Y$ và tính hữu hạn của
$f:X\longrightarrow Y$ được đặt ra hoàn toàn theo cùng một cách nếu, thay cho
đồ thị của một quan hệ tương đương trên $X$, ta chỉ có một tiền đồ thị tương
đương trên $X$, theo nghĩa của đoạn 4.

Vấn đề lấy thương theo một quan hệ tương đương hữu hạn ít nhiều tùy ý xuất
hiện khi dựng các tiền lược đồ bằng cách «dán» những tiền lược đồ $X_i$ đã cho
dọc theo một số tiền lược đồ con đóng; luật dán được biểu thị chính xác bởi
một quan hệ tương đương hữu hạn trên tiền lược đồ tổng $X$ của các $X_i$. Ta
cũng nên chờ đợi rằng lời giải cho các bài toán đặt ra ở đây và cho những biến
thể khác nhau sẽ là một điều kiện sơ bộ để xây dựng một kỹ thuật tổng quát về
các phép dựng không xạ ảnh, theo hướng được mở ra trong [2].

Kết quả tổng quát theo chiều khẳng định duy nhất mà người biên soạn biết là
kết quả sau:

\result{MỆNH ĐỀ}{2.1} Cho $S$ là một tiền lược đồ Noether địa phương, $s$ là
một điểm của $S$, và $\Omega$ là một mở rộng trường đóng đại số của $k(s)$.
Xét
\src{104}«hàm tử thớ» tương ứng $F$, gán cho mọi $S$-lược đồ $X$ hữu hạn
trên $S$ tập các điểm của $X_s$ nhận giá trị trong $\Omega$. Hàm tử này (hiển
nhiên khớp trái) là \emph{khớp phải}, tức giao hoán với các giới hạn quy nạp
hữu hạn, và đặc biệt với các đối hạt nhân của những cặp cấu xạ.

Áp dụng kết quả này cho mọi «điểm hình học» của $S$, ta suy ra rằng phạm trù
«thương» $\mathcal C'$ của $\mathcal C$, thu được bằng cách lập luận «modulo
các cấu xạ radiciel toàn ánh» (tức là ghép thêm một cách hình thức các nghịch
đảo của những cấu xạ đang xét), là một phạm trù «hình học», nghĩa là thỏa mãn
cùng những tính chất «có bản chất hữu hạn» như phạm trù các tập hợp. Đặc biệt,
mọi quan hệ tương đương trong đó đều hữu hiệu. Điều này kéo theo rằng nếu $R$
là một quan hệ tương đương trên một $X$ hữu hạn trên $S$, thì cấu xạ chính tắc
\[
  R\longrightarrow X\times_Y X\qquad (\text{trong đó }Y=X/R)
\]
là \emph{radiciel và toàn ánh} (thực ra là một phép nhúng đóng toàn ánh, vì nó
là một đơn cấu).

\subsection*{3. Trường hợp một nhóm tác tử}
\addcontentsline{toc}{subsection}{3. Trường hợp một nhóm tác tử}

Lại giả sử $\mathcal C$ là một phạm trù tùy ý. Cho $G,X$ là các vật của
$\mathcal C$, và giả sử $G$ là một $\mathcal C$-nhóm tác tử trên vật $X$ của
$\mathcal C$. Điều này có nghĩa (xem [2], A, \S~1) rằng với mọi vật $T$ của
$\mathcal C$, ta đã cho một cấu trúc nhóm trên $G(T)$ và một cấu trúc tập hợp
có nhóm tác tử $G(T)$ trên $X(T)$, sao cho khi $T$ biến thiên, những cấu trúc
đang xét «biến thiên theo hàm tử» theo $T$. Nếu trong $\mathcal C$ tồn tại các
tích $G\times G$ và $G\times X$, thì một cấu trúc như vậy cũng có thể được
xác định bởi một cặp cấu xạ
\[
  G\times G\longrightarrow G,
  \qquad \pi:G\times X\longrightarrow X,
\]
thỏa điều kiện rằng với mọi vật $T$ của $\mathcal C$, các luật hợp thành tương
ứng trên các tập $G(T)$ và $X(T)$ làm cho $G(T)$ thành một nhóm tác động trên
$X(T)$. Việc diễn đạt tiên đề này bằng tính giao hoán của một số biểu đồ trong
$\mathcal C$ là dễ nhưng dài dòng, và thực ra hoàn toàn vô ích trong mọi
trường hợp mà tôi biết.

Giả sử $G\times X$ tồn tại, và xét hai cấu xạ
\[
  p_1,p_2:G\times X\rightrightarrows X
\]
được cho bởi

\src{105}
\[
  p_1=\operatorname{pr}_1,\qquad p_2=\pi.
\]
Ta thấy ngay rằng $(p_1,p_2)$ là một cặp tương đương khi và chỉ khi, với mọi
vật $T$ của $\mathcal C$, ánh xạ
\[
  G(T)\times X(T)\simeq(G\times X)(T)
  \longrightarrow X(T)\times X(T),
\]
do cặp này xác định là đơn ánh, tức là khi nhóm $G(T)$ tác động
\emph{tự do} trên tập $X(T)$, nghĩa là khi $g\in G(T)$, $x\in X(T)$ và
$g\cdot x=x$ kéo theo $g$ là phần tử đơn vị của nhóm $G(T)$. Khi đó, ta nói
rằng $G$ \emph{tác động tự do trên $X$} (hoặc $X$ là một
\emph{$\mathcal C$-không gian chính dưới $G$}). Quan hệ tương đương liên kết
với cặp $(p_1,p_2)$ khi ấy được gọi là \emph{quan hệ tương đương do nhóm $G$
tác động tự do trên $X$ xác định}. Khi $X\times X$ cũng tồn tại và ta xét cấu
xạ
\[
  p:G\times X\longrightarrow X\times X
\]
do cặp $(p_1,p_2)$ xác định, điều kiện $G$ tác động tự do có nghĩa là $p$ là
một \emph{đơn cấu}.

Dĩ nhiên, ngay cả khi $G$ không tác động tự do trên $X$, ta vẫn muốn có các
tiêu chuẩn về sự tồn tại của một thương của $X$ theo $G$, tức đối hạt nhân của
cặp $(p_1,p_2)$ trên.

Đối hạt nhân đang xét thường được ký hiệu là $X/G$, hoặc tốt hơn là
$G\backslash X$ nếu $G$ tác động bên trái (ký hiệu trước được dành cho trường
hợp $G$ tác động bên phải). Lưu ý rằng ngay cả khi «ảnh» của $G\times X$ qua
$p$ tồn tại (chẳng hạn ảnh này được định nghĩa là vật con nhỏ nhất của
$X\times X$ mà qua đó $p$ phân tích được), ký hiệu là $R$, thì nói chung nó
không phải là một quan hệ tương đương trên $X$. Nếu khi ấy ta cố lấy trực tiếp
thương theo $R$ (hay chính xác hơn, theo cặp cấu xạ từ $R$ đến $X$ cảm sinh bởi
hai phép chiếu $\operatorname{pr}_i$), thì ta làm mất không thể khôi phục
những đặc trưng riêng của cặp ban đầu $(p_1,p_2)$. Vì vậy, điều quan trọng là
tìm một sự khái quát hóa của khái niệm quan hệ tương đương, áp dụng trực tiếp
cho cặp được xác định bởi một $\mathcal C$-nhóm tác tử.

\subsection*{4. Tiền quan hệ tương đương}
\addcontentsline{toc}{subsection}{4. Tiền quan hệ tương đương}

Nhắc lại rằng một phạm trù trong đó mọi cấu xạ đều là đẳng cấu được gọi là một
\emph{groupoid}. Mặt khác, một phạm trù phải được định nghĩa là gồm hai tập nền
$(X,R)$, tập các \emph{vật} và tập các \emph{mũi tên}, được trang bị những cấu
trúc sau đây:
\src{106}
\begin{enumerate}
\item[(i)] Một cặp ánh xạ
\[
  p_1,p_2:R\rightrightarrows X
\]
được gọi là \emph{ánh xạ nguồn} và \emph{ánh xạ đích},

\item[(ii)] Một ánh xạ
\[
  \pi:(R,p_2)\times_X(R,p_1)\longrightarrow R
\]
được gọi là \emph{ánh xạ hợp thành}.
\end{enumerate}
Các dữ liệu này phải thỏa mãn những tiên đề quen thuộc, mà ta sẽ không nhắc
lại ở đây, và có thể được diễn đạt bằng tính giao hoán của một số biểu đồ,
cùng với sự tồn tại của một ánh xạ $D:X\to R$ (nhất thiết duy nhất) làm cho
hai biểu đồ khác giao hoán. Ánh xạ $D$ gửi một vật tới mũi tên đồng nhất tương ứng
và thỏa mãn
\[
  p_1D=p_2D=\operatorname{id}_X.
\]
Việc nói phạm trù là một groupoid khi đó có nghĩa là tồn tại một ánh xạ
\[
  s:R\longrightarrow R
\]
(nhất thiết duy nhất), gọi là \emph{phép lấy nghịch đảo} của $R$, biến mỗi
mũi tên thành một mũi tên nghịch đảo. Điều này có thể được diễn đạt bằng tính
giao hoán của bốn biểu đồ khác được tạo từ $s$, $\Delta$ và các dữ liệu trước
đó; hai biểu đồ đầu là
\[
  p_1s=p_2,\qquad p_2s=p_1.
\]

Nhắc lại các điểm này, các định nghĩa tổng quát trong ([2], A, \S~1) đặc biệt
cho thấy thế nào là một cấu trúc $\mathcal C$-\emph{phạm trù}, tương ứng một
cấu trúc $\mathcal C$-\emph{groupoid}, trên một cặp vật $X,R$ của một phạm trù
tùy ý $\mathcal C$: theo định nghĩa, đó là dữ liệu, với mỗi vật $T$ của
$\mathcal C$, của một cấu trúc phạm trù, tương ứng một cấu trúc groupoid, theo
nghĩa tập hợp, có tập các vật là $X(T)$ và tập các mũi tên là $R(T)$, các cấu
trúc này \guillemotleft biến thiên theo hàm tử\guillemotright{} khi $T$ biến
thiên. Do đó, điều này xác định hai cấu xạ
\[
  p_1,p_2:R\rightrightarrows X
\]
được gọi là \emph{cấu xạ nguồn} và \emph{cấu xạ đích}, và khi tích thớ cần
thiết tồn tại, một cấu xạ
\src{107}
\[
  \pi:(R,p_2)\times_X(R,p_1)\longrightarrow R
\]
được gọi là \emph{cấu xạ hợp thành}. Ba dữ liệu này đủ để xác định cấu trúc
phạm trù trên $X,R$; tiên đề đặt lên các dữ liệu ấy là như sau: với mọi $T$,
ba dữ liệu tương ứng cho $X(T),R(T)$ xác định trên cặp tập hợp này một cấu
trúc phạm trù (tương ứng, một cấu trúc groupoid). Khi thích hợp, điều này có
thể được diễn đạt bằng tính giao hoán của một số biểu đồ có chứa một cấu xạ
được xác định duy nhất
\[
  D:X\longrightarrow R
\]
và, trong trường hợp các groupoid, một cấu xạ được xác định duy nhất
\[
  s:R\longrightarrow R;
\]
các biểu đồ được viết tường minh đúng như trong trường hợp «theo nghĩa tập
hợp». May thay, cách diễn giải dài dòng này của các tiên đề là vô ích trong
thực hành. Lợi ích lý thuyết duy nhất của khả năng biểu diễn dữ liệu và tiên
đề bằng các cấu xạ giữa một số tích thớ và các đẳng thức giữa những cấu xạ ấy
là như sau: nếu $F:\mathcal C\to\mathcal C'$ là một hàm tử khớp trái (tức giao
hoán với các tích hữu hạn và các tích thớ), thì nó biến một
$\mathcal C$-phạm trù (tương ứng, một $\mathcal C$-groupoid) thành một
$\mathcal C'$-phạm trù (tương ứng, một $\mathcal C'$-groupoid), với điều kiện
các tích hữu hạn và các tích thớ tồn tại trong $\mathcal C$.

Trong thực hành, điều quan trọng là biết diễn giải các cấu xạ
$p_1,p_2,\pi,D,s$ như những \emph{phép toán đơn hình} trong một vật bán đơn
hình hoặc đơn hình thích hợp của $\mathcal C$, ít nhất khi các tích thớ tồn
tại trong $\mathcal C$. Để cố định thuật ngữ, ta đưa vào \emph{phạm trù
$\mathcal S$ các đơn hình chuẩn} là phạm trù mà các vật là những tập hữu hạn
có dạng
\[
  \Delta_n=[0,n],\qquad\text{trong đó }n\geq 0
\]
(khoảng các số nguyên từ $0$ đến $n$), và các cấu xạ hay mũi tên là
\emph{các ánh xạ tùy ý} giữa những tập hữu hạn ấy. Lưu ý rằng phạm trù
$\mathcal S$ tương đương với phạm trù các tập hữu hạn \emph{khác rỗng}, trong
đó các cấu xạ là những ánh xạ giữa các tập hữu hạn. Trong $\mathcal S$, đối
tích của một họ hữu hạn \emph{khác rỗng} các vật hiển nhiên tồn tại, cũng như
phép đẩy ra của hai vật dưới một vật thứ ba (phép toán đối ngẫu với tích thớ).
Ký hiệu $\mathcal S'$ là phạm trù con của $\mathcal S$ có cùng các vật, nhưng
trong đó các cấu xạ là \emph{các ánh xạ tăng} giữa các $\Delta_n$. Phạm trù
này tương đương với phạm trù các tập hữu hạn khác rỗng được sắp thứ tự toàn
phần. Trong phạm trù này,
\src{108}
đối tích của hai vật không bao giờ tồn tại, và phép đẩy ra của hai vật $A,B$
dưới một vật thứ ba $C$ nói chung không tồn tại (chẳng hạn, lấy
$C=\Delta_0$, $A=B=\Delta_1$, với hai ánh xạ cấu trúc $u:C\to A$ và
$v:C\to B$ bằng nhau). Tuy nhiên, nó tồn tại trong một số trường hợp; chẳng
hạn
\[
  C=\Delta_0,\qquad A=\Delta_m,\qquad B=\Delta_n,\qquad
  u(0)=m,\qquad v(0)=0,
\]
và trong trường hợp này
\[
  \Delta_m\amalg_{\Delta_0}\Delta_n=\Delta_{m+n}.
\]

Một \emph{vật đơn hình} (tương ứng, một \emph{vật bán đơn hình}) trong một
phạm trù $\mathcal C$ theo định nghĩa là một hàm tử phản biến $K$ từ
$\mathcal S$ (tương ứng từ $\mathcal S'$) vào $\mathcal C$. Do đó, một vật
đơn hình xác định một vật bán đơn hình bằng phép hạn chế, nhưng khái niệm thứ
nhất khác khái niệm thứ hai về căn bản bởi sự hiện diện của \emph{các phép
toán đối xứng} trên các $K_n=K(\Delta_n)$, tương ứng với ảnh qua hàm tử $K$
của các phần tử thuộc nhóm đối xứng trên $n+1$ phần tử (được xem là nhóm các
tự đẳng cấu của $\Delta_n$ trong $\mathcal S$).

Với mọi $n$, gọi $\Delta'_n$ (tương ứng $\Delta''_n$) là phạm trù hữu hạn có
tập các vật là $\Delta_n$, được xác định bởi quan hệ thứ tự hỗn loạn (tương
ứng, thứ tự toàn phần tự nhiên) trên $\Delta_n$ (tập các mũi tên là đồ thị của
quan hệ thứ tự ấy). Hiển nhiên $\Delta'_n$ (tương ứng $\Delta''_n$) phụ thuộc
theo hàm tử vào vật $\Delta_n$ của $\mathcal S$ (tương ứng $\mathcal S'$).
Nếu $Z$ là một phạm trù, thì $\operatorname{Hom}(\Delta'_n,Z)$ (tương ứng
$\operatorname{Hom}(\Delta''_n,Z)$), khi $\Delta_n$ biến thiên, là một hàm tử
từ phạm trù $\mathcal S$ (tương ứng $\mathcal S'$) vào phạm trù các tập hợp,
tức một tập đơn hình (tương ứng, một tập bán đơn hình), được gọi là
\emph{liên kết với phạm trù $Z$}; ký hiệu chúng lần lượt là $Z'$ và $Z''$.
Hơn nữa, có một đồng cấu tự nhiên hiển nhiên từ tập bán đơn hình liên kết
với $Z'$ vào
$Z''$, và đồng cấu này là một đẳng cấu khi và chỉ khi $Z$ là một groupoid. Với
các quy ước này:

\result{MỆNH ĐỀ}{4.1} Hàm tử $Z\longmapsto Z''$ từ phạm trù các phạm trù vào
phạm trù các tập bán đơn hình là đầy và trung thành, và xác định một tương
đương giữa phạm trù các \emph{phạm trù} với phạm trù các tập bán đơn hình,
tức các hàm tử phản biến $K$ từ $\mathcal S'$ vào $(\mathrm{Set})$, biến các
phép đẩy ra $A\amalg_C B$ thuộc kiểu được chỉ ra ở trên thành các tích thớ của
các tập hợp. Tương tự, hàm tử $Z\longmapsto Z'$ từ phạm trù các groupoid vào
phạm trù các tập đơn hình là đầy và trung thành, và xác định một tương
đương giữa phạm trù các \emph{groupoid} với phạm trù các
\src{109}
các tập đơn hình, tức các hàm tử phản biến $K$ từ $\mathcal S$ vào
$(\mathrm{Ens})$, biến các phép đẩy ra thành các tích thớ.

Do đó, có thể xem các phạm trù như những tập bán đơn hình đặc biệt, và các
groupoid như những tập đơn hình đặc biệt, dĩ nhiên với điều kiện phải lập luận
«sai khác một đẳng cấu», như bắt buộc khi diễn giải một số cấu trúc bằng những
cấu trúc khác. Phương pháp thông thường quy về trường hợp lý thuyết tập hợp
khi đó kéo theo:

\result{HỆ QUẢ}{4.2} Mệnh đề trên vẫn đúng khi thay các phạm trù, các
groupoid và các tập đơn hình bằng các $\mathcal C$-phạm trù,
$\mathcal C$-groupoid và các $\mathcal C$-vật đơn hình, với điều kiện các tích
thớ tồn tại trong $\mathcal C$.

Vật bán đơn hình $K$ trong $\mathcal C$ liên kết với một phạm trù
$(X,R,\ldots)$ trong $\mathcal C$ có thể được viết tường minh bằng cách diễn
giải thành phần $K_n=K(\Delta_n)$ của $K$ là tích thớ lặp $(n+1)$ lần của
$(R,p_1)$ trên $X$, hoặc tốt hơn, bằng công thức truy hồi
\[
  K_0=R,\qquad
  K_n=(K_{n-1},p_n^{(n-1)})\times_X(R,p_1),
\]
trong đó các $p_i^{(n-1)}$ ($0\leq i\leq n-1$) là các phép chiếu tự nhiên từ
$K_{n-1}$ xuống $X$ (cũng được định nghĩa bằng truy hồi). Theo cách này,
$p_1,p_2,\pi,D,s$ được diễn giải là các phép toán đơn hình lần lượt tương ứng
với các cấu xạ sau trong $\mathcal S$: mặt $0$ của $\Delta_1$, mặt $1$ của
$\Delta_1$, mặt $(0,2)$ của $\Delta_2$, phép suy biến
$\Delta_1\to\Delta_0$, và phép đối xứng của $\Delta_1$. Mọi phép toán bán đơn
hình khác (tương ứng, đơn hình) đều được suy ra một cách hình thức từ bốn
phép toán trước (tương ứng, năm phép toán) bằng phép hợp thành và các tích thớ.

Nay gọi \emph{tiền quan hệ tương đương} trên một vật $X$ của một phạm trù là
việc cho một groupoid mà vật của các vật (nếu có thể nói vậy) là $X$. Vì thế,
dữ liệu ấy bao gồm, ngoài những thứ khác, một vật $R$ và hai cấu xạ
\[
  p_1,p_2:R\rightrightarrows X.
\]
Tuy nhiên, lưu ý rằng riêng các dữ liệu này không xác định cấu trúc đang xét,
khác với trường hợp các cặp tương đương. Trong exposé này, ta quan tâm đến
khái niệm ấy nhằm thu được các tiêu chuẩn để có thể lấy thương, tức là lập đối
hạt nhân của cặp $(p_1,p_2)$. Vì vậy, phát biểu của bài toán không dùng đến
các dữ liệu bổ sung vốn có trong một groupoid. Trong việc chứng minh các kết
quả sẽ nêu, ta sẽ
\src{110}
vẫn sử dụng các dữ liệu bổ sung ấy, và đặc biệt là các phép toán đơn hình
(kể cả các phép toán đối xứng) cho đến chiều $3$ (có liên quan đến tích thớ
lặp bốn lần của $R$ trên $X$).

Một quan hệ tương đương trên một vật $X$ của $\mathcal C$ xác định một tiền
quan hệ tương đương: chỉ cần xét điều này trong trường hợp tập hợp, và khi ấy,
với một quan hệ tương đương trên một tập $X$, ta liên kết groupoid có tập vật
là $X$, còn tập mũi tên là đồ thị của quan hệ tương đương.

Một $\mathcal C$-monoid $G$ tác động lên một vật $X$ của $\mathcal C$ xác định
một $\mathcal C$-phạm trù có các vật cơ sở là $R=G\times X$ và $X$ (với điều
kiện $G\times X$ tồn tại), và là một $\mathcal C$-groupoid khi và chỉ khi $G$
là một nhóm. Một lần nữa, chỉ cần kiểm tra điều này trong trường hợp tập hợp.
Khi đó, phép hợp thành các mũi tên $(g,a)$ và $(g',g\cdot a)$ được định nghĩa
bởi
\[
  (g',g\cdot a)(g,a)=(g'g,a),
\]
tức là, nếu $a,b\in X$, thì theo định nghĩa $\operatorname{Hom}(a,b)$ là tập
chuyển từ $a$ đến $b$, và các cấu xạ được hợp thành nhờ phép hợp thành các phần
tử của $G$.

\medskip
\noindent\textbf{NHẬN XÉT. ---} Có thể tránh các khó khăn logic do một phát
biểu như (4.1) gây ra bằng cách hiểu ngầm rằng mọi vật đang xét đều nằm trong
một «Vũ trụ» cố định (bản thân nó là một tập hợp).

\fsection{5}{Thương theo một quan hệ tương đương hữu hạn và phẳng}

\result{ĐỊNH LÝ}{5.1} Cho $X=\operatorname{Spec}(B)$ là một lược đồ affine và
$R$ là một tiền quan hệ tương đương trên $X$, mà thành phần $R_1$ là affine,
đặt $R_1=\operatorname{Spec}(C)$. Giả sử phép chiếu thứ nhất
$p_1:R_1\to X$ là một cấu xạ hữu hạn và tự do địa phương; nói cách khác, giả
sử đồng cấu vành tương ứng $p_1^*:B\to C$ làm cho $C$ trở thành một
$B$-môđun xạ ảnh hữu hạn sinh. Gọi $A$ là vành con của $B$, là hạt nhân của
cặp đồng cấu
\[
  p_1^*,p_2^*:B\rightrightarrows C
\]
(= tập hợp các phần tử $b$ sao cho $p_1^*(b)=p_2^*(b)$). Đặt
$Y=\operatorname{Spec}(A)$, và gọi $f:X\to Y$ là cấu xạ được xác định bởi phép
nhúng $A$ vào $B$. Trong các điều kiện này:
\begin{enumerate}
\item[(i)] $B$ nguyên trên $A$, tức là $f$ là một cấu xạ nguyên.

\item[(ii)] Cấu xạ $f$ là toàn ánh, và các thớ của nó là các lớp tương đương
\src{111}
theo nghĩa tập hợp $p_2\allowbreak(p_1^{-1}(x))$ trong $X$ modulo $R$, và
tôpô của $Y$ là tôpô thương của tôpô trên $X$.

\item[(iii)] $Y$ là thương của $X$ theo $R$ trong phạm trù các tiền lược đồ.

\item[(iv)] Nếu $R$ xuất phát từ một \emph{quan hệ tương đương}, thì cấu xạ
$f:X\to Y$ là hữu hạn và tự do địa phương (nghĩa là $B$ là một $A$-môđun xạ
ảnh hữu hạn sinh), và quan hệ tương đương là hữu hiệu, nghĩa là
$R_1\to X\times_Y X$ là một đẳng cấu.
\end{enumerate}

Định lý này tổng quát hóa định lý quen thuộc về trường hợp một nhóm hữu hạn
$G$ tác động bằng các tự đẳng cấu lên vành $B$ và về vành bất biến $A$; chứng
minh của nó tương tự chứng minh quen thuộc. Có thể phát biểu chính xác hơn
(iii) như sau:

\result{HỆ QUẢ}{5.2} Cấu xạ chính tắc
$R_1\to X\times_Y X$ là \emph{toàn ánh}.

Vẫn cho $R$ là một tiền quan hệ tương đương «hữu hạn và tự do địa phương»
trên $X$, nhưng bây giờ $X$ là một tiền lược đồ tùy ý. Giả sử có thể tìm được
một tiền lược đồ $Y$ và một cấu xạ $f:X\to Y$ sao cho $fp_1=fp_2$, đồng thời
dãy các đồng cấu giữa các bó vành trên $Y$
\[
  \mathcal O_Y\longrightarrow f_*(\mathcal O_X)
  \rightrightarrows g_*(\mathcal O_R)
\]
là khớp (trong đó $g=fp_i$). Khi đó, từ định lý suy ra các kết luận tương tự
(i) đến (iv) của định lý; đặc biệt, theo (iii), $Y$ là thương của $X$ theo
$R$, và do đó được xác định sai khác một đẳng cấu duy nhất. Trong các điều
kiện ấy, tiền quan hệ tương đương $R$ trên $X$ được gọi là \emph{khả chấp}.
Với định nghĩa này:

\result{ĐỊNH LÝ}{5.3} Cho $X$ là một tiền lược đồ và $R$ là một tiền quan hệ
tương đương trên $X$ sao cho $p_1:R_1\to X$ là một cấu xạ hữu hạn và tự do địa
phương. Khi đó, $R$ khả chấp khi và chỉ khi mọi lớp tương đương theo nghĩa tập
hợp $p_2(p_1^{-1}(x))$ trong $X$ modulo $R$ đều nằm trong một tập mở affine
(điều kiện này luôn được thỏa mãn nếu mọi tập con hữu hạn của $X$ đều nằm
trong một tập mở affine, chẳng hạn nếu $X$ là tựa xạ ảnh trên một lược đồ
affine).

Thật vậy, không khó để chứng minh rằng khi ấy mọi lớp tương đương modulo $R$
trong $X$ đều nằm trong một tập mở affine \emph{ổn định theo $R$}, và thương
$Y$ được dựng bằng cách dán các mảnh thu được nhờ Định lý 5.1.
\src{112}

\result{HỆ QUẢ}{5.4} Giả sử điều kiện ấy được thỏa mãn và ngoài ra $R$ xuất
phát từ một quan hệ tương đương. Khi đó, quan hệ này là hữu hiệu, nghĩa là
$R_1\to X\times_Y X$ là một đẳng cấu, và $f:X\to Y$ là một cấu xạ hữu hạn và
tự do địa phương.

Từ đó, bằng «hạ xuống», ta suy ra ngay:

\result{HỆ QUẢ}{5.5} Dưới các điều kiện của (5.4), $X$ có hạng $n$ tại mọi
điểm trên $Y$ khi và chỉ khi $(R_1,p_1)$ có hạng $n$ tại mọi điểm trên $X$.
Nếu $X,R_1$ là các $Z$-tiền lược đồ và $p_1,p_2$ là các $Z$-cấu xạ, nên $Y$
là một $Z$-tiền lược đồ, thì $X$ phẳng trên $Z$ khi và chỉ khi $Y$ cũng vậy.

Tóm lại:

\result{CHÚ GIẢI}{} Việc cho một cấu xạ hữu hạn, tự do địa phương và toàn ánh
$f:X\to Y$ giữa các tiền lược đồ tương đương với việc cho một tiền lược đồ
$X$ được trang bị một quan hệ tương đương $R$ sao cho $p_1:R\to X$ là hữu hạn
và tự do địa phương, và mọi lớp $p_2(p_1^{-1}(x))$ đều nằm trong một tập mở
affine.

\medskip
\noindent\textbf{CÁC NHẬN XÉT 5.6.}

\smallskip
\noindent $1^\circ$ Không cần đặt giả thiết Noether nào.

\smallskip
\noindent $2^\circ$ Khái niệm lấy thương này bao gồm, như một trường hợp đặc
biệt, «hạ xuống không khả phân» của CARTIER, tương ứng với việc xác định các
cấu xạ hữu hạn và tự do địa phương $f:X\to Y$ sao cho
$f_*(\mathcal O_X)$ có một $p$-cơ sở đối với $\mathcal O_Y$ ($X$ là một tiền
lược đồ đã cho mà bó $\mathcal O_X$ bị số nguyên tố $p>0$ triệt tiêu). Lưu ý
rằng kết quả được phát biểu dễ dàng mà không cần giả thiết chính quy nào về
các vành địa phương và không cần giả sử $X$ là một lược đồ đại số trên một
trường. Lý thuyết JACOBSON--BOURBAKI thu được bằng cách lấy $X$ là phổ của
một trường có đặc số $p$.

\smallskip
\noindent $3^\circ$ Trước đó, GABRIEL đã thu được một trường hợp riêng của
Định lý (5.3) trong lý thuyết lấy thương đối với các nhóm giao hoán hữu hạn
trên một trường $k$. [So sánh (7.4)].

\fsection{6}{Thương theo một quan hệ tương đương riêng và phẳng}

\result{ĐỊNH LÝ}{6.1} Cho $S$ là một tiền lược đồ Noether địa phương, $X$ là
một $S$-lược đồ tựa xạ ảnh, và $R$ là một tiền quan hệ tương đương trên $X$
sao cho:
\src{113}
\begin{enumerate}
\item[(a)] $p_1:R_1\to X$ là riêng và phẳng;
\item[(b)] $R_1\to X\times_S X$ là một cấu xạ hữu hạn, hoặc tương đương theo
(a), có các thớ hữu hạn (điều kiện tự động được thỏa mãn nếu $R$ xuất phát từ
một quan hệ tương đương).
\end{enumerate}
Dưới các điều kiện này:
\begin{enumerate}
\item[(i)] $X/R=Y$ tồn tại và (nếu $S$ là Noether) là tựa xạ ảnh trên $S$.

\item[(ii)] Cấu xạ chính tắc $f:X\to Y$ là toàn ánh, riêng và mở; các thớ của
nó là các lớp tương đương $p_2(p_1^{-1}(x))$ trong $X$ modulo $R$. Do đó,
$Y$ đồng nhất với không gian tôpô thương của $X$ theo quan hệ tương đương theo
nghĩa tập hợp do $R$ xác định. Cuối cùng,
$R_1\to X\times_Y X$ là toàn ánh.

\item[(iii)] Nếu $R$ xuất phát từ một quan hệ tương đương, thì quan hệ ấy là
hữu hiệu, nghĩa là $R_1\to X\times_Y X$ là một đẳng cấu, và hơn nữa
$f:X\to Y$ là phẳng (do đó phẳng trung thành).
\end{enumerate}

Để chứng minh, ta quy về (5.1) nhờ các tựa tiết diện thích hợp của $X$ đối với
$R$; chứng minh tương tự việc dựng các nhóm đại số thương trong Séminaire
Chevalley.

Tóm lại:

\result{CHÚ GIẢI}{} Cho $X$ tựa xạ ảnh trên $S$ Noether địa phương. Việc cho
một cấu xạ riêng, phẳng trung thành và toàn ánh $f:X\to Y$ từ $X$ vào một
$S$-tiền lược đồ $Y$ tương đương với việc cho một quan hệ tương đương $R$ trên
$X$ sao cho $p_1:R\to X$ là riêng và phẳng.

Phương pháp ấy cũng cho kết quả sau:

\result{ĐỊNH LÝ}{6.2} Cho $S$ là một tiền lược đồ Noether, $X$ là một tiền
lược đồ thuộc kiểu hữu hạn trên $S$, và $R$ là một tiền quan hệ tương đương trên
$S$-tiền lược đồ $X$; giả sử rằng:
\begin{enumerate}
\item[(a)] $p_1:R_1\to X$ là phẳng và thuộc kiểu hữu hạn;
\item[(b)] cấu xạ $R_1\to X\times_S X$ là tựa hữu hạn (nghĩa là có các thớ
hữu hạn).
\end{enumerate}
Khi đó, tồn tại một tập mở $U$ trù mật trong $X$, \emph{bão hòa} đối với $R$,
sao cho:
\begin{enumerate}
\item[(i)] Nếu $R_U$ là tiền quan hệ tương đương do $R$ cảm sinh trên $U$,
thì $U/R_U$ tồn tại và thuộc kiểu hữu hạn trên $S$.

\item[(ii)] Cấu xạ chính tắc $U\to U/R_U$ là toàn ánh và mở; các thớ của nó
là các lớp tương đương theo nghĩa tập hợp đối với $R_U$ (do đó $U/R_U$ là
không gian tôpô thương của $U$ theo quan hệ tương đương theo nghĩa tập hợp do
$R_U$ xác định). Cuối cùng, cấu xạ
\[
  (R_U)_1\longrightarrow U\times_{U/R_U}U
\]
là toàn ánh.
\src{114}

\item[(iii)] Nếu $R$ xuất phát từ một quan hệ tương đương, thì có thể giả sử
$U\to U/R_U$ phẳng trung thành và $R_U$ hữu hiệu.
\end{enumerate}

Đây là một kết quả có bản chất chủ yếu là «song hữu tỉ».

\medskip
\noindent\textbf{CÁC NHẬN XÉT 6.3.}

\smallskip
\noindent $1^\circ$ Tôi không biết liệu trong (6.1) và (6.2), giả thiết (b) có
thừa hay không. Khi lấy thương theo các nhóm, giả thiết ấy thực tế buộc ta
giới hạn vào trường hợp mọi nhóm ổn định đều là nhóm hữu hạn.

\smallskip
\noindent $2^\circ$ Có thể hỏi liệu có những kết quả tương tự (6.1) và (6.2)
mà không cần giả thiết phẳng hay không. Tôi không biết phản ví dụ nào theo
hướng này. Mặt khác, ngay cả khi vẫn giữ giả thiết phẳng và chỉ xét các quan
hệ tương đương sao cho $p_1:R\to X$ phẳng và tựa hữu hạn (nhưng không hữu
hạn), với $X$ affine, vẫn có thể xảy ra việc $R$ không hữu hiệu: hãy lấy các
quan hệ tương đương cảm sinh trên các tập mở affine phủ đa tạp Nagata (trên đó
một nhóm gồm hai phần tử tác động theo cách «không khả chấp»).

\fsection{7}{Ứng dụng}

Như đã nói trong phần mở đầu, ứng dụng quan trọng nhất của (6.1) là việc dựng
các lược đồ Picard, cũng như việc tìm nghiệm cho nhiều bài toán moduli khác;
chúng ta sẽ trở lại các bài toán này sau.

Ta thu được một chứng minh đơn giản cho kết quả sau đây của SHIMURA:

\result{MỆNH ĐỀ}{7.1} Cho $A$ là một lược đồ abel được xác định trên một vành
định giá rời rạc $V$ có trường phân thức $K$. Khi đó, mọi lược đồ abel $B'$
trên $K$ đẳng giống với một thương của $A\otimes_V K$ đều «có giảm tốt trên
$V$», tức là đẳng cấu với một $B\otimes_V A$, trong đó $B$ là
một lược đồ abel trên $V$ (về căn bản là duy nhất, như ta đã nhắc).

Có thể giả sử $B'$ là thương của $A_K$ theo một lược đồ con nhóm $C'$.
(LƯU Ý: nói chung, $C'$ sẽ không «rút gọn», tức là các vành địa phương của nó
sẽ có những phần tử lũy linh.) Xét lược đồ con đóng $C$ của $A$ là «bao đóng»
của $C'$, tức là lược đồ con đóng nhỏ nhất của $A$ sao cho $C_K$ chứa $C'$.
Ta có $C_K=C'$, và từ việc $V$ là một vành định giá rời rạc, dễ dàng suy ra
rằng $C$ là một lược đồ con nhóm của $A$ trên $V$. Vì $A$ riêng trên
$\operatorname{Spec}(V)=S$, nên $C$ cũng vậy; mặt khác, $A$ thậm chí còn xạ
ảnh trên $S$. Do đó, ta có thể áp dụng (6.1) để dựng $A/C=B$,
\src{115}
đó là $B$ cần tìm.

Sau cùng, các lập luận về cơ bản đã biết cho phép suy ra từ (6.2) kết quả sau:

\result{ĐỊNH LÝ}{7.2} Cho $S$ là phổ của một vành Artin, $F$ và $G$ là các
lược đồ nhóm thuộc kiểu hữu hạn trên $S$, và $u:F\to G$ là một đồng cấu lược
đồ nhóm trên $S$. Giả sử:
\begin{enumerate}
\item[(i)] $F$ phẳng trên $S$;
\item[(ii)] hạt nhân của $u$ là hữu hạn.
\end{enumerate}
Trong những điều kiện này, lược đồ thương $G/F$ tồn tại; cấu xạ chính tắc
$G\to G/F$ là toàn ánh và mở, và các thớ của nó là các lớp tương đương theo
nghĩa tập hợp được xác định bởi các tác động bên phải của $F$ lên $G$. Sau
cùng, nếu $u$ là một đơn cấu, thì $G\to G/F$ là phẳng và cấu xạ
\[
  G\times F\longrightarrow G\times_{G/F}G
\]
là một đẳng cấu; nói cách khác, $G$ là một không gian thuần nhất chính trên
$G/F$, có nhóm cấu trúc (tác động bên phải) là $F$, hay đúng hơn là
$F\times_S(G/F)$ được xem như một lược đồ nhóm trên $G/F$ (xem [1], B,
\S~6).

\result{HỆ QUẢ}{7.3} Trong những điều kiện này, $G$ phẳng trên $S$ khi và chỉ
khi $G/F$ cũng vậy. Nếu điều kiện này được thỏa mãn, thì việc lấy thương
theo $F$ giao hoán với mọi phép mở rộng cơ sở của $S$; và nếu $F$ là một nhóm
con chuẩn tắc của $G$, thì $G/F$ có thể được trang bị một cấu trúc
\emph{nhóm thương} của $G$ theo $F$.

Tình hình đặc biệt đơn giản khi $S$ là phổ của một trường, vì khi đó mọi
$S$-tiền lược đồ đều đương nhiên phẳng trên $S$. Ta thu được:

\result{HỆ QUẢ}{7.4} Cho $F,G$ là các lược đồ nhóm thuộc kiểu hữu hạn trên một
trường $k$, và cho $u:F\to G$ là một đồng cấu $k$-nhóm. Khi đó $u$ phân tích
thành $F\to F'\to G$, trong đó $F\to F'$ là đồng cấu lấy thương theo nhóm con
đóng $\operatorname{Ker}u$ của $F$, còn $F'\to G$ là một đồng cấu nhóm và đồng
thời là một phép nhúng đóng. Lược đồ thương $G/F=G/F'$ tồn tại. Hình thức luận
thông thường (kiểu các định lý Noether) áp dụng được cho các nhóm đại số trên
$k$.

Kết quả này cho phép xử lý một cách thống nhất việc lấy thương trong các nhóm
đại số theo nghĩa cổ điển (tức là những nhóm bất khả quy trên $k$ và trơn trên
$k$), và việc lấy thương theo các nhóm con «vô cùng bé»
\src{116}
được CARTIER xét. Có lợi khi xem các «hyperalgebra» do CARTIER đưa ra, tiếp nối
các công trình của DIEUDONNÉ về các nhóm hình thức, như những nhóm trong phạm
trù các lược đồ hình thức trên $k$, và khi thích hợp (nếu chúng tương ứng với
các hyperalgebra có hạng hữu hạn trên $k$), như các nhóm đại số hữu hạn trên
$k$.

\fsection{8}{Một giả thuyết}

Nó đáp ứng nhu cầu lấy thương theo nhóm xạ ảnh tác động trên một số lược đồ con
của các «lược đồ Hilbert» (trong lý thuyết lược đồ, các lược đồ sau thay thế
các đa tạp Chow).

Cho $S$ là một tiền lược đồ và $n$ là một số nguyên. Với mỗi tiền lược đồ $S'$
trên $S$, ta cho tương ứng nhóm
$\operatorname{Gl}(n,\Gamma(S',\mathcal O_{S'}))$ gồm các ma trận khả nghịch có
$n$ hàng và $n$ cột, với các hệ số trong vành các tiết diện của
$\mathcal O_{S'}$. Ta thu được một hàm tử phản biến theo $S'$; dễ chứng minh
rằng nó biểu diễn được, nên tương ứng với một lược đồ nhóm trên $S$, hơn nữa
affine trên $S$, ký hiệu $\operatorname{Gl}(n)_S$. Phép dựng này tương thích
với đổi cơ sở, vì vậy trên thực tế tất cả xuất phát từ một lược đồ nhóm trên
$\mathbf Z$, ký hiệu $\operatorname{Gl}(n)$. Nhóm $\operatorname{Gl}(1)$, gọi
là nhóm nhân và thường ký hiệu $\mathbf G_m$, tương ứng với hàm tử
$S\longmapsto\Gamma(S,\mathcal O_S)^*$, nhóm các «đơn vị» trên $S$. Có một
đồng cấu hiển nhiên $\operatorname{Gl}(1)\to\operatorname{Gl}(n)$, và dễ dàng
dựng được nhóm thương, ký hiệu $\operatorname{GP}(n-1)$ và gọi là
\emph{nhóm xạ ảnh bậc $n-1$ trên $\mathbf Z$}. Nó biểu diễn hàm tử gán cho
$S$ nhóm các tiết diện của bó
$\widetilde{\operatorname{Gl}(n)}_S/
\widetilde{\operatorname{Gl}(1)}_S$, trong đó
$\widetilde{\operatorname{Gl}(n)}_S$ ký hiệu bó các mầm tiết diện của
$\operatorname{Gl}(n)_S$ trên $S$. (Cần lưu ý rằng các tiết diện của
$\operatorname{GP}(n-1)_S$ trên $S$ nói chung không xuất phát từ các tiết diện
của $\operatorname{Gl}(n)_S$ trên $S$!) Lưu ý rằng có thể chứng minh
$\operatorname{GP}(n-1)$ cũng biểu diễn hàm tử
$S\longmapsto\operatorname{Aut}_S(\mathbf P_S^{n-1})$, trong đó
$\mathbf P_S^{n-1}$ là không gian xạ ảnh chuẩn có chiều tương đối $n-1$ trên
$S$, ít nhất khi $S$ là Noether. Chính theo cách này mà nó xuất hiện trong lý
thuyết moduli.

Cho $S$ là một lược đồ Noether, nếu muốn có thể giả sử affine, và cho $X$ là
một $S$-tiền lược đồ tựa xạ ảnh, được trang bị một bó khả nghịch $\mathcal L$
rất rộng tương đối với $S$. Giả sử nhóm $G=\operatorname{GP}(n)_S$ tác động
trên $X$ và đồng thời trên $\mathcal L$ (tương thích với tác động của nó trên
$X$), và tác động \emph{tự do} trên $S$.
\src{117}

\result{GIẢ THUYẾT}{8.1} Trong các điều kiện trên:

\smallskip
\noindent $1^\circ$ \emph{Quan hệ tương đương do $G$ xác định là hữu hiệu,
thương $Y=X/G$ thuộc kiểu hữu hạn trên $S$, và cấu xạ chính tắc $f:X\to Y$ là
phẳng và toàn ánh} (do đó $X$ trở thành một không gian thuần nhất chính trên
$Y$, với nhóm $G\times_S Y=\operatorname{GP}(n)_Y$).

\smallskip
\noindent $2^\circ$ Cho $\mathcal L'$ là bó khả nghịch trên $Y$ được suy ra
từ $\mathcal L$ bằng «hạ xuống phẳng trung thành» theo $f$ (xem [1], B,
\S~th. 1). Khi đó $\mathcal L'$ là \emph{tiền rộng trên $Y$ tương đối với
$S$}, nghĩa là tồn tại một số nguyên $m$ và một cấu xạ tựa hữu hạn từ $Y$ vào
một không gian xạ ảnh chuẩn thích hợp $\mathbf P_S^N$, sao cho
$(\mathcal L')^{\otimes m}$ đẳng cấu với ảnh ngược của
$\mathcal O_{\mathbf P_S^N}(1)$.

Lưu ý rằng ngay cả khi $X$ tách trên $S$, vẫn có thể xảy ra việc $Y$ không
tách trên $S$ (một tình huống gặp trong những bài toán moduli hoàn toàn không
bệnh lý). Nếu $1^\circ$ đúng, thì $Y$ tách khi và chỉ khi quan hệ tương đương
do $G$ xác định có đồ thị đóng, nghĩa là khi $G\times X\to X\times X$ có ảnh
đóng (khi ấy nó là một phép nhúng đóng). Nếu $Y$ tách, thì $\mathcal L'$ tiền
rộng trên $Y$ tương đối với $S$ khi và chỉ khi nó rộng, nghĩa là khi một lũy
thừa tensor thích hợp xác định một phép nhúng xạ ảnh. Trong các bài toán moduli
được nhắc đến ở phần mở đầu, có thể chứng minh rằng quan hệ tương đương thu
được quả thật có đồ thị đóng.

\medskip
\noindent\textbf{NHẬN XÉT 8.1. ---} Ta đã giả sử
$G=\operatorname{GP}(n)_S$ để cố định ý tưởng và vì hiện nay đây là trường hợp
quan trọng nhất trong thực hành. Giả thiết hợp lý hơn đối với $G$ dường như là
$G$ thuộc một trong các «dạng» trên $S$ của một trong các nhóm Tohokû (phép
dựng các nhóm này trên các số nguyên cũng do CHEVALLEY thực hiện). Kết quả tích
cực duy nhất mà tôi biết theo hướng của giả thuyết trên là như sau: cho $X$ là
\emph{một lược đồ affine trên một trường $k$ có đặc trưng $0$, trên đó nhóm
$\operatorname{Gl}(n)_k$ hoặc $\operatorname{GP}(n-1)_k$ tác động tự do. Khi
đó quan hệ tương đương do $G$ xác định là hữu hiệu, thương $X/G$ cũng affine,
và cấu xạ $X\to X/G$ là phẳng và toàn ánh}. Chứng minh sử dụng sự kiện sau
(hiện chỉ được chứng minh trong đặc trưng $0$): khi cho $G$ tác động trên vành
tọa độ affine $A$ của $G$, xem như một không gian vectơ trên $k$, biểu diễn tầm
thường của $G$ chỉ xuất hiện một lần (trong một dãy hợp thành của một không
gian vectơ con hữu hạn chiều trên $k$, ổn định dưới $G$). Có vẻ như việc sử
dụng có hệ thống
\src{118}
lý thuyết biểu diễn tuyến tính của $G$ rốt cuộc sẽ cho một chứng minh của
giả thuyết, ít nhất khi cơ sở là một trường. Khi cơ sở không còn là một trường,
diễn giả không biết gì cả.

\section*{TÀI LIỆU THAM KHẢO}
\addcontentsline{toc}{section}{Tài liệu tham khảo}

\begin{description}
\item[{[1]}] GROTHENDIECK (Alexander). --- Techniques de descente et théorèmes
d'existence en géométrie algébrique, I: Généralités, Séminaire Bourbaki,
1959/60, n\textsuperscript{o} 190, 29 p.

\item[{[2]}] GROTHENDIECK (Alexander). --- Techniques de descente et théorèmes
d'existence en géométrie algébrique, II: Le théorème d'existence en théorie
formelle des modules, Séminaire Bourbaki, 1959/60,
n\textsuperscript{o} 195, 22 p.

\item[{[3]}] GROTHENDIECK (A.) và DIEUDONNÉ (J.). --- Éléments de géométrie
algébrique, I: Le langage des schémas. --- Paris, Presses universitaires de
France, 1960 (Institut des hautes Études scientifiques, Publications
mathématiques, 4).
\end{description}

\par

\clearpage
\chapter*{Đính chính cho Báo cáo số~212}
\addcontentsline{toc}{chapter}{Đính chính cho Báo cáo số~212}
\setcounter{footnote}{2}
\sid{e212}
\markboth{Đính chính 212}{Đính chính 212}
\src{301}

\noindent\textbf{5. TDTE III: Các tiền lược đồ thương.}

\begin{center}
[Séminaire Bourbaki, tập 13, 1960/61, số 212, 20 trang.]
\end{center}

\noindent\textbf{Trang 15, khẳng định (i).} --- Việc $Y=X/R$ là tựa xạ ảnh
trên $S$ dường như, cho đến nay, chỉ được chứng minh trong trường hợp $R$
phát sinh từ một quan hệ tương đương.

\noindent\textbf{Trang 18 và 19.} --- Như chúng tôi chỉ ra ở cuối báo cáo kế
tiếp, phỏng đoán đang xét hoàn toàn sai. ``Sự kiện khẳng định'' nêu trong Nhận
xét 8.1 dường như đã được nhiều tác giả chứng minh đồng thời (NAGATA,
ROSENLICHT, GROTHENDIECK, \ldots).

Xin lưu ý rằng việc áp dụng trong TDTE V lý thuyết chuyển sang thương được
phát triển ở đây vào việc xây dựng các lược đồ Picard cũng có thể được thay
thế bằng việc sử dụng thích hợp các lược đồ Hilbert (xem SHMT). Như đã nói
trong §~8, thiếu sót quan trọng nhất của lý thuyết được trình bày

\src{302}
ở đây là chưa có một tiêu chuẩn tồn tại cho các thương theo một quan hệ tương
đương không proper, chẳng hạn các quan hệ tương đương phát sinh từ một số tác
động của nhóm xạ ảnh. Một định lý quan trọng theo hướng này đã được MUMFORD
chứng minh\footnote{MUMFORD (David). --- An elementary theorem in geometric
invariant theory, \emph{Bull. Amer. Math. Soc.}, tập 67, 1961, tr. 483--486.}.
Về một sự tinh chỉnh kết quả của ông và nhiều ứng dụng vào lý thuyết, xem
SHMT.

\par

\unit{Thuyết trình 221 -- Các lược đồ Hilbert}
\sid{221}
\markboth{Thuyết trình 221}{Thuyết trình 221}
\src{249}
\seminar{Năm thứ 13, 1960/61, số 221\\Tháng 5 năm 1961}
\paperhead{CÁC KỸ THUẬT XÂY DỰNG VÀ CÁC ĐỊNH LÝ TỒN TẠI TRONG HÌNH HỌC ĐẠI SỐ\\
IV. CÁC LƯỢC ĐỒ HILBERT}%
  {Alexander GROTHENDIECK}

\subsection*{Mở đầu}
\addcontentsline{toc}{subsection}{Mở đầu}

Các kỹ thuật được trình bày trong [2], I và II, về cơ bản không phụ thuộc vào
bất kỳ giả thiết xạ ảnh nào đối với các lược đồ được xét. Tiếc rằng hiện nay
chúng chưa cho phép giải các bài toán tồn tại được đặt ra trong Exposé II của
[2]. Trong exposé này và exposé kế tiếp, chúng tôi sẽ giải các bài toán ấy dưới
các giả thiết xạ ảnh. Các kỹ thuật được sử dụng mang tính xạ ảnh điển hình và
hầu như không dùng đến các kết quả của I và II trong [2]. Ở đây, chúng tôi sẽ
xây dựng các «lược đồ Hilbert», nhằm thay thế việc sử dụng các tọa độ Chow, như
đã nói trong [2], II, số 2. Trong exposé kế tiếp, lý thuyết phép lấy thương đối
với các lược đồ được phát triển trong III, kết hợp với lý thuyết các lược đồ
Hilbert, sẽ cho phép chúng tôi, chẳng hạn, xây dựng các lược đồ Picard (được
định nghĩa trong [2], II, số 3) dưới những điều kiện khá tổng quát.

Tóm lại, có thể nói rằng hiện nay ta đã có một kỹ thuật xây dựng xạ ảnh khá
thỏa đáng, ngoại trừ việc vẫn còn thiếu\footnote{Xem phần bổ sung ở cuối
exposé.} một định lý về phép lấy thương theo các nhóm như nhóm xạ ảnh, tác động
«không có điểm bất động» (xem [2], III, số 8). Tình hình dường như còn thuận lợi
hơn đôi chút trong hình học giải tích (khi ta chỉ xét các không gian giải tích
«xạ ảnh» trên một không gian giải tích đã cho), bởi vì đối với các không gian
giải tích, khó khăn của phép lấy thương theo một nhóm tác động tốt sẽ biến mất.
Mặt khác, trong hình học đại số cũng như trong hình học giải tích, vẫn còn phải
phát triển một kỹ thuật xây dựng có hiệu lực mà không cần các giả thiết xạ ảnh.

\fsection{1}{Các họ bị chặn của các bó: các tính chất ổn định}

Cho $k$ là một trường và $X$ là một $k$-tiền lược đồ, giả sử thuộc kiểu hữu hạn
để đơn giản. Với mỗi mở rộng $K/k$, ta do đó thu được một $K$-tiền lược đồ
$X_K=X\otimes_k K$. Nếu $\mathcal F$ là một bó coherent trên $X_K$ và nếu
$K'$ là một mở rộng của $K$, thì
$\mathcal F\otimes_K K'=\mathcal F_{K'}$ là một bó quasi-coherent trên
$X_K\otimes_K K'=X_{K'}$. Như vậy, nếu $K$ và $K'$ là hai mở rộng tùy ý của
$k$, $\mathcal F$ là một bó quasi-coherent trên $X_K$ và $\mathcal F'$
\src{250}
là một bó quasi-coherent trên $X_{K'}$, ta nói $\mathcal F$ và
$\mathcal F'$ là \emph{tương đương} nếu tồn tại các $k$-đồng cấu từ $K$,
$K'$ vào cùng một mở rộng $K''$ của $k$ sao cho $\mathcal F_{K''}$ và
$\mathcal F'_{K''}$ đẳng cấu trên $X_{K''}$. Đây là một quan hệ tương đương;
ta sẽ quan tâm đến các lớp tương đương của các bó theo quan hệ này và đến các
tập hợp lớp tương đương. Lưu ý rằng, nếu $X_0$ thuộc kiểu hữu hạn trên
$k$, thì mỗi lớp bó coherent có thể được xác định bởi một bó coherent trên
một $X_K$, trong đó $K$ là một mở rộng hữu hạn sinh của $k$. Vì vậy, trong
định nghĩa các lớp bó coherent, ta có thể chỉ xét các mở rộng đóng đại số của
$k$; cũng có thể chỉ xét một mở rộng đóng đại số cố định $\Omega$ của $k$, có
bậc siêu việt vô hạn. Khi ấy, hai bó coherent $\mathcal F$ và $\mathcal F'$
trên $X_\Omega$ tương đương khi và chỉ khi tồn tại một $k$-tự đẳng cấu $\sigma$
của $\Omega$ sao cho $\mathcal F\otimes_\Omega(\Omega,\sigma)$ đẳng cấu với
$\mathcal F'$. Do đó có một tương ứng song ánh giữa các lớp bó coherent theo
một trong hai định nghĩa.

Cho $E$, $E'$ là hai tập hợp các lớp bó coherent trên $X$. Xét các lớp của mọi
bó có dạng $\mathcal F\otimes\mathcal F'$, trong đó $\mathcal F$ và
$\mathcal F'$ là các bó coherent trên cùng một $X_K$, lớp của $\mathcal F$
thuộc $E$ và lớp của $\mathcal F'$ thuộc $E'$. Ta thu được một tập hợp các lớp
bó coherent, ký hiệu là $E\otimes E'$. Tương tự, ta định nghĩa các
$\operatorname{Tor}_i(E,E')$, v.v. Tổng quát hơn, với mỗi hàm $U$ gán cho mọi
dãy $\mathcal F_1,\ldots,\mathcal F_n$ gồm $n$ bó coherent trên cùng một $X_K$
một tập hợp $U(\mathcal F_1,\ldots,\mathcal F_n)$ các bó coherent trên $X_K$,
và có tính chất tương thích hiển nhiên với các đẳng cấu bó và các ảnh ngược qua
phép đổi cơ sở, ta liên kết một hàm, cũng ký hiệu là $U$, gán cho mọi dãy gồm
$n$ tập hợp $E_1,\ldots,E_n$ các lớp bó coherent một tập hợp
$U(E_1,\ldots,E_n)$ các lớp bó coherent.

Mục đích của chúng ta trong mục này là định nghĩa một số tập hợp lớp bó, gọi là
\emph{bị chặn}, và chứng minh rằng các phép toán $U$ thông dụng nhất, khi áp
dụng cho các tập hợp bị chặn, vẫn cho các tập hợp bị chặn.

Cho $X$ là một tiền lược đồ thuộc kiểu hữu hạn trên một lược đồ Noether $S$.
Với mỗi $s\in S$, thớ $X_s$ là một tiền lược đồ thuộc kiểu hữu hạn trên
$k(s)$, và ta sẽ xét các lớp bó coherent trên $X_s$ theo nghĩa trên. Điều này
mang lại nghĩa cho cụm từ: lớp bó coherent trên một thớ của $X/S$, cũng như
cho các cụm từ tương tự. Cũng vậy, khi tiến hành riêng trên từng thớ, ta vẫn có
thể xét các phép toán như $E\otimes E'$, v.v., gán cho mỗi hệ gồm những tập hợp
\src{251}
các lớp bó coherent trên các thớ của $X/S$ một tập hợp khác gồm các lớp bó
coherent trên các thớ của $X/S$.

\result{ĐỊNH NGHĨA}{1.1} Cho $E$ là một tập hợp các lớp bó coherent trên các
thớ của $X/S$. Ta nói $E$ là \emph{bị chặn} nếu tồn tại một tiền lược đồ $S'$
thuộc kiểu hữu hạn trên $S$ và một bó coherent $\mathcal F'$ trên
$X'=X\times_S S'$ sao cho $E$ được chứa trong tập hợp các lớp bó trên các thớ
của $X/S$ do $\mathcal F'$ xác định.

Theo định nghĩa, tập hợp cuối cùng này gán cho mỗi $s\in S$ các lớp của những bó
$\mathcal F'\otimes_{S'}k(s')$, trong đó $s'$ chạy qua các điểm của $S'$ nằm
trên $s$ (do đó $k(s')$ là một mở rộng của $k(s)$, và
$X_s\otimes_{k(s)}k(s')$ được đồng nhất với thớ
$X'\otimes_{S'}k(s')=X'_{s'}$ của $X'$ tại $s'$). Có thể nói rằng các họ bị
chặn là những họ được chứa trong một \emph{họ đại số} các bó coherent, được
tham số hóa bởi một $S'$ thuộc kiểu hữu hạn trên $S$.

Một hợp hữu hạn của các họ bị chặn là bị chặn (lấy tiền lược đồ tổng rời nhau
của các tiền lược đồ tham số $S_i$ xác định những họ đại số bao chứa). Một phép
đổi cơ sở $T\longrightarrow S$ biến một họ bị chặn tương đối với $X/S$ thành
một họ bị chặn tương đối với $X_T/T$; mệnh đề đảo đúng nếu
$T\longrightarrow S$ là toàn ánh (hoặc, tổng quát hơn, nếu ảnh của nó chứa các
$s$ thực sự xuất hiện trong họ $E$ đã cho đối với $X/S$). Về lý thuyết, điều
này quy việc xác định các họ bị chặn về trường hợp $S$ là phổ của một đại số
thuộc kiểu hữu hạn trên vành số nguyên $\mathbf Z$.

Nếu $E$ và $E'$ là các họ bị chặn gồm các lớp bó tương đối với $X/S$, thì
$E\otimes E'$ cũng bị chặn. Thật vậy, giả sử $E$ và $E'$ được chứa tương ứng
trong các họ đại số xác định bởi $T\longrightarrow S$ và $\mathcal F$ trên
$X_T$, bởi $T'\longrightarrow S$ và $\mathcal F'$ trên $X_{T'}$. Khi đó,
$E\otimes E'$ được chứa trong họ đại số xác định bởi
$T''=T\times_S T'\longrightarrow S$ và bó $\mathcal F''$ trên $X_{T''}$, là
tích tensor của các ảnh ngược của $\mathcal F$ và $\mathcal F'$ trên
$X_{T''}$. Lập luận này chỉ đúng bởi vì hàm tử
$\mathcal F\otimes\mathcal F'$ theo các biến $\mathcal F$, $\mathcal F'$ là
khớp phải, nên giao hoán với mở rộng cơ sở (và đặc biệt với phép chuyển sang
các thớ). Nó không áp dụng trực tiếp cho các phép toán địa phương như
$\operatorname{Tor}_i(E,E')$, $\operatorname{Hom}(E,E')$,
$\operatorname{Ext}^i(E,E')$. Tuy nhiên, có thể chứng minh rằng các phép toán
này vẫn biến các tập hợp bị chặn thành các tập hợp bị chặn, bằng cách tiến hành
như đối với $E\otimes E'$, nhưng sử dụng thêm các kết quả thuộc kiểu sau đây
(tất cả đều có trong [3], IV, 6.11): một họ bị chặn $E$ luôn được chứa trong
một họ đại số xác định bởi một bó coherent $\mathcal F$ trên một $X_T$ ($T$
thuộc kiểu hữu hạn trên $S$), phẳng trên $T$. (Ta «cắt thành từng mảnh» không
gian tham số ban đầu.) Những
\src{252}
tính chất phẳng của các bó thích hợp quả thực bảo đảm rằng các phép toán như
$\operatorname{Tor}_i(\mathcal F,\mathcal F')$ giao hoán với mọi phép đổi cơ
sở. Cùng phương pháp ấy áp dụng cho các phép toán mang tính toàn cục: ảnh trực
tiếp và ảnh trực tiếp dẫn xuất của các bó coherent qua các cấu xạ riêng, các
$\operatorname{Ext}$ toàn cục tương đối với các cấu xạ riêng (xem [5], III,
\S~6), v.v.; mọi phép toán này đều biến các họ bị chặn của những bó thành các
họ bị chặn của những bó. (Lưu ý. -- Ở đây những tiền lược đồ mà trên đó ta lấy
các bó khác nhau có thể thay đổi qua các phép toán thuộc kiểu đang xét.)

Hai phát biểu sau được chứng minh về cơ bản bằng cùng kỹ thuật về tính phẳng;
đối với sự phân tích nguyên sơ trên các thớ của một cấu xạ thuộc kiểu hữu hạn,
xem riêng [5], IV.

\result{MỆNH ĐỀ}{1.2} Cho $E$, $E'$ là những tập hợp bị chặn gồm các lớp bó
trên các thớ của $X/S$, giả sử $X$ riêng trên $S$. Khi đó:

\medskip
\noindent (i) Họ các hạt nhân, đối hạt nhân và ảnh của những đồng cấu
$\mathcal F\longrightarrow\mathcal F'$, trong đó lớp của $\mathcal F$ thuộc
$E$ và lớp của $\mathcal F'$ thuộc $E'$, là bị chặn.

\smallskip
\noindent (ii) Họ các bó $\mathcal F''$ là những mở rộng của một $\mathcal F$
bởi một $\mathcal F'$, trong đó lớp của $\mathcal F$ thuộc $E$ và lớp của
$\mathcal F'$ thuộc $E'$, là bị chặn.

Sau một phép đổi cơ sở thích hợp, có thể giả sử $E$ và $E'$ lần lượt được xác
định bởi một bó coherent $\mathcal G$ và một bó coherent $\mathcal G'$ trên
một $X_T/T$, với $T$ thuộc kiểu hữu hạn trên $S$. Hơn nữa, có thể giả sử một
số giả thiết về tính phẳng được thỏa mãn, kéo theo việc tạo thành các bó
$f_{T*}(\operatorname{Hom}_{X_T}(\mathcal G,\mathcal G'))$ và
$\operatorname{Ext}^{1}_{f_T}(\mathcal G,\mathcal G')$ giao hoán với phép đổi
cơ sở theo một cấu xạ tùy ý $T'\longrightarrow T$. Cũng có thể giả sử các bó
coherent nói trên là tự do địa phương trên $T$. Khi đó, gọi $T_0$ và $T_1$
là các phân thớ vectơ trên $T$ mà bó các mầm tiết diện của chúng lần lượt là
các bó nói trên. Khi ấy ta xác định chính tắc một đồng cấu
$\mathcal G_{T_0}\longrightarrow\mathcal G'_{T_0}$ giữa các bó coherent trên
$X_{T_0}$, và một mở rộng
\[
0\longrightarrow\mathcal G_{T_1}\longrightarrow\mathcal G''
\longrightarrow\mathcal G'_{T_1}\longrightarrow0
\]
gồm các bó coherent trên $X_{T_1}$, có một tính chất phổ dụng hiển nhiên. Bó
thứ hai này xác định một họ đại số bao chứa họ đang xét trong (ii). Điều tương
tự cũng đúng với hạt nhân, đối hạt nhân và ảnh của đồng cấu
\src{253}
trước đó, mà việc xét nó chứng minh (i) (miễn là giả sử đối hạt nhân phẳng trên
$T_0$, trường hợp mà ta lại quy về bằng cách cắt $T_0$ thành từng
mảnh\ldots).

\result{MỆNH ĐỀ}{1.3} Cho $E$ là một họ bị chặn gồm các lớp bó trên các thớ
của $X/S$. Khi đó, các lớp của các bó cấu trúc của
$(\operatorname{supp}\mathcal F)_{\mathrm{red}}$, trong đó $\mathcal F$ là một
bó coherent trên một $X_K$, với $K$ đóng đại số, và lớp của nó thuộc $E$, tạo
thành một họ bị chặn.

Ở đây, $(\operatorname{supp}\mathcal F)_{\mathrm{red}}$ ký hiệu cho giá của
$\mathcal F$ được trang bị cấu trúc rút gọn cảm sinh; tức là bó cấu trúc của nó
là thương của $\mathcal O_{X_K}$ theo bó iđêan lớn nhất xác định
$\operatorname{supp}\mathcal F$. Có thể chứng minh kết quả tương tự (1.3) cho
những bó được suy ra chính tắc từ $\mathcal F$ bằng lý thuyết phân tích nguyên
sơ, chẳng hạn các $\mathcal F/\mathcal F_i$, trong đó các $\mathcal F_i$ là
những bó con của $\mathcal F$ nguyên sơ ứng với các thành phần của giá của
$\mathcal F$ và tối tiểu đối với tính chất này; hoặc các
$\mathcal O_{X_K}/\mathfrak p$, trong đó $\mathfrak p$ là một bó iđêan nguyên
tố liên kết với $\mathcal F$; hoặc các $\mathcal O_{X_K}/\mathfrak q$, trong đó
$\mathfrak q$ là một bó iđêan nguyên sơ liên kết với một thành phần của giá của
$\mathcal F$ (trường tham chiếu được giả sử đóng đại số).

\fsection{2}{Các họ bị chặn và đa thức Hilbert}

Trong phần tiếp theo, ta giả sử $X$ xạ ảnh trên $S$ và được trang bị một bó rất
rộng, ký hiệu $\mathcal O_X(1)$. Với mọi mở rộng $K$ của trường thặng dư $k(s)$
tại một điểm $s$ của $S$, ta xét trên $X_K$ bó tương ứng
$\mathcal O_{X_K}(1)$, bó này vẫn rất rộng.

Với mỗi bó coherent $\mathcal F$ trên $X_K$, ta gán hàm
\[
P_{\mathcal F}(n)=\text{đặc trưng Euler--Poincaré của $\mathcal F(n)$
trên $X_K$},
\]
là một đa thức theo số nguyên $n$, gọi là \emph{đa thức Hilbert} của
$\mathcal F$. Với các giá trị đủ lớn của $n$, $P(n)$ không gì khác hơn là số
chiều của $H^0(X_K,\mathcal F(n))$ trên $K$, vì các
$H^i(X_K,\mathcal F(n))$ bằng không khi $i>0$ và $n$ đủ lớn.

Bây giờ, nếu $\mathcal F$ là một bó coherent trên $X$, phẳng trên $S$, thì các
đa thức Hilbert của những bó $\mathcal F_s$ cảm sinh trên các thớ $X_s$ thuộc
cùng một thành phần liên thông của $S$ đều bằng nhau [5], III, \S~7. Do đó
(không cần giả thiết phẳng), tập hợp các đa thức Hilbert của những bó
$\mathcal F_s$, $s\in S$, là hữu hạn đối với mọi bó coherent $\mathcal F$ trên
$X$.
\src{254}
Mặt khác, nhắc lại rằng nếu $\mathcal F$ là một bó coherent trên $X$, thì nó
đẳng cấu với một bó thương của bó $\mathcal O_X(-n)^N$, với $n$ và $N$ đủ
lớn. Do đó, các bó $\mathcal F_s$ cảm sinh trên các thớ cũng là thương của bó
$\mathcal O(-n)$ trên thớ.

Từ hai nhận xét này, ta suy ra chiều «chỉ khi» của định lý sau:

\result{ĐỊNH LÝ}{2.1} Cho $X$ xạ ảnh trên $S$ Noether, với
$\mathcal O_X(1)$ rất rộng trên $X$ tương đối với $S$. Cho $E$ là một tập hợp
các lớp bó trên các thớ của $X/S$. Để $E$ bị chặn, điều kiện cần là nó thỏa mãn các điều kiện sau:

\medskip
\noindent (a) Tồn tại một bó coherent $\mathcal L$ trên $X$ (có thể giả sử có
dạng $\mathcal O_X(-n)^N$) sao cho $E$ được chứa trong họ các lớp của những bó
coherent là thương của các bó dạng $\mathcal L_K$.

\smallskip
\noindent (b) Các đa thức Hilbert $P_{\mathcal F}$ của những bó $\mathcal F$
có lớp thuộc $E$ đều thuộc cùng một tập hợp hữu hạn các đa thức.

Còn phải chứng minh chiều «nếu», đây sẽ là một trường hợp riêng của một kết quả
chính xác hơn. Với mọi môđun coherent $\mathcal F$ trên một tiền lược đồ thuộc
kiểu hữu hạn trên một trường $K$, và mọi số nguyên $r$, gọi $N_r$ là môđun con
của $\mathcal F$ mà các tiết diện trên một tập mở là những tiết diện của
$\mathcal F$ trên tập mở đó có giá mang chiều $<r$. Khi đó
$N_r=\mathcal F$ nếu $r>\dim\operatorname{supp}\mathcal F$, còn $N_r=0$ nếu
$r\leq0$, và ta thu được một lọc tăng hữu hạn của $\mathcal F$ mà các thương
$N_r/N_{r+1}$ có các chu trình nguyên tố
liên kết đúng bằng những chu trình nguyên tố liên kết với $\mathcal F$ mang
chiều $r$. Ta đặt
\[
\mathcal F_{(r)}=\mathcal F/N_r,
\]
thì các chu trình nguyên tố liên kết với $\mathcal F_{(r)}$ đúng bằng những chu
trình nguyên tố liên kết với $\mathcal F$ mang chiều $\geq r$; đặc biệt, nó
bằng $\mathcal F$ khi và chỉ khi các chu trình nguyên tố liên kết với
$\mathcal F$ mang chiều $\geq r$. Với ký hiệu này:

\result{ĐỊNH LÝ}{2.2} Dưới các giả thiết của (2.1), cho $s$ là một số nguyên
và giả sử $E$ thỏa mãn điều kiện (a) cùng dạng yếu hơn sau đây của (b):

\medskip
\noindent $(b_s)$ Các đa thức Poincaré $P_{\mathcal F}$ của những bó
$\mathcal F$ có lớp thuộc $E$ có các hệ số ở bậc $\leq s-1$ được chặn.

Dưới các điều kiện này, các bó $\mathcal F_{(s)}$ (khi lớp của $\mathcal F$
chạy trong $E$) tạo thành một họ bị chặn. Hơn nữa, các hệ số bậc $s-2$ của
các $P_{\mathcal F}$ bị chặn dưới.
\src{255}
Do đó:

\result{HỆ QUẢ}{2.3} Giả sử những bó $\mathcal F$ có lớp thuộc $E$ sao cho
mọi chu trình nguyên tố liên kết của chúng đều mang chiều $d$ thỏa mãn
$s\leq d\leq r$. Khi đó, trong (2.1) (b), chỉ cần xét các hệ số của
$P_{\mathcal F}$ ở các bậc từ $s-1$ đến $r$.

Phần còn lại của mục này dành cho việc phác thảo chứng minh (2.2). Các bổ đề
then chốt là hai bổ đề sau, trong đó bổ đề thứ nhất đã quen thuộc (và tóm tắt
nội dung toán học hữu dụng của các tọa độ Chow):

\result{BỔ ĐỀ 2.4 (CHOW)}{} Xét các bó cấu trúc của những lược đồ con $Y$ của
các thớ $X_K$ ($K$ chạy qua các mở rộng đóng đại số của những trường thặng dư
của $S$), với $Y$ rút gọn và mọi thành phần của nó cùng mang chiều $r$
($\mathcal O_Y$ được xem như một bó thương của $\mathcal O_X$). Nếu các bậc
của $Y$ bị chặn, thì các $Y$ tạo thành một họ bị chặn.

Ở đây, bậc $a$ của $Y$ có thể được định nghĩa thuận tiện nhất bằng hệ số của
số hạng đầu của
\[
P_{\mathcal O_Y}=a\frac{n^r}{r!}+\cdots.
\]

\result{BỔ ĐỀ}{2.5} Cho $\mathcal L$ là một bó coherent trên $X$ và $E$ là
một tập hợp các lớp bó thương $\mathcal F$ của những bó $\mathcal L_K$ ($K$
là một mở rộng trường thặng dư của $S$). Giả sử các thớ của $X$ trên $S$ mang
chiều $\leq r$, và đặt
\[
P_{\mathcal F}(n)=a_{\mathcal F}\frac{n^r}{r!}
+b_{\mathcal F}\frac{n^{r-1}}{(r-1)!}
+\text{các số hạng bậc $<r-1$}.
\]
Khi đó, hệ số $a_{\mathcal F}$ bị chặn, còn $b_{\mathcal F}$ bị chặn dưới.
Nếu $b_{\mathcal F}$ bị chặn, thì họ các $\mathcal F_{(r)}$ bị chặn.

\emph{Chứng minh.} -- Sau khi thay $S$, nếu cần, bằng một hợp các lược đồ con
của $S$ phủ $S$, ta có thể giả sử tồn tại một cấu xạ hữu hạn
$f:X\longrightarrow\mathbf P^r_S$ sao cho $\mathcal O_X(1)$ đẳng cấu với ảnh
ngược của $\mathcal O_{\mathbf P^r_S}(1)$; do đó, với mọi bó coherent
$\mathcal F$ trên $X$, ta có
$P_{\mathcal F}=P_{f_*(\mathcal F)}$. Mặt khác, dễ kiểm tra (bằng kỹ thuật của
mục trước) rằng một tập hợp các bó $\mathcal F$ trên $X$ bị chặn khi và chỉ khi
tập hợp các $f_*(\mathcal F)$ bị chặn. Cuối cùng, ta có
\[
f_*(\mathcal F)_{(r)}=f_*(\mathcal F_{(r)}).
\]
\src{256}
Như vậy, ta quy về trường hợp $X=\mathbf P^r_S$. Hơn nữa, có thể giả sử
$\mathcal L=\mathcal O_{\mathbf P^r_S}(k)^s$ với $k$, $s$ thích hợp. Hệ số
$a_{\mathcal F}$ thỏa mãn
\[
0\leq a_{\mathcal F}\leq s
\]
và do đó bị chặn. Mặt khác, việc các $P_{\mathcal F}(n)$ có hệ số của
$n^{r-1}$ bị chặn dưới (tương ứng, bị chặn) tương đương với việc khẳng định
điều đó cho
$P_{\mathcal F}(n-k)=P_{\mathcal F(-k)}(n)$. Như vậy, ta quy về trường hợp
\[
\mathcal L=\mathcal O_{\mathbf P^r_S}^{s}.
\]

Xét dãy khớp
\[
0\longrightarrow N_r\longrightarrow\mathcal F
\longrightarrow\mathcal F_{(r)}\longrightarrow0,
\]
từ đó
\[
P_{\mathcal F}=P_{\mathcal F_{(r)}}+P_{N_r},
\]
và vì hệ số của $n^{r-1}$ trong $P_{N_r}$ không âm (do
$\dim\operatorname{supp}N_r\leq r-1$), ta có
\[
b_{\mathcal F_{(r)}}\leq b_{\mathcal F}.
\]
Điều này cho phép ta, để chứng minh các khẳng định của bổ đề, thay
$\mathcal F$ bằng $\mathcal F_{(r)}$, tức là giả sử những bó thương
$\mathcal F$ của $\mathcal L$ đang xét là không xoắn.

Vì $\mathbf P^r_K$ chuẩn tắc, suy ra $\mathcal F$ tự do địa phương hạng
$a=a_{\mathcal F}$ trên một tập mở $U=\mathbf P^r_K-Y$, trong đó $Y$ có đối
chiều $\geq2$. Do đó, $\bigwedge^a\mathcal F$ là một bó trên
$\mathbf P^r_K$ mà hạn chế của nó lên $U$ là khả nghịch; vì vậy (do
$\mathbf P^r_K$ chính quy và $Y$ có đối chiều $\geq2$), hạn chế này đẳng cấu
với hạn chế của một bó khả nghịch trên $\mathbf P^r_K$, được xác định sai khác
một đẳng cấu. Bó sau có dạng $\mathcal O_{\mathbf P^r_K}(d)$ với một số nguyên
$d$ được xác định duy nhất. Vì $\bigwedge^a\mathcal F$ là thương của
$\bigwedge^a\mathcal O_{\mathbf P^r_K}^{n}\simeq
\mathcal O_{\mathbf P^r_K}^{N}$, với $N=\binom{s}{a}$, nó có $N$ tiết diện chính tắc, do đó xác
định các tiết diện của
\src{257}
$\mathcal O_{\mathbf P^r_K}(d)$ trên $U$, là những hạn chế của các tiết diện
$s_i$ ($1\leq i\leq N$) của $\mathcal O_{\mathbf P^r_K}(d)$ (vì
$\mathbf P^r_K$ chuẩn tắc và $Y$ có đối chiều $\geq2$). Các tiết diện này sinh
$\mathcal O_{\mathbf P^r_K}(d)$ tại các điểm của $U$, nên chúng không đồng thời
bằng không; điều này kéo theo $d\geq0$. Mặt khác, một phép tính đơn giản cho
\[
b_{\mathcal F}=a_{\mathcal F}(r+1)/2+d.
\]
Đặc biệt, điều này cho thấy $b_{\mathcal F}>0$, nên $b_{\mathcal F}$ bị chặn
dưới. Nó bị chặn khi và chỉ khi $d$ bị chặn; ta chứng minh rằng khi đó
$\mathcal F$ thuộc một họ bị chặn. Có thể giả sử $a_{\mathcal F}$ và
$b_{\mathcal F}$ cố định, lần lượt bằng $a$ và $b$, nên $d$ cố định. Việc biết
$N$ tiết diện $s_i$ của $\mathcal O_{\mathbf P^r_K}(d)$, tức là biết một đồng
cấu
\[
s:\bigwedge^a\mathcal L_K\longrightarrow\mathcal O_{\mathbf P^r_K}(d),
\]
cho phép khôi phục $\mathcal F$ như đối ảnh của đồng cấu hợp thành tương ứng:
\[
\mathcal L_K\longrightarrow
\operatorname{Hom}\!\left(\bigwedge^{a-1}\mathcal L_K,
\bigwedge^a\mathcal L_K\right)\longrightarrow
\operatorname{Hom}\!\left(\bigwedge^{a-1}\mathcal L_K,
\mathcal O_{\mathbf P^r_K}(d)\right),
\]
trong đó đồng cấu thứ nhất là đồng cấu chính tắc sinh bởi tích ngoài, còn đồng
cấu thứ hai được suy ra từ $s$. Khi ấy, kết luận suy ra từ (1.2) (i).

Kết hợp hai bổ đề trước, ta chứng minh được:

\result{BỔ ĐỀ}{2.6} Dưới các giả thiết sơ bộ của (2.1), giả sử với mọi
$\mathcal F$ ta có
\[
P_{\mathcal F}(n)=a_{\mathcal F}\frac{n^r}{r!}
+b_{\mathcal F}\frac{n^{r-1}}{(r-1)!}
+\text{các số hạng bậc $<r-1$},
\]
và các hệ số $a_{\mathcal F}$ bị chặn. Khi đó, các hệ số $b_{\mathcal F}$ bị
chặn dưới. Nếu các $b_{\mathcal F}$ bị chặn, thì các $\mathcal F_{(r)}$ tạo
thành một họ bị chặn.

Có thể giả sử các trường cơ sở $K$ của những bó $\mathcal F$ đều đóng đại số.
Trang bị cho mỗi $\operatorname{supp}\mathcal F_{(r)}$, là hợp của các thành
phần mang bậc $r$
của $\operatorname{supp}\mathcal F$, cấu trúc rút gọn cảm sinh. Khi đó, các
bậc của $\operatorname{supp}\mathcal F_{(r)}$ bị chặn trên bởi $a$; do đó,
theo (2.4), các $\operatorname{supp}\mathcal F_{(r)}$ tạo thành một tập hợp bị
chặn. Hơn nữa, với mỗi thành phần của
$\operatorname{supp}\mathcal F_{(r)}$, độ dài của $\mathcal F_{(r)}$ dọc theo
thành phần đó là $\leq a$; vì vậy, nếu $\mathcal J_{\mathcal F}$ là iđêan xác
định $\operatorname{supp}\mathcal F_{(r)}$, thì $\mathcal F_{(r)}$
\src{258}
có thể được xem như một môđun trên lược đồ con $Y_{\mathcal F}$ của $X$ được
xác định bởi $\mathcal J_{\mathcal F}^{a}$. Như trong bổ đề trước, ta cũng quy
về trường hợp $\mathcal F=\mathcal F_{(r)}$, do đó $\mathcal F$ xuất phát từ
một môđun trên $Y_{\mathcal F}$. Các $Y_{\mathcal F}$ tương ứng với một họ bị
chặn gồm các môđun thương của các $\mathcal O_{X_K}$, và vì thế xuất phát từ
một lược đồ con đóng $Y$ của một lược đồ $X\times_S T$. Khi đó, ta có thể áp
dụng (2.5) cho $Y/T$ và $\mathcal L\otimes_XY$, suy ra kết luận.

Bây giờ ta có thể chứng minh (2.2) bằng quy nạp theo cận trên $r$ của các số
$\dim\operatorname{supp}\mathcal F$. Mệnh đề là hiển nhiên với $r<0$; vì vậy,
giả sử $r\geq0$ và mệnh đề đã được chứng minh với mọi $r'<r$. Theo (2.6), các
$\mathcal F_{(r)}$ tạo thành một tập hợp bị chặn; do đó, theo (1.2) (i), các
hạt nhân của những đồng cấu
$\mathcal L_K\longrightarrow\mathcal F_{(r)}$ cũng tạo thành một tập hợp bị
chặn. Vì thế, tồn tại một môđun coherent $\mathcal L'$ trên $X$ sao cho các hạt
nhân đang xét, và do đó cả các
\[
N_r(\mathcal F)=\operatorname{Ker}(\mathcal F\longrightarrow\mathcal F_{(r)}),
\]
đều là những môđun thương của các môđun $\mathcal L'_K$. Vì các
$\mathcal F_{(r)}$ bị chặn, các $P_{\mathcal F_{(r)}}$ vẫn bị chặn, và công
thức
\[
P_{\mathcal F}=P_{\mathcal F_{(r)}}+P_{N_r}
\]
cho thấy các $P_{N_r}$ thỏa mãn cùng điều kiện $(b_s)$ như các
$P_{\mathcal F}$. Do đó, các $N_r$ thỏa mãn các điều kiện (a) và $(b_s)$, và
theo giả thiết quy nạp, các $(N_r)_{(s)}$ vẫn bị chặn. Mặt khác,
$\mathcal F_{(s)}$ là một mở rộng của $\mathcal F_{(r)}$ bởi
$(N_r)_{(s)}$; vì thế, theo (1.2) (ii), các $\mathcal F_{(s)}$ vẫn bị chặn.
Đối với khẳng định cuối cùng của (2.2), ta lưu ý rằng các hạt nhân $N_s$ của
$\mathcal F\longrightarrow\mathcal F_{(s)}$ vẫn bị chặn theo (1.2) (i), và hệ
số của số hạng chứa $n^{s-1}$ trong $P_{N_s}$ vẫn bị chặn; khi đó Bổ đề 2.6
chứng minh rằng hệ số của số hạng kế tiếp bị chặn dưới. Điều này kết thúc chứng
minh.

\fsection{3}{Các lược đồ Hilbert: định nghĩa, định lý tồn tại}

Cho $X$ là một tiền lược đồ trên một tiền lược đồ khác $S$, và $\mathcal F$ là
một môđun quasi-coherent trên $X$. Ký hiệu bởi
\[
\operatorname{Quot}(\mathcal F/X/S)
\]
tập hợp các môđun quasi-coherent thương của $\mathcal F$ phẳng trên $S$. Bây
giờ, cho $S'\longrightarrow S$ là một cấu xạ đổi cơ sở; đặt
$X'=X\times_S S'$ và
$\mathcal F'=\mathcal F\otimes_{\mathcal O_S}\mathcal O_{S'}$; do đó, $X'$
là một tiền lược đồ trên $S'$ được trang bị một môđun
\src{259}
quasi-coherent $\mathcal F'$, và ta có thể xét
$\operatorname{Quot}(\mathcal F'/X'/S')$. Ta đặt
\[
\underline{\operatorname{Quot}}_{\mathcal F/X/S}(S')
=\operatorname{Quot}(\mathcal F'/X'/S')
\qquad\text{(trong đó $X'=X\times_SS'$)}.
\]
Bây giờ, nếu $S''\longrightarrow S'$ là một $S$-cấu xạ, thì
$X''=X\times_SS''$ đẳng cấu với $X'\times_{S'}S''$, còn $\mathcal F''$ đẳng
cấu với $\mathcal F'\otimes_{\mathcal O_{S'}}\mathcal O_{S''}$. Vì hàm tử ảnh
ngược
\[
\mathcal G'\longmapsto
\mathcal G'\otimes_{\mathcal O_{S'}}\mathcal O_{S''}
\]
từ phạm trù các môđun quasi-coherent trên $X'$ vào phạm trù các môđun
quasi-coherent trên $X''$ là khớp phải và biến các môđun phẳng trên $S'$ thành
các môđun phẳng trên $S''$, ta thu được một ánh xạ tự nhiên
\[
\operatorname{Quot}(\mathcal F'/X'/S')\longrightarrow
\operatorname{Quot}(\mathcal F''/X''/S''),
\]
do đó $\underline{\operatorname{Quot}}_{\mathcal F/X/S}(S')$ là một hàm tử
phản biến theo $S'$ (một tiền lược đồ trên $S$), nhận giá trị trong phạm trù
các tập hợp. Từ đây về sau, ta giả sử $X$ xạ ảnh trên một $S$ Noether,
$\mathcal F$ coherent, và để đơn giản, ta chỉ xét các $S'$ trên $S$ là
Noether địa phương.

\result{ĐỊNH LÝ}{3.1} Dưới các điều kiện này, hàm tử phản biến
$\underline{\operatorname{Quot}}_{\mathcal F/X/S}$ trên phạm trù các
$S$-tiền lược đồ Noether địa phương được biểu diễn bởi một $S$-tiền lược đồ
$\operatorname{Quot}_{\mathcal F/X/S}$, là hợp rời nhau của một dãy các
$S$-lược đồ xạ ảnh (đặc biệt,
$\operatorname{Quot}_{\mathcal F/X/S}$ thuộc kiểu hữu hạn địa phương trên
$S$).

Ta sẽ thu được một phân rã như vậy theo cách sau. Cho $\mathcal O_X(1)$ là một
bó khả nghịch trên $X$, rất rộng tương đối với $S$. Với mỗi đa thức $P(n)$ có
các hệ số hữu tỉ, ký hiệu $\operatorname{Quot}^P(\mathcal F/X/S)$ là tập con
của $\operatorname{Quot}(\mathcal F/X/S)$ gồm các môđun thương coherent
$\mathcal G$ của $\mathcal F$, phẳng trên $S$, và có đa thức Hilbert tại mọi
$s\in S$ bằng $P$. Khi đó ta đặt
\[
\underline{\operatorname{Quot}}^P_{\mathcal F/X/S}(S')
=\operatorname{Quot}^P(\mathcal F'/X'/S')
\]
và như vậy thu được một hàm tử con của
$\underline{\operatorname{Quot}}_{\mathcal F/X/S}$. Tính bất biến của các đa
thức Hilbert được nhắc lại trong mục 2 kéo theo điều sau: để
$\underline{\operatorname{Quot}}_{\mathcal F/X/S}$ biểu diễn được, điều kiện
cần và đủ là các $\underline{\operatorname{Quot}}^P_{\mathcal F/X/S}$ đều
biểu diễn được; khi đó, $S$-tiền lược đồ
$\operatorname{Quot}_{\mathcal F/X/S}$ biểu diễn nó đẳng cấu với tiền lược đồ
là hợp rời nhau của các $\operatorname{Quot}^P_{\mathcal F/X/S}$ biểu diễn
những hàm tử $\underline{\operatorname{Quot}}^P_{\mathcal F/X/S}$. Vì vậy,
(3.1) là một hệ quả của định lý sau:
\src{260}
\result{ĐỊNH LÝ}{3.2} Với các ký hiệu trên, hàm tử
$\underline{\operatorname{Quot}}^P_{\mathcal F/X/S}$ được biểu diễn bởi một
$S$-tiền lược đồ xạ ảnh $\operatorname{Quot}^P_{\mathcal F/X/S}$.

Phần còn lại của mục này dành cho việc chứng minh (3.2).

Cho $\nu$ là một số nguyên. Với mọi $S'$ trên $S$, ta ký hiệu $A_\nu(S')$ là
tập hợp các môđun thương $\mathcal G=\mathcal F'/\mathcal K$ của
$\mathcal F'=\mathcal F\otimes_{\mathcal O_S}\mathcal O_{S'}$, coherent,
phẳng trên $S'$, và thỏa mãn các điều kiện sau:
\begin{align*}
\text{a.}\quad &R^if'_*(\mathcal G(n))=0 &&\text{với $i>0$ và $n\geq\nu$},\\
\text{b.}\quad &R^if'_*(\mathcal K(n))=0 &&\text{với $i>0$ và $n\geq\nu$},\\
\text{c.}\quad &f'_*(\mathcal K(\nu+k))
=\mathcal S'_k f'_*(\mathcal K(\nu)) &&\text{với $k\geq0$}.
\end{align*}
Đối với quan hệ cuối cùng này, ta giả sử $X$ được viết thành phổ nguyên tố
thuần nhất của một đại số phân bậc quasi-coherent $\mathcal S_*$ trên $S$, có
các bậc dương và được sinh bởi $\mathcal S_1$, rồi đặt
$\mathcal S'_* =\mathcal S_*\otimes_{\mathcal O_S}\mathcal O_{S'}$, sao cho
$X'$ là phổ nguyên tố thuần nhất của $\mathcal S'_*$. Để chứng minh (3.2), ta
dễ dàng quy về trường hợp $X=\mathbf P^r_S$ (vì $S$ là hợp của các tập mở $U$
sao cho $X|U$ là một lược đồ con đóng của $\mathbf P^r_U$ và
$\mathcal O_X(1)$ được cảm sinh bởi $\mathcal O_{\mathbf P^r_U}(1)$), và
trường hợp $\mathcal F$ có dạng $\mathcal O_{\mathbf P^r_S}^{N}$, do đó phẳng
trên $S$. Khi ấy, trong phần trên, các bó $\mathcal K$ cũng phẳng trên $S'$.
Từ các quan hệ Künneth ([5], III, \S~7), suy ra các điều kiện (a) và (b) ổn
định qua phép đổi cơ sở, đồng thời kéo theo rằng với $n\geq\nu$, việc tạo thành
$f'_*(\mathcal G(n))$ và $f'_*(\mathcal K(n))$ giao hoán với phép mở rộng cơ
sở (loc. cit.); do đó, (c) cũng ổn định qua phép mở rộng cơ sở.

Nói cách khác, $A_\nu(S')$ là một hàm tử phản biến theo $S'$, chính xác hơn là
một hàm tử con của
$A(S')=\underline{\operatorname{Quot}}^P_{\mathcal F/X/S}(S')$. Khi $\nu$
biến thiên, ta thu được một dãy tăng các tập con $A_\nu(S')$ của $A(S')$, có
hợp bằng $A(S')$ theo một định lý quen thuộc của Serre [5], III, \S~2. Bây giờ,
giả sử $\mathcal G$ là một môđun thương coherent của $\mathcal F'$, phẳng trên
$S'$, và $s$ là một điểm của $S'$ sao cho phép đổi cơ sở
$\operatorname{Spec}(k(s))\longrightarrow S'$ cho một môđun thương
$\mathcal G_s$ của $\mathcal F_s$ thỏa mãn các điều kiện (a), (b), (c), tức là
thuộc $A_\nu(\operatorname{Spec}(k(s)))$. Khi đó tồn tại một lân cận mở $U$
của $s$ sao cho các điều kiện ấy cũng được thỏa mãn đối với
$\mathcal G|f'^{-1}(U)$, tức là môđun thương này thuộc $A_\nu(U)$; đối với (a)
và (b), điều này suy ra từ «định lý về các hàm chỉnh hình» [5], III, \S~7, còn
(c) suy ra từ bổ đề Nakayama và từ việc trong mọi trường hợp ta biết rằng
$f'_*(\mathcal K(n+k))=\mathcal S'_k f'_*(\mathcal K(n))$
\src{261}
với $n$ đủ lớn và $k\geq0$ (xem [5], III, \S~2). Từ những nhận xét này, ta
suy ra điều sau (so sánh [4], IV). Để hàm tử $A$ biểu diễn được, điều kiện cần
và đủ là các hàm tử $A_\nu$ biểu diễn được; khi đó, $S$-tiền lược đồ $Q$ biểu
diễn $A$ là hợp tăng của các tiền lược đồ con mở $Q_\nu$ biểu diễn các
$A_\nu$.

Cho
\[
M_* =\bigoplus_{n\geq0}f_*(\mathcal F(n))=\mathcal S_*^{N},
\]
do đó
\[
M'_* =M_*\otimes_{\mathcal O_S}\mathcal O_{S'}
=\bigoplus_{n\geq0}f'_*(\mathcal F'(n))=\mathcal S'_*{}^{N}.
\]
Từ (a) suy ra:

\noindent (a') $f'_*(\mathcal G(n))$ tự do địa phương hạng $P(n)$ với
$n\geq\nu$, [5], III, \S~7, và từ (b), với $i=1$, suy ra:

\noindent (a'') $f'_*(\mathcal G(n))$ là một môđun thương của $M'_n$.

Hơn nữa, việc biết môđun thương này khi $n=\nu$, theo (c), xác định các môđun
con $f'_*(\mathcal K(n))$ của $M'_n$ với $n\geq\nu$, do đó xác định
$\mathcal K$ và tiếp theo là $\mathcal G$. Như vậy, ta thu được một ánh xạ đơn
\[
A_\nu(S')\longrightarrow\operatorname{Grass}_{P(\nu)}(M'_\nu)
\]
từ $A_\nu(S')$ vào tập hợp các môđun thương tự do địa phương hạng $P(\nu)$ của
$M'_\nu$, và do đó một đồng cấu theo hàm tử
\[
i_\nu:A_\nu(S')\longrightarrow
\underline{\operatorname{Grass}}_{P(\nu)}(M_\nu)(S'),
\]
trong đó hàm tử ở vế phải được biểu diễn bởi lược đồ Grassmann
$\operatorname{Grass}_{P(\nu)}(M_\nu)$ (so sánh [4], V), là một lược đồ xạ
ảnh trên $S$. Ta khẳng định:

\result{BỔ ĐỀ}{3.3} $A_\nu(S')$ là một hàm tử biểu diễn được, và cấu xạ
$Q_\nu\longrightarrow\operatorname{Grass}_{P(\nu)}(M_\nu)$ biểu diễn đồng cấu
$i_\nu$ là một phép nhúng (điều này kéo theo $Q_\nu$ tựa xạ ảnh trên $S$).

Khẳng định này tương đương với khẳng định sau (so sánh [4], IV): giả sử cho
một môđun thương $N$ của $M'_\nu$, tự do địa phương hạng $P(\nu)$; khi đó tồn
tại một tiền lược đồ con $Z$ của $S'$ sao cho với mọi tiền lược đồ Noether địa
phương $T'$
\src{262}
trên $S'$, ảnh ngược của $N$ trên $T'$ thuộc
$\operatorname{Im}A_\nu(T')$ khi và chỉ khi $T'\longrightarrow S'$ phân tích
qua tiền lược đồ con $Z$. Thay đổi ký hiệu, ta có thể giả sử $S'=S$, tức là ta
có một môđun thương $N_\nu$ của $M_\nu$ theo một môđun con $R_\nu$. Để môđun
thương này phát sinh từ một phần tử của $A_\nu(S)$, điều kiện cần và đủ là nó
thỏa mãn hai điều kiện sau:

\medskip
\noindent (i) $M_{\nu+k}/\mathcal S_kR_\nu$ tự do địa phương hạng
$P(\nu+k)$ với $k\geq0$.

\smallskip
\noindent (ii) Bó con $\mathcal K$ của $\mathcal F$, được xác định bởi môđun
con phân bậc
\[
R_* =\sum_{k\geq0}\mathcal S_kR_\nu
\]
của $M_*$ ([5], II, \S~3), và môđun thương
$\mathcal G=\mathcal F/\mathcal K$ thỏa mãn các điều kiện (a) và (b) ở trên.

Các điều kiện này hiển nhiên là cần thiết; ngược lại, nếu chúng được thỏa mãn,
thì bó $\mathcal G$ xác định trong (ii), do đẳng cấu với bó liên kết với
$\mathcal S_*$-môđun phân bậc $N_*$ là tổng của các
$M_{\nu+k}/\mathcal S_kR_\nu$, phẳng trên $S$ (vì các thớ của nó là những hạng
tử trực tiếp của các môđun thu được bằng cách địa phương hóa $N$ tại các
iđêan nguyên tố thuần nhất của $\mathcal S_*$), và có đa thức Hilbert $P$ theo
(i). Xét đến (ii), ta
thấy khi ấy (a') và (a'') được thỏa mãn, do đó, với $n\geq\nu$, đồng cấu tự
nhiên $N_n\longrightarrow f_*(\mathcal G(n))$ là một đồng cấu toàn ánh giữa
các môđun tự do địa phương có cùng hạng, nên là một đẳng cấu; vì vậy
\[
f_*(\mathcal K(\nu+k))=\mathcal S_kR_\nu
\quad\text{với mọi $k\geq0$},
\]
điều này chứng minh rằng $\mathcal G\in A_\nu(S)$ và $R_\nu$ là phần tử của
$\operatorname{Grass}_{P(\nu)}(M_\nu)$ được xác định bởi $\mathcal G$.

Tiêu chuẩn (i)--(ii) này cũng áp dụng cho tình huống thu được sau một phép đổi
cơ sở $S'\longrightarrow S$. Trước hết, ta sẽ chứng minh rằng việc điều kiện
(i) được thỏa mãn sau phép đổi cơ sở $S'\longrightarrow S$ có thể biểu thị
bằng điều kiện rằng $S'\longrightarrow S$ phân tích qua một tiền lược đồ con
$Z$ nào đó của $S$; sau khi có kết quả này, ta quy về (bằng cách thay $S$ bởi
$Z$) trường hợp điều kiện (i) đã được thỏa mãn trên $S$, và vì nó ổn định qua
phép đổi cơ sở nên chỉ còn phải biểu thị điều kiện (ii). Nhưng khi ấy, nếu $U$
ký hiệu tập hợp các $s\in S$ sao cho đối đồng điều của các bó cảm sinh trên thớ
$X_s$ bởi $\mathcal G(n)$ và $\mathcal K(n)$ triệt tiêu ở bậc $>0$ với
$n\geq\nu$, ta đã lưu ý rằng $U$ là mở, và điều kiện (ii) sẽ được thỏa mãn sau
một phép đổi cơ sở $S'\longrightarrow S$ khi và chỉ khi
$S'\longrightarrow S$ phân tích qua $U$; điều này chứng minh Bổ đề 3.3. Do đó,
còn phải chứng minh bổ đề sau:

\result{BỔ ĐỀ}{3.4} Cho $S$ là một tiền lược đồ Noether địa phương, được trang
bị một đại số phân bậc quasi-coherent $\mathcal S_*$ có các bậc dương, được
sinh bởi $\mathcal S_1$; cho $M_*$ là một $\mathcal S_*$-môđun phân bậc quasi-coherent
hữu hạn sinh, $P$ là một đa thức có hệ số hữu tỉ, và $\nu$ là một số nguyên.
Khi đó tồn tại một tiền lược đồ con $Z$ của $S$ (hiển nhiên là duy nhất)
\src{263}
có tính chất sau: với mọi tiền lược đồ $S'$ trên $S$, điều kiện cần và đủ để
$M_n\otimes_{\mathcal O_S}\mathcal O_{S'}$ tự do địa phương hạng $P(n)$ với
mọi $n\geq\nu$ là $S'\longrightarrow S$ phân tích qua $Z$.

Ta hiển nhiên có thể giả sử $S$ affine, do đó là Noether. Khi ấy:

\result{BỔ ĐỀ}{3.5} Với mọi số nguyên $N\geq\nu$, đặt $U_N$ là tập mở của $S$
gồm các $s\in S$ sao cho
\[
\operatorname{rank}_{k(s)}
M_{ns}\otimes_{\mathcal O_{S,s}}k(s)\leq P(n)
\quad\text{với mọi $\nu\leq n\leq N$}.
\]
Khi đó dãy giảm các tập mở $U_N$ là dừng.

Ta biết [3], IV, rằng $S$ có một phân hoạch hữu hạn thành các tiền lược đồ con
rút gọn $S_i$, sao cho mỗi $M\otimes_{\mathcal O_S}\mathcal O_{S_i}$ phẳng
trên $S_i$. Do đó, ta có thể giả sử $M$ phẳng, nên các $M_n$ phẳng. Cuối cùng,
ta hiển nhiên có thể giả sử $S$ liên thông. Nhưng khi ấy tồn tại một số nguyên
$n_0$ và một đa thức $Q$ sao cho
\[
\operatorname{rank}_{k(s)}
M_{ns}\otimes_{\mathcal O_{S,s}}k(s)=Q(n)
\quad\text{với $n\geq n_0$},
\]
[5], III, \S~7. Trước hết giả sử $P\neq Q$; khi đó
$P(n)\neq Q(n)$ với $n$ lớn, nên hiển nhiên $U_N=\varnothing$ với $N$ lớn;
huống chi dãy các $U_N$ là dừng. Ngược lại, nếu $P=Q$, thì
$U_N=U_{n_0}$ với $N\geq n_0$, và dãy các $U_N$ vẫn là dừng.

Đặc biệt, tập $U_\infty$ gồm các $s\in S$ sao cho
\[
\tag{*}
\operatorname{rank}_{k(s)}
M_{ns}\otimes_{\mathcal O_{S,s}}k(s)\leq P(n)
\quad\text{với mọi $n\geq\nu$},
\]
là mở vì nó là giao của các $U_N$. Khi ấy, để chứng minh (3.4), ta có thể thay
$S$ bởi tập mở $U$, nhờ đó quy về trường hợp bất đẳng thức (*) được thỏa mãn
tại mọi $s\in S$. Hơn nữa:

\result{BỔ ĐỀ}{3.6} Cho $M$ là một môđun trên một tiền lược đồ Noether địa
phương $S$, và $r$ là một số nguyên. Khi đó tồn tại một tiền lược đồ con $Z$
của $S$ (hiển nhiên là duy nhất) có tính chất sau: với mọi $S'$ trên $S$, điều
kiện cần và đủ để $M\otimes_{\mathcal O_S}\mathcal O_{S'}$ tự do địa phương
hạng $r$ là $S'\longrightarrow S$ phân tích qua $Z$. Nếu
\[
\operatorname{rank}_{k(s)}
M_s\otimes_{\mathcal O_{S,s}}k(s)\leq r
\quad\text{với mọi $s$},
\]
thì $Z$ là một tiền lược đồ con đóng của $S$ (giả sử $M$ coherent).
\src{264}
Thật vậy, lập luận trên quy ta về trường hợp bất đẳng thức (*) đúng với mọi
$s\in S$ (nếu cần, thay $S$ bằng phần mở gồm các $s$ tại đó bất đẳng thức
được thỏa mãn). Khi ấy ta có thể giả sử $M$ nằm trong một dãy khớp
\[
\mathcal O_S^q\longrightarrow\mathcal O_S^r\longrightarrow M\longrightarrow0,
\]
và điều kiện đang xét đối với các $S'$ trên $S$ cũng có nghĩa là, trong dãy
khớp tương ứng $\mathcal O_{S'}^q\longrightarrow\mathcal O_{S'}^r
\longrightarrow M'\longrightarrow0$, mũi tên thứ hai là một đẳng cấu, tức là
mũi tên thứ nhất bằng không. Khi đó, tiền lược đồ con đóng $Z$ của $S$ được
xác định bởi iđêan sinh bởi các hệ số của ma trận xác định đồng cấu
$\mathcal O_S^q\longrightarrow\mathcal O_S^r$ thỏa mãn điều kiện mong muốn.

Trở lại chứng minh (3.4) tại chỗ ta đã dừng, gọi $Z_n$ là lược đồ con đóng
của $S$ được liên kết, theo (3.6), với môđun $M_n$ và số nguyên $r=P(n)$, và
gọi $Z'_N$ là $\operatorname{Inf}$ của các $Z_n$ với
$\nu\leq n\leq N$. Khi ấy các $Z'_N$ tạo thành một dãy giảm các lược đồ con
đóng của $Z$, nên tất yếu dừng. Gọi $Z$ là giá trị không đổi của các
$Z'_N$ khi $N$ lớn. Đó là $Z$ cần tìm trong (3.4). Điều này kết thúc chứng
minh (3.4), và do đó cả chứng minh (3.3).

Vậy ta đã chứng minh rằng
$\underline{\operatorname{Quot}}^P_{\mathcal F/X/S}$ được biểu diễn bởi một
$S$-tiền lược đồ $Q$ là hợp tăng của các tiền lược đồ con mở $Q_\nu$ tựa xạ
ảnh trên $S$. Để tiến thêm, phải viện dẫn Định lý 2.1, từ đó dễ dàng suy ra
$Q$ là quasi-compact (vì nó là ảnh của tiền lược đồ $S'$ thuộc kiểu hữu hạn
trên $S$ tham số hóa họ các bó thương của $\mathcal F_K$ có đa thức Hilbert
$P$). Do đó $Q$ bằng một trong các $Q_\nu$, vì thế tựa xạ ảnh trên $S$. Để
chứng minh nó xạ ảnh trên $S$, chỉ còn phải chứng minh nó riêng trên $S$; để
làm điều này, chỉ cần viện dẫn tiêu chuẩn định giá cho tính riêng dưới dạng
[5], II, 7.3.8. Chỉ cần kiểm tra điều sau:

\result{BỔ ĐỀ}{3.7} Cho $S$ là phổ của một vành định giá rời rạc, $s$ là điểm
tổng quát của nó, $X$ là một tiền lược đồ trên $S$, $\mathcal F$ là một môđun
quasi-coherent trên $X$, và $\mathcal G_s$ là một môđun thương quasi-coherent
của $\mathcal F_s=\mathcal F\otimes_{\mathcal O_S}k(s)$ trên $X_s$. Khi đó
tồn tại duy nhất một môđun thương quasi-coherent $\mathcal G$ của
$\mathcal F$, phẳng trên $S$, sao cho hạn chế của nó lên $X_s$ là
$\mathcal G_s$.
\src{265}
Thật vậy, nếu $\mathcal G_s=\mathcal F_s/\mathcal K_s$, chỉ cần xét bó con lớn
nhất $\mathcal K$ của $\mathcal F$ cảm sinh $\mathcal K_s$, [5], I, 9.4.2, và
lấy $\mathcal G=\mathcal F/\mathcal K$. Dễ dàng kiểm tra rằng bó này thỏa mãn
yêu cầu.

Định lý 3.2, và do đó (3.1), được chứng minh hoàn toàn.

Đồng thời, chứng minh cũng cho ta điều sau:

\result{MỆNH ĐỀ}{3.8} Trong các điều kiện của (3.2), đặt
$Q=\operatorname{Quot}^P_{\mathcal F/X/S}$, $X_Q=X\times_SQ$,
$\mathcal F_Q=\mathcal F\otimes_{\mathcal O_S}\mathcal O_Q$, và cho
$\mathcal G$ là môđun thương coherent của $\mathcal F_Q$, phẳng trên $Q$, có
đa thức Hilbert tương đối $P$, sao cho $(Q,\mathcal G)$ biểu diễn hàm tử
$\underline{\operatorname{Quot}}^P_{\mathcal F/X/S}$. Tồn tại một số nguyên
$\nu$ sao cho, với $n\geq\nu$, $(f_Q)_*(\mathcal G(n))$ là một môđun trên $Q$
tự do địa phương hạng $P(n)$ và rất rộng tương đối với $S$, tức là xác định
một phép nhúng của $Q$ vào một lược đồ Grassmann
$\operatorname{Grass}_{P(n)}(M)$ trên $S$. Đặc biệt, với $n\geq\nu$, bó
\[
\bigwedge^{P(n)}(f_Q)_*(\mathcal G(n))
\]
trên $Q$ là khả nghịch và rất rộng tương đối với $S$.

Thật vậy, như đối với (3.2), ta có thể quy về trường hợp $\mathcal F$ phẳng
trên $S$, và khi đó chỉ cần lấy một số nguyên $\nu$ sao cho $A_\nu=A$ (với các
ký hiệu trước).

Trường hợp áp dụng quan trọng nhất của (3.2) là khi
$\mathcal F=\mathcal O_X$. Khi ấy ta viết
\[
\operatorname{Quot}_{\mathcal O_X/X/S}=\operatorname{Hilb}_{X/S},
\qquad
\operatorname{Quot}^P_{\mathcal O_X/X/S}=\operatorname{Hilb}^P_{X/S},
\]
do đó có một phân rã
\[
\operatorname{Hilb}_{X/S}=\coprod_P\operatorname{Hilb}^P_{X/S}.
\]
Theo định nghĩa, $\operatorname{Hilb}_{X/S}$ biểu diễn hàm tử
$\underline{\operatorname{Hilb}}_{X/S}(S')$, là tập hợp các tiền lược đồ con
đóng của $X'=X\times_SS'$ phẳng trên $S'$, và
$\operatorname{Hilb}^P_{X/S}$ biểu diễn hàm tử con tương ứng với các tiền lược
đồ con đóng có một đa thức Hilbert cho trước $P$. Các tiền lược đồ này lần lượt
còn được gọi là \emph{tiền lược đồ Hilbert} của $X$ trên $S$ và
\emph{tiền lược đồ Hilbert chỉ số $P$}. Thuật ngữ này được biện minh bởi vai
trò của các đa thức Hilbert trong lý thuyết. Về bản chất, sự khác biệt của
chúng với các đa tạp Chow cổ điển (dùng để tham số hóa các chu trình, chứ không
phải các đa tạp) tương tự như sự khác biệt giữa
\src{266}
vành Chow của các lớp chu trình trên một đa tạp và vành các lớp bó trên đa
tạp đó (như vành này xuất hiện trong định lý Riemann--Roch [1]); thật vậy,
lưu ý rằng khi $X=\mathbf P^r_S$, với $S$ là phổ của một trường, việc biết đa
thức Hilbert của một môđun coherent $\mathcal F$ trên $X$ cũng tương đương với
việc biết các lớp Chern của $\mathcal F$, hay lớp của $\mathcal F$ trong vành
các lớp bó coherent trên $X$.

\result{CÁC NHẬN XÉT}{3.9} Cũng cần lưu ý rằng việc xây dựng các
$\operatorname{Quot}_{\mathcal F/X/S}$ và
$\operatorname{Quot}^P_{\mathcal F/X/S}$ đã được quy về trường hợp
$X=\mathbf P^r_S$, $\mathcal F=\mathcal O_X^N$, trong đó
$\mathcal O_X(1)$ là bó rất rộng thông thường; chính xác hơn, các
$\operatorname{Quot}_{\mathcal F/X/S}$ tổng quát là các tiền lược đồ con đóng
của những tiền lược đồ vừa nêu. Vì việc tạo thành các
$\operatorname{Quot}^P_{\mathcal F/X/S}$ hiển nhiên tương thích với các phép
đổi cơ sở $S'\longrightarrow S$, ta thấy có thể quy tiếp về trường hợp
$S=\operatorname{Spec}(\mathbf Z)$. Sau cùng, ta được quy về việc nghiên
cứu các lược đồ xạ ảnh sau đây trên $\mathbf Z$:
\[
\mathfrak Q^P_{r,N}=
\operatorname{Quot}^{P}_{(\mathcal O_{\mathbf P^r_{\mathbf Z}})^N/
\mathbf P^r_{\mathbf Z}/\operatorname{Spec}(\mathbf Z)},
\]
và đặc biệt hơn là các lược đồ Hilbert tuyệt đối:
\[
\operatorname{Hilb}^P_r=\mathfrak Q^P_{r,1}.
\]
Sẽ rất đáng quan tâm nếu nghiên cứu chi tiết hơn các lược đồ này, bắt đầu bằng
việc xác định các thành phần liên thông của chúng (các lược đồ này có liên
thông không?), và các thành phần bất khả quy của
chúng (theo Serre [7], có thể có những thành phần bất khả quy nằm hoàn toàn
phía trên một số nguyên tố $p\neq0$). Nhắc lại câu hỏi của Weil: liệu các thành
phần bất khả quy của các thớ của $\operatorname{Hilb}^P_r$ trên các
$s\in\operatorname{Spec}(\mathbf Z)$ có tương ứng với các mở rộng
``chính quy'' của trường nguyên tố hay không, tức là liệu chúng có ``liên thông
tương đối'' hay không. Không loại trừ khả năng những câu hỏi này dễ tiếp cận
đối với các lược đồ Hilbert hơn là đối với các ``đa tạp Chow''.

\fsection{4}{Các biến thể}

\noindent a. Trong các điều kiện của (3.1), cho $U$ là một tập mở trong $X$,
và ký hiệu $A'$ là hàm tử con của
$A=\underline{\operatorname{Quot}}_{\mathcal F/X/S}$ sao cho $A'(S')$ là tập
các môđun thương $\mathcal G$ của $\mathcal F'$, phẳng trên $S'$, có giá chứa
trong $U'$. Ta thấy ngay rằng $A'$ được biểu diễn bởi một tập mở của tiền lược
đồ $\operatorname{Quot}_{\mathcal F/X/S}$ biểu diễn $A$. Do đó các định lý
(3.1) và (3.2)
\src{267}
vẫn còn đúng khi giả sử $X$ tựa xạ ảnh trên $S$ thay vì xạ ảnh trên $S$, đồng
thời trong các kết luận cũng thay từ ``xạ ảnh'' bằng ``tựa xạ ảnh'', và giờ
đây ký hiệu
$\underline{\operatorname{Quot}}_{\mathcal F/X/S}(S')$ là tập các môđun
thương coherent $\mathcal G$ của $\mathcal F'$, phẳng trên $S'$, có giá riêng
trên $S'$.

\medskip
\noindent b. Nói chung, có thể áp đặt lên các môđun thương $\mathcal G$ của
$\mathcal F'=\mathcal F\otimes_{\mathcal O_S}\mathcal O_{S'}$, phẳng trên
$S'$, mọi loại điều kiện tự nhiên bổ sung ổn định qua phép đổi cơ sở, qua đó thu
được các hàm tử con của
$\underline{\operatorname{Quot}}_{\mathcal F/X/S}$ mà ta muốn biểu diễn. Trong
nhiều trường hợp, tiêu chuẩn thông thường cho phép chứng minh rằng các hàm tử
thu được vẫn được biểu diễn bởi các tập mở của
$\operatorname{Quot}_{\mathcal F/X/S}$. Đặc biệt, điều này đúng khi áp đặt một
trong các tính chất bổ sung sau:

\smallskip
\noindent $1^\circ$ Chiều của mỗi chu trình nguyên tố liên kết với các môđun
$\mathcal G_{s'}$ ($s'\in S'$), được cảm sinh trên các thớ $X'_{s'}$, thuộc
một tập số nguyên đã cho.

\smallskip
\noindent $2^\circ$ (Khi $\mathcal F=\mathcal O_X$, do đó $\mathcal G$ tương
ứng với một tiền lược đồ con đóng $Y$ của $X'$); $Y$ là một tiền lược đồ trơn
([3], IV) trên $S$, hoặc chuẩn tắc trên $S$ (tức là các thớ $Y_{s'}$ chuẩn
tắc ``trên $k(s')$'', nghĩa là vẫn chuẩn tắc sau mọi mở rộng của trường cơ
sở), hoặc (khi $X$ phẳng trên $S$) là một giao đầy đủ địa phương theo cơ sở
$S$ trong $X$, xét tương đối trên $S$ (tức là các thớ $Y_{s'}$ là các giao
đầy đủ địa phương trong các $X_{s'}$).

Những điều kiện khác sẽ liên quan đến các tính chất đối đồng điều
của các môđun $\mathcal G_{s'}$ được cảm sinh trên các $X'_{s'}$, v.v. Dĩ
nhiên, hội của các điều kiện mà mỗi điều kiện được biểu diễn bởi một tập mở
$U_i$ của $\operatorname{Quot}_{\mathcal F/X/S}$ được biểu diễn bởi tập mở thu
được bằng cách lấy giao các tập ấy. Chẳng hạn, với mỗi $S'$ trên $S$, xét tập
các tiền lược đồ con đóng $Y$ của $X'=X\times_SS'$ mà mỗi tiền lược đồ là một
phủ étale ([3], I) hạng cho trước $r$ của $S'$, ta thu được một hàm tử phản biến
biểu diễn được theo biến $S'$.

\medskip
\noindent c. Các tiền lược đồ $\operatorname{Hom}_S(X,Y)$,
$\Pi_{X/S}(Z/X)$, $\operatorname{Isom}_S(X,Y)$, được định nghĩa trong [2], II,
C, n\textsuperscript{o} 2, tồn tại dưới các giả thiết xạ ảnh thích hợp và được
hiện thực hóa dưới dạng các tập mở trong những tiền lược đồ Hilbert thích hợp.
Vì ta có
\[
\operatorname{Hom}_S(X,Y)=\Pi_{X/S}((X\times_SY)/X),
\]
trường hợp $\operatorname{Hom}_S(X,Y)$ được quy về trường hợp
$\Pi_{X/S}(Z/X)$. Khi đó, lưu ý rằng với mỗi $S'$ trên $S$, tập các tiết diện
của $Z'=Z\times_SS'$ trên $X'=X\times_SS'$ tương ứng song ánh với tập các tiền
lược đồ con $\Gamma$ của $Z'$ (nhất thiết đóng nếu $Z$ tách trên $X$) sao cho
cấu xạ $\Gamma\longrightarrow X'$ cảm sinh bởi $Z'\longrightarrow X'$ là một
đẳng cấu. Bằng cách này,
\src{268}
khi $X$ phẳng và riêng trên $S$, còn $Z$ tựa xạ ảnh trên $S$,
$\Pi_{X/S}(Z/X)$ tồn tại và được hiện thực hóa dưới dạng một tiền lược đồ con mở của
$\operatorname{Hilb}_{Z/S}$. Do đó, khi $X$ xạ ảnh và phẳng trên $S$, còn $Y$
tựa xạ ảnh trên $S$, $\operatorname{Hom}_S(X,Y)$ tồn tại và được hiện thực hóa
dưới dạng một tiền lược đồ con mở của
$\operatorname{Hilb}_{(X\times_SY)/S}$. Khi $X$ và $Y$ đều xạ ảnh trên $S$,
suy ra ngay $\operatorname{Isom}_S(X,Y)$ cũng tồn tại và được biểu diễn bởi
một tập mở của $\operatorname{Hom}_S(X,Y)$. Tương tự, khi $X$ phẳng và xạ ảnh
trên $S$, còn $Y$ tựa xạ ảnh trên $S$, $S$-tiền lược đồ
$\operatorname{Imm}_S(X/Y)$, tương ứng với hàm tử con của hàm tử được
$\operatorname{Hom}_S(X,Y)$ biểu diễn, gồm các $S'$-cấu xạ
$X'\longrightarrow Y'$ là các phép nhúng, cũng được biểu diễn bởi một tập mở
của $\operatorname{Hom}_S(X,Y)$.

Cho $\mathcal L$ (tương ứng $\mathcal M$) là một bó khả nghịch trên $X$ (tương
ứng trên $Y$), rất rộng tương đối với $S$; từ đó có một bó rất rộng
$\mathcal L\otimes_{\mathcal O_S}\mathcal M$ trên $X\times_SY$ tương đối với
$S$. Vì vậy, với mọi đa thức $P$ có hệ số hữu tỉ,
$\operatorname{Hilb}^P_{(X\times_SY)/S}$ được xác định và là một tiền lược đồ
tựa xạ ảnh trên $S$. Do đó nó cảm sinh trên $\operatorname{Hom}_S(X,Y)$ một
tập vừa mở vừa đóng, tựa xạ ảnh trên $S$, mà ta ký hiệu là
$\operatorname{Hom}_S(X,Y)^P$. Vậy các tiết diện của
$\operatorname{Hom}_S(X,Y)^P$ trên $S$ là các $S$-cấu xạ
$g:X\longrightarrow Y$ sao cho với mọi số nguyên $n$, ta có
\[
\chi\!\left(\left(\mathcal L\otimes_{\mathcal O_X}g^*(\mathcal M)\right)^{
\otimes n}\right)=P(n).
\]

Theo cách đó, ta thu được những tổng quát hóa của định lý Matsusaka, khẳng định
rằng các tự đẳng cấu của một đa tạp xạ ảnh ``phân cực'' tạo thành một nhóm đại
số. Ở đây, khẳng định này hiển nhiên mang một ý nghĩa chính xác hơn, vì ta có
một định nghĩa của nhóm ấy như nghiệm của một bài toán phổ dụng. Cũng cần lưu ý
rằng, trên một trường đóng đại số, nhóm các tự đẳng cấu được xét trước đây là
nhóm thu được từ nhóm ``thật'' được định nghĩa ở đây bằng cách lấy thương theo
iđêan các phần tử lũy linh; điều này giải thích vì sao khó có thể thực hiện các
phép dựng cũ trên một trường cơ sở không hoàn hảo: iđêan các phần
tử lũy linh xuất hiện sau một phép mở rộng trường cơ sở không nhất thiết
``được xác định trên $k$''. Nhận xét này cũng áp dụng cho phần lớn các phép dựng
theo lối cũ.

\fsection{5}{Nghiên cứu vi phân các lược đồ Hilbert}

Nghiên cứu này suy ra từ kết quả sau:
\src{269}
\result{MỆNH ĐỀ}{5.1} Cho $S$ là một tiền lược đồ, $S_0$ là một tiền lược đồ con
được xác định bởi một iđêan quasi-coherent $\mathcal J$ có bình phương bằng
không, $X$ là một $S$-tiền lược đồ, $\mathcal F$ là một môđun quasi-coherent
trên $X$, đặt $X_0=X\times_SS_0$ và
$\mathcal F_0=\mathcal F\otimes_{\mathcal O_S}\mathcal O_{S_0}$, và cuối cùng
cho $\mathcal G_0=\mathcal F_0/\mathcal K_0$ là một môđun thương
quasi-coherent của $\mathcal F_0$, phẳng trên $S_0$. Với mỗi tập mở $U$ của
$X$, gọi $\mathcal E(U)$ là tập các môđun thương quasi-coherent $\mathcal G$
của $\mathcal F|U$, phẳng trên $S$, và thỏa mãn
$\mathcal G\otimes_{\mathcal O_S}\mathcal O_{S_0}=\mathcal G_0$; do đó, khi
$U$ thay đổi, các $\mathcal E(U)$ là các tiết diện của một bó $\mathcal E$
trên $U$. Khi ấy, bó các nhóm
\[
\mathcal A=\operatorname{Hom}_{\mathcal O_{X_0}}
(\mathcal K_0,\mathcal G_0\otimes_{\mathcal O_{S_0}}\mathcal J)
\]
tác động tự nhiên lên $\mathcal E$, nhờ đó bó này trở thành một bó
``thuần nhất chính một cách hình thức dưới $\mathcal A$'' (nghĩa là, với mọi
tập mở $U$ trong $X$, $\mathcal E(U)$ hoặc là rỗng, hoặc là một tập thuần nhất
chính dưới $\mathcal A(U)$).

Từ đó suy ra:

\result{HỆ QUẢ}{5.2} Giả sử địa phương trên $X$ tồn tại một phép nâng
$\mathcal G$ của $\mathcal G_0$ thành một môđun thương của $\mathcal F$ phẳng
trên $S$ (nghĩa là các thớ của bó $\mathcal E$ không rỗng). Khi đó tồn tại một
lớp cản trở chính tắc
\[
c(\mathcal G_0)\in H^1(X,\mathcal A),
\]
mà sự triệt tiêu của nó là điều kiện cần và đủ để tồn tại một phép nâng toàn
cục $\mathcal G$ của $\mathcal G_0$ thành một môđun thương của $\mathcal F$
phẳng trên $S$. Nếu lớp này bằng không, thì tập $\mathcal E(X)$ gồm mọi phép
nâng khả dĩ là một tập thuần nhất chính dưới
\[
\mathcal A(X)=\operatorname{Hom}_{\mathcal O_X}
(\mathcal K_0,\mathcal G_0\otimes_{\mathcal O_{S_0}}\mathcal J).
\]

Đặc biệt, sự tồn tại của một phép nâng toàn cục được bảo đảm nếu
$H^1(X,\mathcal A)=0$.

\result{HỆ QUẢ}{5.3} Giả sử
$Q=\operatorname{Quot}_{\mathcal F/X/S}$ tồn tại (xem 4 (a)) -- chẳng hạn,
$X$ tựa xạ ảnh trên tiền lược đồ Noether địa phương $S$ và $\mathcal F$ coherent. Cho
$x\in Q$ tương ứng với một mở rộng trường thặng dư $K=k(x)$ của $k(s)$
($s\in S$):
\src{270}
do đó, $x$ được xác định bởi một môđun thương coherent
$\mathcal G_0=\mathcal F_0/\mathcal K_0$ của môđun
$\mathcal F_0=\mathcal F_K$ trên $K$-tiền lược đồ $X_K$. Gọi $\mathcal A$ là
bó coherent trên $X_K$ được xác định bởi
\[
\mathcal A=\operatorname{Hom}_{\mathcal O_{X_0}}(\mathcal K_0,\mathcal G_0).
\]
Khi đó, không gian tiếp tuyến Zariski của thớ $Q_s$ tại điểm $x$ (đối ngẫu trên
$K$ của $\mathfrak m/\mathfrak m^2$, trong đó $\mathfrak m$ là iđêan cực đại
của $\mathcal O_{Q_K,x}$) đẳng cấu chính tắc với $H^0(X_K,\mathcal A)$.

Mô tả không gian tiếp tuyến Zariski được tổng quát hóa như sau. Với một
$S$-cấu xạ cho trước $g:S'\longrightarrow Q$, tức là một tiết diện $g'$ của
$Q'=Q\times_SS'$ trên $S'$, môđun
\[
\Omega=g^*(\Omega^1_{Q/S})=g'^*(\mathcal J/\mathcal J^2)
\]
(trong đó $\mathcal J$ là iđêan trên $Q'$ được xác định bởi tiết diện $g'$ của
$Q'$ trên $S'$) được đặc trưng bởi công thức sau, có tính hàm tử theo môđun coherent
$\mathcal M$ trên $S'$:
\[
\operatorname{Hom}_{\mathcal O_{S'}}(\Omega,\mathcal M)
\simeq H^0(X',\mathcal A),
\]
trong đó $\mathcal A$ vẫn là môđun trên $X'=X\times_SS'$ được xác định bởi
\[
\mathcal A=\operatorname{Hom}_{\mathcal O_{X'}}
(\mathcal K,\mathcal G\otimes_{\mathcal O_S}\mathcal M),
\]

($\mathcal G=\mathcal F'/\mathcal K$ là môđun thương của
$\mathcal F'=\mathcal F\otimes_{\mathcal O_S}\mathcal O_{S'}$ tương ứng với
$g$). Thật vậy, chỉ cần áp dụng (5.1), trong đó thay $S_0$ bởi $S'$, thay $S$
bởi tiền lược đồ $D(\mathcal M)=(S',\mathcal O_{S'}+\mathcal M)$, với
$\mathcal M$ được xem như một iđêan có bình phương bằng không.

Khi $\mathcal F=\mathcal O_X$ trong (5.1), dữ liệu của
$\mathcal G_0$ tương ứng với một tiền lược đồ con đóng $Y_0$ của
$X_0$, phẳng trên $S_0$ và được xác định bởi iđêan
$\mathcal J_0=\mathcal K_0$; khi đó (*) cho
\[
\mathcal A=\operatorname{Hom}_{\mathcal O_{X_0}}
(\mathcal J_0/\mathcal J_0^2,
\mathcal O_{Y_0}\otimes_{\mathcal O_{S_0}}\mathcal J),
\]
trong đó $\mathcal J/\mathcal J^2$ được hiểu là bó đối pháp tuyến của $Y_0$
trong $X_0$, cũng được ký hiệu là
\src{271}
$\check{\mathcal N}_{Y_0/X_0}$; khi đó nên xem $\mathcal A$ như một môđun
trên $Y_0$ và tính $H^0$ và $H^1$ trên $Y_0$. Hơn nữa, nếu $Y_0$ là một giao
đầy đủ địa phương trong $X_0$, với $X$ phẳng trên $S$, thì trong (5.1) khả
năng nâng địa phương được bảo đảm; mặt khác,
$\mathcal J/\mathcal J^2$ tự do địa phương trên $Y_0$ và ta có thể viết
\[
\mathcal A=\mathcal N_{Y_0/X_0}\otimes_{\mathcal O_{S_0}}\mathcal J,
\]
trong đó thừa số thứ nhất ở vế phải là bó pháp tuyến của $Y_0$ trong $X_0$.
Áp dụng tiêu chuẩn cơ bản về tính trơn ([3], III, 3.1), chẳng hạn ta thu được:

\result{HỆ QUẢ}{5.4} Trong các điều kiện của (5.3), giả sử rằng
$\mathcal F=\mathcal O_X$, rằng $X$ phẳng trên $S$, và rằng tiền lược đồ con
đóng $Y_0$ của $X_0$ tương ứng với $\mathcal G_0$ là một giao đầy đủ địa
phương. Khi đó, không gian tiếp tuyến Zariski của $Q_s$ tại $x$ đẳng cấu chính
tắc với $H^0(Y_0,\mathcal N_{Y_0/X_0})$. Nếu
$H^1(Y_0,\mathcal N_{Y_0/X_0})=0$, thì tiền lược đồ Hilbert
$\operatorname{Hilb}_{X/S}$ trơn trên $S$ tại điểm $x$.
($\mathcal N_{Y_0/X_0}$ là bó pháp tuyến của $Y_0$ trong $X_0$.)

\result{CHÚ GIẢI}{5.5} Phát biểu này áp dụng đặc biệt khi $Y_0$ là một giao
đầy đủ trong $X_0$ được xác định bởi một phương trình, tức là một ``ước
Cartier'' hiệu dụng. Khi đó, $\mathcal N_{Y_0/X_0}$ đẳng cấu với bó trên $Y_0$
cảm sinh bởi bó khả nghịch $\mathcal J^{-1}$ trên $X_0$ được xác định bởi
ước $Y_0$. Tình huống này xuất hiện đặc biệt khi nghiên cứu các họ ước hiệu
dụng trên một đa tạp xạ ảnh không kỳ dị $X_0$. Đẳng cấu giữa không gian tiếp
tuyến Zariski tại điểm $x$ của $Q$ (hoặc, nếu muốn, của tập mở $D$ của $Q$
tương ứng với các ước) và $H^0(Y_0,\mathcal N_{Y_0/X_0})$ từng được biết trong
hình học đại số cổ điển dưới tên ``đồng cấu đặc trưng'' từ không gian thứ nhất
đến không gian thứ hai. Đồng cấu này chỉ được xác định khi $x$ là một điểm trơn
của đa tạp tham số $T$ của một ``họ liên tục đầy đủ'' các ước; theo quan điểm
của chúng ta, đa tạp tham số ấy là một thành phần bất khả quy của lược đồ $D$
với cấu trúc rút gọn cảm sinh. Khi đó, không gian tiếp tuyến của $T$ tại $x$ là
một không gian con của không gian tiếp tuyến của $D$ tại $x$; do đó đồng cấu
đặc trưng cổ điển quả là đơn ánh, nhưng chỉ toàn ánh dưới các điều kiện bổ
sung, chẳng hạn nếu $D$ là nguyên tại $x$. Thật vậy, Zappa [8] đã xây dựng
một ví dụ
\src{272}
(với $X$ là một mặt xạ ảnh không kỳ dị trên trường số phức) trong đó, ngay cả
tại điểm tổng quát của $T$, đồng cấu đặc trưng cũng không toàn ánh. Do đó, điều
này có nghĩa là $D$ không nguyên ngay cả tại điểm tổng quát của thành phần bất
khả quy đang xét. Điều này cho thấy một cách đặc biệt nổi bật rằng các đa tạp
có phần tử lũy linh là cần thiết để hiểu những hiện tượng thuộc lý thuyết cổ
điển nhất về các mặt.


\subsection*{5.6}

Trong (5.4), ta đã nêu một tiêu chuẩn về tính trơn, áp dụng đặc biệt cho các
lược đồ ước. Trong [6], Kodaira đã nêu một tiêu chuẩn khác, cụ thể là sự triệt
tiêu của $H^1(X_0,\mathcal L)$, trong đó $\mathcal L=\mathcal J_0^{-1}$ là bó
khả nghịch trên $X_0$ được xác định bởi ước $Y_0$; tiêu chuẩn này đúng khi $S$
là phổ của một trường có đặc số 0, và được chứng minh trong [6] bằng phương
pháp siêu việt khi trường cơ sở là $\mathbf C$. Ta lưu ý rằng, nói chung, khi
$S$ lại là tùy ý, điều kiện của Kodaira là một điều kiện đủ để cấu xạ chính tắc
\[
D\longrightarrow\operatorname{Pic}_{X/S}
\]
từ tiền lược đồ các ước đến tiền lược đồ Picard của $X/S$ trơn tại điểm $x$
đang xét (điều này dễ dàng được kiểm tra bằng tiêu chuẩn thông thường về tính
trơn, một khi đã biết sự tồn tại của $\operatorname{Pic}_{X/S}$). Do đó, nếu
thêm vào đó $\operatorname{Pic}_{X/S}$ trơn trên $S$ tại điểm ảnh của $x$
(chẳng hạn nếu $\operatorname{Pic}_{X/S}$ trơn trên $S$), thì $D$ trơn trên
$S$ tại $x$. Mặt khác, Cartier đã chứng minh rằng mọi tiền lược đồ nhóm thuộc
kiểu hữu hạn địa phương trên một trường $k$ có đặc số 0 đều trơn trên $k$. Kết
hợp hai kết quả này, ta thu lại kết quả của Kodaira. Cũng cần lưu ý rằng từ các
nhận xét này suy ra rằng, trên một trường $K$ có đặc số $p>0$, khi
$\operatorname{Pic}_{X/S}$ không trơn trên $k$ (đó là trường hợp khi $X$ là
mặt Igusa), điều kiện $H^1(X_0,\mathcal L)=0$ trái lại kéo theo rằng $D$ không
trơn tại $x$, và thậm chí không rút gọn tại $x$ nếu $K$ đóng đại số.

Cuối cùng, ta còn nêu kết quả sau đây, vốn đóng một vai trò quan trọng trong
nghiên cứu vi phân các không gian phân thớ:

\result{MỆNH ĐỀ}{5.7} Cho $X$ là một tiền lược đồ hữu hạn và phẳng trên tiền
lược đồ Noether địa phương $S$, và cho $Z$ là một tiền lược đồ trên $S$ sao
cho $\Pi_{X/S}(Z/X)$ tồn tại (điều này đúng nếu $Z$ tựa xạ ảnh trên $X$). Khi
đó, nếu $Z$ trơn trên $X$, thì $\Pi_{X/S}(Z/X)$ trơn trên $S$.

Đây là một hệ quả trực tiếp của định nghĩa và của tiêu chuẩn thông thường về
tính trơn ([3], III, 3.1). Ta lưu ý rằng khi $X$ hữu hạn và phẳng trên $S$,
vấn đề về sự tồn tại của
$Z/K$ có thể
được xử lý rất sơ cấp, không cần
\src{273}
sử dụng lý thuyết lược đồ Hilbert. Chẳng hạn, ta thấy rằng nếu $X$ là
radiciel trên $S$, thì $\Pi_{X/S}(Z/X)$ tồn tại mà không có bất kỳ hạn chế nào
đối với $Z$. Một ví dụ khác, cho $T$ là một $S$-tiền lược đồ, và cho $T_n$ là
``lân cận vô cùng bé bậc $n$'' của đường chéo của $T\times_ST$ trong
$T\times_ST$, được trang bị các cấu xạ $p_1,p_2:T_n\longrightarrow T$ cảm
sinh bởi hai phép chiếu. Ta xem $T_n$ như một tiền lược đồ hữu hạn trên $T$
qua $p_1$, và giả sử thêm rằng nó phẳng trên $T$ (điều này đúng nếu $T$ trơn
trên $S$). Với mọi tiền lược đồ $X$ trên $T$, đặt
\[
(X/T/S)^{(n)}=
\Pi_{T_n/S}\bigl(p_2^*(X/T)/T_n\bigr),
\]

đó là một tiền lược đồ trên $T$ được gọi là \emph{phân thớ các mầm tiết diện
bậc $n$ của $X$ trên $T$} (tương đối với $S$). Nó phụ thuộc theo hàm tử vào
$X$, và trơn trên $T$ nếu $X$ cũng vậy.

\fsection{6}{Quan hệ với khái niệm chuẩn và các tích đối xứng}

Cho $S$ là một tiền lược đồ, $X$ và $Y$ là các $S$-tiền lược đồ,
\[
u:(X/S)^n\longrightarrow Y
\]
là một $S$-cấu xạ đối xứng từ lũy thừa Descartes thứ $n$ của $X/S$ vào $Y$.
Để đơn giản, ta giả sử $S$ là Noether địa phương, và $X$ cùng $Y$ thuộc kiểu
hữu hạn trên $S$. Khi đó, với mỗi môđun coherent $\mathcal F$ trên $X$, có giá
hữu hạn trên $S$, phẳng trên $S$ và có hạng tương đối trên $S$ bằng $n$ (tức
là $f_*(\mathcal F)$ là một môđun tự do địa phương hạng $n$ trên $S$), ta có
thể liên kết một cách tự nhiên một tiết diện của $Y$ trên $S$:
\[
\pi^u_{X/S}(\mathcal F)\in\Gamma(Y/S).
\]
Ở đây ta sẽ không đưa ra định nghĩa hình thức, mà chỉ lưu ý rằng hình thức
luận thu được là một sự tổng quát hóa tự nhiên của hình thức luận thông thường
về chuẩn và vết. Khi lũy thừa đối xứng thứ $n$ của $X$ trên $S$ tồn tại (chẳng
hạn khi các quỹ đạo của nhóm đối xứng $\mathfrak S_n$ tác động trên $(X/S)^n$
được chứa trong các tập mở affine), ta có thể lấy
$Y$ bằng $\operatorname{Symm}^n_S(X)$, và thu được một phần tử chính tắc
\[
\pi_{X/S}(\mathcal F)\in\Gamma(\operatorname{Symm}^n_S(X)/S),
\]
\src{274}
cho phép khôi phục các $\pi^u_{X/S}(\mathcal F)$. Một trường hợp quan trọng
khác là khi $X$ là một monoid giao hoán trên $S$, và $X=Y$, còn cấu xạ $u$
xuất phát từ luật hợp thành của $X$. Khi đó ta sẽ chỉ viết $\pi(\mathcal F)$
cho tiết diện của $X$ trên $S$ liên kết với môđun $\mathcal F$ trên $X$.

Bây giờ giả sử ta có một môđun coherent $\mathcal F$ trên $X$ sao cho
$\operatorname{Quot}_{\mathcal F/X/S}$ tồn tại, hoặc ít nhất sao cho hàm tử
$\underline{\operatorname{Quot}}^n_{\mathcal F/X/S}$ gán cho mỗi $S'$ trên
$S$ tập hợp các bó thương coherent $\mathcal M$ của
$\mathcal F'=\mathcal F\otimes_{\mathcal O_S}\mathcal O_{S'}$ phẳng trên
$S'$ và có hạng tương đối $n$, được biểu diễn bởi một $S$-tiền lược đồ
$\operatorname{Quot}^n_{\mathcal F/X/S}$. (Khi $X$ quasi-xạ ảnh trên $S$,
thì $\operatorname{Quot}^n_{\mathcal F/X/S}$ quả thực tồn tại và, với ký hiệu
của mục 3, chính là $\operatorname{Quot}^P_{\mathcal F/X/S}$, trong đó $P$ là
đa thức chỉ còn số hạng hằng $n$.) Vì phép dựng
$\pi^u_{X/S}(\mathcal M)$ tương thích với phép đổi cơ sở, ta thu được một cấu
xạ chính tắc
\[
\pi^u_{X/S}:\operatorname{Quot}^n_{\mathcal F/X/S}\longrightarrow Y,
\]
đặc biệt, nếu lũy thừa đối xứng thứ $n$ của $X$ trên $S$ tồn tại:
\[
\pi_{X/S}:\operatorname{Quot}^n_{\mathcal F/X/S}
\longrightarrow\operatorname{Symm}^n_S(X).
\]
Trường hợp quan trọng nhất là $\mathcal F=\mathcal O_X$, cho ta một cấu xạ:
\[
\pi_{X/S}:\operatorname{Hilb}^n_{X/S}
\longrightarrow\operatorname{Symm}^n_S(X).
\]
Hiển nhiên đây là một đẳng cấu khi $n=0$ hoặc $n=1$. Nhưng khi
$n\geq1$, ngay cả nếu $S$ là phổ của một trường $k$ và $X$ trơn trên
$S$, nói chung đây không phải là một đẳng cấu, thậm chí cũng không phải là một
cấu xạ đơn ánh, vì một lược đồ con chiều 0 của $X$ (chẳng hạn tương ứng với
một iđêan $\mathfrak i$ nguyên sơ đối với iđêan cực đại trong vành địa phương
$\mathcal O_{X,x}$, với $x$ là một điểm đóng của $X$) không được xác định khi
chỉ biết chu trình do nó xác định (trong trường hợp này, khi biết đối chiều
trên $k$ của $\mathfrak i$ trong $\mathcal O_{X,x}$). Ta chỉ có thể nói điều
sau đây ($S$ lại là tùy ý):

\medskip
\noindent a. Nếu $X$ trơn trên $S$, thì cấu xạ chuẩn xác định một đẳng cấu từ
tập mở của $\operatorname{Hilb}^n_{X/S}$ phân loại các phủ étale hạng $n$ nằm
trong $X$ (xem số \textsuperscript{o} 4, (b)) lên tập mở của
$\operatorname{Symm}^n_S(X)$ tương ứng với các $n$-chu trình không có thành
phần bội.
\src{275}

\medskip
\noindent b. Nếu thêm vào đó $X$ có chiều tương đối 1 trên $S$, thì cấu xạ
chuẩn thậm chí xác định một đẳng cấu từ $\operatorname{Hilb}^n_{X/S}$ lên
$\operatorname{Symm}^n_{X/S}$.

Sở dĩ như vậy là vì một lược đồ con chiều 0 của một đường cong đại số không
kỳ dị được xác định hoàn toàn bởi chu trình tương ứng. Cùng nhận xét ấy áp
dụng tổng quát hơn cho các ước Cartier hiệu dụng trên một lược đồ đại số không
kỳ dị (và không loại trừ khả năng rằng, trong trường hợp rất đặc biệt này, đa
tạp Chow cho cùng một đối tượng như lược đồ Hilbert).

\fsection{7}{Các bổ sung và câu hỏi}

Như J.-P. Serre đã nhận xét, từ một ví dụ nổi tiếng của Nagata suy ra rằng có
thể tìm được một lược đồ $S$, là phổ của một trường $k$, một $S$-lược đồ $S'$,
là phổ của một mở rộng bậc hai $k'$ của $k$, và cuối cùng một $S'$-lược đồ $X$
riêng và trơn (nhưng không xạ ảnh) chiều 3, sao cho $\Pi_{S'/S}(X/S')$ không
tồn tại. Điều này càng kéo theo rằng lược đồ Hilbert
$\operatorname{Hilb}^2_{X/S}$ không tồn tại (thậm chí cả $k$-lược đồ biểu diễn
các phủ étale hạng 2 của $S$ nằm trong $X$ cũng không tồn tại, và càng không có
bình phương đối xứng của $X$; xem mục trước). Do đó, khả năng thực hiện các
phép dựng không xạ ảnh trong hình học đại số bị đặt dưới những hạn chế nghiêm
trọng. (Tuy nhiên, có thể tin rằng những hạn chế như vậy sẽ không xuất hiện
trong hình học giải tích, cũng như chúng không xuất hiện trong hình học hình
thức (xem [2], II).) Mặt khác, nếu $X$ là một lược đồ riêng trên phổ $S$ của
một trường $k$, và nếu $Z$ tựa xạ ảnh trên $X$, thì $\Pi_{X/S}(Z/X)$ tồn tại
và là một lược đồ, hợp rời nhau của một dãy các lược đồ tựa xạ ảnh trên $S$
(như trong trường hợp xạ ảnh (3.1)). Thật vậy, ta quy về trường hợp bản thân
$X$ là xạ ảnh bằng cách để $X$ được trội bởi một $S$-lược đồ xạ ảnh $X'$; ở đây
chúng tôi không trình bày chi tiết chứng minh, vốn còn sử dụng kết quả về việc
phân tích một cấu xạ hữu hạn thành hợp được nêu trong [2], I, A,
n\textsuperscript{o} 2 (b). Sự thành công của phương pháp nằm ở chỗ, vì $S$ là
phổ của một trường, $X'$ xuất hiện trong bổ đề Chow tự động phẳng trên $S$.

Tôi không biết kết quả có còn đúng nếu không đặt giả thiết nào lên $S$ hay
không, khi chỉ giả sử $X$ riêng và phẳng trên $S$, còn $Z$ tựa xạ ảnh trên
$X$. Một trường hợp quan trọng trong ứng dụng là khi $Z$ là một lược đồ con
đóng của $X$; nếu khi ấy $\Pi_{X/S}(Z/X)$ tồn tại, thì nó nhất thiết là một
lược đồ con đóng của $S$. Ta có thể
\src{276}
dựng nó trực tiếp khá đơn giản khi $X$ xạ ảnh trên $S$, mà không dùng lý
thuyết lược đồ Hilbert; tổng quát hơn, phương pháp được sử dụng cho thấy rằng
nếu $Z$ affine trên $X$, thì $\Pi_{X/S}(Z/X)$ tồn tại và affine trên $S$. Nó
cũng cho thấy rằng nếu $X$ riêng và phẳng trên $S$ (không nhất thiết xạ ảnh
trên $S$), thì với mọi phân thớ vectơ $Z$ tầm thường địa phương trên $X$,
$\Pi_{X/S}(Z/X)$ tồn tại và là một phân thớ vectơ trên $S$. Rất nên xem xét
lại và thống nhất các kết quả này.

\section*{TÀI LIỆU THAM KHẢO}
\addcontentsline{toc}{section}{Tài liệu tham khảo}

\begin{description}
\item[{[1]}] BOREL (A.) và SERRE (J.-P.). --- Le théorème de Riemann--Roch,
\emph{Bull. Soc. math. France}, tập 86, 1958, tr. 97--136.

\item[{[2]}] GROTHENDIECK (Alexander). --- Technique de construction et
théorèmes d'existence en géométrie algébrique, I, II, III, \emph{Séminaire
Bourbaki}, tập 13, 1959/60, n\textsuperscript{o} 190, 29 tr.,
n\textsuperscript{o} 195, 22 tr.; và tập 14, 1960/61,
n\textsuperscript{o} 212, 20 tr.

\item[{[3]}] GROTHENDIECK (Alexander). --- Séminaire de Géométrie algébrique,
I, II, III, IV. --- Paris, Institut des hautes Études scientifiques, 1960/61
(in rô-nê-ô).

\item[{[4]}] GROTHENDIECK (Alexander). --- Technique de construction en
géométrie analytique, IV, V, \emph{Séminaire Cartan}, tập 13, 1960/61,
n\textsuperscript{o} 11 và 12.

\item[{[5]}] GROTHENDIECK (A.) và DIEUDONNÉ (J.). --- \emph{Éléments de
géométrie algébrique}, I: Le langage des schémas. --- Paris, Presses
universitaires de France, 1960 (Institut des hautes Études scientifiques,
Publications mathématiques, 4); II, III, IV (sắp xuất bản).

\item[{[6]}] KODAIRA (K.). --- Characteristic linear systems of complete
continuous systems, \emph{Amer. J. of Math.}, tập 78, 1956, tr. 716--744.

\item[{[7]}] SERRE (Jean-Pierre). --- Exemples de variétés projectives en
caractéristique $p$ non relevables en caractéristique 0, \emph{Proc. Nat. Acad.
Sc. U.S.A.}, tập 47, 1961, tr. 108--109.

\item[{[8]}] ZAPPA (G.). --- Sull'esistenza, sopra le superficie algebriche, di
sistemi continui completi infiniti, la cui curva è a serie caratteristica
incompleta, \emph{Pont. Acad. Sc. Acta}, tập 9, 1945, tr. 91--93.
\end{description}

\section*{BỔ SUNG}
\addcontentsline{toc}{section}{Bổ sung}
\begin{center}
{\small [bổ sung khi sửa bản in thử]}
\end{center}

Giờ đây có thể thấy rằng các phỏng đoán của [2], III,
n\textsuperscript{o} 8 là sai, ngay cả đối với các đa tạp không kỳ dị trên một
trường đặc trưng 0, cả về sự tồn tại lẫn tính tựa xạ ảnh của thương, và ngay cả
khi $G$ tác động với đồ thị đóng.

\par

\clearpage
\chapter*{Đính chính cho Chuyên đề số 221}
\addcontentsline{toc}{chapter}{Đính chính cho Chuyên đề số 221}
\setcounter{footnote}{3}
\sid{e221}
\markboth{Đính chính 221}{Đính chính 221}
\src{302}

\noindent\textbf{6. TDTE IV: Các lược đồ Hilbert.}

\begin{center}
[Hội thảo Bourbaki, tập 13, 1960/61, số 221, 28 tr.]
\end{center}

\noindent\textbf{Trang 6, Định lý 2.1, dòng thứ 3.} --- Thay vì
«cần», đọc là «cần và đủ».

\noindent\textbf{Trang 15, dòng 13.} --- Thay vì «$P\neq Q$», đọc là
«$P<Q$».

\noindent\textbf{Dòng 15.} --- Thay vì «Nếu trái lại $P=Q$», đọc là
«Trong trường hợp còn lại, ta có $P(n)\geq Q(n)$ khi $n$ đủ lớn».

\noindent\textbf{Dòng 18, công thức (*).} --- Thay vì «$<$» và «$>$»,
đọc là «$\leq$» và «$\geq$».

\noindent\textbf{Dòng thứ 5 tính từ dưới lên.} --- Thay vì «$r$», đọc là
«$Z$».

\noindent\textbf{Trang 18, Nhận xét 3.9.} --- Xin lưu ý rằng Hartshorne đã
nghiên cứu các thành phần liên thông của những lược đồ Hilbert trên một trường
đóng đại số. Ông chứng minh rằng các $\operatorname{Hilb}^{P}_{r}$ là liên
thông và xác định những cặp $(r,P)$ sao cho
$\operatorname{Hilb}^{P}_{r}\neq\varnothing$\footnote{HARTSHORNE
(R.). --- Connectedness of the Hilbert scheme (Luận án, Harvard, 1962).}.

\noindent\textbf{Trang 23, dòng thứ 5 tính từ dưới lên.} --- Thay vì «một
thành phần», đọc là «của một thành phần».

\noindent\textbf{Trang 23 và 24, Nhận xét 5.5.} --- Liên quan đến ví dụ của
Zappa, xin lưu ý rằng Mumford thậm chí vừa xây dựng một thành phần bất khả quy
của lược đồ Hilbert của $\mathbf P^3_{\mathbf C}$ (các điểm tổng quát của nó
biểu diễn những đường cong không kỳ dị bậc 14, giống 24); thành phần này không
rút gọn tại điểm tổng quát của nó. Bằng cách thổi lên dọc theo một trong các
đường cong thu được, ông cũng thu được một lược đồ xạ ảnh chính quy chiều 3
trên $\mathbf C$, mà lược đồ moduli hình thức của nó không rút gọn tại điểm
tổng quát; hay tương đương, đa tạp moduli địa phương của nó, theo nghĩa của
hình học giải tích trên $\mathbf C$ (xem Hội thảo Cartan, tập 13, 1960/61),
không rút gọn tại mọi điểm.

\par

\unit{Thuyết trình 232 -- Các lược đồ Picard: các định lý tồn tại}
\sid{232}
\markboth{Thuyết trình 232}{Thuyết trình 232}
\src{143}
\seminar{Năm thứ 14\textsuperscript{}, 1961/62, số 232\textsuperscript{}\\Tháng Hai 1962}
\paperhead{KỸ THUẬT HẠ XUỐNG VÀ CÁC ĐỊNH LÝ TỒN TẠI TRONG HÌNH HỌC ĐẠI SỐ\\
V. CÁC LƯỢC ĐỒ PICARD: CÁC ĐỊNH LÝ TỒN TẠI}%
  {do Alexander GROTHENDIECK}

\fsection{1}{Các nhóm và hàm tử Picard tương đối}

Với mọi tiền lược đồ (tổng quát hơn, mọi không gian có vành) $X$, ta gọi
\emph{nhóm Picard} (tuyệt đối) của $X$, và ký hiệu là
$\operatorname{Pic}(X)$, nhóm các lớp đẳng cấu của các Môđun khả nghịch
(nghĩa là đẳng cấu địa phương với $\mathcal O_X$) trên $X$. Do đó ta có một
đẳng cấu chính tắc
\begin{equation}
\operatorname{Pic}(X)\xrightarrow{\sim}H^1(X,\mathcal O_X^*),
\tag{1.1}
\end{equation}
trong đó $\mathcal O_X^*$ chỉ bó các đơn vị của $\mathcal O_X$ (quả thật bó
này đồng nhất với bó các tự đẳng cấu của Môđun khả nghịch chuẩn
$\mathcal O_X$). Lưu ý rằng $X\longmapsto\operatorname{Pic}(X)$ hiển nhiên là
một hàm tử phản biến theo $X$, và đẳng cấu (1.1) có tính hàm tử.

Nếu $X$ là một tiền lược đồ trên một tiền lược đồ khác $S$, thì khi $S'$ biến
thiên trong phạm trù $(\mathrm{Sch})/S$ các tiền lược đồ trên $S$, từ điều trên
ta có hàm tử phản biến
$S'\longmapsto\operatorname{Pic}(X\times_S S')$. Hàm tử này hoàn toàn không
thể «biểu diễn được» ([4], II, A), vì do tồn tại những tự đẳng cấu không tầm
thường của các Môđun khả nghịch cần phân loại, hàm tử này không «có bản chất
địa phương» ([5], IV, 5.4). Vì vậy cần «làm cho nó thành địa phương» bằng cách,
nói chung, với mọi tiền lược đồ tương đối $X/S$, đưa vào một nhóm có bản chất
tương đối,
\begin{equation}
\operatorname{Pic}'(X/S)=H^0\bigl(S,R^1f_*(\mathcal O_X^*)\bigr),
\tag{1.2}
\end{equation}
(trong đó $f:X\longrightarrow S$ là cấu xạ cấu trúc) (so sánh [4], II,
C 3). Ở chỗ đã dẫn, nhóm này được gọi là nhóm Picard tương đối; ở đây nên gọi
nó là \emph{nhóm Picard tương đối hạn chế} của $X/S$, vì những lý do sẽ xuất
hiện sau đây. Khi $S'$ biến thiên trong $(\mathrm{Sch})/S$,
$S'\longmapsto\operatorname{Pic}'(X\times_S S'/S')$ là một hàm tử phản biến
theo $S'$, cũng được ký hiệu là $\operatorname{Pic}'_{X/S}$, do đó về cơ bản
được định nghĩa bởi công thức
\begin{equation}
\operatorname{Pic}'_{X/S}(S')=
\operatorname{Pic}'(X\times_S S'/S').
\tag{1.3}
\end{equation}
\src{144}
Hàm tử này nay «có bản chất địa phương», vì ta đã làm điều cần thiết cho mục
đích ấy. Một cách trực giác, vế phải của (1.3) được hiểu là tập các «họ đại số»
của những lớp bó khả nghịch trên (các thớ của) $X/S$, được đánh chỉ số bởi tiền
lược đồ tham số $S'/S$. Khi hàm tử $\operatorname{Pic}'$ biểu diễn được, tiền
lược đồ trên $S$ biểu diễn nó được ký hiệu là $\operatorname{Pic}'_{X/S}$ và
được gọi là \emph{tiền lược đồ Picard} của $X$ trên $S$; do đó khi ấy ta có
\begin{equation}
\operatorname{Hom}_S(S',\operatorname{Pic}'_{X/S})
\simeq\operatorname{Pic}'_{X/S}(S')
=\operatorname{Pic}'(X\times_S S'/S').
\tag{1.4}
\end{equation}

Tuy nhiên, có những trường hợp quan trọng trong đó
$\operatorname{Pic}'_{X/S}$ không biểu diễn được (ví dụ: một đa tạp
«Brauer--Severi» trên một trường $k$, không có điểm hữu tỉ trên $k$), nhưng vẫn
có một định nghĩa tự nhiên cho tiền lược đồ Picard tương đối. Điều này là do
trong định nghĩa hàm tử $\operatorname{Pic}'$ từ các nhóm Picard tuyệt đối
$\operatorname{Pic}(X\times_S S'/S')$, ta vẫn chưa địa phương hóa đủ; chính
xác hơn, nói chung $\operatorname{Pic}'$ không «tương thích với hạ phẳng trung
thành». Ta làm rõ điều này.

Gọi $\mathfrak M$ là lớp các cấu xạ tiền lược đồ \emph{phẳng trung thành và tựa
compact}; lớp này ổn định qua đổi cơ sở và qua hợp thành. Cho $P$ là một hàm tử
phản biến từ $(\mathrm{Sch})/S$ đến phạm trù các tập hợp, và với mọi
$S$-cấu xạ $u:T'\longrightarrow T$, $u\in\mathfrak M$, xét biểu đồ
\begin{equation}
P(T)\longrightarrow P(T')\rightrightarrows P(T'\times_TT')
\tag{1.5}
\end{equation}
là ảnh qua $P$ của biểu đồ
\[
T\longleftarrow T'\mathrel{\substack{\xleftarrow{\,\mathrm{pr}_1\,}\\[-0.35ex]
\xleftarrow[\,\mathrm{pr}_2\,]{}}}T'\times_TT'.
\]
Nếu $P$ biểu diễn được, thì từ lý thuyết hạ xuống ([4], I, B, định lý 2) suy ra
biểu đồ (1.5) là khớp với mọi $u\in\mathfrak M$. Ta diễn đạt sự kiện này bằng
cách nói $P$ tương thích với $\mathfrak M$, trong trường hợp này là $P$
«tương thích với hạ phẳng trung thành», hay «tiền bó» $P$ trên
$(\mathrm{Sch})/S$ là một «bó» đối với khái niệm địa phương hóa do lớp
$\mathfrak M$ cung cấp. Khi $P$ tùy ý, một quy trình chuẩn, quen thuộc trong
trường hợp địa phương hóa tôpô thông thường, cho phép liên kết với nó một «bó»
$\mathcal P$ và một đồng cấu hàm tử $P\longrightarrow\mathcal P$, phổ dụng
theo nghĩa hiển nhiên. Có thể tính $\mathcal P$ một cách tường minh như sau:
để định nghĩa $\mathcal P(T)$, với mọi $T'$ trên $T$ sao cho cấu xạ
$u:T'\longrightarrow T$ thuộc $\mathfrak M$, ta ký hiệu bởi
\src{145}
$\overline H^0(T'/T,P)$ tập con của $P(T')$ gồm các phần tử $\xi$ sao cho hai
ảnh $\xi_1,\xi_2$ của chúng trong $P(T'\times_TT')$ thỏa điều kiện: tồn tại
một cấu xạ $v:T''\longrightarrow T'\times_TT'$, $v\in\mathfrak M$, sao cho
$\xi_1$ và $\xi_2$ có cùng ảnh trong $P(T'')$. [Ghi chú. -- Tập
$\overline H^0$ được định nghĩa như vậy lớn hơn tập $H^0(T'/T,P)$ đã đưa vào
trong [4], I, A, 4(a).] Khi $T'$ biến thiên trên $T$ cố định (vẫn với
$u\in\mathfrak M$), các $\overline H^0(T'/T,P)$ tạo thành một hệ quy nạp (khi
tập các $T'$ được trang bị tiền thứ tự xác định bởi quan hệ trội), và ta đặt
\begin{equation}
\mathcal P(T)=\varinjlim_{T'}\overline H^0(T'/T,P).
\tag{1.6}
\end{equation}
Quy luật hàm tử theo $T$ của biểu thức này được tường minh hóa một cách hiển
nhiên.

Khi
\[
P(T)=\operatorname{Pic}(X\times_ST),
\]
hàm tử phản biến trên $(\mathrm{Sch})/S$ được định nghĩa bởi (1.6) được gọi là
\emph{hàm tử Picard tương đối} của $X$ trên $S$, và được ký hiệu là
$\operatorname{Pic}_{X/S}$; ta gọi \emph{nhóm Picard tương đối} của $X$ trên
$S$, và ký hiệu là $\operatorname{Pic}(X/S)$, nhóm
$\operatorname{Pic}_{X/S}(S)$. Khi đó ta có song ánh hiển nhiên
\begin{equation}
\operatorname{Pic}_{X/S}(T)\simeq\operatorname{Pic}(X\times_ST/T).
\tag{1.7}
\end{equation}

Do đó một phần tử của $\operatorname{Pic}(X/S)$ được định nghĩa nhờ một phần tử
$\xi'$ của một nhóm $\operatorname{Pic}(X\times_SS')$ (trong đó
$S'\longrightarrow S$ phẳng trung thành và tựa compact), sao cho có thể tìm
được một cấu xạ phẳng trung thành và tựa compact
$S''\longrightarrow S'\times_SS'$ sao cho hai ảnh ngược của $\xi'$ trong
$\operatorname{Pic}(X\times_SS'')$ bằng nhau. Một phần tử $\xi'$ của
$\operatorname{Pic}(X\times_SS')$ và $\xi_1$ của
$\operatorname{Pic}(X\times_SS_1)$ (thỏa các điều kiện vừa nêu) xác định cùng
một phần tử của $\operatorname{Pic}(X/S)$ khi và chỉ khi tồn tại một cấu xạ
phẳng trung thành và tựa compact $S'_1\longrightarrow S'\times_SS_1$ sao cho
ảnh của hai phần tử đang xét trong $\operatorname{Pic}(X\times_SS'_1)$ bằng
nhau. Thường vẫn tiện khi làm việc với hàm tử
$P'=\operatorname{Pic}'_{X/S}$ đã đưa vào ở trên; ta thấy ngay rằng cấu xạ
chính tắc $P\longrightarrow P'$ xác định một đẳng cấu
\begin{equation}
\mathcal P\xrightarrow{\sim}\mathcal P',
\tag{1.8}
\end{equation}
điều này cho một mô tả của $\operatorname{Pic}_{X/S}$ theo
$\operatorname{Pic}'_{X/S}=P'$, thường tiện hơn. Theo 2.3 dưới đây, khi trong
mô tả vừa nêu của $\operatorname{Pic}(X/S)$ ta thay $P$ bởi $P'$, ta thực ra
có thể
\src{146}
lấy $S''=S'\times_SS'$, $S'_1=S'\times_SS_1$, ít nhất dưới các điều kiện được
nêu rõ ở chỗ đã dẫn.

Khi hàm tử $\operatorname{Pic}_{X/S}$ biểu diễn được, ta nói $X/S$ có một tiền
lược đồ Picard, và tiền lược đồ trên $S$ biểu diễn hàm tử được gọi là
\emph{tiền lược đồ Picard} của $X$ trên $S$, vẫn được ký hiệu là
$\operatorname{Pic}_{X/S}$. Hiển nhiên, chỉ cần
$P'=\operatorname{Pic}'_{X/S}$ biểu diễn được, vì khi ấy $P'$ đã là một «bó»,
và công thức (1.8) chứng tỏ cấu xạ $P'\longrightarrow\mathcal P'$ đồng nhất
với cấu xạ chính tắc
\begin{equation}
\operatorname{Pic}'_{X/S}\longrightarrow\operatorname{Pic}_{X/S},
\tag{1.9}
\end{equation}
mà khi ấy là một đẳng cấu. Do đó thuật ngữ của ta tương thích với thuật ngữ đã
đưa vào ở trên cùng với (1.4). Nói chung, khi $\operatorname{Pic}_{X/S}$ tồn
tại, nó được định nghĩa bởi đẳng cấu có tính hàm tử
\begin{equation}
\operatorname{Hom}_S(S',\operatorname{Pic}_{X/S})\xrightarrow{\sim}
\operatorname{Pic}(X\times_SS'/S').
\tag{1.10}
\end{equation}

\fsection{2}{Quan hệ giữa các nhóm Picard tương đối và tuyệt đối khác nhau}

\result{MỆNH ĐỀ}{2.1} Cho $f:X\longrightarrow S$ là một cấu xạ sao cho
$\mathcal O_S\xrightarrow{\sim}f_*(\mathcal O_X)$. Khi đó ta có một dãy khớp
\[
0\longrightarrow\operatorname{Pic}(S)\longrightarrow\operatorname{Pic}(X)
\longrightarrow\operatorname{Pic}'(X/S).
\]
Khi $X$ có một tiết diện trên $S$, cấu xạ cuối là toàn ánh, nghĩa là ta có một
đẳng cấu
\[
\operatorname{Pic}'(X/S)\xrightarrow{\sim}
\operatorname{Pic}(X)/\operatorname{Pic}(S).
\]

Dãy khớp này có thể được xem là dãy khớp ở các bậc thấp tương ứng với dãy phổ
Leray cho $f$ và $\mathcal O_X^*$. Khẳng định thứ hai cũng có tính hình thức.

\result{MỆNH ĐỀ}{2.2} Cho $f:X\longrightarrow S$ là một cấu xạ tựa compact và
tách sao cho $\mathcal O_S\xrightarrow{\sim}f_*(\mathcal O_X)$, và cho
$S'\longrightarrow S$ là một cấu xạ phẳng trung thành và tựa compact. Khi đó:

\noindent (i) $\operatorname{Pic}'(X/S)\longrightarrow
\operatorname{Pic}'(X\times_SS'/S')$ là đơn ánh;

\noindent (ii) nếu $X$ có một tiết diện địa phương trên $S$ (nghĩa là mỗi
$s\in S$ có một lân cận mở $U$ sao cho $X|U$ có một tiết diện trên $U$), thì
biểu đồ
\[
\operatorname{Pic}'(X/S)\longrightarrow
\operatorname{Pic}'(X\times_SS'/S')\rightrightarrows
\operatorname{Pic}'(X\times_SS''/S'')
\]
\src{147}
(trong đó $S''=S'\times_SS'$) là khớp.

Khẳng định thứ nhất suy ra, nhờ các tính chất sơ cấp của hạ phẳng trung thành,
từ nhận xét tổng quát sau đây. Nếu $f:X\longrightarrow S$ là một cấu xạ sao cho
$\mathcal O_S\xrightarrow{\sim}f_*(\mathcal O_X)$, thì hàm tử
$\mathcal F\longmapsto f^*(\mathcal F)$, từ phạm trù các Môđun tự do địa
phương kiểu hữu hạn trên $S$ đến phạm trù các Môđun tự do địa phương kiểu hữu
hạn trên $X$, là đầy đủ và trung thành, và ảnh cốt yếu gồm các Môđun
$\mathcal G$ trên $X$ sao cho $f_*(\mathcal G)$ tự do địa phương và đồng cấu
chính tắc
\[
f^*\bigl(f_*(\mathcal G)\bigr)\longrightarrow\mathcal G
\]
là một đẳng cấu. Khẳng định thứ hai đã được chứng minh bằng lý thuyết hạ xuống
trong [4], I, B, 4.

Các kết quả của 2.2 cũng có thể được phát biểu như sau:

\result{HỆ QUẢ}{2.3} Dưới các điều kiện của 2.2, đồng cấu chính tắc
(1.9) $\operatorname{Pic}'_{X/S}\longrightarrow\operatorname{Pic}_{X/S}$ là
đơn ánh, và thậm chí song ánh nếu $X$ có một tiết diện địa phương trên $S$.
(Do đó, trong trường hợp sau, nhóm Picard tương đối
$\operatorname{Pic}(X/S)$ và nhóm Picard tương đối hạn chế
$\operatorname{Pic}'(X/S)$ đồng nhất với nhau.)

Kết hợp với 2.1, ta thu được:

\result{HỆ QUẢ}{2.4} Dưới các điều kiện của 2.2, ta có một dãy khớp
\[
0\longrightarrow\operatorname{Pic}(S)\longrightarrow\operatorname{Pic}(X)
\longrightarrow\operatorname{Pic}(X/S).
\]
Khi $X$ có một tiết diện trên $S$, đồng cấu cuối là toàn ánh, nghĩa là khi ấy
ta có một đẳng cấu
\[
\operatorname{Pic}(X/S)\xrightarrow{\sim}
\operatorname{Pic}(X)/\operatorname{Pic}(S).
\]

\result{NHẬN XÉT}{2.5} Cho $f:X\longrightarrow S$ là một cấu xạ sao cho
$\mathcal O_S\xrightarrow{\sim}f_*(\mathcal O_X)$, và cho $g$ là một tiết
diện của $X$ trên $S$. Cho $\mathcal L$ là một Môđun khả nghịch trên $X$; gọi
$g$-rigid hóa của $\mathcal L$ một đẳng cấu
$\mathcal O_S\xrightarrow{\sim}g^*(\mathcal L)$, và gọi
\emph{Môđun khả nghịch được $g$-rigid hóa} một Môđun khả nghịch $\mathcal L$
trên $X$ được trang bị một phép $g$-rigid hóa. Mọi tự đẳng cấu của một cấu
trúc như vậy đều tầm thường, và
$\operatorname{Pic}'(X/S)$ đồng nhất với tập các lớp đẳng cấu của những Môđun
khả nghịch được $g$-rigid hóa trên $S$. (Chính sự kiện này đã cho phép dùng lý thuyết hạ
xuống để chứng minh 2.2(ii).) Điều này cho một cách hiểu mới về
$\operatorname{Pic}(X/S)$, ít nhất khi $f$ còn tựa compact và tách, do đó
$\operatorname{Pic}(X/S)\simeq\operatorname{Pic}'(X/S)$ theo 2.3.
\src{148}
\result{NHẬN XÉT}{2.6} Cho $f:X\longrightarrow S$ là một cấu xạ như trong
2.2, và cho $S'\longrightarrow S$ là một cấu xạ phẳng trung thành và tựa
compact sao cho tồn tại một $S$-cấu xạ $S'\longrightarrow X$, nghĩa là sao cho
tồn tại một tiết diện của $X'=X\times_SS'$ trên $S'$. Đặt
$S''=S'\times_SS'$, $X''=X\times_SS''$, và xét dãy khớp
\[
\operatorname{Pic}(X/S)\longrightarrow\operatorname{Pic}(X'/S')
\rightrightarrows\operatorname{Pic}(X''/S'').
\]
Áp dụng 2.4, ta thu được một dãy khớp
\[
\operatorname{Pic}(X/S)\longrightarrow
\operatorname{Pic}(X')/\operatorname{Pic}(S')
\rightrightarrows
\operatorname{Pic}(X'')/\operatorname{Pic}(S''),
\]
đặc biệt, mọi phần tử của nhóm Picard tương đối đã «đến từ» một phần tử của
$\operatorname{Pic}(X')$. Như vậy điều này đơn giản hóa đáng kể mô tả nhóm
Picard tương đối đã cho trong mục trước, và tương tự đối với hàm tử Picard của
$X$ trên $S$, vì với mọi $T$ trên $S$, ta có thể áp dụng điều trên cho
$X\times_ST/T$ và cấu xạ $T'=S'\times_ST\longrightarrow T$. Chẳng hạn, nếu
chính $f$ phẳng trung thành, ta có thể lấy $S'=X$; điều này cho phép, khi $f$
còn thuộc kiểu hữu hạn (tương ứng trơn, v.v.), trong mô tả hàm tử Picard tương
đối $T\longmapsto\operatorname{Pic}(X\times_ST/T)$, chỉ cần xét các phép đổi
cơ sở $T'\longrightarrow T$ thuộc kiểu hữu hạn (tương ứng trơn, v.v.). Khi
$f$ xạ ảnh và phẳng, $S$ Noether địa phương, ta chứng minh được rằng có thể
chọn $S'\longrightarrow S$ ở trên sao cho $S'$ là hợp rời của các phủ phẳng
$S'_i$ của các tập mở $S_i$ của $S$ phủ $S$; nếu $f$ còn khả phân, có thể lấy $S'_i$
étale trên $S_i$.

\fsection{3}{Định lý tồn tại chính: phát biểu}

Ta chưa có, ngay cả dưới dạng phỏng đoán, một phát biểu về sự tồn tại của các
tiền lược đồ Picard bao quát mọi trường hợp đã biết. Một điều kiện «hầu như
cần thiết», nếu có thể nói như vậy, là $f:X\longrightarrow S$ phải riêng (bảo
đảm các tính chất hữu hạn thiết yếu) và phẳng. Các điều kiện này không đủ,
ngay cả khi $S$ là phổ của đại số các số đối ngẫu $k[t]/(t^2)$ trên một trường
$k$ (chẳng hạn trường các số phức $\mathbf C$), và $X$ có chiều 1. Vào thời
điểm viết thuyết trình này, các định lý tồn tại quan trọng nhất cho tiền lược
đồ Picard được suy ra từ định lý sau:

\result{ĐỊNH LÝ}{3.1} Cho $f:X\longrightarrow S$ là một cấu xạ giữa các tiền
lược đồ Noether địa phương. Giả sử:

\noindent (i) $f$ xạ ảnh;
\src{149}
\noindent (ii) $f$ phẳng;

\noindent (iii) các thớ hình học của $f$ là nguyên.

Dưới các điều kiện này, $\operatorname{Pic}_{X/S}$ tồn tại.

Chứng minh, sẽ được phác thảo trong hai mục tiếp theo, đồng thời sẽ chỉ ra điều
sau: cho $\xi$ là tiết diện của $\operatorname{Pic}_{X/S}$ tương ứng với một
bó rất rộng $\mathcal O_X(1)$ trên $X/S$ (nghĩa là được cảm sinh bởi một phép
nhúng xạ ảnh $X\longrightarrow\mathbf P(\mathcal E)$); khi đó tồn tại một tập
con mở $U$ của $\operatorname{Pic}_{X/S}$, là hợp rời của các tập con mở tựa
xạ ảnh trên $S$, sao cho $U$ ổn định dưới phép tịnh tiến bởi $\xi$, và
$\operatorname{Pic}_{X/S}$ là hợp tăng của các tập con mở $U-n\xi$ (tất cả
đều đẳng cấu với $U$). Đặc biệt suy ra rằng dưới các điều kiện của 3.1,
$\operatorname{Pic}_{X/S}$ tách trên $S$.

\result{NHẬN XÉT}{3.2} Các ví dụ cho thấy (chẳng hạn với $S$ là phổ của một
vành định giá rời rạc, và $X$ có chiều tương đối 1 trên $S$) rằng nếu trong
3.1 ta bỏ giả thiết (iii), thay nó bằng giả thiết yếu hơn rằng với mọi
$s\in S$, đồng cấu
\[
k(s)\longrightarrow H^0(X_s,\mathcal O_{X_s})
\]
là một đẳng cấu, thì $\operatorname{Pic}_{X/S}$ không nhất thiết tách trên
$S$; điều này xảy ra cả trong trường hợp các thớ hình học của $f$ rút gọn,
nhưng một thớ tổng quát hình học nguyên «vỡ» do chuyên biệt hóa thành hai thành
phần bất khả quy, lẫn trong trường hợp các thớ hình học của $f$ bất khả quy,
nhưng một thớ tổng quát hình học nguyên chuyên biệt hóa thành một «thớ bội».
Trường hợp thứ nhất xảy ra chẳng hạn với một đường conic suy biến thành hai
đường thẳng cắt nhau; một ví dụ của trường hợp thứ hai do D. Mumford cung cấp,
với một đường cong elliptic suy biến thành một đường cong elliptic kép. Các ví
dụ này đúng trong mọi đặc số.

\result{NHẬN XÉT}{3.3} Dưới các điều kiện của 3.1, tôi không biết
$\operatorname{Pic}_{X/S}$ có phải là một hợp rời của các tập con mở thuộc
kiểu hữu hạn, do đó tựa xạ ảnh, trên $S$ hay không. Lưu ý rằng việc xét các
đa thức Hilbert $Q\in\mathbf Q[t]$ cho phép, như trong trường hợp các lược đồ
Hilbert ([4], IV), phân tích $\operatorname{Pic}_{X/S}$ thành một hợp rời của
các tập con mở $\operatorname{Pic}^{Q}_{X/S}$, và có vẻ hợp lý rằng các tập
sau là tựa xạ ảnh trên $S$; ít nhất ta sẽ thấy điều đó trong thuyết trình kế
tiếp khi $f$ là một cấu xạ trơn. Cần lưu ý rằng nếu thay giả thiết (i) bằng
giả thiết: $X$ xạ ảnh địa phương trên $S$ (điều này đủ để 3.1 đúng, vì bài toán
tồn tại của $\operatorname{Pic}_{X/S}$ hiển nhiên có tính địa phương trên
$S$), thì ngược lại, dễ dàng
\src{150}
đưa ra những ví dụ trong đó $\operatorname{Pic}_{X/S}$ chứa các thành phần
liên thông không thuộc kiểu hữu hạn trên $S$. Chẳng hạn, cho $X_0$ là một đa
tạp đại số xạ ảnh không kỳ dị trên một trường đóng đại số $k$, được trang bị
một tự đẳng cấu $u$ và một phần tử $\xi$ trong nhóm Néron--Severi của $X_0$,
sao cho các $u^n(\xi)$ đôi một khác nhau. Chẳng hạn có thể lấy $X_0$ là tích
của một đường cong elliptic $E$ với chính nó, và lấy $u$ là tự đẳng cấu
\[
(x,y)\longmapsto(x,y+x)
\]
của $E\times E$. Cho $S$ là hợp của hai đường cong bất khả quy không kỳ dị,
cắt nhau tại hai điểm $a$ và $b$. Trên $S$ có một phủ chính liên thông $P$ với
nhóm $\mathbf Z$; dùng các tác động của $\mathbf Z$ trên $X_0$ được xác định
bởi $u$, ta suy ra một phân thớ liên kết trên $S$ với thớ $X_0$ (tầm thường
trên $S-a$ và $S-b$), và trong trường hợp riêng đang xét, thực ra là một lược
đồ Abel trên $S$. Dễ thấy rằng $\operatorname{Pic}_{X/S}$, cũng là phân thớ
liên kết với $P$ và với các tác động của $\mathbf Z$ trên
$\operatorname{Pic}_{X_0/k}$ qua $u$, chứa một thành phần liên thông đẳng cấu
với
\[
P\times\operatorname{Pic}^{0}_{X_0/k},
\]
(trong đó $\operatorname{Pic}^{0}$ chỉ thành phần liên thông của phần tử đơn
vị trong $\operatorname{Pic}$), và thành phần này không thuộc kiểu hữu hạn
trên $S$. [Những hiện tượng tương tự cũng xảy ra trong nhiều trường hợp tiền
lược đồ Picard không tách trên $S$, như đã xét trong 3.3.]

\fsection{4}{Các ước Cartier tương đối và các phân thớ xạ ảnh}

Ta sẽ chỉ cần dùng các ước dương, và trong phần còn lại của mục này sẽ bỏ qua
tính từ bổ sung ấy.

Cho $X$ là một tiền lược đồ. Một \emph{ước Cartier} trên $X$, hay để đơn giản
chỉ gọi là ước trên $X$, là một tiền lược đồ con đóng $D$ của $X$, được xác
định bởi một Iđêan $\mathcal J$ là một Môđun khả nghịch, nghĩa là được sinh địa
phương bởi một tiết diện không là ước của không của $\mathcal O_X$. Với $D$,
ta liên kết một Môđun khả nghịch, cụ thể là
\[
\mathcal L(D)=\mathcal J^{-1},
\]
và đơn ánh chính tắc $\mathcal J\longrightarrow\mathcal O_X$ cho một đồng cấu
chính tắc
\[
s_D:\mathcal O_X\longrightarrow\mathcal J^{-1}=\mathcal L(D),
\qquad\text{nghĩa là}\quad s_D\in\Gamma(X,\mathcal L(D)).
\]
Hơn nữa, dữ liệu của một ước về cơ bản tương đương với dữ liệu của một Môđun
khả nghịch $\mathcal L$ trên $X$, được trang bị một tiết diện $s$ không là ước
của không ở mọi nơi; ngược lại, với một cặp như vậy $(\mathcal L,s)$, ta liên
kết «ước» của $s$, ký hiệu là $\operatorname{div}(s)$. Với một
$\mathcal L$ khả nghịch đã cho trên $X$, tập các
\src{151}
ước $D$ xác định $\mathcal L$ tương ứng song ánh với tập thương
\[
\Gamma(X,\mathcal L)^*/\Gamma(X,\mathcal O_X^*),
\]
trong đó $\Gamma(X,\mathcal L)^*$ chỉ bộ phận của
$\Gamma(X,\mathcal L)$ gồm các tiết diện không là ước của không ở mọi nơi.

Bây giờ giả sử ta có một cấu xạ thuộc kiểu hữu hạn địa phương
$f:X\longrightarrow S$, và để đơn giản giả sử $S$ Noether địa phương. Cho
$\mathcal J$ là một Iđêan coherent trên $X$, $D$ là lược đồ con của $X$ do nó
xác định, $x\in X$, và $s=f(x)$. Ta chứng minh được rằng các điều kiện sau
tương đương:

\noindent (i) $\mathcal J$ khả nghịch tại $x$ (nghĩa là $\mathcal J_x$ được
sinh bởi một phần tử chính quy của $\mathcal O_{X,x}$) và $D$ phẳng trên $S$
tại $x$;

\noindent (ii) $X$ và $D$ phẳng trên $S$ tại $x$, và $D_s$ là một ước Cartier
trên thớ $X_s$ tại điểm $x$;

\noindent (iii) $X$ phẳng trên $S$ tại $x$, và $\mathcal J_x$ được sinh bởi
một phần tử $f_x$ cảm sinh trên $X_s$ một mầm không là ước của không.

Khi đó ta nói $D$ là một \emph{ước Cartier tương đối} trên $X/S$ (hay chỉ là
ước tương đối trên $X/S$) tại điểm đang xét. Từ (i), lưu ý rằng khi đó $D$
cũng là một ước tương đối tại các điểm lân cận của $x$; do đó nếu $X$ và $D$
phẳng trên $S$, còn $D$ riêng trên $S$, thì tập các $s\in S$ sao cho $D_s$ là
một ước Cartier trong $X_s$ (nghĩa là sao cho $D$ là một ước Cartier tương đối
tại các điểm của $X_s$) là một bộ phận mở của $S$. Mặt khác, định nghĩa trên
đã được đặt ra chính xác để khái niệm ước Cartier tương đối ổn định qua một
phép đổi cơ sở tùy ý $S'\longrightarrow S$. Bây giờ xét tập
$\operatorname{Div}(X/S)$ các ước tương đối trên $X/S$, rồi hàm tử phản biến
theo $S'$ biến thiên trên $S$ được định nghĩa bởi
\[
\operatorname{Div}_{X/S}(S')=
\operatorname{Div}(X\times_SS'/S').
\]

Giả sử $X$ phẳng và riêng trên $S$. Khi đó, theo đặc trưng (ii) của các ước
Cartier tương đối, $\operatorname{Div}_{X/S}$ có thể được xem là một hàm tử
con của hàm tử $\operatorname{Hilb}_{X/S}$ được định nghĩa trong [4], IV, và
cấu xạ bao hàm
\[
\operatorname{Div}_{X/S}\longrightarrow\operatorname{Hilb}_{X/S}
\]
«biểu diễn được bởi các phép nhúng mở» (xem [5], IV, 3.13) theo các nhận xét
trên. Dùng định lý tồn tại chính của [4], IV, ta thu được:
\src{152}
\result{MỆNH ĐỀ}{4.1} Giả sử $f:X\longrightarrow S$ xạ ảnh và phẳng. Khi đó
hàm tử $\operatorname{Div}_{X/S}$ biểu diễn được, và chính xác hơn, được biểu
diễn bởi một tập mở của $\operatorname{Hilb}_{X/S}$.

Dùng (đối với một bó $\mathcal O_X(1)$ rất rộng đã cho trên $X/S$) phân tích
chính tắc của $\operatorname{Hilb}_{X/S}$ thành tổng các tập mở
$\operatorname{Hilb}^{Q}_{X/S}$ tương ứng với các đa thức Hilbert
$Q\in\mathbf Q[t]$, ta cũng suy ra một phân tích tương tự
\[
\operatorname{Div}_{X/S}=
\coprod_{Q\in\mathbf Q[t]}\operatorname{Div}^{Q}_{X/S}
\]
thành tổng các tập mở rời nhau, tựa xạ ảnh trên $S$.

Mặt khác, dùng ánh xạ $D\longmapsto\mathcal L(D)$, ta thu được một đồng cấu có
tính hàm tử
\begin{equation}
\operatorname{Div}_{X/S}\longrightarrow\operatorname{Pic}_{X/S}
\tag{+}
\end{equation}
mà ta sẽ nghiên cứu; dưới những điều kiện khá tổng quát, nó sẽ tỏ ra biểu diễn
được tương đối ([5], IV, 3). Vì vậy, hãy xuất phát với một phần tử $\xi$ của
$\operatorname{Pic}_{X/S}(S')$, để đơn giản ký hiệu giả sử $S'=S$, và chứng
minh rằng hàm tử con tương ứng của $\operatorname{Div}_{X/S}$ biểu diễn được.
Trước hết, xét trường hợp $\xi$ được xác định bởi một Môđun khả nghịch
$\mathcal L$ trên $X$. Ta sẽ giả sử $X$ riêng và phẳng trên $S$, và các thớ
hình học của $X$ trên $S$ là nguyên; điều này cũng kéo theo ([2], III, đoạn 7)
rằng ta có $\mathcal O_S\xrightarrow{\sim}f_*(\mathcal O_X)$ và quan hệ này
vẫn đúng sau mọi phép đổi cơ sở $S'\longrightarrow S$. Khi đó, các ước Cartier
tương đối $D$ trên $X/S$ sao cho $\mathcal L(D)$ và $\mathcal L$ xác định cùng
một phần tử của $\operatorname{Pic}(X/S)=\operatorname{Pic}_{X/S}(S)$, nghĩa
là theo 2.4, sao cho $\mathcal L(D)$ và $\mathcal L$ đẳng cấu địa phương trên
$S$, tương ứng song ánh với các tiết diện của bó thương
$f_*(\mathcal L)^*/\mathcal O_S^*$. Sự tương ứng này tương thích với các phép
đổi cơ sở. Mặt khác, các xét đoán tổng quát kiểu «Künneth» ở chỗ đã dẫn (xem
thêm [6]) cho thấy tính riêng của $X/S$ và tính phẳng của $\mathcal L$ trên
$S$ kéo theo sự tồn tại của một Môđun coherent $\mathcal Q$ trên $S$, được
xác định sai khác một đẳng cấu duy nhất, và một đẳng cấu các bó
\[
f_*(\mathcal L)\xrightarrow{\sim}
\operatorname{Hom}_{\mathcal O_X}(\mathcal Q,\mathcal O_S),
\]
 hơn nữa, việc tạo thành $\mathcal Q$ tương
thích với các phép đổi cơ sở. Ở đây $f_*(\mathcal L)^*$ chỉ bó con các tập hợp
của $f_*(\mathcal L)$ mà các tiết diện trên $U$ là những tiết diện của
$\mathcal L$ trên $f^{-1}(U)$ xác định các ước Cartier tương đối trên
\src{153}
$f^{-1}(U)/U$, nghĩa là cảm sinh các tiết diện không là ước của không trên các
$X_s$ ($s\in U$). Dùng giả thiết rằng các thớ $X_s$ là nguyên, điều này đơn
giản có nghĩa là các tiết diện cảm sinh trên các thớ $X_s$ không đồng nhất
bằng không, hay nói theo các đồng cấu địa phương
$\mathcal Q\longrightarrow\mathcal O_S$, rằng các đồng cấu ấy là toàn ánh
(Nakayama). Điều này chứng tỏ tập các tiết diện của
$f_*(\mathcal L)^*/\mathcal O_S^*$ tương ứng song ánh với tập các Môđun thương
khả nghịch của $\mathcal Q$, hay cũng vậy, theo định nghĩa của phân thớ xạ ảnh
$\mathbf P(\mathcal Q)$ liên kết với Môđun coherent $\mathcal Q$ (xem [5], V,
2), với tập các tiết diện của $\mathbf P(\mathcal Q)$ trên $S$. Mô tả này
tương thích với việc tạo thành các ảnh ngược, và do đó ta thu được định lý sau.

\result{ĐỊNH LÝ}{4.3} Cho $f:X\longrightarrow S$ là một cấu xạ riêng và phẳng
có các thớ hình học nguyên, với $S$ Noether địa phương, và cho $\mathcal L$ là
một Môđun khả nghịch trên $X$. Với mọi $S'$ trên $S$, gọi $T(S')$ là tập các
ước tương đối $D$ trên $X\times_SS'/S'$ sao cho $\mathcal L(D)$ đẳng cấu với
$\mathcal L\otimes_{\mathcal O_S}\mathcal O_{S'}$ địa phương trên $S'$ (nghĩa
là sao cho $\mathcal L(D)$ và
$\mathcal L\otimes_{\mathcal O_S}\mathcal O_{S'}$ xác định cùng một phần tử
của $\operatorname{Pic}(X\times_SS'/S')$). Khi đó tồn tại một Môđun coherent
$\mathcal Q$ trên $S$, được xác định sai khác một đẳng cấu duy nhất, sao cho
hàm tử $T$ được biểu diễn bởi phân thớ xạ ảnh $\mathbf P(\mathcal Q)$.

\result{HỆ QUẢ}{4.4} Nếu còn giả sử $f$ xạ ảnh, thì đồng cấu có tính hàm tử
$\operatorname{Div}_{X/S}\longrightarrow\operatorname{Pic}_{X/S}$ biểu diễn
được bởi các cấu xạ xạ ảnh.

Thật vậy, khi $X$ có một tiết diện (tương ứng, có một tiết diện địa phương)
trên $S$, thì đồng cấu trên được biểu diễn bởi các phân thớ xạ ảnh (tương ứng,
bởi các phân thớ xạ ảnh địa phương), nhờ 4.3 và 2.1. Trong trường hợp $f$ tựa
xạ ảnh, ta dễ dàng quy về trường hợp trên bằng một phương pháp hạ xuống, dùng
các tựa tiết diện hữu hạn và phẳng của $X$ địa phương trên $S$.

\result{NHẬN XÉT}{4.5} Dưới các điều kiện của 4.3, Môđun $\mathcal Q$ trên
$S$ nói chung không tự do địa phương; điều này được biểu thị bởi sự kiện rằng
chiều của các thớ rút gọn của $\mathcal Q$, nghĩa là chiều của
$H^0(X_s,\mathcal L_s)$ khi $s\in S$ biến thiên, có thể nhảy. Cho một Môđun
coherent $\mathcal Q$ trên tiền lược đồ Noether địa phương $S$, dễ kiểm tra
rằng với một $s\in S$ đã cho, $\mathcal Q$ tự do tại $s$ khi và chỉ khi
$\mathbf P(\mathcal Q)$ phẳng trên $S$ tại các điểm nằm trên $s$ (trong trường
hợp đó, nó thậm chí trơn trên $S$ tại các điểm nằm trên $s$); trong tình huống
đang xét, vì $\mathcal Q$ được định nghĩa theo $\mathcal L$ như trên, điều này
cũng có nghĩa là việc tạo thành ảnh trực tiếp $f_*(\mathcal L)$ «giao hoán với
phép đổi cơ sở trong
\src{154}
một lân cận của $s$», hay cũng vậy,
$f_*(\mathcal L)_s\longrightarrow H^0(X_s,\mathcal L_s)$ là toàn ánh. Chẳng
hạn điều này xảy ra nếu $H^1(X_s,\mathcal L_s)=0$. Với điều kiện các tiền lược
đồ đang xét tồn tại, các tiêu chuẩn này đặc biệt áp dụng cho tình huống phổ
dụng $\operatorname{Div}_{X/S}\longrightarrow\operatorname{Pic}_{X/S}$, và
cho một điều kiện cần và đủ, tương ứng một điều kiện đủ, để cấu xạ này trơn tại
một điểm đã cho của $\operatorname{Div}_{X/S}$.

\fsection{5}{Chứng minh định lý tồn tại chính}

Dưới các điều kiện của 3.1, chọn một Môđun $\mathcal O_X(1)$ rất rộng trên
$X/S$, và gọi $\xi$ là phần tử tương ứng của $\operatorname{Pic}(X/S)$. Để
viết gọn, đặt
\[
P(S')=\operatorname{Pic}(X\times_SS'/S'),
\]
và trước hết, để đơn giản, giả sử $X/S$ có một tiết diện. Gọi $P^+(S')$ là bộ
phận của $P(S')$ gồm các lớp của những $\mathcal L$ khả nghịch trên
$X\times_SS'$ sao cho
\[
R^if'_*(\mathcal L(n))=0\quad\text{với $i>0$ và mọi $n\geq0$},
\qquad
f'_*(\mathcal L(n))\neq0\quad\text{với mọi $n\geq0$}.
\]
Đó là những điều kiện ổn định qua đổi cơ sở, nên xác định một hàm tử con
$P^+$ của $P$, hiển nhiên ổn định dưới phép tịnh tiến bởi $\xi$. Dùng «các
định lý A và B» của Serre ([2], III, 2) và các kết quả tổng quát ([5], IV, 5),
dễ thấy rằng $P$ biểu diễn được khi và chỉ khi $P^+$ biểu diễn được, và khi ấy
$P^+$ sẽ được biểu diễn bởi một tập mở $U$ của tiền lược đồ
$\operatorname{Pic}_{X/S}$ biểu diễn $P$, còn tiền lược đồ này sẽ là hợp tăng
của các tập mở $U-n\xi$.

Để viết gọn, đặt $D=\operatorname{Div}_{X/S}$, và gọi $D^+$ là ảnh ngược của
$P^+$ qua cấu xạ chính tắc $D\longrightarrow P$. Do đó ta có một cấu xạ
\begin{equation}
D^+\longrightarrow P^+,
\tag{+}
\end{equation}
và ta đã biết $D^+$ được biểu diễn bởi một tập mở $D^+$ của tiền lược đồ
$D=\operatorname{Div}_{X/S}$ (sự tồn tại của nó được bảo đảm bởi 4.1). Từ 4.3
dễ suy ra cấu xạ (+) được biểu diễn bởi các cấu xạ trơn, xạ ảnh và toàn ánh
(chính xác hơn, bởi các phân thớ xạ ảnh liên kết với các Môđun tự do địa
phương không triệt tiêu ở mọi nơi): điều này là do khi $\mathcal L$ trên
$X\times_SS'$ như ở đầu mục này, thì $f'_*(\mathcal L)$ là một Môđun tự do
địa phương $\neq0$, mà việc tạo thành nó giao hoán với đổi cơ sở; với các ký
hiệu của 4.3, khi đó $\mathcal Q$ đẳng cấu với đối ngẫu của
$f'_*(\mathcal L)$. Dùng tiêu chuẩn tổng quát ([5], IV, 4.7), ta sẽ suy ra
tính biểu diễn được của $P^+$. Ở chỗ đã dẫn, ta sẽ lấy $\mathfrak C$ là lớp
các cấu xạ tiền lược đồ phẳng trung thành
\src{155}
và tựa compact (quả thật là các toàn cấu hiệu dụng theo [4], I B). Điều kiện
(a) ở chỗ đã dẫn, tức là (+) biểu diễn được bởi các cấu xạ là phần tử của
$\mathfrak C$, đã được kiểm chứng như vừa thấy; điều kiện (b) cũng vậy, nó có
nghĩa là hàm tử $P^+$ tương thích với hạ phẳng trung thành tựa compact, điều
này là hiển nhiên. Còn phải chứng minh điều kiện (c) ở chỗ đã dẫn, tức là quan
hệ tương đương $R$ trong tiền lược đồ $D^+$ suy ra từ cấu xạ (+) biểu diễn được
bởi $\mathfrak C$ là $\mathfrak C$-hiệu dụng, nghĩa là hiệu dụng và sao cho
$D^+\longrightarrow D^+/R$ thuộc $\mathfrak C$. Để làm điều này, trước hết
lưu ý rằng các tập mở $D^{+Q}$ của $D^+$ tương ứng với các đa thức Hilbert ảo
khác nhau $Q\in\mathbf Q[t]$ ổn định dưới $R$ (vì các thớ của $R$ liên thông);
điều này quy ta về việc chứng minh rằng với mọi $Q$, quan hệ tương đương $R^Q$
cảm sinh trong $D^{+Q}$ là $\mathfrak C$-hiệu dụng. Nhưng giờ đây $D^{+Q}$ tựa
xạ ảnh, còn quan hệ tương đương $R^Q$ xạ ảnh và phẳng. Do đó ta ở trong các
điều kiện áp dụng của [4], III, 6.1, và kết quả cần chứng minh suy ra.

Trong trường hợp tổng quát, khi $X/S$ không nhất thiết có một tiết diện, ta dễ
dàng quy về trường hợp trên bằng kỹ thuật hạ xuống, hoặc lặp lại chứng minh
trên với sửa đổi cần thiết trong định nghĩa của $P^+$.

\result{CÁC NHẬN XÉT}{5.1} Phương pháp được dùng về cơ bản là phương pháp của
Matsusaka để xây dựng xạ ảnh các đa tạp Picard. Kết quả được viện dẫn từ [4],
III về khả năng chuyển sang thương hiệu dụng cũng có thể dễ dàng suy ra từ
định lý tồn tại các lược đồ Hilbert (xem chẳng hạn [6]) (theo cách cổ điển, các
thương này được xây dựng bằng cách dùng các tọa độ Chow). Lưu ý rằng việc tạo
thành tập mở $\operatorname{Pic}^{+}_{X/S}$ của
$\operatorname{Pic}_{X/S}$ và sự phân tích nó thành các tập mở
$\operatorname{Pic}^{+Q}_{X/S}$ tựa xạ ảnh trên $S$, theo các đa thức Hilbert
của các ước xác định những Môđun khả nghịch đang xét, tương thích với đổi cơ
sở (điều này cho phép áp dụng kỹ thuật hạ xuống).

\result{CÁC NHẬN XÉT}{5.2} Không loại trừ khả năng
$\operatorname{Pic}_{X/S}$ tồn tại mỗi khi $f:X\longrightarrow S$ riêng và
phẳng, và các đồng cấu
\[
k(s)\longrightarrow H^0(X_s,\mathcal O_{X_s})\qquad(s\in S)
\]
là những đẳng cấu (điều kiện cuối này khi đó cũng có nghĩa là
$\mathcal O_S\xrightarrow{\sim}f_*(\mathcal O_X)$ và quan hệ này vẫn đúng sau
mọi phép đổi cơ sở $S'\longrightarrow S$). Ít nhất đây là điều ta chứng minh
được trong khuôn khổ các không gian giải tích khi $f$ còn xạ ảnh, bằng một
phương pháp khác (phương pháp của Chow, nếu không nhầm) được giải thích trong
[5], IX, 3.1. Trong phương pháp này, việc chuyển sang
\src{156}
thương theo một quan hệ tương đương riêng và phẳng trong một tập mở của lược
đồ các ước được thay bằng việc lấy thương theo nhóm xạ ảnh trong lược đồ các
phép nhúng của $X$ vào $\mathbf P^r_S$. Phương pháp này có lẽ có thể thích nghi
với trường hợp các lược đồ, bằng cách dùng các kết quả của Mumford về việc lấy
thương theo nhóm xạ ảnh [8]; hiện nay chỉ có chứng minh thành văn khi $X$ có
«nhiều tiết diện địa phương» trên $S$, chẳng hạn khi $X$ khả ly trên một vành
địa phương đầy đủ. Về nguyên tắc, phương pháp đang xét
có phạm vi tổng quát hơn, vì nó cũng cho sự tồn tại của các tiền lược đồ Picard
trong những trường hợp chúng không tách, và do đó phương pháp thứ nhất tất yếu
thất bại. (Về mặt kỹ thuật, khó khăn xuất phát từ sự kiện rằng khi các thớ hình
học của $f$ không còn nguyên, thì hàm tử được xét trong 4.3 không còn được biểu
diễn bởi chính phân thớ xạ ảnh $\mathbf P(\mathcal Q)$, mà bởi một tập mở của
nó; điều này dẫn đến bài toán tế nhị về việc lấy thương theo một quan hệ tương
đương phẳng nhưng không riêng.)

\result{NHẬN XÉT}{5.3} Lưu ý rằng chứng minh được đưa ra ở đây không viện đến
phép dựng trước các đa tạp Jacobi của các đường cong hay các họ đường cong,
cũng không viện đến lý thuyết các đa tạp Abel hay các lược đồ Abel, và do đó
khác về bản chất với những trình bày truyền thống như trong sách của Lang [7]
hay thuyết trình của Chevalley [1], vốn đi theo con đường do A. Weil phác thảo.
Ngay cả trong trường hợp các đa tạp Jacobi của những đường cong không kỳ dị
trên một trường đóng đại số (chẳng hạn trường các số phức), phép dựng đa tạp
Jacobi được đưa ra ở đây là phép dựng duy nhất đã biết cho những tính chất rất
mạnh mà ta đã lấy làm định nghĩa ở mục 1 (về cơ bản là các tính chất của
Chevalley, nhưng có tính đến các «đa tạp tham số» với phần tử lũy linh). Việc
phép dựng các lược đồ Picard phải đi trước chứ không đi sau lý thuyết các đa
tạp Abel là rõ ràng a priori, vì nói chung các lược đồ Picard không phải là các
lược đồ Abel cũng không quy được về chúng; điều này đã thấy trong trường hợp
các đường cong kỳ dị trên một trường đóng đại số, nơi xuất hiện các «đa tạp
Jacobi suy rộng» của Rosenlicht, vốn không phải là các đa tạp Abel. Hơn nữa, lý
thuyết các đa tạp Abel, và tổng quát hơn lý thuyết các lược đồ Abel, trở nên
đơn giản hơn nhiều khi đã có một lý thuyết tổng quát về các lược đồ Picard. Đặc
biệt, lý thuyết đối ngẫu cho các lược đồ Abel, nhất là các kết quả kiểu Cartier,
khi đó trở nên gần như hình thức (xem chẳng hạn [10]).
\src{157}
\result{CÁC NHẬN XÉT}{5.4} «Nguyên lý tương thích» của Igusa đối với đa tạp
Jacobi của một đường cong suy biến thành một đường cong kỳ dị chỉ có thể được
hiểu đúng như một định lý tồn tại của lược đồ Picard của một lược đồ đường cong
tương đối $X/S$, không nhất thiết trơn trên $S$. Vì vậy, đó là một trường hợp
riêng của định lý tồn tại chính 3.1 khi đường cong chuyên biệt là nguyên (nghĩa
là, theo ngôn ngữ cổ điển, bất khả quy với bội 1). Lưu ý rằng hiện nay trường
hợp một đường cong đặc biệt khả quy (ngay cả khi các thành phần có bội 1, nghĩa
là khi đường cong đặc biệt khả ly trên trường thặng dư) vẫn nằm ngoài các định
lý tồn tại đã biết, trừ trường hợp ta ở trên một vành định giá rời rạc đầy đủ
có trường thặng dư đóng đại số; so sánh 5.2. Bài toán tồn tại này chắc chắn sẽ
đặt ra khi xây dựng bằng hình học, trong lý thuyết lược đồ, các «phép compact
hóa» Baily--Satake của các lược đồ môđuli đường cong giống $g$. (Hiện nay phép
compact hóa này chỉ được biết với $g=1$, nhờ công trình của Igusa.)

\fsection{6}{Các định lý tồn tại tương đối}

Ở đây ta sẽ phác thảo một số trường hợp hữu ích trong đó sự tồn tại của một số
lược đồ Picard kéo theo sự tồn tại của một số lược đồ khác, nhờ đó có thể suy
ra từ định lý chính 3.1 nhiều định lý tồn tại khác.

\result{MỆNH ĐỀ}{6.1} Cho $f:X\longrightarrow S$ là một cấu xạ phẳng và xạ ảnh
sao cho trong phân tích Stein $f=f''f'$, cấu xạ
$f':X\longrightarrow S'$ phẳng và có các thớ hình học nguyên (do đó thỏa các
giả thiết của 3.1), còn cấu xạ hữu hạn $f'':S'\longrightarrow S$ phẳng. Khi đó
$\operatorname{Pic}_{X/S}$ tồn tại và (với các ký hiệu được đưa vào trong [4],
II, C 2) ta có một đẳng cấu chính tắc
\[
\operatorname{Pic}_{X/S}\xrightarrow{\sim}
\Pi_{S'/S}\operatorname{Pic}_{X/S'}.
\]

Để chứng minh, trước hết ta thiết lập một đẳng cấu các hàm tử
\[
\operatorname{Pic}_{X/S}\xrightarrow{\sim}
\Pi_{S'/S}\operatorname{Pic}_{X/S'},
\]
rồi dùng 3.1, vốn kéo theo rằng $\operatorname{Pic}_{X/S'}$ biểu diễn được, và
dùng cấu trúc đã nêu ở mục số\textsuperscript{} 3 của
$\operatorname{Pic}_{X/S'}$ (kéo theo rằng mọi bộ phận hữu hạn của một thớ của
$\operatorname{Pic}_{X/S'}$ trên $S$ đều nằm trong một tập mở afin) để có sự
tồn tại của $\Pi_{S'/S}\operatorname{Pic}_{X/S'}$.
\src{158}
Chẳng hạn, khi $X$ là một lược đồ tổng của các lược đồ $X_i$ trên $S$ thỏa các
điều kiện của 3.1, phát biểu 6.1 quy về phát biểu hiển nhiên
\[
\operatorname{Pic}_{X/S}\xrightarrow{\sim}
\prod_i\operatorname{Pic}_{X_i/S}.
\]

\result{HỆ QUẢ}{6.2} Cho $f:X\longrightarrow S$ là một cấu xạ xạ ảnh và phẳng
có các thớ hình học nguyên địa phương (chẳng hạn một cấu xạ xạ ảnh và chuẩn),
thì $\operatorname{Pic}_{X/S}$ tồn tại.

Hơn nữa, trong trường hợp này $S'$ là một phủ étale của $S$ (đúng mỗi khi $f$
khả ly, nghĩa là phẳng với các thớ hình học rút gọn), và ta kiểm tra được rằng
định lý cấu trúc phát biểu ở mục số\textsuperscript{} 3 cho
$\operatorname{Pic}_{X/S}$ vẫn đúng, nhờ cấu trúc tương tự của
$\operatorname{Pic}_{X/S'}$.

Mặt khác, một thủ tục hạ xuống cho một định lý tồn tại tương đối, mà phạm vi
của nó phụ thuộc vào lời giải cho các bài toán hạ không phẳng được nêu trong
[4], I, A (c), và ở đây ta chỉ trình bày tường minh một trường hợp riêng:

\result{MỆNH ĐỀ}{6.3} Cho $f:X\longrightarrow S$ là một cấu xạ riêng, và cho
$X_1$, $X_2$ là hai tiền lược đồ con của $X$ phẳng trên $S$, được xác định bởi
hai Iđêan coherent $\mathcal I_1$ và $\mathcal I_2$ sao cho
$\mathcal I_1\cap\mathcal I_2=(0)$ và
$\mathcal O_X/(\mathcal I_1+\mathcal I_2)$ phẳng trên $S$ (nghĩa là tiền lược
đồ con của $X$ vốn là hợp của $X_1$ và $X_2$ thì bằng $X$, còn giao $X$ của chúng
phẳng trên $S$). Giả sử thêm rằng với mọi $s\in S$,
các đồng cấu
\[
k(s)\longrightarrow H^0(X_{i,s},\mathcal O_{X_{i,s}}),\qquad i=1,2,
\]
là song ánh. Khi đó đồng cấu tự nhiên của các hàm tử
\[
\operatorname{Pic}_{X/S}\longrightarrow
\operatorname{Pic}_{X_1/S}\times\operatorname{Pic}_{X_2/S}
\]
biểu diễn được bởi các cấu xạ afin; do đó, nếu
$\operatorname{Pic}_{X_1/S}$ và $\operatorname{Pic}_{X_2/S}$ tồn tại, thì
$\operatorname{Pic}_{X/S}$ cũng tồn tại, và cấu xạ chính tắc
\[
\operatorname{Pic}_{X/S}\longrightarrow
\operatorname{Pic}_{X_1/S}\times\operatorname{Pic}_{X_2/S}
\]
là afin.

Bằng hạ phẳng trung thành, ta quy về trường hợp $Z$ có một tiết diện trên $S$,
do đó xác định các tiết diện của $X$, $X_1$, $X_2$ trên $S$, cho phép loại bỏ
các tự đẳng cấu trong những cấu trúc đang xét như đã giải thích trong 2.5.
Chứng minh khi đó chỉ cần lưu ý rằng dữ liệu của một Môđun khả nghịch
\src{159}
«được rigid hóa» $\mathcal L$ trên $X$ tương đương với dữ liệu của một bộ ba
$(\mathcal L_1,\mathcal L_2,u)$, trong đó $\mathcal L_i$ là một Môđun khả
nghịch «được rigid hóa» trên $X_i$, còn $u$ là một đẳng cấu từ
$\mathcal L_1|Z$ đến $\mathcal L_2|Z$, tương thích với các phép rigid hóa. Chỉ
cần kiểm tra rằng với $\mathcal L_1$, $\mathcal L_2$ cố định, dữ liệu của $u$
được biểu thị bởi một tiết diện của một lược đồ thích hợp trên $S$, afin trên
$S$, điều này là dễ. Từ 6.3, ta dễ dàng suy ra:

\result{HỆ QUẢ}{6.4} Cho $X$ là một lược đồ riêng và khả ly trên một trường
$k$, và cho $X_i$ là các thành phần bất khả quy của $X$. Nếu các
$\operatorname{Pic}_{X_i/k}$ tồn tại, thì $\operatorname{Pic}_{X/k}$ cũng tồn
tại, và cấu xạ chính tắc
\[
\operatorname{Pic}_{X/k}\longrightarrow\prod_i\operatorname{Pic}_{X_i/k}
\]
là afin.

Kết hợp với 6.2, điều này chẳng hạn chứng tỏ sự tồn tại của
$\operatorname{Pic}_{X/k}$ mỗi khi $X$ là một lược đồ xạ ảnh và khả ly trên
một trường $k$. Khi $X$ không còn khả ly trên $k$, ta cũng có một kết quả quy
giản, dùng một lập luận của Oort [11]. Phương pháp cũng áp dụng trên một lược đồ
cơ sở tùy ý (trường hợp hữu ích chẳng hạn để chứng minh trong thuyết trình tiếp
theo kết quả hữu hạn đã thông báo ở 3.3). Để tránh một phát biểu quá kỹ thuật,
ta chỉ xét trường hợp cơ sở là một trường:

\result{MỆNH ĐỀ}{6.5} Cho $X$ là một lược đồ riêng trên một trường $k$, và cho
$X_0$ là một lược đồ con có cùng tập hợp nền (do đó được xác định bởi một Iđêan
lũy linh trên $X$). Khi đó cấu xạ có tính hàm tử
\[
\operatorname{Pic}_{X/k}\longrightarrow\operatorname{Pic}_{X_0/k}
\]
biểu diễn được bởi các cấu xạ afin. Đặc biệt, nếu
$\operatorname{Pic}_{X_0/k}$ tồn tại, thì $\operatorname{Pic}_{X/k}$ cũng tồn
tại, và cấu xạ
$\operatorname{Pic}_{X/k}\longrightarrow\operatorname{Pic}_{X_0/k}$ là afin.

Kết hợp với 6.4, ta dễ dàng kết luận:

\result{HỆ QUẢ}{6.6} Cho $X$ là một tiền lược đồ xạ ảnh trên một trường $k$,
thì $\operatorname{Pic}_{X/k}$ tồn tại.

\result{NHẬN XÉT}{6.6} Rất có cơ sở để tin rằng với mọi lược đồ $X$ riêng trên
một trường $k$, $\operatorname{Pic}_{X/k}$ tồn tại. Các kết quả trên
quy bài toán này về trường hợp $k$ đóng đại số và $X$ nguyên. Khi đó ta biết
tồn tại một lược đồ nguyên $X'$ xạ ảnh trên $k$ và một cấu xạ trội
$g:X'\longrightarrow X$ (bổ đề Chow). Do đó chỉ cần chứng minh rằng
\src{160}
cấu xạ có tính hàm tử tương ứng
\[
\operatorname{Pic}_{X/k}\longrightarrow\operatorname{Pic}_{X'/k}
\]
biểu diễn được (và hẳn là biểu diễn được bởi các cấu xạ afin), vì mặt khác ta
biết $\operatorname{Pic}_{X'/k}$ biểu diễn được. Điều này đặt ra những bài toán
hạ không phẳng hiện còn chưa giải quyết. Lưu ý rằng khi ta chỉ xét hạn chế của
hàm tử $\operatorname{Pic}_{X/k}$ lên các tiền lược đồ rút gọn (với $X$ riêng
và nguyên trên $k$ đóng đại số), ta quả thực thu được một hàm tử biểu diễn
được, như Chevalley [1] đã chứng minh khi $X$ chuẩn, và Seshadri [12] chứng minh
bằng một phương pháp hạ xuống trong trường hợp tổng quát. Nhưng với các ký hiệu
của ta, lược đồ do các tác giả ấy dựng không phải là
$\operatorname{Pic}_{X/k}$, mà là
$(\operatorname{Pic}_{X/k})_{\mathrm{réd}}$, lược đồ rút gọn tương ứng với
$\operatorname{Pic}_{X/k}$.

\result{NHẬN XÉT}{6.7} Tổng quát hơn, cho $f:X'\longrightarrow X$ là một cấu
xạ toàn ánh của các tiền lược đồ riêng trên $k$. Khi đó, các xét đoán về hạ
không phẳng dẫn đến phỏng đoán rằng
\[
\operatorname{Pic}_{X/k}\longrightarrow\operatorname{Pic}_{X'/k}
\]
là một cấu xạ afin, điều này đặc biệt kéo theo (sau khi chia cho các thành phần liên
thông của các phần tử đơn vị) rằng đồng cấu tương ứng trên các nhóm
Néron--Severi là đơn ánh modulo xoắn. Đây là điều có thể kiểm tra bằng lý thuyết
giao khi $X$ và $X'$ không kỳ dị. Câu trả lời dường như chưa được biết trong
bất kỳ trường hợp nào khác.

\result{NHẬN XÉT}{6.8} Trái với điều người ta có thể nghĩ, việc xét các lược đồ
Picard của những lược đồ đại số có phần tử lũy linh là hữu ích, thậm chí không
thể thiếu, trong nhiều bài toán. Chẳng hạn, khi $X$ là một lược đồ xạ ảnh, giả
sử chính quy, và $Y$ là một tiết diện siêu phẳng, cần xét các «lân cận vô cùng
bé» mọi cấp $X_n$ của $Y$, và các lược đồ Picard
$\operatorname{Pic}_{X_n/k}$; khi $X$ bất khả quy có chiều $\geq4$ (tương ứng
$\geq3$), cấu xạ chính tắc
\[
\operatorname{Pic}_{X/k}\longrightarrow\operatorname{Pic}_{X_n/k}
\qquad(n\text{ lớn})
\]
là một đẳng cấu (tương ứng, cảm sinh một đẳng cấu trên các ảnh ngược của những
nhóm con xoắn của Néron--Severi), một kết quả sẽ hữu ích cho ta trong nghiên
cứu định tính các lược đồ Picard ở thuyết trình tiếp theo. Tương tự, việc xét
các lược đồ Picard của một số đường cong có phần tử lũy linh và các định lý nền
tảng của hình học hình thức [3] cho phép kiểm chứng, trong trường hợp đồng đặc
trưng, một phỏng đoán của Mumford, cụ thể là với mọi vành địa phương Noether
chuẩn đầy đủ $A$ có chiều 2, nhóm các lớp ước của $A$ có thể được xem là tập
các điểm hữu tỉ trên $k$ của một
\src{161}
nhóm đại số $G$ trên trường thặng dư $k$ (hơn nữa, $G$ được xác định chính tắc
khi ta chọn một trường hệ số trong $A$). Khi $A$ có chiều tùy ý, rất có cơ sở
để tin rằng tồn tại một nhóm pro-đại số trên $k$ đóng vai trò như $G$ ở trên;
trong trường hợp có thể «giải kỳ dị» $\operatorname{Spec}(A)$, nhóm này được
dựng như một giới hạn xạ ảnh của các lược đồ Picard của những lược đồ xạ ảnh
thích hợp (có phần tử lũy linh) trên $k$.

\section*{TÀI LIỆU THAM KHẢO}
\addcontentsline{toc}{section}{Tài liệu tham khảo}

\begin{description}
\item[{[1]}] CHEVALLEY (Claude). -- Sur la théorie de Picard,
\emph{Amer. J. of Math.}, tập 82, 1960, tr. 435--490.

\item[{[2]}] GROTHENDIECK (A.) và DIEUDONNÉ (J.). -- Éléments de géométrie
algébrique, Chapitre I et suivants. -- Paris, Presses universitaires de France,
1960--1961 (Institut des Hautes Études Scientifiques, Publications
mathématiques 4, 8, \ldots).

\item[{[3]}] GROTHENDIECK (Alexander). -- Géométrie formelle et géométrie
algébrique, \emph{Séminaire Bourbaki}, tập 11, 1958/59,
số\textsuperscript{} 182, 28 tr.

\item[{[4]}] GROTHENDIECK (Alexander). -- Techniques de descente et théorèmes
d'existence en géométrie algébrique, I--IV, \emph{Séminaire Bourbaki}, tập 12,
1959/60, số\textsuperscript{} 190, 29 tr. và số\textsuperscript{} 195, 22 tr.;
tập 13, 1960/61, số\textsuperscript{} 212, 20 tr. và
số\textsuperscript{} 221, 28 tr.

\item[{[5]}] GROTHENDIECK (Alexander). -- Techniques de construction en
géométrie analytique, I--X, \emph{Séminaire Cartan}, tập 13, 1960/61:
Déformation des structures analytiques complexes, số\textsuperscript{} 7--17.

\item[{[6]}] GROTHENDIECK (Alexander). -- Séminaire de géométrie algébrique,
do Lichtenbaum biên soạn, Harvard University, 1961 (sắp xuất bản).

\item[{[7]}] LANG (Serge). -- \emph{Abelian varieties}. -- New York,
Interscience Publishers, 1959 (Interscience Tracts in pure and applied
Mathematics, 7).

\item[{[8]}] MUMFORD (D.). -- An elementary theorem in geometric invariant
theory, \emph{Bull. Amer. math. Soc.}, tập 67, 1961, tr. 483--486.

\item[{[9]}] MUMFORD (D.). -- The topology of normal singularities of an
algebraic surface and a criterion for simplicity. -- Paris, Presses
universitaires de France, 1961 (Institut des Hautes Études Scientifiques,
Publications mathématiques, 9).

\item[{[10]}] MUMFORD (D.) và TATE (J.). -- Séminaire de géométrie algébrique,
Harvard University, học kỳ Xuân 1962 (sắp xuất bản).

\item[{[11]}] OORT (F.). -- Sur le schéma de Picard,
\emph{Bull. Soc. math. France}, tập 90, 1962 (sắp xuất bản).

\item[{[12]}] SESHADRI (C. S.). -- Luận án sắp xuất bản trong
\emph{Annali di Matematica}.
\end{description}

\par

\clearpage
\chapter*{Đính chính cho Báo cáo số~232}
\addcontentsline{toc}{chapter}{Đính chính cho Báo cáo số~232}
\sid{e232}
\markboth{Đính chính 232}{Đính chính 232}
\src{303}

\noindent\textbf{7. TDTE V: Các lược đồ Picard. Các định lý tồn tại.}

\begin{center}
[Séminaire Bourbaki, tập 14, 1961/62, số 232, 19 trang.]
\end{center}

\noindent\textbf{Trang 7, Nhận xét 3.3.} --- Vấn đề nêu ra ở đây đã được
MUMFORD giải đáp khẳng định (xem phần Bình luận cho báo cáo kế tiếp).

\noindent\textbf{Trang 13, Nhận xét 5.2.} --- Như tôi chỉ ra ở đầu báo cáo kế
tiếp, phỏng đoán tồn tại được nêu ở đây là sai; tuy nhiên, bằng các phương pháp
của mình, MUMFORD chứng minh được một định lý yếu hơn đôi chút.

\noindent\textbf{Trang 14, dòng 7.} --- Sau ``đầy đủ'', thêm ``có trường thặng
dư đóng đại số''. Hơn nữa, phản ví dụ của MUMFORD vừa nêu cho thấy hạn chế này
là thiết yếu.

\noindent\textbf{Trang 17, Nhận xét 6.6.} --- Như chúng tôi chỉ ra ở đầu báo
cáo kế tiếp, vấn đề nêu ra ở đây vừa được MURRE giải đáp khẳng định.

\noindent\textbf{Trang 18, Nhận xét 6.7.} --- Vấn đề nêu ra ở đây có lời giải
khẳng định (xem đoạn cuối của phần Bình luận cho báo cáo kế tiếp).

\noindent\textbf{Trang 18, Nhận xét 6.8, dòng thứ 4.} --- Thay ``chính quy''
bằng ``đơn trên $k$''.

\par

\unit{Thuyết trình 236 -- Các lược đồ Picard: các tính chất tổng quát}
\sid{236}
\markboth{Thuyết trình 236}{Thuyết trình 236}
\src{221}
\seminar{Năm thứ 14\textsuperscript{}, 1961/62, số 236\textsuperscript{}\\Tháng Năm 1962}
\paperhead{KỸ THUẬT HẠ XUỐNG VÀ CÁC ĐỊNH LÝ TỒN TẠI TRONG HÌNH HỌC ĐẠI SỐ\\
VI. CÁC LƯỢC ĐỒ PICARD: CÁC TÍNH CHẤT TỔNG QUÁT}%
  {do Alexander GROTHENDIECK}

\fsection{0}{Các bổ sung cho thuyết trình trước ([3], V)\texorpdfstring{\textsuperscript{(1)}}{(1)}}

Đã có một số tiến bộ liên quan đến các vấn đề tồn tại của những tiền lược đồ
Picard được nêu ra trong [3], V:

\medskip
\noindent a. (Mumford). Nói chung, không đúng rằng khi
$f:X\longrightarrow S$ là một cấu xạ xạ ảnh và khả ly (= phẳng với các thớ
khả ly), tiền lược đồ Picard $\operatorname{Pic}_{X/S}$ tồn tại, ngay cả khi
các thớ của $f$ có chiều 1 và $S$ là phổ của một vành định giá rời rạc đầy đủ.
Một ví dụ thu được bằng cách lấy $S=\operatorname{Spec}\mathbf R[[t]]$ và lấy
$X$ là lược đồ con của $\mathbf P^2_S$ được xác định bởi phương trình (trong đó
$x,y,z$ là các biến thuần nhất)
\[
x^2+y^2=tz^2.
\]
Như vậy, đây là một conic suy biến về mặt hình học thành hai đường thẳng giao
nhau, nhưng thớ đặc biệt trên trường $\mathbf R$ vẫn bất khả quy (nó được cho
bởi phương trình $x^2+y^2=0$ trên $\mathbf R$). Dễ thấy rằng, sau mở rộng étale
$S'\longrightarrow S$, với $S'=\operatorname{Spec}\mathbf C[[t]]$, tiền lược
đồ Picard của $X'/S'$ tồn tại và có một mô tả tường minh như tổng rời nhau của
các bản sao $\widetilde S$, trong đó $\widetilde S$ thu được từ $S$ bằng
cách tách đôi gốc vô hạn lần. Ta dễ dàng nhận thấy dữ liệu hạ xuống trên
$\operatorname{Pic}_{X'/S'}$ đối với $S'\longrightarrow S$ (ở đây được cho
bởi tác động của nhóm Galois $\mathbf Z/2\mathbf Z$ của $S'$ trên $S$) không
hữu hiệu: nhóm này hoán đổi một số điểm đã tách đôi, nên có những quỹ đạo không
nằm trong một tập mở afin. Tuy nhiên, Mumford có thể chứng minh rằng nếu
$f:X\longrightarrow S$ là một cấu xạ xạ ảnh và khả ly sao cho, với mọi
$s\in S$, các thành phần bất khả quy của $X_s$ là bất khả quy hình học trên
$k(s)$, thì $\operatorname{Pic}_{X/S}$ tồn tại; chứng minh của ông dựa trên
một sự tinh chỉnh định lý thương của mình, xem [7]. Mặt khác, cần lưu ý rằng,
ngay cả khi không có giả thiết nào về các thành phần bất khả quy của các thớ
$X_s$, lược đồ $\operatorname{Pic}^{\tau}_{X/S}$ (sẽ được giới thiệu dưới
đây) vẫn có thể tồn tại.
\src{222}
\noindent b. (Murre). Cho $X$ là một lược đồ riêng trên một trường; khi đó
$\operatorname{Pic}_{X/k}$ tồn tại. Chứng minh dựa một phần trên chứng minh
của Chevalley [1] và sử dụng cốt yếu cấu trúc nhóm của hàm tử Picard.

Đối với một số bổ sung về lý thuyết các lược đồ Picard, đặc biệt trong mối
liên hệ với các lược đồ abel, nên tham khảo [7]. Cuối cùng, một thiếu sót đáng
kể của thuyết trình này là chưa có các «tiêu chuẩn tương đương» cho phép so
sánh lược đồ Picard của một lược đồ xạ ảnh với các lược đồ Picard của những
tiết diện siêu phẳng của nó; các định lý then chốt để xây dựng những tiêu
chuẩn như vậy nằm trong [5], và cần ghép thêm các định lý tồn tại của các
lược đồ Picard.

\fsection{1}{Các tính chất tôpô của các tiền lược đồ nhóm giao hoán}

Trước hết, cho $k$ là một trường và $G$ là một tiền lược đồ nhóm trên $k$.
Vì phần tử đơn vị $e$, do hữu tỷ trên $k$, nhất thiết là đóng, nên đường chéo
của $G\times_kG$ lập tức là đóng, và do đó $G$ là tách: mọi tiền lược đồ nhóm
trên một trường đều tách. Ta ký hiệu bởi $G^\circ$ thành phần liên thông của
phần tử đơn vị $e$. Vì $e$ hữu tỷ trên $k$, $G^\circ$ thực ra liên thông hình
học, nghĩa là phép tạo $G^\circ$ tương thích với mọi đổi trường cơ sở. Cũng
suy ra rằng $G^\circ$ ổn định dưới phép nhân khi xét như một tập; và khi $G$
Noether địa phương, nên $G^\circ$ mở, ta có thể xem $G^\circ$ như một nhóm
con mở của $G$. Sau đây ta giả sử $G$ thuộc kiểu hữu hạn địa phương trên $k$;
khi đó $G^\circ$ bất khả quy hình học và thuộc kiểu hữu hạn trên $k$. Thật
vậy, có thể giả sử $k$ đóng đại số, rồi giả sử thêm $G$ rút gọn (vì khi ấy
$G_{\mathrm{réd}}$ là một nhóm con của $G$, do
$G_{\mathrm{réd}}\times_kG_{\mathrm{réd}}$ vẫn rút gọn); do đó trơn trên
$k$ tại một tập mở không rỗng, và vì tịnh tiến tập mở ấy nên trơn ở mọi nơi.
Nhưng khi đó $G$ bất khả quy địa phương, nên các thành phần bất khả quy của nó
trùng với các thành phần liên thông; do đó $G^\circ$ bất khả quy. Cho $U$ là
một lân cận afin của $e$ trong $G^\circ$; dùng tính bất khả quy của $G^\circ$,
ta thấy ngay
$U\cdot U=G^\circ$, điều này chứng minh $G^\circ$ tựa compact, nên thuộc kiểu
hữu hạn trên $k$.

Để đơn giản, giả sử $G$ giao hoán. Với mỗi số nguyên $n>0$, cho $G^{(n)}$ là
ảnh ngược của $G^\circ$ qua đồng cấu $\varphi_n$ lũy thừa bậc $n$ trong
$G$; do đó $G^{(n)}$ là một nhóm con mở của $G$. Ta đặt:
\[
G^\tau=\bigcup_nG^{(n)}.
\]
\src{223}
\[
G^\sigma=\bigcup_{(n,p)=1}G^{(n)},
\qquad
G^\rho=\bigcup_hG^{(p^h)},
\]
trong đó $p$ là số mũ đặc trưng của trường $k$. Như vậy ta thu được các nhóm
con mở của $G$, thỏa mãn
\[
G^\sigma\cap G^\rho=G^\circ,
\qquad
G^\sigma G^\rho=G^\tau.
\]

\result{NHẬN XÉT}{} Ta có thể dựng lược đồ nhóm thương
$G/G^\circ=N$ (xem [3], IV), rồi định nghĩa $G^\tau$, $G^\sigma$, $G^\rho$
là các ảnh ngược trong $G$, tương ứng, của nhóm con xoắn của $N$, của thành
phần $p$-nguyên sơ, và của tổng các thành phần $q$-nguyên sơ với $q$ là số
nguyên tố $\neq p$. Cần lưu ý rằng $N$ là một lược đồ nhóm rời rạc khả ly trên
$k$; do đó, sau khi chọn một bao đóng đại số $\overline k$ của $k$, sinh ra
một nhóm Galois $\pi$, nó được đồng nhất với một nhóm rời rạc thông thường mà
$\pi$ tác động bằng các tự đẳng cấu. Chính bằng cách này ta có thể diễn giải
hiển nhiên việc dựng nhóm con xoắn và sự phân tích nhóm ấy thành các thành
phần $q$-nguyên sơ. Khi $G$ là lược đồ Picard của một lược đồ $X$ riêng trên
$k$, có thể gọi $N$ là \emph{lược đồ Néron--Severi} (rút gọn) của $X$ trên
$k$. Khi $G^\circ_{\mathrm{réd}}$ là một lược đồ nhóm con của $G^\circ$, điều
này đặc biệt xảy ra mỗi khi $k$ hoàn hảo hoặc $G^\circ$ riêng trên $k$ (chẳng
hạn khi $X$ chuẩn hình học), cũng nên đưa vào thương
$N'=G/G^\circ_{\mathrm{réd}}$; thương này có khuynh hướng ứng xử tốt hơn $N$
về phương diện chuyên biệt hóa, nghĩa là khi $X$ biến thiên trong một họ các
lược đồ đại số.

Bây giờ cho $S$ là một tiền lược đồ Noether địa phương và $G$ là một tiền
lược đồ nhóm trên $S$, thuộc kiểu hữu hạn địa phương trên $S$. Ta không giả
sử $G$ thuộc kiểu hữu hạn trên $S$, cũng không giả sử là tách trên $S$. Ta
đặt
\[
G^\circ=\bigcup_{s\in S}(G_s)^\circ,
\]
và, khi $G$ giao hoán,
\[
G^\tau=\bigcup_{s\in S}(G_s)^\tau,
\qquad
G^\sigma=\bigcup_{s\in S}(G_s)^\sigma,
\qquad
G^\rho=\bigcup_{s\in S}(G_s)^\rho.
\]
\src{224}
Đó là các tập con của $G$, ổn định dưới phép nhân của $G$; điều này hiển nhiên
không kéo theo rằng chúng có thể được định nghĩa bằng các tiền lược đồ nhóm
con của $G$. Đặc biệt, dường như nói chung không tồn tại tiền lược đồ nhóm
con nào của $G$ có tập nền là $G^\circ$. Dĩ nhiên, nếu một trong các tập ấy
mở thì, với cấu trúc cảm sinh, nó là một tiền lược đồ nhóm con mở của $G$. Ta
sẽ thấy điều này luôn đúng với $G^\tau$; vì thế, từ quan điểm các hàm tử biểu
diễn được, đặc biệt từ quan điểm «chuyên biệt hóa», tương đương số ứng xử thỏa
đáng hơn tương đương đại số. Sau đây là các tính chất tổng quát chính của
những tập vừa được định nghĩa:

\result{ĐỊNH LÝ}{1.1} $G^\circ$, $G^\tau$, $G^\sigma$, $G^\rho$ đều kiến
thiết được địa phương. Hơn nữa:

\noindent (i) $G^\circ$ tựa compact trên $S$. Khi các $G_s^\circ$ riêng và
$G$ tách trên $S$, thì $G^\circ$ riêng trên $S$, do đó đóng trong $G$.

\noindent (ii) $G^\tau$ mở. Nếu $G^\circ$ đóng thì $G^\tau$ cũng đóng.

\noindent (iii) Nếu $G^\circ$ đóng thì $G^\sigma$ cũng đóng, miễn là ta ở
trong đồng đặc trưng, nghĩa là mọi trường thặng dư của $S$ có cùng đặc trưng.
Nếu $G^\circ$ đóng và $G\longrightarrow S$ mở phổ dụng tại các điểm của
$G^\sigma$ (xem Hệ quả 1.5 dưới đây), thì
$G^\sigma\longrightarrow S$ mở phổ dụng.

\noindent (iv) Nếu $G^\circ$ đóng thì $G^\rho$ cũng đóng. Giả sử ta ở trong
đồng đặc trưng và, với mọi số nguyên $n>0$ sao cho $(n,p)=1$, đồng cấu lũy
thừa bậc $n$ trong $G$ là mở; khi đó $G^\rho$ mở.

Ta nêu vài chỉ dẫn cho chứng minh. Việc $G^\circ$ kiến thiết được địa phương
và tựa compact trên $S$ nằm trong bổ đề sau:

\result{BỔ ĐỀ}{1.2} Giả sử $S\neq\varnothing$. Khi đó tồn tại một tập mở
không rỗng $U$ trong $S$, một lược đồ nhóm $H$ thuộc kiểu hữu hạn trên $U$ có
các thớ liên thông, và một đồng cấu lược đồ nhóm
$H\longrightarrow G|U$ là một phép nhúng mở có ảnh $G^\circ|U$.

Để chứng minh bổ đề, có thể giả sử $S$ bất khả quy; gọi $\eta$ là điểm tổng
quát của nó và thực hiện đổi cơ sở
$S'=\operatorname{Spec}(\mathcal O_{S,\eta})\longrightarrow S$. Ta thu được
một lược đồ nhóm $G'$ trên vành Artin địa phương
$\mathcal O_{S,\eta}=A$, trong đó có một lược đồ nhóm con mở
$G'{}^\circ$ thuộc kiểu hữu hạn trên $A$, như đã nói ở trên. Vì vậy nó thu
được từ một lược đồ nhóm $H$ thuộc kiểu hữu hạn trên một lân cận mở $U$ của
$\eta$, và phép nhúng chính tắc $G'{}^\circ\longrightarrow G'$ thu được từ
một phép nhúng
\src{225}
mở $H\longrightarrow G|U$, và phép nhúng này sẽ là một đồng cấu lược đồ nhóm
khi $U$ đủ nhỏ. Các thớ của $H$ liên thông nếu lấy $U$ đủ nhỏ và đều có
cùng chiều, bằng chiều của các thớ của $G$ khi $U$ đủ nhỏ; do đó, với mọi
$s\in U$, ảnh của
$H_s$ trong $G_s$ chính là $G_s^\circ$ (sau khi lại thu nhỏ $U$), chứng minh
1.2.

Khẳng định thứ hai trong (i) nằm trong bổ đề sau (áp dụng cho một lân cận mở
tựa compact của $G^\circ$ trong $G$):

\result{BỔ ĐỀ}{1.3} Cho $X$ là một tiền lược đồ thuộc kiểu hữu hạn và tách
trên $S$ Noether địa phương, $g$ là một tiết diện của $X$ trên $S$, và
$X^\circ$ là hợp của các thành phần liên thông của $g(s)$ trong các $X_s$.
Cho $s\in S$ sao cho $X_s^\circ$ riêng trên $k(s)$. Khi đó tồn tại một lân
cận mở $U$ của $s$ sao cho $X^\circ|U$ riêng trên $U$, và do đó đóng trong
$X|U$.

Bằng hạ phẳng trung thành trên cơ sở, ta quy về trường hợp $S$ là phổ của một
vành địa phương đầy đủ và $s$ là điểm đóng của nó. Áp dụng [2], III, 5.5.1,
ta thấy $X$ phân tích thành tổng của hai tập mở rời nhau $X'$ và $X''$, trong
đó tập thứ nhất riêng trên $S$ và thỏa $X'_s=X_s^\circ$. Điều này quy ta về
trường hợp $X=X'$, nghĩa là $X$ riêng trên $S$. Trong trường hợp ấy, ta dùng
một chứng minh chuẩn dựa trên tiêu chuẩn định giá cho tính riêng của một bộ
phận (bị bỏ sót trong [2], chương II).

Ta chứng minh $G^\tau$ mở; hay tương đương, vì phép tạo $G^\tau$ (cũng như
phép tạo $G^\circ$, $G^\sigma$, $G^\rho$) giao hoán với mở rộng cơ sở: với
mọi tiết diện $g$ của $G$ trên $S$, $g^{-1}(G^\tau)$ mở. Điều này có nghĩa là
hai điều:

\medskip
\noindent a. Cho $y\in S$ sao cho $g(y)\in G_y^\tau$; khi đó, với mọi
$y'\in\overline{\{y\}}$ đủ gần $y$, ta có $g(y')\in G_{y'}^\tau$.

\medskip
\noindent b. Cho $y'\in S$ sao cho $g(y')\in G_{y'}^\tau$; khi đó, với mọi
tổng quát hóa $y$ của $y'$, ta có $g(y)\in G_y^\tau$.

Đối với (a), lưu ý rằng tồn tại một số nguyên $n>0$ sao cho
$g^n(y)\in G_y^\circ$; vì $G^\circ$ kiến thiết được,
$(g^n)^{-1}(G^\circ)$ kiến thiết được, nên
$g^n(y')\in G_{y'}^\circ$ với $y'\in\overline{\{y\}}$ đủ gần $y$. Đối với
(b), lưu ý rằng các $g(y')^n=g^n(y')$ nằm trong một tập mở tựa compact của
$G_{y'}$ (vì chúng thuộc một số hữu hạn lớp modulo $G_{y'}^\circ$); do đó tồn
tại một tập mở tựa compact $U$ trong $G$ chứa các $g^n(y')$, và vì thế cũng
chứa các tổng quát hóa $g^n(y)$ của chúng. Vậy các lũy thừa của $g(y)$ nằm
trong một tập mở tựa compact của $G_y$, điều này dễ dàng kéo theo
$g(y)\in G_y^\tau$.
\src{226}
Giả sử $G^\circ$ đóng; ta chứng minh $G^\tau$ cũng đóng. Vì đã biết
$G^\tau$ mở, nên kiến thiết được địa phương, chỉ còn phải chứng minh nó ổn
định dưới phép chuyên biệt hóa. Điều này suy ra từ việc nó là hợp của các tập
đóng, cụ thể là các ảnh ngược, qua các đồng cấu $\varphi_n$ lũy thừa bậc
$n$, của tập đóng $G^\circ$.

Cũng lập luận ấy sẽ chứng minh $G^\sigma$ và $G^\rho$ đóng nếu $G^\circ$
đóng (trong trường hợp thứ nhất, với điều kiện ta ở trong đồng đặc trưng), một
khi đã chứng minh $G^\sigma$ và $G^\rho$ kiến thiết được địa phương. Thật vậy,
với $x\in G^\tau$, gọi $\nu(x)$ là số nguyên nhỏ nhất $n>0$ sao cho đồng cấu
$\varphi_n$ lũy thừa bậc $n$ đưa $x$ vào $G^\circ$. Khi đó $G^\sigma$
(tương ứng, $G^\rho$) gồm các $x\in X$ sao cho $\nu(x)$ nguyên tố cùng nhau với $p$ (tương ứng,
là một lũy thừa của $p$), và khẳng định của ta về tính kiến thiết được suy ra
từ khẳng định chính xác hơn sau đây:

\result{BỔ ĐỀ}{1.4} Hàm $\nu$ trên $G^\tau$ kiến thiết được địa phương.

Thật vậy, điều này có nghĩa là, với mọi số nguyên $n>0$, tập các
$x\in G^\tau$ sao cho $\nu(x)=n$ kiến thiết được địa phương; nhưng đó là hiệu
của $\varphi_n^{-1}(G^\circ)$ và hợp của các
$\varphi_d^{-1}(G^\circ)$, trong đó $d$ chạy qua các ước thực sự của $n$. Vì
$G^\circ$ kiến thiết được địa phương, các $\varphi_d^{-1}(G^\circ)$ cũng vậy,
và do đó hiệu nói trên cũng vậy.

Giả sử $G^\circ$ đóng và $G\longrightarrow S$ mở phổ dụng tại các điểm của
$G^\sigma$; ta chứng minh $G^\sigma\longrightarrow S$ là mở, nghĩa là biến
một lân cận trong $G^\sigma$ của $x\in G^\sigma$ thành một lân cận của
$y=f(x)$. Vì $G^\sigma$ kiến thiết được địa phương, ta quy về việc chứng minh
rằng, với mọi tổng quát hóa $y'$ của $y$, tồn tại một tổng quát hóa $x'$ của
$x$ trong $G^\sigma$ nằm trên $y'$. Bằng đổi cơ sở, điều này quy ta về trường
hợp $S$ là phổ của một vành định giá rời rạc, còn $y,y'$ lần lượt là điểm đóng
và điểm tổng quát. Dùng việc $G\longrightarrow S$ mở tại $x$ (do đó tồn tại
một tổng quát hóa $x_1$ của $x$ trong $G$ nằm trên $y'$), ta có thể giả sử
thêm rằng tồn tại một tiết diện $g$ của $G$ trên $S$ sao cho $x=g(y)$ (sau
khi lại đổi cơ sở nếu cần). Nếu $k(y')$ có đặc trưng 0, chỉ cần lấy một tổng
quát hóa bất kỳ $x'$ của $x$ trong $G$ nằm trên $y'$; nó thuộc $G^\tau$ vì
$G^\tau$ mở, và do đó thuộc $G^\sigma$ vì
$G_{y'}^\sigma=G_{y'}^\tau$. Nếu $k(y')$ có đặc trưng $p>0$, đặt
\[
\nu(g(y'))=p^hm\qquad\text{với }(m,p)=1 ;
\]
\src{227}
chọn hai số nguyên $a$ và $b$ sao cho $ap^h+bm=1$, và đặt
\[
g_1=g^{ap^h},\qquad g_2=g^{bm},\qquad\text{do đó}\qquad g=g_1g_2 ;
\]
theo cách dựng, ta có $g_1(y')\in G^\sigma$ và
$g_2(y')\in G^\rho$. Vì $G^\circ$ đóng, suy ra
$g_1(y)\in G^\sigma$ và $g_2(y)\in G^\rho$; mặt khác, từ
\[
g(y)=g_1(y)g_2(y)\in G^\sigma,
\]
ta cũng có $g_2(y)\in G^\sigma$, nên
\[
g_2(y)\in G_y^\sigma\cap G_y^\rho=G_y^\circ.
\]
Từ giả thiết và việc $S$ là phổ của một vành định giá rời rạc, suy ra
$G-(G_y-G_y^\circ)$ là một tập mở của $G$ mà trên đó
$G\longrightarrow S$ cảm sinh một cấu xạ mở; do đó, qua mỗi điểm của
$G_y^\circ$ có một «tựa tiết diện». Vì vậy, sau khi lại thực hiện một mở rộng
trên cơ sở $S$ nếu cần, ta có thể giả sử tồn tại một tiết diện $g'_2$ của
$G^\circ$ trên $S$ sao cho $g'_2(y)=g_2(y)$. Đặt $g'=g_1g'_2$; theo cách dựng,
$g'(y)=g(y)=x$, và
\[
g'(y')=g_1(y')g'_2(y')\in G^\sigma,
\]
do đó $g(y')=x'$ là một tổng quát hóa
của $x$ trong $G^\sigma$ nằm trên $y'$, chứng minh (iii).

Cuối cùng, ta chứng minh khẳng định cuối của (iv). Ta quy về việc chứng minh
rằng, nếu $x\in G^\rho$, thì mọi tổng quát hóa $x'$ của $x$ đều thuộc
$G^\rho$. Có thể giả sử (sau khi lấy ảnh qua $\varphi_{n^h}$ với $h$ thích
hợp nếu cần) rằng thậm chí $x\in G^\circ$. Khi đó, với mọi số nguyên $n$
nguyên tố cùng nhau với đặc trưng, $x$ thuộc ảnh của $\varphi_n$ (vì lũy thừa
bậc $n$ trong một nhóm liên thông thuộc kiểu hữu hạn trên một trường có đặc
trưng nguyên tố cùng nhau với $n$ là toàn ánh). Vì $\varphi_n$ mở, suy ra $x'$
cũng thuộc ảnh của $\varphi_n$. Chính xác hơn, cho $U$ là một tập mở tựa
compact của $G^\tau$ chứa $G_y^\circ$; khi đó
$x'\in\varphi_n(U)$ với mọi $n$ nguyên tố cùng nhau với $p$. Lấy $n$ là một
bội chung của các thừa số nguyên tố cùng nhau với $p$ của $\nu(z)$, với
$z\in U$, ta thu được $x'\in G^\rho$.

Khẳng định (1.1), (iii), được bổ sung như sau:

\result{HỆ QUẢ}{1.5} Cho $n>1$ là một số nguyên sao cho đồng cấu lũy thừa bậc
$n$, $\varphi_n:G\longrightarrow G$, mở phổ dụng (chẳng hạn étale), và đặt
\[
G^{(n)}=\varphi_n^{-1}(G^\circ),
\]
Giả sử các thớ liên thông $G_s^\circ$ không «chứa
\src{228}
thành phần cộng tính» (nghĩa là nhóm thu được sau khi mở rộng trường $k(s)$
đến bao đóng đại số không chứa nhóm con đẳng cấu với $\mathbf G_a$). Khi đó
$G\longrightarrow S$ mở phổ dụng tại các điểm của các $G^{(n^h)}$. Đặc biệt,
nếu các $G_s^\circ$ không chứa thành phần cộng tính và, với mọi số nguyên
$n>1$, đồng cấu $\varphi_n$ mở phổ dụng tại các điểm có đặc trưng thặng dư
nguyên tố cùng nhau với $p$, thì $G\longrightarrow S$ mở phổ dụng tại các điểm của $G^\sigma$, và
do đó, theo (1.1), (iii), nếu $G^\circ$ đóng thì
$G^\sigma\longrightarrow S$ mở phổ dụng. Trong các trường hợp này,
$s\longmapsto\dim G_s$ là một hàm hằng địa phương trên $S$.

Thật vậy, từ giả thiết suy ra hạt nhân ${}_nG$ của $\varphi_n$ mở phổ dụng
trên $S$, nên $G\longrightarrow S$ mở phổ dụng tại các điểm của ${}_nG$, và
do đó (thay $n$ bởi $n^h$) tại các điểm của ${}_{n^h}G$. Giả thiết về các thớ
$G_s^\circ$ được dùng chính xác để bảo đảm rằng các điểm có cấp là một lũy
thừa của $n$ trong $G_s$ trù mật trong hợp của các $G^{(n^h)}$; từ đó dễ suy
ra $G\longrightarrow S$ mở phổ dụng tại mọi điểm của hợp này, đặc biệt dọc
theo tiết diện đơn vị, và rồi dễ suy ra $s\longmapsto\dim G_s$ là một hàm
hằng địa phương.

\result{NHẬN XÉT}{1.6} Nhắc lại rằng một cấu xạ $X\longrightarrow S$ được gọi
là mở phổ dụng nếu nó biến mọi tập mở thành một tập mở và vẫn giữ tính chất
đáng khen ấy sau mọi phép đổi cơ sở $S'\longrightarrow S$. Điều này cũng có
nghĩa (khi $S$ Noether địa phương và $X\longrightarrow S$ thuộc kiểu hữu hạn
địa phương) rằng mọi thành phần bất khả quy của $X$ trội lên $S$, và tính chất
này được bảo toàn sau mọi phép đổi cơ sở $S'\longrightarrow S$. Trong cả hai
khẳng định ấy, hơn nữa, chỉ cần (với các giả thiết hữu hạn ở trên) kiểm tra
bằng các phép đổi cơ sở $S'\longrightarrow S$, trong đó $S'$ là phổ của một
vành định giá rời rạc (đầy đủ với trường thặng dư đóng đại số nếu muốn~...).
Định nghĩa mở rộng hiển nhiên cho trường hợp một bộ phận $Z$ của $X$ (chẳng
hạn bộ phận $G^\sigma$ của $G$). Hoàn toàn bình thường là, ngay cả khi bắt đầu
từ một cấu xạ xạ ảnh trơn $X\longrightarrow S$ có các thớ hình học liên thông
(chẳng hạn tích thớ của họ môđun các đường cong elliptic trên
$S=\operatorname{Spec}\mathbf C[j]$ với chính nó),
$\operatorname{Pic}_{X/S}$ vẫn không mở phổ dụng trên $S$, nghĩa là có thể có
những thành phần bất khả quy của $\operatorname{Pic}_{X/S}$ nằm hoàn toàn
trên một điểm duy nhất của $S$: điều này liên quan đến việc hạng của nhóm
Néron--Severi của các thớ của $X/S$ có thể nhảy lên (các hiện tượng «nhân
phức»). Trái lại, (1.5) bảo đảm rằng, trong những trường hợp thuận lợi,
$\operatorname{Pic}_{X/S}^\sigma$ (và dường như thường ngay cả
$\operatorname{Pic}_{X/S}^\tau$) mở phổ dụng trên $S$.
\src{229}
Cuối cùng, sau đây là một trường hợp hữu ích trong đó, một cách ngoại lệ,
$G^\circ$ lại là mở:

\result{HỆ QUẢ}{1.7} Giả sử $G\longrightarrow S$ mở phổ dụng tại các điểm của
$G^\circ$ (xem (1.5)), và các thớ $G_s$ khả ly, do đó trơn trên $k(s)$ (điều
kiện sau tự động được thỏa mãn khi đặc trưng thặng dư bằng không, theo một kết
quả của Cartier). Khi đó $G^\circ$ mở trong $G$. Nếu thêm nữa $S$ rút gọn, thì
$G^\circ$ trơn (đặc biệt là phẳng) trên $S$.

Khẳng định thứ nhất còn có thể được làm chính xác hơn như sau. Nếu, với một
$s\in S$ cho trước, $G\longrightarrow S$ mở phổ dụng tại các điểm của
$G_s^\circ$ và $G_s^\circ$ khả ly trên $k(s)$, thì $G^\circ$ là một lân cận
của $G_s^\circ$ (hơn nữa, giả thiết đặt tại $s$ vẫn sẽ đúng tại các điểm lân
cận). Mệnh đề này, vốn độc lập với mọi cấu trúc nhóm, sẽ được chứng minh trong
[2], IV, đoạn 7. Khẳng định cuối của (1.7), cũng độc lập với mọi cấu trúc nhóm
và mọi vấn đề về thành phần liên thông, là một trường hợp riêng của một tiêu
chuẩn phẳng được cho trong [IV], đoạn 5, và tiêu chuẩn này tổng quát hơn còn
kéo theo:

\result{HỆ QUẢ}{1.8} Cho $U$ là một bộ phận mở của một thớ $G_s$, sao cho
$G\longrightarrow S$ mở phổ dụng tại các điểm của $U$ (xem (1.5)). Nếu $G_s$
khả ly trên $k(s)$ và $S$ rút gọn tại $s$, thì $G$ phẳng trên $S$ tại các điểm
của $U$ (do đó, trong trường hợp này, trơn trên $S$ tại các điểm của $U$).

\result{CÁC NHẬN XÉT}{1.9} Tôi không biết, khi $G$ tách trên $S$, liệu
$G^\circ$ hay $G^\tau$ có luôn đóng trong $G$ hay không; điều này có vẻ khó xảy
ra. Dù sao, ta có những phản ví dụ hiển nhiên nếu bỏ giả thiết tách, chẳng hạn
lấy đường thẳng afin mà gốc được tách đôi vô hạn đếm được lần. Như vậy ta thu
được một lược đồ nhóm $G$ trên đường thẳng afin, mọi thớ của nó, trừ một thớ
bằng $\mathbf Z$, đều rút về nhóm chỉ gồm phần tử đơn vị. Lược đồ nhóm này là
con mở, cụ thể là bao đóng của tiết diện đơn vị, trong tiền lược đồ Picard của
$S$-lược đồ $X$ tương ứng với một họ conic suy biến thành hai đường thẳng giao
nhau.

Hơn nữa, ngay cả khi bắt đầu với một lược đồ nhóm hữu hạn và phẳng $G$ trên
phổ $S$ của một vành định giá rời rạc $V$, sao cho $G^\circ$ chỉ còn là tiết
diện đơn vị và do đó là đóng, nhiều kết luận vẫn trở thành sai nếu bỏ một số
giả thiết. Giả sử $V$ có đồng đặc trưng $p>0$, và cho $G$ là hạt nhân của đồng
cấu
$\mathbf G_a\longrightarrow\mathbf G_a$ được xác định bởi đồng cấu các hàm tử
\[
f\longmapsto f^p-tf,
\]
trong đó $t$ là một phần tử đều hóa của $V$.
\src{230}
(Nhắc lại rằng, theo định nghĩa, «nhóm cộng» $\mathbf G_a$ trên $S$ biểu diễn
hàm tử $S'\longmapsto\Gamma(S',\mathcal O_{S'})$.) Khi đó, với tư cách là một
$S$-lược đồ,
\[
G=\operatorname{Spec}V[x]/(x^p-tx)
\]
là hợp của $p$ «đường thẳng» giao nhau, mô tả một nhóm
$\mathbf Z/p\mathbf Z$ suy biến thành một nhóm vô cùng bé kiểu cộng. Trong ví
dụ này, $G^\circ=G^\sigma$ (và đó là một trong $p$ đường thẳng), nên không mở
trong $G$. Trong trường hợp khác đặc trưng, với đặc trưng thặng dư $p>0$, ta
bắt đầu từ lược đồ nhóm $H=\symbf{\mu}_p$, là hạt nhân của phép lấy lũy thừa
bậc $p$ trong $\mathbf G_m$, do đó
\[
H=\operatorname{Spec}V[x]/(x^p-1) ;
\]
đây lại là hợp của $p$ đường thẳng giao nhau, mô tả một nhóm
$\mathbf Z/p\mathbf Z$ (trong đặc trưng 0) suy biến thành một nhóm vô cùng bé
kiểu nhân trong đặc trưng $p$. Cho $H'$ là «lược đồ nhóm hằng» được xác định
bởi nhóm hữu hạn thông thường $\mathbf Z/p\mathbf Z$, tức hợp rời của
$\mathbf Z/p\mathbf Z$ bản sao của $S$, hay cũng là
$\operatorname{Spec}V^{\mathbf Z/p\mathbf Z}$. Khi đó $G=H\times_S H'$ mô tả
một nhóm $(\mathbf Z/p\mathbf Z)^2$ trong đặc trưng 0, suy biến thành tích của
một nhóm vô cùng bé với $\mathbf Z/p\mathbf Z$ trong đặc trưng $p$. Ở đây
$G^\rho$ là hợp của tiết diện đơn vị và thớ đặc biệt, nên $G^\rho$ không mở,
trong khi $G^\sigma$ là hợp của tiết diện đơn vị và thớ tổng quát, nên không
đóng, trái với điều được khẳng định trong các trường hợp đồng đặc trưng ở
(1.1), (iii) và (iv). Dĩ nhiên, đây là những hiện tượng gắn với đặc trưng
$p>0$. Thật vậy, các kết quả trên cho:

\result{HỆ QUẢ}{1.10} Giả sử mọi đặc trưng thặng dư của $S$ đều bằng không,
do đó $G^\tau=G^\sigma$ và $G^\circ=G^\rho$. Khi đó $G^\tau=G^\sigma$ là mở,
thậm chí vừa mở vừa đóng nếu $G$ tách trên $S$ và các $G_s^\circ$ riêng; dưới
cùng giả thiết ấy, $G^\circ=G^\rho$ riêng trên $S$, nên đóng trong $G$. Cuối
cùng, giả sử với mọi số nguyên $n>1$, đồng cấu lũy thừa bậc $n$ trong $G$ mở
phổ dụng; khi đó $G^\circ$ mở. Nếu thêm nữa các $G_s^\circ$ không có thành
phần cộng (chẳng hạn các $G_s^\circ$ riêng, như trên), thì
$G^\tau\longrightarrow S$ mở phổ dụng, và thậm chí trơn nếu $S$ rút gọn.

Cuối cùng, ta ghi nhận kết quả dễ sau đây:

\result{MỆNH ĐỀ}{1.11} Tồn tại một bộ phận mở $U$ của $S$ sao cho tập $G'$ gồm
các điểm của $G$ tại đó $G$ trơn (tương ứng phẳng) trên $S$ là tập nền của một
lược đồ nhóm con mở cảm sinh từ $G|U$. Hơn nữa, mọi tiết diện của $G$ trên $U$
đều là một tiết diện của $G'$ trên $U$.
\src{231}
\result{HỆ QUẢ}{1.12} Nếu $G$ trơn (tương ứng phẳng) trên $S$ tại các điểm của
tiết diện đơn vị, thì nó cũng như vậy tại các điểm của mọi tiết diện của $G$
trên $S$, và tại mọi điểm của $G^\circ$. Nếu thêm nữa, với mọi số nguyên
$n>0$, đồng cấu lũy thừa bậc $n$
$\varphi_n:G\longrightarrow G$ là étale tại các điểm có đặc trưng nguyên tố
cùng nhau với $n$, thì $G$ trơn (tương ứng phẳng) trên $S$ tại mọi điểm của
$G^\sigma$.

\fsection{2}{Ứng dụng vào các tính chất địa phương của các lược đồ Picard}

\result{ĐỊNH LÝ}{2.1}

\noindent(i) Cho $f:X\longrightarrow S$ là một cấu xạ riêng và trơn, và giả
sử $\operatorname{Pic}_{X/S}$ tồn tại (chẳng hạn $f$ là một cấu xạ xạ ảnh).
Khi đó $\operatorname{Pic}_{X/S}$ tách trên $S$, và với mọi bộ phận đóng $Z$
của $\operatorname{Pic}_{X/S}$ thuộc kiểu hữu hạn trên $S$, $Z$ riêng trên
$S$.

\noindent(ii) Cho $X$ là một tiền lược đồ trên một trường $k$, riêng và chuẩn
hình học. Khi đó $\operatorname{Pic}_{X/S}^\circ$ riêng trên $k$.

\begin{proof}
(i) Các tiêu chuẩn định giá ([2], II, đoạn 7) quy ta về việc chứng minh điều
sau: nếu $S$ là phổ của một vành định giá rời rạc đầy đủ, $U$ là tập mở chỉ
gồm điểm tổng quát của $S$, thì mọi tiết diện hữu tỷ của
$\operatorname{Pic}_{X/S}$ trên $S$, nghĩa là mọi tiết diện trên $U$, đều kéo
dài duy nhất thành một tiết diện trên $S$. Theo định nghĩa của
$\operatorname{Pic}_{X/S}$, điều này tương đương với mệnh đề sau: với mọi
Môđun khả nghịch $\mathcal L$ trên $V=f^{-1}(U)$, tồn tại một Môđun khả nghịch
trên $X$ kéo dài $\mathcal L$, duy nhất sai khác một đẳng cấu. Điều này dễ dàng
suy ra từ mô tả các Môđun khả nghịch trên $V$, tương ứng trên $X$, bằng các
lớp «ước Cartier», vì các vành địa phương của $X$ là chính quy ($X$ trơn trên
lược đồ chính quy $S$), nên là các vành phân tích duy nhất theo Auslander; do
đó mọi ước trên $S$ đều là một ước
Cartier. Thật vậy, mọi ước trên $V$ đều có thể được kéo dài thành một ước trên
$X$ bằng cách lấy «bao đóng» của nó.
\end{proof}

\result{NHẬN XÉT}{2.2} Chứng minh vẫn đúng nếu chỉ giả sử $f$ phẳng và các thớ
$X_s$ là các giao đầy đủ địa phương và trơn trên $k(s)$ ở đối chiều $\leq2$,
theo kết quả sau được chứng minh trong [5]: một vành địa phương Noether là giao
đầy đủ và chính quy ở đối chiều $\leq3$ thì là vành phân tích duy nhất («phỏng
đoán Samuel»). Cần lưu ý rằng
\src{232}
kết quả trở thành sai nếu thay «đối chiều $\leq2$» bằng «đối chiều
$\leq1$», nghĩa là bằng giả thiết «chuẩn», như có thể thấy từ ví dụ một họ các
siêu mặt bậc hai không kỳ dị suy biến thành một nón bậc hai.

(ii) Dùng bổ đề Chow, ta quy về trường hợp $X$ xạ ảnh, do đó được nhúng vào
$\mathbf P_k^n$; ta có thể giả sử $X$ liên thông. Nếu $\dim X=1$, thì $X$ trơn
trên $k$, và ta áp dụng (i). Nếu $\dim X\geq2$, ta có thể dùng một biến thể của
các «tiêu chuẩn tương đương» đã biết, từ đó suy ra rằng tồn tại hữu hạn đường
cong $Y_i$ trơn nằm trong $X$ (thu được bằng cách giao $X$ với các không gian
con tuyến tính thích hợp của $\mathbf P_k^n$), sao cho
\[
\operatorname{Pic}_{X/k}^\tau\longrightarrow
\prod_i\operatorname{Pic}_{Y_i/k}^\tau
\]
là một đơn cấu, và do đó cảm sinh một đơn cấu trên các thành phần liên thông.
Vì thành phần liên thông của vế phải riêng trên $k$ theo kết quả trên, và vì
đồng cấu đang xét là một đồng cấu các lược đồ nhóm, tất yếu là một phép nhúng
đóng, nên $\operatorname{Pic}_{X/k}^\circ$ cũng riêng trên $k$. Có thể tránh
việc dùng các tiêu chuẩn tương đương tinh tế bằng cách sử dụng cấu trúc của các
nhóm đại số giao hoán trên một trường đóng đại số (do Chevalley--Borel); ta quy
về việc chứng minh rằng mọi cấu xạ từ đường thẳng afin vào
$\operatorname{Pic}_{X/k}^\tau$ là hằng. Điều này tương đương với việc mọi
Môđun khả nghịch trên $X[t]$ đều đến từ một Môđun khả
nghịch trên $X$. Kết quả quen thuộc và rất sơ cấp này thậm chí không dùng đến
việc $X$ riêng trên $k$ (giả thiết $X$ chuẩn cho phép lập tức quy về trường hợp
$X$ chính quy).

\result{HỆ QUẢ}{2.3} Cho $f:X\longrightarrow S$ là một cấu xạ riêng và chuẩn
(nghĩa là phẳng với các thớ hình học chuẩn), và giả sử
$\operatorname{Pic}_{X/S}$ tồn tại. Khi đó $\operatorname{Pic}_{X/S}^\circ$
riêng trên $S$, nên đóng trong $\operatorname{Pic}_{X/S}$; hơn nữa,
$\operatorname{Pic}_{X/S}^\tau$ và $\operatorname{Pic}_{X/S}^\rho$ cũng đóng,
và $\operatorname{Pic}_{X/S}^\sigma$ cũng vậy trong «đồng đặc trưng».

Ta áp dụng (1.1) và (2.1), (ii).

\result{HỆ QUẢ}{2.4} Cho $f:X\longrightarrow S$ là một cấu xạ riêng và trơn,
sao cho $\operatorname{Pic}_{X/S}$ tồn tại và là hợp rời của các lược đồ
$P^{(i)}$ thuộc kiểu hữu hạn trên $S$ (xem [3], V, 4.1). Khi đó mỗi $P^{(i)}$
riêng trên $S$.

Điều này suy ra từ (2.1), (i). Như đã nêu trong (2.2), kết quả có thể được tổng
quát hóa bằng cách đặt các giả thiết ít hạn chế hơn lên các thớ của $f$, nhưng
trở thành sai nếu ta chỉ giả sử $f$ chuẩn. Trong trường hợp này, tôi cũng không
biết liệu $\operatorname{Pic}_{X/S}^\tau$ có vẫn riêng trên $S$ hay không,
ngay cả khi giả sử nó thuộc kiểu hữu hạn trên $S$.
\src{233}
\result{ĐỊNH LÝ}{2.5} Cho $f:X\longrightarrow S$ là một cấu xạ riêng và
phẳng, sao cho $\operatorname{Pic}_{X/S}$ tồn tại, và với mỗi số nguyên $n$,
cho $\varphi_n$ là đồng cấu lũy thừa bậc $n$ trong tiền lược đồ nhóm này. Khi
đó $\varphi_n$ là étale tại mọi điểm
$x\in X$ có đặc trưng
thặng dư nguyên tố cùng nhau với $n$.

Theo đặc trưng hóa vi phân của các cấu xạ étale, khẳng định này tương đương với
khẳng định sau:

\result{BỔ ĐỀ}{2.6} Giả sử $S$ là phổ của một vành Artin địa phương $A$ mà
iđêan cực đại $\mathfrak m$ có lũy thừa bậc $(\nu+1)$ bằng không, và đặt
\[
A_{\nu-1}=A/\mathfrak m^\nu,
\qquad X_{\nu-1}=X\otimes_A A_{\nu-1},
\]
$\mathcal L$ là một Môđun khả nghịch trên $X$, và cuối cùng
$\mathcal L'_{\nu-1}$ là một Môđun khả nghịch trên $X_{\nu-1}$ mà lũy thừa
tensor bậc $n$ đẳng cấu với
\[
\mathcal L_{\nu-1}=\mathcal L\otimes_A A_{\nu-1}.
\]
Khi đó tồn tại một Môđun khả nghịch $\mathcal L'$ trên $X$ mà lũy thừa tensor
bậc $n$ đẳng cấu với $\mathcal L$ (nếu $n$ nguyên tố cùng nhau với đặc trưng
thặng dư của $k=k(A)$).

Ta chứng minh bổ đề. Đặt
\[
V=\mathfrak m^\nu=\mathfrak m^\nu/\mathfrak m^{\nu+1} ;
\]
đây là một không gian vectơ trên $k=k(A)$. Trước hết ta tìm cách kéo dài
$\mathcal L'_{\nu-1}$ thành một Môđun khả nghịch tùy ý $\mathcal L'$ trên
$X$. Trở ngại cho việc này nằm trong
\[
H^2(X_0,\mathcal O_{X_0})\otimes_k V,
\]
nhưng theo giả thiết về $\mathcal L'_{\nu-1}$ và sự kiện
$\mathcal L_{\nu-1}$ có thể được kéo dài, ta thấy tích của trở ngại này với
$n$ bằng không; do đó chính nó bằng không vì $n$ nguyên tố cùng nhau với đặc
trưng. Hơn nữa, độ tùy ý trong phép kéo dài đã chọn nằm trong
\[
H^1(X_0,\mathcal O_{X_0})\otimes_k V,
\]
và sai lệch $\xi$ của $\mathcal L'^{\otimes n}$ so với $\mathcal L$ nằm trong
cùng Môđun ấy. Nếu ta tìm cách hiệu chỉnh $\mathcal L'$ để làm sai lệch này
bằng không, thì bài toán quy về việc tìm một $\eta$ trong Môđun nói trên sao
cho $n\eta=\xi$. Điều này lại khả thi vì $n$ nguyên tố cùng nhau với đặc
trưng.

\result{HỆ QUẢ}{} Dưới các điều kiện của (2.5), giả sử thêm rằng các lược đồ
Picard của các thớ $X_s$ không chứa thành phần cộng (chẳng hạn các $X_s$ chuẩn
hình học, xem (2.1), (ii)). Khi đó
$\operatorname{Pic}_{X/S}\longrightarrow S$ mở phổ dụng tại các điểm của
$\operatorname{Pic}_{X/S}^\circ$. Nếu $\operatorname{Pic}_{X/S}^\circ$ đóng
(chẳng hạn các $X_s$ chuẩn hình học, xem (2.3)), thì chính
$\operatorname{Pic}_{X/S}^\sigma$ mở phổ dụng trên $S$. Cuối cùng, trong
trường hợp đồng đặc trưng, $\operatorname{Pic}_{X/S}^\rho\longrightarrow S$
mở phổ dụng.

Ta áp dụng (1.5) và (1.1).
\src{234}
\result{HỆ QUẢ}{2.7} Cho $f:X\longrightarrow Y$ là một cấu xạ riêng và phẳng, sao cho $\operatorname{Pic}_{X/S}$ tồn
tại. Khi đó hàm số
\[
s\longmapsto\dim\operatorname{Pic}_{X_s/k(s)}
\]
trên $S$ là nửa liên tục trên (nghĩa là nó có thể nhảy lên, nhưng không thể
nhảy xuống), và thậm chí liên tục (nghĩa là hằng địa phương) nếu các
$\operatorname{Pic}_{X_s/k(s)}$ không chứa thành phần cộng.

Khẳng định thứ nhất đúng một cách tầm thường, hoặc gần như vậy, đối với mọi
tiền lược đồ nhóm thuộc kiểu hữu hạn địa phương trên một cơ sở Noether địa
phương, vì chỉ cần xét dọc theo tiết diện đơn vị. Khẳng định thứ hai suy ra từ
(2.5).

\result{NHẬN XÉT}{2.8} Cho $s,s'\in S$ sao cho $s$ là một chuyên biệt hóa của
$s'$. Khi đó (2.7) tương đương với một bất đẳng thức (tương ứng một đẳng thức)
giữa các chiều của các lược đồ Picard của $X_{s'}$ và của «chuyên biệt hóa»
$X_s$ của nó. Hơn nữa, ngay từ trước khi các lược đồ Picard được xây dựng,
Serre đã nhận xét rằng tính bất biến của chiều các lược đồ Picard của $X_s$ trong
trường hợp một cấu xạ trơn $f:X\longrightarrow S$ là hệ quả hình thức của lý
thuyết chuyên biệt hóa nhóm cơ bản ([4], X), và của các quan hệ Kummer cổ điển
giữa các điểm có cấp hữu hạn trên đa tạp Picard với nhóm cơ bản Abel hóa ([4],
XI). Nếu nói chung ta ký hiệu bởi $\alpha,\mu,\lambda$ các chiều của phần Abel,
phần nhân và phần cộng của $\operatorname{Pic}_{X_s/k(s)}^\circ$, và định
nghĩa tương tự $\alpha',\mu',\lambda'$, thì các quan hệ đã biết được biểu thị
bởi những bất đẳng thức sau:
\[
\tag{*}\alpha+\mu+\lambda\geq\alpha'+\mu'+\lambda'
\]
(đúng miễn là $\operatorname{Pic}_{X/S}$ tồn tại, do đó có lẽ đúng trong mọi
trường hợp), và khi $\lambda=\lambda'=0$, bất đẳng thức này rút về một đẳng
thức, đúng dưới cùng các giả thiết tồn tại:
\[
\alpha+\mu=\alpha'+\mu' ;
\]
mặt khác ta có
\[
\tag{**}2\alpha+\mu\leq2\alpha'+\mu'
\]
\src{235}
theo lập luận của Serre, nếu các $X_s$ khả ly (thậm chí không cần giả sử
$\operatorname{Pic}_{X/S}^\tau$ tồn tại), hoặc nếu các
${}_n\operatorname{Pic}_{X/S}$ (các hạt nhân của $\varphi_n$ trong hàm tử
Picard) tách trên $S$ (vì chúng étale trên $S$ theo (2.5)). Ta có xu hướng
phỏng đoán rằng (*) luôn là một đẳng thức, hoặc ít nhất khi các $X_s$ khả ly,
và rằng ta có các bất đẳng thức
\[
\tag{***}\alpha\leq\alpha',\qquad\lambda\geq\lambda',
\]
những bất đẳng thức này phải đúng mỗi khi ta có một tiền lược đồ nhóm thuộc
kiểu hữu hạn địa phương trên $S$ Noether địa phương, trong đó chiều của các
thớ là hằng (xem (1.3) về một kết quả khẳng định theo hướng này).

\result{NHẬN XÉT}{2.9} Trong mọi trường hợp đã biết,
$\operatorname{Pic}_{X/S}^\tau$ mở phổ dụng trên $S$, nhưng có lẽ không nên vì
thế mà nuôi quá nhiều ảo tưởng, ngay cả khi $f:X\longrightarrow S$ trơn. Dù
sao, Mumford đã xây dựng một ví dụ (đúng là với $S$ không rút gọn, thực ra $S$
là phổ của một vành Artin) trong đó $\operatorname{Pic}_{X/S}^\tau$ không
phẳng trên $S$, bằng cách làm biến thiên vô cùng bé mặt Igusa. Hơn nữa, điểm
được Mumford xét nằm trong $\operatorname{Pic}_{X/S}^\rho$, và vẫn có thể xảy
ra rằng (với $f:X\longrightarrow S$ trơn) $\operatorname{Pic}_{X/S}^\tau$
phẳng trên $S$ tại các điểm của $\operatorname{Pic}_{X/S}^\sigma$; tuy nhiên,
diễn giả nghi ngờ rằng điều này luôn đúng, ngay cả khi giới hạn ở các điểm của
$\operatorname{Pic}_{X/S}^\circ$ và giả sử $S$ là phổ của một vành định giá
rời rạc. Câu hỏi này còn gắn với việc nghiên cứu các điểm bất động của một
lược đồ Abel dưới tác động của một nhóm hữu hạn các tự đẳng cấu, một tình
huống mà dường như ta còn thiếu ví dụ.

Dường như, ngay cả khi chỉ xét $f:X\longrightarrow S$ trơn và xạ ảnh, các kết
quả về tính chính quy địa phương của $\operatorname{Pic}_{X/S}$ được nêu trong
mục này, cùng với các phỏng đoán được chỉ ra trong (2.8), bao quát gần như tất
cả những gì có thể nói về vấn đề này nếu không có các giả thiết cụ thể hơn về bản chất
các thớ của $f$. Tuy nhiên, nhắc lại rằng nếu các thớ hình học của
$\operatorname{Pic}_{X/S}$ rút gọn và không có thành phần cộng, thì từ (1.8)
và (2.5) suy ra rằng $\operatorname{Pic}_{X/S}$ trơn trên $S$ tại các điểm của
$\operatorname{Pic}_{X/S}^\sigma$ khi $S$ rút gọn; kết quả này thậm chí vẫn
đúng, nếu $f:X\longrightarrow S$ chuẩn, mà không có giả thiết nào về $S$, như
ta sẽ thấy trong (3.5). Liên quan đến điều này, ta ghi nhận:

\result{MỆNH ĐỀ}{2.10}

\noindent(i) Nếu $\operatorname{Pic}_{X/S}$ trơn (tương ứng phẳng) trên $S$
tại các điểm của tiết diện đơn vị, thì nó cũng như vậy tại mọi điểm của
$\operatorname{Pic}_{X/S}^\sigma$, và tại các điểm của mọi tiết diện của
$\operatorname{Pic}_{X/S}$ trên $S$.
\src{236}
\noindent(ii) Cho $s\in S$ sao cho
\[
H^2(X_s,\mathcal O_{X_s})=0,
\]
khi đó tồn tại một lân cận mở $U$ của $s$ sao cho
$\operatorname{Pic}_{X/S}|U$ trơn trên $U$.

\noindent(iii) Cho $X$ là một lược đồ riêng trên một trường, khi đó
\[
\dim\operatorname{Pic}_{X/k}\leq\dim H^1(X,\mathcal O_X),
\]
dấu bằng xảy ra khi và chỉ khi $\operatorname{Pic}_{X/k}$ trơn trên $k$; điều
này luôn đúng nếu $k$ có đặc số không.

(i) suy ra từ (2.5) và (1.12), còn (ii) suy ra từ tiêu chuẩn vi phân vô cùng
đối với các cấu xạ trơn và một phép tính chướng ngại quen thuộc, vì (theo
``định lý nửa liên tục'') giả thiết đặt tại $s$ vẫn đúng tại các điểm lân cận.
Cuối cùng, (iii) suy ra từ sự kiện $H^1(X,\mathcal O_X)$ đẳng cấu với không
gian tiếp xúc Zariski của $\operatorname{Pic}_{X/k}$ tại phần tử đơn vị; khẳng
định cuối cùng là một trường hợp riêng của một định lý của Cartier, nói rằng
một ``nhóm hình thức'' trong đặc số 0 là trơn hình thức trên $k$.

\fsection{3}{Lược đồ con Abel chính tắc của
\texorpdfstring{$\operatorname{Pic}_{X/S}$}{Pic(X/S)} và lược đồ Albanese}

\result{MỆNH ĐỀ}{3.1} Cho $k$ là một trường và $G$ là một lược đồ nhóm giao
hoán kiểu hữu hạn trên $k$, ``không có thành phần cộng''. Khi đó
$G_{\mathrm{réd}}^\circ$ khả ly trên $k$, và do đó nó là một lược đồ nhóm con
trơn trên $k$.

Điều này hiển nhiên nếu $k$ hoàn hảo, đặc biệt đối với $G_{\bar k}$, trong đó
$\bar k$ là bao đóng đại số của $k$; vì thế chỉ cần chứng minh rằng
$(G_{\bar k}^\circ)_{\mathrm{réd}}$ được hạ xuống từ một lược đồ con của $G$.
Bây giờ, giả thiết trên $G_{\bar k}$ (không có thành phần cộng) dễ dàng suy ra rằng tồn tại một số nguyên $m$ sao cho
$(G_{\bar k}^\circ)_{\mathrm{réd}}$ là ảnh theo nghĩa lược đồ của đồng cấu lũy
thừa bậc $m$ trong $G_{\bar k}$. Vì đồng cấu cuối này đến từ đồng cấu tương tự
đối với $G$, ảnh theo nghĩa lược đồ của đồng cấu sau đáp ứng yêu cầu.

\result{HỆ QUẢ}{3.2} Cho $X$ là một lược đồ chuẩn và riêng trên $k$ sao cho
$\operatorname{Pic}_{X/k}$ tồn tại. Khi đó tồn tại một lược đồ con Abel $A$
của $\operatorname{Pic}_{X/k}$ có tập nền là
$\operatorname{Pic}_{X/k}^\circ$.

Thật vậy, theo (2.1), (ii), $\operatorname{Pic}_{X/k}^\circ$ riêng trên $k$
nên thỏa mãn điều kiện của (3.1).
\src{237}
Vậy kết quả trước cho thấy rằng, trong một số trường hợp, ``đa tạp Picard''
cổ điển (mà trong lý thuyết hiện nay sẽ được ký hiệu là
$(\operatorname{Pic}_{X/k})_{\mathrm{réd}}^\circ$) ``được xác định trên
$k$'', mà không cần giả sử trường $k$ là hoàn hảo.

Bây giờ, cho $f:X\longrightarrow S$ là một lược đồ tương đối riêng và phẳng,
với
\[
\mathcal O_S\xrightarrow{\sim}f_*(\mathcal O_X)
\]
để đơn giản, sao cho $\operatorname{Pic}_{X/S}$ tồn tại và
$\operatorname{Pic}_{X/S}^\circ$ riêng trên $S$. Ngoài ra, đối với (3.3),
(ii), ta giả sử rằng có một tập mở của $\operatorname{Pic}_{X/S}$ chứa
$\operatorname{Pic}_{X/S}^\circ$ và tựa xạ ảnh trên $S$; như đã thấy, điều
kiện này được thỏa mãn nếu $f$ là xạ ảnh và có các thớ hình học khả ly và bất khả
quy. Nhắc lại rằng một lược đồ Abel trên $S$ là một lược đồ nhóm trên $S$,
riêng và trơn trên $S$, có các thớ liên thông hình học. Ta xét xem có tồn tại
một lược đồ nhóm con $A$ của $\operatorname{Pic}_{X/S}$, vừa là một lược đồ
Abel vừa có tập nền là $\operatorname{Pic}_{X/S}^\circ$, hay không. Ta vừa
thấy rằng một lược đồ như thế luôn tồn tại nếu $S$ là phổ của một trường. Sau
đây là điều có thể khẳng định trong trường hợp tổng quát đang xét:

\result{ĐỊNH LÝ}{3.3} Dưới các điều kiện trên:

\noindent(i) Nếu tồn tại một lược đồ con Abel của
$\operatorname{Pic}_{X/S}$ có tập nền là
$\operatorname{Pic}_{X/S}^\circ$, thì nó là duy nhất. Do đó, phép tạo thành
nó tương thích với các phép đổi cơ sở.

\noindent(ii) Để một lược đồ con Abel như thế tồn tại, chỉ cần nó tồn tại sau
mọi phép đổi cơ sở $S'\longrightarrow S$, trong đó $S'$ là Artin địa phương;
nếu $S$ là phổ của một vành địa phương, thì thậm chí chỉ cần kiểm tra với
$S'=\operatorname{Spec}(A_n)$, trong đó
$A_n=A/\mathfrak m^{n+1}$. Nếu $S$ rút gọn, cũng chỉ cần kiểm tra với các phép
đổi cơ sở $S'\longrightarrow S$, trong đó $S'$ là phổ của một vành định giá
rời rạc (đầy đủ với trường thặng dư đóng đại số, nếu muốn).

\noindent(iii) Giả sử $A$ tồn tại, và cho
$B=\operatorname{Alb}^0(X/S)$ là lược đồ Abel đối ngẫu (tức là
$B=\operatorname{Pic}_{A/S}^\circ$ [7]). Khi đó, ta có thể xây dựng chính tắc
một không gian thuần nhất chính $P=\operatorname{Alb}^1(X/S)$ dưới $B$, và
một $S$-cấu xạ $X\longrightarrow P$ phổ dụng đối với các $S$-cấu xạ từ $X$
vào các lược đồ para-Abel (tức là các không gian thuần nhất chính dưới các
lược đồ Abel). Phép tạo thành $\operatorname{Alb}^0(X/S)$,
$\operatorname{Alb}^1(X/S)$ và
$X\longrightarrow\operatorname{Alb}^1(X/S)$ giao hoán với phép đổi cơ sở.

Ta phác thảo chứng minh:

(i) là một tính chất cứng tổng quát của các lược đồ con Abel trong các lược đồ
nhóm giao hoán: nếu hai lược đồ con như thế trùng nhau theo nghĩa tập hợp tại
một điểm $s\in S$, thì chúng trùng nhau trên toàn bộ thành phần
\src{238}
liên thông của $s$ [7]. (Kết quả này khái quát một định lý cổ điển của Chow.)

(ii) Dùng các lược đồ Hilbert, ta thấy hàm tử gán cho mỗi $S'$ trên $S$ tập
(chỉ gồm 1 hoặc 0 phần tử) các lược đồ con Abel chính tắc của
$(\operatorname{Pic}_{X/S})\times_S S'$ được biểu diễn bởi một lược đồ $T$
kiểu hữu hạn trên $S$. Theo (i), $T\longrightarrow S$ là một đơn cấu, và theo
(3.2), nó là toàn ánh. Nói rằng tồn tại một lược đồ con Abel chính tắc của
$\operatorname{Pic}_{X/S}$ có nghĩa là $T$ có một tiết diện trên $S$, hay
tương đương, $T\longrightarrow S$ là một đẳng cấu. Điều này cũng tương đương
với việc $T\longrightarrow S$ là étale, hoặc, trong trường hợp $S$ rút gọn,
với việc $T\longrightarrow S$ là riêng. Từ đó suy ra ngay (ii).

(iii) Chỉ cần dùng định nghĩa của $\operatorname{Pic}_{X/S}$, ta thấy rằng,
đối với mọi lược đồ Abel $C$ trên $S$, việc cho một $S$-cấu xạ từ $X$ vào một
không gian thuần nhất chính dưới $C$ tương đương với việc cho một đồng cấu nhóm
$C'\longrightarrow\operatorname{Pic}_{X/S}$, trong đó $C'$ là lược đồ Abel
đối ngẫu của $C$. Mà nếu lược đồ con Abel chính tắc $A$ của
$\operatorname{Pic}_{X/S}$ tồn tại, thì các đồng cấu này nhất thiết phân tích
qua $A$ (cũng có thể thấy điều này bằng cách dùng các điểm có cấp hữu hạn). Từ
đó suy ra ngay (iii).

\result{NHẬN XÉT}{3.4} Ta sẽ ký hiệu bởi
$\operatorname{Pic}_{X/S}^{\circ\circ}$ lược đồ con Abel chính tắc của
$\operatorname{Pic}_{X/S}$, nếu nó tồn tại. Đáng tiếc là điều này không phải
lúc nào cũng xảy ra, như có thể thấy khi làm biến thiên vô cùng bé mặt Igusa
(bằng biến dạng modular bậc nhất). Tuy nhiên, có thể
$\operatorname{Pic}_{X/S}^{\circ\circ}$ ít nhất tồn tại nếu $S$ rút gọn,
hoặc, theo (ii), tương đương với điều đó, nếu $S$ là phổ của một vành định giá
rời rạc. Khi đó, cho $X_0,X_1$ là thớ đặc biệt và tổng quát của $X$, và đặt
\[
A_1=\operatorname{Pic}_{X_1/K}^\circ,
\]
trong đó $K$ là trường phân thức của vành định giá $V$. Theo Koizumi, tồn tại
một lược đồ Abel $A$ trên $S$, duy nhất về căn bản, có thớ tổng quát là $A_1$,
và dễ dàng thấy như trong (2.1), (i) (từ đây giả sử $X$ trơn trên $S$), rằng
cấu xạ đồng nhất của $A_1$ mở rộng thành một cấu xạ
\[
A\longrightarrow\operatorname{Pic}_{X/S}.
\]
Do đó có một đồng cấu
\[
\tag{*}A_0\longrightarrow\operatorname{Pic}_{X_0/k}^{\circ\circ},
\]
mà ta dễ dàng chứng minh là toàn ánh, có hạt nhân là một $p$-nhóm hữu hạn,
trong đó $p$ là số mũ đặc số của trường thặng dư $k$
\src{239}
(vẫn bằng cách dùng các điểm có cấp hữu hạn). Với quy ước đó, các điều kiện
sau đây trên $X/S$ là tương đương:

\medskip
\noindent a. Đồng cấu (*) ở trên là một đẳng cấu (điều mà Shimura sẽ diễn đạt
bằng cách nói rằng phép tạo thành ``đa tạp Picard'' là ``tương thích với phép
chuyên biệt hóa'').

\noindent b. $\operatorname{Pic}_{X/S}^{\circ\circ}$ tồn tại (khi đó nó
không gì khác ngoài $A$).

\noindent c. (Để ghi nhớ) Các
$\operatorname{Pic}_{X_n/S}^{\circ\circ}$ tồn tại.

Theo nhận xét đã nêu về hạt nhân của (*), điều kiện (a) được thỏa mãn nếu đặc
số thặng dư bằng không, nhưng kết quả này sẽ được khái quát đáng kể trong
(3.5).

Dĩ nhiên, nếu $\operatorname{Pic}_{X/S}$ trơn trên $S$ tại các điểm của
$\operatorname{Pic}_{X/S}^\circ$, thì tập sau là mở trong
$\operatorname{Pic}_{X/S}$ (xem (1.7)) và do đó, với cấu trúc cảm sinh, nó là
một lược đồ con Abel của $\operatorname{Pic}_{X/S}$, nên bằng
$\operatorname{Pic}_{X/S}^{\circ\circ}$, vốn tồn tại trong trường hợp này.
Nhưng ta còn có kết quả mạnh hơn nhiều:

\result{ĐỊNH LÝ}{3.5} Dưới các điều kiện của (3.3), cho $s\in S$ sao cho
$\operatorname{Pic}_{X_s/k(s)}$ trơn trên $k(s)$ (hay, tương đương,
\[
\dim\operatorname{Pic}_{X_s/k(s)}=\dim H^1(X_s,\mathcal O_{X_s})).
\]
Khi đó tồn tại một lân cận mở $U$ của $s$ sao cho
$\operatorname{Pic}_{X/S}$ trơn trên $S$ trong các điểm của $\operatorname{Pic}_{X/S}^\circ|U$, và do đó đây là một lược
đồ con Abel mở trong $\operatorname{Pic}_{X/S}|U$. Đặc biệt,
$\operatorname{Pic}_{X|U/U}^{\circ\circ}$ tồn tại.

Ta nêu nguyên lý của chứng minh. Những điều trên đưa ta về trường hợp $S$ là
phổ của một vành artin địa phương $A$, và ta quy nạp theo cấp vô cùng bé của
$A$. Do đó, có thể giả sử $\operatorname{Pic}_{X_n/A_n}^\circ$ trơn trên
$A_n$, và chỉ cần chứng minh rằng
$\operatorname{Pic}_{X_{n+1}/A_{n+1}}^\circ$ trơn trên $A_{n+1}$. Ta thấy
rằng để làm điều này, chỉ cần xây dựng một lược đồ Abel $P_{n+1}$ trên
$A_{n+1}$ mở rộng
\[
P_n=\operatorname{Pic}_{X_n/A_n}^\circ,
\]
đồng thời xây dựng một Môđun khả nghịch $\mathcal L_{n+1}$ trên
$X_{n+1}\times_{A_{n+1}}P_{n+1}$ mở rộng Môđun khả nghịch $\mathcal L_n$
trên $X_n\times_{A_n}P_n$ xuất hiện trong định nghĩa của lược đồ Picard
$\operatorname{Pic}_{X_n/A_n}$ như nghiệm của một bài toán phổ dụng. (N.B. ---
Ta có thể giả sử $X$ có một tiết diện trên $S$~...) Để thực hiện phép xây dựng,
phải dùng kết quả then chốt sau: mọi lược đồ Abel xác định trên một vành thương
của một vành artin địa phương đều có thể mở rộng (nói cách khác, ``lược đồ
hình thức các Môđun'' tuyệt đối ([3], II) của một lược đồ Abel trên một trường
đóng đại số là trơn trên vành các
\src{240}
vectơ Witt của $k$); bản thân kết quả này thu được bằng cách dùng các tính chất
hình thức tổng quát của chướng ngại đối với phép nâng và các phép toán nhóm.
Chấp nhận kết quả này, trước hết ta mở rộng $P_n$ một cách tùy ý thành
$P_{n+1}$; khi đó gặp một chướng ngại đối với việc nâng $\mathcal L_n$, nằm
trong
\[
H^2(X_0\times P_0,\mathcal O_{X_0\times P_0})\otimes_k V
\qquad(\text{trong đó }V=\mathfrak m^{n+1}/\mathfrak m^{n+2}),
\]
và thực ra nằm trong không gian con
\[
H^1(X_0,\mathcal O_{X_0})\otimes
H^1(P_0,\mathcal O_{P_0})\otimes_k V
\]
(do hạn chế của $\mathcal L_n$ lên hai nhân tử $X_n$ và $P_n$ là tầm thường).
Mà không gian sau chính là
\[
H^1(P_0,\mathcal T_{P_0/k})\otimes V,
\]
trong đó $\mathcal T_{P_0/k}$ là bó tiếp xúc của $P_0/k$, vì vậy nó cũng là
không gian biểu thị tính bất định trong phép nâng $P_n$ thành $P_{n+1}$ ([3],
II). Do đó, có thể hiệu chỉnh phép nâng này (theo một và chỉ một cách, như phải
vậy) để triệt tiêu chướng ngại đối với việc nâng $\mathcal L_n$.

\result{HỆ QUẢ}{3.6} Dưới các điều kiện của (3.5),
$R^1f_*(\mathcal O_X)$ là một Môđun tự do địa phương trên $S$ trong một lân
cận của $s$, và phép tạo thành nó giao hoán với phép đổi cơ sở.

Thật vậy, Môđun này chính là Môđun tiếp xúc của
$\operatorname{Pic}_{X/S}$ dọc theo tiết diện đơn vị.

\result{NHẬN XÉT}{3.7} Dùng cùng lập luận như đối với (3.5), ta có thể chứng
minh rằng nếu $S'$ là một lược đồ con của $S$ được xác định bởi một Iđêan
coherent lũy linh, và nếu chỉ giả sử $\operatorname{Pic}_{X'/S'}^\circ$ tồn
tại và có các thớ trơn, thì $\operatorname{Pic}_{X/S}^\circ$ nhất thiết tồn
tại và là một lược đồ Abel trên $S$. Điều này cho phép xây dựng các lược đồ
Picard trong một số trường hợp mặc dù không có giả thiết xạ ảnh; chẳng hạn,
lược đồ Abel đối ngẫu của một lược đồ Abel trên một vành artin luôn tồn tại.
Mặt khác, dùng (3.6) trong trường hợp $X$ là một lược đồ Abel trên $S$, và dùng
cấu trúc đã biết của $H^*(X_s,\mathcal O_{X_s})$ như một đại số ngoài trên
$H^1(X_s,\mathcal O_{X_s})$ (Rosenlicht--Serre), ta thấy rằng với mọi $i$,
$R^if_*(\mathcal O_X)$ tự do địa phương, và chính xác hơn, đẳng cấu với lũy
thừa ngoài thứ $i$ của $R^1f_*(\mathcal O_X)$.

\result{NHẬN XÉT}{3.8} Trong trường hợp một cấu xạ xạ ảnh trơn
$f:X\longrightarrow S$, với $S$ rút gọn và có các đặc số thặng dư bằng không,
kết quả (3.6) đã được biết bằng phương pháp siêu việt, như một hệ quả của lý
thuyết Hodge. Thật ra, tất cả các $R^pf_*(\Omega^q_{X/S})$ khi đó đều tự do
địa phương. Trái lại, có các phản ví dụ
\src{241}
trong trường hợp đặc số hỗn hợp đối với tính tự do địa phương của
$R^1f_*(\mathcal O_X)$, bằng cách dùng các đa tạp Serre ([8], tr.~50). Dường
như không có phản ví dụ nào trong trường hợp cùng đặc số.

\result{HỆ QUẢ}{3.8} Dưới các điều kiện của (3.5),
$\operatorname{Pic}_{X/S}|U$ trơn trên $U$ tại mọi điểm của
$\operatorname{Pic}_{X/S}^\sigma|U$.

Áp dụng (2.10), (i). Đặc biệt, có tính đến (2.10), (iii):

\result{HỆ QUẢ}{3.9} Giả sử $S$ có các đặc số thặng dư bằng không. Khi đó
$\operatorname{Pic}_{X/S}^\tau$ trơn trên $S$.

\result{NHẬN XÉT}{3.10} Từ đó suy ra, chẳng hạn, nếu
$\operatorname{Pic}_{X/S}$ cũng riêng trên $S$, thì nhóm con xoắn của các
nhóm Néron--Severi của các thớ hình học của $f$ không đổi trên mỗi thành phần
liên thông của $S$ (điều này cũng hiển nhiên bằng phương pháp siêu việt khi
$f$ trơn và xạ ảnh). Ta lưu ý rằng việc áp dụng trực tiếp (2.5) cho phép chứng
minh tổng quát hơn rằng, nếu $\operatorname{Pic}_{X/S}$ riêng trên $S$ (chẳng
hạn $f:X\longrightarrow S$ trơn và xạ ảnh) và nếu $q$ là một số nguyên tố
khác các đặc số thặng dư của $S$, thì thành phần $q$-nguyên sơ của các nhóm
con xoắn Néron--Severi của các thớ hình học của $X/S$ không đổi trên mỗi thành
phần liên thông của $S$. Tuy nhiên, trong đặc số $p>0$, điều này không còn đúng
đối với thành phần $p$-nguyên sơ của phần xoắn. Dù vậy, vẫn có thể xảy ra rằng
hạng trên $k(s)$ của
\[
\operatorname{Pic}_{X_s/k(s)}^\tau/
\operatorname{Pic}_{X_s/k(s)}^{\circ\circ}=T_{X_s/k(s)}
\]
là hằng địa phương; khi $S$ rút gọn, có thể chứng minh rằng điều này tương
đương với việc $\operatorname{Pic}_{X/S}^{\circ\circ}$ tồn tại và
$\operatorname{Pic}_{X/S}^\tau$ phẳng trên $S$, và chỉ cần kiểm tra câu hỏi
khi $S$ là phổ của một vành định giá rời rạc. Đó là điều tôi đã kiểm tra trong
một vài ví dụ đã xét; nhưng vì phát biểu tương ứng với $S$ artin là sai (xem
các nhận xét (2.9) và (3.4)), có lẽ không nên nuôi quá nhiều ảo tưởng.

\fsection{4}{Định lý về tính hữu hạn đối với lược đồ Picard}

Cho $f:X\longrightarrow S$ là một cấu xạ xạ ảnh và phẳng sao cho
$\operatorname{Pic}_{X/S}$ tồn tại; khi đó, việc xét các ``đa thức Hilbert''
$Q\in\mathbf Q[t]$ cho phép phân tích $\operatorname{Pic}_{X/S}$ thành tổng
rời nhau của các tập mở $\operatorname{Pic}_{X/S}^Q$. Nếu không đặt thêm các
giả thiết, chẳng hạn bảo đảm $\operatorname{Pic}_{X/S}$ tách trên $S$, thì
nói chung các tập mở này không nhất thiết thuộc kiểu hữu hạn trên $S$; ta có
một phản ví dụ khi $X$ là ``một đường conic suy biến thành hai đường thẳng''.
Điều này
\src{242}
có thể vẫn đúng nếu $f$ khả ly và có các thớ hình học bất khả quy. Câu hỏi này
liên quan đến việc $\operatorname{Pic}_{X/S}^\tau$ có thuộc kiểu hữu hạn trên
$S$ hay không, điều có lẽ đúng mà không cần giả thiết nào về các thớ của
$X/S$. Khi $f:X\longrightarrow S$ trơn, lưu ý rằng
$\operatorname{Pic}_{X/S}^\tau$ nằm trong một trong các
$\operatorname{Pic}_{X/S}^Q$ (nhờ sự kiện rằng, trên một đa tạp xạ ảnh không
kỳ dị, tương đương ``theo xoắn'' tinh hơn tương đương số đối với các ước, và
thật ra, theo Matsusaka, là ``đồng nhất với'' tương đương số), nên nó thuộc
kiểu hữu hạn trên $S$ nếu các $\operatorname{Pic}_{X/S}^Q$ cũng vậy. Lưu ý
rằng các câu hỏi về tính hữu hạn vừa nêu vẫn có một ý nghĩa hiển nhiên ngay cả khi
không giả sử $\operatorname{Pic}_{X/S}$ tồn tại, vì chúng được diễn đạt bằng
cách nói rằng một số họ Môđun khả nghịch là ``bị chặn'' theo nghĩa của [3],
IV: với mỗi trường đóng đại số $k$, xét các lược đồ con nguyên của
$\mathbf P_k^r$, có chiều $n$ và bậc $d$, cùng các Môđun khả nghịch trên các
tiền lược đồ như thế, có đa thức Hilbert $Q$ (trong đó $r,n,d,Q$ được cho
trước); cần chứng minh rằng họ các Môđun này (được xem là các Môđun coherent
trên các $\mathbf P_k^r$) bị chặn, tức là có thể được tham số hóa bởi một lược
đồ kiểu hữu hạn trên $\mathbf Z$.

Dùng phương pháp của Matsusaka [6], một chứng minh khá kỹ thuật (viện đến các
``tiêu chuẩn tương đương'' dưới một dạng thích hợp) cho phép trả lời khẳng định
khi ta giới hạn ở các đa tạp con không kỳ dị của $\mathbf P_k^r$. Chính xác
hơn, ta được kết quả sau:

\result{ĐỊNH LÝ}{4.1} Cho: $X\longrightarrow S$ là một cấu xạ trơn và xạ ảnh có các thớ liên thông hình học, với $S$
noether, $\mathcal O_X(1)$ là một Môđun rất rộng trên $X$, tương đối với $S$,
$E$ là một tập con của $\operatorname{Pic}_{X/S}$ tương ứng với một họ
$(\mathcal L_i)$ các Môđun khả nghịch trên các thớ hình học
$\overline X_{s_i}$ của $X/S$, $D_i$ là một ước (không nhất thiết dương) trên
$\overline X_{s_i}$ xác định $\mathcal L_i$,
\[
a_n^{(i)}t^n+\cdots+a_0^{(i)}
\]
là đa thức Hilbert của $\mathcal L_i$, $\xi=\xi_i$ là một ước xác định
$\mathcal O_X(1)$, tức là một tiết diện siêu phẳng. Các điều kiện sau là tương
đương:

\noindent a. $Q$ là tựa
compact, tức là nằm trong một tập mở kiểu hữu hạn trên $S$, tức là họ các
$(\mathcal L_i)$ bị chặn.

\noindent b. $E$ nằm trong hợp của một số hữu hạn các tập
$\operatorname{Pic}_{X/S}^Q$, $Q\in\mathbf Q[t]$, tức là tập các đa thức
Hilbert của các $\mathcal L_i$ là hữu hạn.

\noindent b'. (Nếu các thớ của $X/S$ đều có cùng chiều $n$.) Hai hệ số
$a_{n-1}^{(i)}$ và $a_{n-2}^{(i)}$ của các đa thức Hilbert của
$\mathcal L_i$ nằm trong một tập hữu hạn.
\src{243}
\noindent b''. (Nếu các thớ của $X/S$ đều có cùng chiều $n$.) Các số nguyên
\[
\xi^{n-1}D_i\qquad\text{và}\qquad\xi^{n-2}D_i^2
\]
(các bậc xạ ảnh của $D_i$ và $D_i^2$) bị chặn trên.

\result{HỆ QUẢ}{4.2} Cho $f:X\longrightarrow S$ là một cấu xạ trơn và xạ ảnh
có các thớ liên thông hình học; khi đó các lược đồ
$\operatorname{Pic}_{X/S}^Q$ và $\operatorname{Pic}_{X/S}^\tau$
($Q\in\mathbf Q[t]$) là xạ ảnh trên $S$.

Vì chúng thuộc kiểu hữu hạn trên $S$ theo (4.1), ta có thể áp dụng (2.4).

\section*{TÀI LIỆU THAM KHẢO}
\addcontentsline{toc}{section}{Tài liệu tham khảo}

\begin{description}
\item[{[1]}] CHEVALLEY (Claude). -- Sur la théorie de Picard,
\emph{Amer. J. of Math.}, t. 82, 1960, tr. 435--490.

\item[{[2]}] GROTHENDIECK (A.) và DIEUDONNÉ (J.). -- Éléments de géométrie
algébrique, Chương I và các chương tiếp theo. -- Paris, Presses universitaires
de France, 1960--1961 (Institut des Hautes Études Scientifiques, 4, 8, 11,
\ldots).

\item[{[3]}] GROTHENDIECK (Alexander). -- Technique de descente et théorèmes
d'existence en géométrie algébrique, I--V, \emph{Séminaire Bourbaki}, t. 12,
1959/60, n\textsuperscript{o} 190, 29 tr. và n\textsuperscript{o} 195, 22 tr.;
t. 13, 1960/61, n\textsuperscript{o} 212, 20 tr. và
n\textsuperscript{o} 221, 28 tr.; t. 14, 1961/62,
n\textsuperscript{o} 232, 19 tr.

\item[{[4]}] GROTHENDIECK (Alexander). -- Séminaire de géométrie algébrique,
1960/61. -- Paris, Institut des Hautes Études Scientifiques, 1961.

\item[{[5]}] GROTHENDIECK (Alexander). -- Séminaire de géométrie algébrique,
1961/62, do một nhóm thính giả biên soạn. -- Paris, Institut des Hautes Études
Scientifiques (sắp xuất bản).

\item[{[6]}] MATSUSAKA (T.). -- The criteria for algebraic equivalence and the
torsion group, \emph{Amer. J. Math.}, t. 79, 1957, tr. 53--66.

\item[{[7]}] MUMFORD (D.) và TATE (J.). -- Séminaire de géométrie algébrique,
Harvard University, học kỳ mùa xuân 1962 (sắp xuất bản).

\item[{[8]}] SERRE (Jean-Pierre). -- Sur la topologie des variétés algébriques
en caractéristique $p$, \emph{Symposium internacional de topología algebraica}
[1956, Mexico], tr. 24--53. -- Mexico, Universidad nacional autónoma, 1958.
\end{description}

\par

\clearpage
\chapter*{Đính chính và bổ sung cho Báo cáo số~236}
\addcontentsline{toc}{chapter}{Đính chính và bổ sung cho Báo cáo số~236}
\sid{e236}
\markboth{Đính chính 236}{Đính chính 236}
\src{303}

\noindent\textbf{8. TDTE VI: Các lược đồ Picard. Các tính chất tổng quát.}

\begin{center}
[Séminaire Bourbaki, tập 14, 1961/62, số 236, 23 trang.]
\end{center}

\noindent\textbf{Trang 12, dòng 16.} --- Thay ``đường thẳng afin'' bằng
``đường thẳng afin bỏ gốc tọa độ''.

\noindent\textbf{Dòng 18.} --- Thay ``$X[t]$'' bằng
``$X[t,t^{-1}]$''.

\noindent\textbf{Trang 16, Mệnh đề 3.1, dòng thứ 6.} --- Thay ``Bây giờ giả
thiết ...'' bằng ``Bây giờ, từ giả thiết ...''.

\noindent\textbf{Trang 21, §~4: Định lý hữu hạn cho lược đồ Picard.} --- Các
vấn đề hữu hạn thuộc loại được nêu trong mục này hầu như đã được giải quyết
hoàn toàn kể từ khi báo cáo hiện tại được biên soạn. Chúng tôi nêu các sự kiện
chính hiện đã biết theo hướng này. Để đơn giản hóa các phát biểu, chúng tôi
ngầm giả sử mọi tiền lược đồ Picard xuất hiện trong các phát biểu đều tồn tại,
mặc dù một sửa đổi hiển nhiên của các phát biểu cho phép loại bỏ mọi giả
thiết tồn tại. Trong phần sau, $S$ chỉ một lược đồ Noether, còn $X,Y$ là các
lược đồ proper trên $S$.

\medskip
\noindent(i) Cho $f:X\longrightarrow Y$ là một cấu xạ toàn ánh. Khi đó
\[
f^*: \operatorname{Pic}_{Y/S}\longrightarrow\operatorname{Pic}_{X/S}
\]
là một cấu xạ kiểu hữu hạn.

\src{304}
\noindent(ii) Giả sử $Y$ xạ ảnh trên $S$, được trang bị một môđun khả nghịch
ample tương đối trên $S$, và cho $X$ là lược đồ không của một tiết diện tùy ý
của môđun này, $f:X\longrightarrow Y$ là phép nhúng chính tắc. Cuối cùng, giả
sử các thành phần bất khả quy của các thớ của $Y/S$ có chiều $\geq3$. Khi đó
\[
f^*: \operatorname{Pic}_{Y/S}\longrightarrow\operatorname{Pic}_{X/S}
\]
là kiểu hữu hạn.

\noindent(iii) Giả sử $X$ xạ ảnh trên $S$ và mọi thớ hình học của nó đều
nguyên, có chiều $n$. Cho $\mathcal O_X(1)$ là một môđun khả nghịch ample trên
$X$, cho phép xác định các đa thức Hilbert. Cho $M$ là một tập con của
$\operatorname{Pic}_{X/S}$; khi đó $M$ tựa compact khi và chỉ khi, trong các
đa thức Hilbert
\[
a_0x^n+a_1x^{n-1}+a_2x^{n-2}+\cdots
\]
của các phần tử thuộc $M$, các hệ số $a_1$ và $a_2$ vẫn bị chặn.

\noindent(iv) Với mọi số nguyên $n\neq0$, đồng cấu lũy thừa bậc $n$ trong
tiền lược đồ nhóm $\operatorname{Pic}_{X/S}$ là một cấu xạ kiểu hữu hạn.

Lưu ý rằng (i) và (ii) cũng có nghĩa là (dưới các giả thiết tương ứng trong
hai phát biểu ấy) một tập con $M$ của $\operatorname{Pic}_{Y/S}$ tựa compact
khi và chỉ khi ảnh của nó trong $\operatorname{Pic}_{X/S}$ tựa compact. Suy
ra một môđun khả nghịch $\mathcal L$ trên $Y$ là $\tau$-tương đương với không
khi và chỉ khi ảnh ngược của nó trên $X$ có tính chất ấy; nói cách khác,
$\operatorname{Pic}_{Y/S}^\tau$ là ảnh ngược của
$\operatorname{Pic}_{X/S}^\tau$. Đặc biệt, để chứng minh tiền lược đồ thứ
nhất thuộc kiểu hữu hạn trên $S$, chỉ cần chứng minh điều đó cho tiền lược đồ
thứ hai, vì $\operatorname{Pic}_{Y/S}\longrightarrow
\operatorname{Pic}_{X/S}$ thuộc kiểu hữu hạn. Khi đó, dùng (i), Bổ đề Chow và
(iii), ta thu được:

\medskip
\noindent(v) $\operatorname{Pic}_{X/S}^\tau$ thuộc kiểu hữu hạn trên $S$.

Tổng quát hơn, việc kết hợp (i) cho một cấu xạ hữu hạn với (ii) cho phép quy
hầu hết các vấn đề hữu hạn về trường hợp $X/S$ có các thớ hình học nguyên và
chuẩn, chiều $\leq2$; thường thì, bằng cách áp dụng (i) cho một cấu xạ toàn
ánh không nhất thiết hữu hạn và dùng phép giải kỳ dị của các mặt đại số (do
Abhyankar thực hiện trong đặc số tùy ý), ta thậm chí quy được về trường hợp
$X/S$ trơn, do đó có các thớ hình học không kỳ dị chiều 2. Kết hợp điều này
với (v) và bất đẳng thức Picard--Igusa chặn hạng của nhóm Néron--Severi của
một mặt xạ ảnh không kỳ dị, ta chứng minh được sự khái quát hóa sau đây của
định lý hữu hạn Néron:

\medskip
\noindent(vi) Cho $X/S$ proper trên $S$ và ngoài ra tùy ý. Khi đó các nhóm
Néron--Severi
\[
\operatorname{Pic}_{X_i/k_i}/\operatorname{Pic}_{X_i/k_i}^\circ
\]
của các thớ hình học $X_i/k_i$ của $X/S$

\src{305}
đều thuộc kiểu hữu hạn, còn hạng của chúng và cấp của các nhóm con xoắn của
chúng vẫn bị chặn trên.

Cùng một phương pháp quy về trường hợp các mặt không kỳ dị, kết hợp với các
định lý đã biết trong trường hợp đó (cụ thể là định lý hữu hạn Néron và tính
không suy biến của dạng giao trên nhóm Néron--Severi), kéo theo:

\medskip
\noindent(vii) Cho $X$ là một lược đồ proper trên một trường đóng đại số. Khi
đó tồn tại hữu hạn đường cong đóng nguyên $C_i$ ($1\leq i\leq r$) trong $X$
có tính chất sau: với mọi tập con $M$ của $\operatorname{Pic}_{X/k}$, $M$ tựa
compact khi và chỉ khi các số nguyên
\[
\deg_{C_i'}\mathcal L\qquad(\mathcal L\in M)
\]
vẫn bị chặn (trong đó $C_i'$ chỉ chuẩn hóa của $C_i$).

Ở đây có thể lấy $r$ bằng hạng của nhóm Néron--Severi. Một khi đã biết nhóm
này thuộc kiểu hữu hạn, (vii) quy về sự kiện rằng các dạng tuyến tính trên
nhóm Néron--Severi được xác định bởi các đường cong $C$ trong $X$ chỉ đồng
thời triệt tiêu trên các phần tử xoắn của nhóm Néron--Severi. Trong trường hợp
$X$ xạ ảnh và không kỳ dị, kết quả này, cũng như (v), là do Matsusaka. Dùng
(vii), ta dễ dàng thu được đặc trưng sau đây của các môđun khả nghịch
$\tau$-tương đương với không trên $X$:

\medskip
\noindent(viii) Cho $X/k$ là một lược đồ proper trên một trường và
$\mathcal L$ là một môđun khả nghịch trên $X$. Các điều kiện sau là tương
đương:

\noindent a. $\mathcal L$ là $\tau$-tương đương với không;

\noindent b. Với mọi môđun coherent $F$, ta có
\[
\chi(F\otimes\mathcal L)=\chi(F),
\]
trong đó $\chi$ chỉ đặc trưng Euler--Poincaré;

\noindent b'. Như (b), nhưng với $F=\mathcal O_Y$, trong đó $Y$ là một lược
đồ con đóng nguyên chiều 1 trong $X$;

\noindent c. Với mọi $Y$ như trên, ký hiệu $Y'$ là đường cong chuẩn hóa, ta
có
\[
\deg\mathcal L_{Y'}=0.
\]

Nếu $X/k$ xạ ảnh và được trang bị một môđun khả nghịch ample
$\mathcal O_X(1)$, thì các điều kiện trước cũng tương đương với các điều kiện
sau:

\noindent d. Với mọi số nguyên $m$, $\mathcal L^{\otimes m}(1)$ là ample.

\noindent e. (Nếu $X$ nguyên.) Với mọi cặp số nguyên $m,n$, ta có
\[
\chi(\mathcal L^{\otimes m}(n))=\chi(\mathcal O_X(n))
\]

\src{306}
(nghĩa là (b) đúng khi lấy $F=\mathcal L^{\otimes m}(n)$).

Để chứng minh tính đủ của điều kiện cuối này, lưu ý rằng nó có nghĩa là các
đa thức Hilbert của $\mathcal L^{\otimes m}$ đều bằng nhau; do đó, theo tiêu
chuẩn Mumford (iii), các $\mathcal L^{\otimes m}$ vẫn nằm trong một tập tựa
compact của $\operatorname{Pic}_{X/k}$, nghĩa là ta có (a). Các điều kiện (b),
(b') và (c) phải được xem là những biến thể (trên một lược đồ proper tùy ý)
của khái niệm tương đương số, vốn thường được định nghĩa trên các đa tạp xạ
ảnh không kỳ dị. Đối với các đa tạp ấy, sự tương đương giữa (a) và (c) dĩ
nhiên đã được biết rõ (Matsusaka).

Tiêu chuẩn (e) cũng kéo theo kết quả sau:

\medskip
\noindent(ix) Cho $f:X\longrightarrow S$ là một cấu xạ xạ ảnh và phẳng có
các thớ hình học nguyên. Khi đó $\operatorname{Pic}_{X/S}^\tau$ vừa mở vừa
đóng trong $\operatorname{Pic}_{X/S}$.

Chúng tôi chỉ giới hạn ở một vài bình luận về chứng minh của các kết quả then
chốt (i), (ii), (iii) (kết quả (iv) hơi tách biệt với các kết quả khác và
được chứng minh chỉ bằng cách dùng (i) cho các cấu xạ hữu hạn, toàn ánh và
radiciel, cụ thể là cho một cấu xạ Frobenius). Đối với (i), ta sử dụng một
cách thiết yếu các ý tưởng của descent không phẳng (xem đặc biệt TDTE I,
trang 9). Do chỉ nhắm đến các kết quả hữu hạn, việc thiếu các tiêu chuẩn hữu
hiệu cho dữ liệu descent là vô hại. Mặt khác, Mumford gần đây đã chứng minh
một dạng hơi yếu hơn của (iii), cụ thể là tiêu chuẩn có sự tham gia của mọi
hệ số trong đa thức Hilbert. Lập luận của ông cực kỳ đơn giản và được gợi ý từ
chứng minh của một tiêu chuẩn ample của Nakai (do Nakai phát biểu cho các mặt
không kỳ dị và được Mumford khái quát hóa cho các cấu xạ xạ ảnh tùy ý). Tuy
nhiên, theo tôi, lập luận này chỉ đúng khi đặt thêm một hạn chế nhẹ lên các thớ
của $X/S$, tức điều kiện $(S_2)$ của Serre, vốn được thỏa mãn nếu các thớ hình
học là chuẩn. Sau đó ta dùng tiêu chuẩn bị hạn chế này trong chứng minh (ii):
tiêu chuẩn (i) thực sự cho phép quy về trường hợp $Y/S$ phẳng, có các thớ hình
học nguyên và chuẩn; rồi áp dụng tiêu chuẩn Mumford, ta dễ dàng quy về trường
hợp $X/S$ thỏa mãn các điều kiện tương tự. Từ giả thiết về chiều suy ra các
thớ hình học của $Y$ và $X$ có độ sâu $\geq2$ tại các điểm đóng của chúng;
điều này cho phép áp dụng các ``tiêu chuẩn tương đương'' dưới dạng được nêu
trong SGA 1962 (Báo cáo XII) và hoàn tất chứng minh của (ii). Một khi đã có
(i) và (ii), không khó để loại bỏ khỏi tiêu chuẩn Mumford mọi giả thiết về
tính chuẩn của các thớ và cũng chứng minh tiêu chuẩn ấy dưới dạng mạnh hơn
được nêu trong (iii).

\src{307}
Cuối cùng, lưu ý rằng chứng minh của (i) và (ii) cũng cho thấy rằng, trong
trường hợp $S$ là phổ của một trường $k$, cấu xạ
\[
\operatorname{Pic}_{Y/k}\longrightarrow\operatorname{Pic}_{X/k}
\]
là afin (chứ không chỉ thuộc kiểu hữu hạn).

\par

\unit{Bình luận năm 1962}
\sid{com}
\markboth{Bình luận năm 1962}{Bình luận năm 1962}
\src{297}
\seminar{Năm thứ 14, 1961/62, Phần bổ sung\\Tháng 5 năm 1962}
\paperhead{NỀN TẢNG CỦA HÌNH HỌC ĐẠI SỐ.\\BÌNH LUẬN}%
  {bởi Alexander GROTHENDIECK}

\section*{Lời báo trước}
\addcontentsline{toc}{section}{Lời báo trước}

Từ năm 1957 đến năm 1962, chúng tôi đã trình bày tại Séminaire Bourbaki tám
báo cáo liên quan đến nền tảng của hình học đại số. Ngoại trừ báo cáo đầu
tiên, các báo cáo này được viết trong khuôn khổ lý thuyết lược đồ. Mọi kết quả
được phát biểu sẽ tìm thấy vị trí của mình trong bộ ``Éléments de géométrie
algébrique'' của Jean DIEUDONNÉ và Alexander GROTHENDIECK. Tuy nhiên, nội dung
của không một báo cáo nào trong số này hiện được bao quát bởi bất kỳ chương
nào đã xuất bản hoặc đang biên soạn, cũng không bởi bất kỳ cuốn sách hay bài
báo nào khác, và hẳn trong vài năm tới vẫn chưa được bao quát. Đó là lý do
chính khiến chúng tôi tập hợp các báo cáo này, qua đó cung cấp cho người đọc
một số khái niệm và kết quả then chốt của lý thuyết lược đồ trong khi chờ một
bản trình bày hoàn chỉnh. Hơn nữa, việc đọc các văn bản này sẽ giúp người đọc
nhanh chóng làm quen với các kết quả và khái niệm nói trên mà không bị cản trở
bởi những chi tiết tất yếu rườm rà của một chuyên luận có hệ thống, đồng thời
cung cấp các động cơ thiết yếu cho việc nghiên cứu chuyên luận như vậy.

Để tiện cho người đọc, chúng tôi tập hợp tại đây một số bình luận và đính
chính theo từng báo cáo; chúng đặc biệt ghi nhận những tiến bộ đạt được kể từ
khi văn bản được biên soạn và chỉ ra một số tài liệu tham khảo bổ sung.

Một số kết quả xuất hiện trong các báo cáo này đã được nghiên cứu chi tiết tại
Séminaire de géométrie algébrique của IHES, cũng như trong hai seminar liên
tiếp ở Harvard năm 1961/62, seminar thứ nhất do chính tôi trình bày và seminar
thứ hai do MUMFORD--TATE trình bày; các ghi chú in rô-nê-ô cho hai seminar ấy
(do ông LICHTENBAUM biên soạn) hiện đang được chuẩn bị.

\section*{Các chữ viết tắt}
\addcontentsline{toc}{section}{Các chữ viết tắt}

\begin{description}
\item[SGA :] Séminaire de géométrie algébrique, Institut des Hautes Études
Scientifiques, 1960--1962.

\src{298}
\item[SHG :] Séminaire GROTHENDIECK, Harvard University, 1961/62.

\item[SHMT :] Séminaire MUMFORD--TATE, Harvard University, 1961/62.

\item[EGA :] GROTHENDIECK (A.) và DIEUDONNÉ (J.). --- Éléments de géométrie
algébrique. --- Paris, Presses universitaires de France, 1960, 1961, …
(Institut des Hautes Études Scientifiques, Publications mathématiques, 4, 8,
11, …).

\item[TDTE :] GROTHENDIECK (Alexander). --- Technique de descente et théorèmes
d'existence en géométrie algébrique, I--VI, \emph{Séminaire Bourbaki}, tập 12,
1959/60, số 190 và 195; tập 13, 1960/61, số 212 và 221; tập 14, 1961/62, số
232 và 236.
\end{description}

\par


\end{document}
